\documentclass[11pt,a4paper]{article}
\usepackage[T1]{fontenc}
\usepackage[utf8]{inputenc}
\usepackage{lmodern,amsmath,amssymb,amsthm,mathtools}
\usepackage[margin=26mm,headheight=15pt]{geometry}
\usepackage{microtype,booktabs,longtable,array,enumitem,needspace}
\usepackage[dvipsnames]{xcolor}
\usepackage{listings,fancyhdr,xurl}
\usepackage[colorlinks=true,linkcolor=MidnightBlue,citecolor=MidnightBlue,urlcolor=MidnightBlue,breaklinks=true]{hyperref}
\hypersetup{pdftitle={Prime-Representing Constants and Related Structures},pdfauthor={Keio},pdfsubject={Final consolidated research edition: prime constants, algebraic power sequences, and modular obstructions},bookmarksnumbered=true}
\newtheorem{theorem}{Theorem}[section]
\newtheorem{lemma}[theorem]{Lemma}
\newtheorem{proposition}[theorem]{Proposition}
\newtheorem{corollary}[theorem]{Corollary}
\newtheorem{criterion}[theorem]{Criterion}
\theoremstyle{definition}
\newtheorem{definition}[theorem]{Definition}
\newtheorem{example}[theorem]{Example}
\theoremstyle{remark}
\newtheorem{remark}[theorem]{Remark}
\numberwithin{equation}{section}
\newcommand{\N}{\mathbb N}\newcommand{\Z}{\mathbb Z}
\newcommand{\Q}{\mathbb Q}\newcommand{\R}{\mathbb R}\newcommand{\C}{\mathbb C}
\DeclareMathOperator{\Tr}{Tr}\DeclareMathOperator{\Norm}{N}
\DeclareMathOperator{\ord}{ord}\DeclareMathOperator{\lcm}{lcm}
\DeclareMathOperator{\rad}{rad}\DeclareMathOperator{\agcd}{agcd}
\DeclareMathOperator{\asdim}{asdim}\DeclareMathOperator{\np}{nextprime}
\newcommand{\floor}[1]{\left\lfloor#1\right\rfloor}
\newcommand{\ceil}[1]{\left\lceil#1\right\rceil}
\setlist{itemsep=2pt,topsep=4pt}
\setlength{\emergencystretch}{3em}
\setlength{\parindent}{1.2em}
\setlength{\parskip}{2pt}
\raggedbottom
\pagestyle{fancy}\fancyhf{}
\fancyhead[L]{\small Prime-representing constants and related structures}
\fancyhead[R]{\small Keio}
\fancyfoot[C]{\thepage}
\renewcommand{\headrulewidth}{0.3pt}
\lstset{language=Python,basicstyle=\ttfamily\footnotesize,
 keywordstyle=\bfseries,commentstyle=\color{black!65},
 stringstyle=\color{black!80},showstringspaces=false,columns=fullflexible,
 keepspaces=true,breaklines=true,breakatwhitespace=false,
 numbers=left,numberstyle=\tiny\color{black!50},numbersep=8pt,
 xleftmargin=14pt,frame=none,aboveskip=10pt,belowskip=10pt,
 upquote=true,tabsize=4}
\setcounter{tocdepth}{1}
\title{\bfseries Prime-Representing Constants\\and Related Structures\\[7pt]
\large Encodings, Arithmetic, Geometry,\\and Modular Obstructions}
\author{Keio}
\date{18 October 2026\\\small Final consolidated research edition}
\begin{document}
\maketitle
\begin{abstract}
We study constants whose prescribed powers have prime integer parts,
together with related algebraic power sequences. Exact prime cells yield
certified minima, effective independent perfect sets, and reproducible
large-prime certificates. Geometric arguments give full Hausdorff
dimension for unbounded power schedules and iterated-exponential prime
orbits. Modular return and local minimality give transcendence criteria
for perturbed exponential and factorial schedules. Further results concern
polynomial growth constants, ceiling endpoints, optimal finite prime
covers, and local computability. We extend support-sensitive independence
to all left boundary minima, obtain density-one cubic compositeness
results at fixed norm, and lower a Type-C transcendence threshold.
The final arithmetic analysis proves infinitely many composite outputs
for non-totally-real cubic Pisot numbers, exact parity-dependent
$3$-adic trace depths in a cubic-unit branch, and a short-interval prime
deficit forced by a hypothetical rational square-clock minimum.
An exact finite-state counterexample limits a proposed prime-divisor
return mechanism. Related sections treat polynomial patterns,
palindromes, quotient sums, and dimension profiles. Proofs and standalone
integer verification programs are included in one self-contained source.
The unconditional transcendence of Mills' constant, the irrationality of
$\xi(2^k)$, the universal totally real cubic Pisot problem, and the
twin-prime conjecture are not resolved.
\end{abstract}
\noindent\textbf{Keywords:} prime-representing constants; Pisot numbers;
transcendence; Hausdorff dimension; modular dynamics; finite prime covers.
\clearpage
\begingroup\small\setlength{\parskip}{0pt}
\tableofcontents
\endgroup
\clearpage

\renewcommand{\theHtable}{intro.\arabic{table}}
\section{Objects, principal results, and logical scope}
\label{sec:overview}
For a sequence $C=(C_k)$ of positive real exponents, write
\[
 W(C)=\{A>1:\lfloor A^{C_k}\rfloor\text{ is prime for every }k\ge1\},
 \qquad \xi(C)=\min W(C)
\]
whenever the set is nonempty. Write $c_1=C_1$ and
$c_{k+1}=C_{k+1}/C_k$. The classical Mills choice is $C_k=3^k$
\cite{Mills1947}; the arithmetic and geometry of more general choices
are developed in Saito and Saito--Takeda
\cite{SaitoDimension2021,SaitoTakeda2022,SaitoMills,SaitoVariants}.
This manuscript combines exact constructions with extensions of those
methods and with related problems from the same research programme.

\subsection{The objects that must be distinguished}
The set $W(C)$, its least element, a selected branch inside it, and the
least element after deleting initial conditions are different objects.
The following conventions remain in force throughout.
\begin{longtable}{@{}p{0.25\textwidth}p{0.70\textwidth}@{}}
\toprule
Object & Meaning and scope\\ \midrule
$W(C)$ & All permitted bases for one fixed exponent sequence.
For successive ratios at least two its Hausdorff dimension is one,
by Theorem~\ref{fol:thm:unbounded-dimension} and
Corollary~\ref{fol:cor:arbitrary-initial-scale}.\\[4pt]
$\xi(C)$ & The global least permitted base. Statements about this
specific number require a minimality argument.\\[4pt]
Local left minimum & A point $A\in W(C)$ with an interval immediately
to its left containing no other point of $W(C)$. The localized
transcendence criteria concern these points.\\[4pt]
Selected perfect set $K$ & A constructed subset of $W(C)$. Some such
sets have upper box dimension zero and algebraically independent points.
Their dimension need not equal the dimension of $W(C)$.\\[4pt]
Tail minimum $\xi_N$ & The least base subject only to the conditions
with index at least $N$. It generally differs from $\xi(C)$.\\[4pt]
Adaptive exponent sequence & A sequence chosen during a construction.
Its flexibility does not imply the same result for a previously fixed
sequence or its global minimum.\\
\bottomrule
\end{longtable}

The ordinary factorial clock $C_k=k!$, the perturbations $k!+k$ and
$k!+k^2$, and the sparse recursion $E_{n+1}=E_nE_n!$ are also distinct.
The enormous decimal-prefix estimates in Part~\ref{part:encoding}
belong to the explicitly defined sparse schedules there. They are not
asserted for every factorial perturbation.

\subsection{Main conclusions and their dependencies}
Part~\ref{part:encoding} begins with the finite exponent list
$(4,6,7,12)$ and proves its exact minimum $29^{1/12}$.
A certified prime seed extends this finite problem to explicit infinite
schedules. The construction supplies algebraically independent perfect
sets, preserves Turing degrees, and gives deterministic certificates for
a $1966$-digit prime and two finite twin-prime pairs. The twin-output
window equivalence is kept separate from a proof that its windows are
nonempty infinitely often. Spectral applications transfer the constructed
algebraic independence to finite matrices.

Part~\ref{part:geometry} proves full dimension for arbitrary unbounded
multiplier sequences with successive ratios at least two, and for Wright
prime orbits in every real base $b>1$. It then establishes local dimension,
support-sensitive separation, and relative independence of logarithms and
finite exponential towers. The complete independent binary tree in
Part~\ref{part:encoding} contains its greedy minimum. A later analytic
selection need not retain that minimum; in particular it cannot retain
it when independence is required over a field already containing it.
Likewise, the field generated by a selected thin set is distinct from
the ring generated by the full spectrum.

Part~\ref{part:arithmetic} develops modular return and local minimality.
In particular, the full minima for
\[
 3^k-k,\quad 3^k-2k,\quad 3^k-2k+1,\quad k!+k,\quad k!+k^2
 \qquad(k\ge1)
\]
are transcendental; these are consequences of Saito's criteria and the
arithmetic verifications given here, not a new proof of those external
criteria. The more general family $P(a^n)+Q(n)$ is treated with its exact
growth and nonvanishing assumptions; its general conclusion concerns
sufficiently late minima. For the ordinary Mills constant $A$, the
degree conclusion is conditional:
\[
 A\in\overline{\mathbb Q}\quad\Longrightarrow\quad
 [\mathbb Q(A):\mathbb Q]=3^{t+1},\qquad
 \operatorname{MinPoly}_A(X)=f(X^{3^t})
\]
for some integer $t\ge0$ and a monic irreducible cubic $f\in\mathbb Z[X]$.
It is a necessary condition, not a contradiction. The same distinction
applies to the cubic prime-divisor sieve and bounded-degree certificates.

Part~\ref{part:patterns} treats selected related questions in additive
combinatorics and digital structure. Its statements are collected with
their exact scopes in the following table.
\begin{longtable}{@{}p{0.31\textwidth}p{0.64\textwidth}@{}}
\toprule
Source question or extension & Conclusion proved here\\ \midrule
Saito, dimension Questions~1.4--1.5 in arXiv:2102.04038v1
& Full dimension for unbounded power schedules and for Wright prime
orbits; Theorems~\ref{fol:thm:unbounded-dimension} and
\ref{fol:thm:wright-full-dimension}.\\[4pt]
Pongsriiam--Saito, palindrome Question~27
& The coloring and positive-relative-density assertions hold exactly
for integer bases $b\ge3$; Theorem~\ref{sur:thm:palindrome-classification}.\\[4pt]
Saito--Yoshida, polynomial-pattern Question~6.1
& The stated order of magnitude holds under positive lower density;
Theorem~\ref{sur:thm:ps-density}.\\[4pt]
Saito--Yoshida, three-term waiting-time conjecture
& $L_{\alpha,3,r}(x)=O_{\alpha,r}(x^{1-\alpha/2})$ for
$1<\alpha<2$ and fixed $r$; Theorem~\ref{sur:thm:ps-waiting}.\\[4pt]
Positive balanced equations
& An exceptional-dimension upper bound and almost-everywhere finiteness;
Theorem~\ref{sur:thm:ps-balanced}. No exact dimension is claimed.\\[4pt]
Normality and lattice dimension profiles
& A sparse analytic equivalence criterion and all permissible numerical
lattice profiles. No specified algebraic irrational is proved normal.\\
\bottomrule
\end{longtable}

\subsection{Attribution, notation, and verification}
Question numbers always refer to the source version explicitly identified
at the point of use. Some consequences settle the literal statements of
questions in those versions; historical priority has not been established
by a comprehensive literature search. The paper is a research manuscript,
not a claim of completed external peer review. Existing theorems used as
inputs are attributed where needed. The conditional results are not
counted as unconditional resolutions.

All logarithms are natural unless a base is displayed. A prime is a
positive rational prime. Algebraic independence means that every nonzero
polynomial in finitely many distinct specified entries is nonzero.
Hausdorff dimension is denoted by $\dim_H$; computability is understood
in the Cauchy-name sense and carries no efficiency assertion. Arithmetic
progressions are nonconstant unless stated otherwise. In the final
pattern sections, ``almost every'' means Lebesgue almost every.
Appendix~\ref{enc:app:verification-code} gives executable exact integer
checks for the finite numerical certificates; they were rerun during
consolidation. Proofs of infinite assertions remain separate from these
finite checks.


\subsection{Extensions incorporated in this edition}
Alkauskas--Dubickas constructed transcendental prime-representing
constants for suitable rapidly growing schedules
\cite{AlkauskasDubickas2004}. Saito--Takeda's results already include
transcendence of a least factorial constant \cite{SaitoTakeda2022}.
These antecedents distinguish existence of transcendental members from
arithmetic conclusions about a specified least member.
Parts~\ref{part:classifications} and \ref{part:extensions} incorporate the
two subsequent research manuscripts, preserving their substantive proofs
and finite certificates while unifying the notation and bibliography.
Part~\ref{part:classifications} treats polynomial growth, factorial
return, ceiling geometry, finite covers, and computability under integer
successive multipliers at least two with infinitely many at least three.
Part~\ref{part:extensions} develops left-boundary independence,
fixed-norm cubic sieves, square-clock criteria, non-reversible windows,
the refined Type-C cutoff, and the two density applications.
Part~\ref{part:focused} incorporates the subsequent focused analysis of
the three principal remaining arithmetic problems. In particular,
the non-totally-real cubic theorem is included with its full proof;
no irrational-angle equidistribution hypothesis is used.

\begin{center}
\begin{tabular}{@{}p{0.40\linewidth}p{0.54\linewidth}@{}}
\toprule
Target & Status of this manuscript\\ \midrule
Unconditional transcendence of Mills' least constant
& Necessary arithmetic conditions and finite exclusions; not proved.\\[3pt]
Irrationality of $\xi(2^k)$
& Exact conditional prime-deficit theorem; not proved.\\[3pt]
Infinitely many composite cubic Pisot outputs at $3^n$
& Non-totally-real case proved; individual totally real exceptions remain.\\[3pt]
Infinite twin primes
& Exact output windows and finite certificates only.\\
\bottomrule
\end{tabular}
\end{center}
The word ``final'' identifies the consolidated edition of this research
programme. It does not certify external peer review or historical priority.
All conclusions retain the hypotheses stated in their individual results.

\setcounter{table}{0}
\renewcommand{\theHtable}{\arabic{table}}

\clearpage
\part{Exact prime encodings and arithmetic certificates}\label{part:encoding}
This part constructs specified constants and selected sets from exact prime cells. It establishes their minimality or branch status before deriving arithmetic and computational consequences.

\section{Prime cells and the exact seed minimum}
The constructions below distinguish the global least admissible constant
for a fixed exponent sequence, a selected branch inside an admissible
cell, and a constant for a modified exponent sequence. They build on the
prime-representing framework of Mills~\cite{Mills1947} and the algebraic
independence results of Saito and Takeda~\cite{SaitoTakeda2022}.
The explicit constructions and the consequences proved here do not imply
infinitely many twin primes or a general primality-testing speedup.

Write $\mathbb P$ for the positive primes and $\np(x)$ for the least prime strictly greater than the real number $x$. For a positive integer $E$ and a prime $s$, the half-open interval
\[
J_E(s)=[s^{1/E},(s+1)^{1/E})
\]
is exactly the set of positive bases whose $E$th power has integer part $s$. We also use its closed hull $I_E(s)=[s^{1/E},(s+1)^{1/E}]$. The upper endpoint of $I_E(s)$ must not be treated as admissible without an additional argument.

\begin{lemma}[Strict nesting]\label{enc:lem:nesting}
Let $E_{n+1}=E_nR_n$, where $R_n$ is an integer greater than one. Suppose primes $s_n$ satisfy
\[
s_n^{R_n}<s_{n+1},\qquad s_{n+1}+1<(s_n+1)^{R_n}.
\]
Then $I_{E_{n+1}}(s_{n+1})$ is contained in the interior of $I_{E_n}(s_n)$. If $E_n\to\infty$, the intervals have a unique common point $\alpha$, and $\lfloor\alpha^{E_n}\rfloor=s_n$ for every $n$.
\end{lemma}
\begin{proof}
Taking positive $E_{n+1}$th roots proves strict containment. The mean value theorem gives $|I_E(s)|\le1/E$. Hence the nested compact intervals have a unique intersection. The next interval is strictly inside each given parent, so the intersection point avoids both endpoints of every parent. This proves the floor identities.
\end{proof}

For a finite exponent list $A=(a_1,\ldots,a_t)$, let
\[
\eta(A)=\min\{x>1:\lfloor x^{a_i}\rfloor\in\mathbb P\text{ for all }i\},
\]
whenever the set is nonempty. A finite lower-bound certificate can prove such a minimum exactly. If every admissible $x$ is at least $u$ and $\lfloor u^c\rfloor=m$ is composite, then every admissible $x$ must be at least $\np(m)^{1/c}$. The jump is forced because no intervening integer is prime.

\begin{proposition}[An exact four-condition minimum]\label{enc:prop:finite-minimum}
For $A=(4,6,7,12)$,
\[
\eta(A)=29^{1/12},\qquad
(\lfloor\eta^4\rfloor,\lfloor\eta^6\rfloor,
\lfloor\eta^7\rfloor,\lfloor\eta^{12}\rfloor)=(3,5,7,29).
\]
Every $x\in[29^{1/12},30^{1/12})$ has these four outputs.
\end{proposition}
\begin{proof}
The first condition forces $x\ge2^{1/4}$. Table~\ref{enc:tab:jumps} applies the preceding forced-jump principle. Each entry is an exact comparison of integer powers, not a floating-point floor computation. At the final lower bound all four outputs are prime. Monotonicity, together with the values at the two endpoints, verifies the entire half-open interval.
\end{proof}

\begin{table}[htbp]
\centering
\caption{An exact certificate of the finite global minimum.}\label{enc:tab:jumps}
\begin{tabular}{@{}lll@{}}\toprule
Current lower bound & Outputs at $(4,6,7,12)$ & Next forced lower bound\\\midrule
$2^{1/4}$ & $(2,2,3,8)$ & $11^{1/12}$\\
$11^{1/12}$ & $(2,3,4,11)$ & $5^{1/7}$\\
$5^{1/7}$ & $(2,3,5,15)$ & $17^{1/12}$\\
$17^{1/12}$ & $(2,4,5,17)$ & $5^{1/6}$\\
$5^{1/6}$ & $(2,5,6,25)$ & $7^{1/7}$\\
$7^{1/7}$ & $(3,5,7,28)$ & $29^{1/12}$\\
$29^{1/12}$ & $(3,5,7,29)$ & all conditions satisfied\\\bottomrule
\end{tabular}
\end{table}

Define the seed prime
\begin{equation}\label{enc:eq:seedprime}
p_0=5132842708382182842735812571653
=29^{21}+24=\np(29^{21}).
\end{equation}
The exact certificates in Section~\ref{enc:sec:certificates} and Appendix~\ref{enc:app:verification-code} prove both primality and the absence of an intervening prime. In particular this is not an assumed probable prime. Put
\[
b=p_0^{1/252}.
\]
Since $p_0+1<2\cdot29^{21}<30^{21}$,
\begin{equation}\label{enc:eq:seed-cell}
I_{252}(p_0)\subset(29^{1/12},30^{1/12})\subset(1,4/3).
\end{equation}

\section{The least constant for the factorial tail}
The exponent sequence used first is
\begin{equation}\label{enc:eq:tail}
C=(4,6,7,12,E_0,E_1,\ldots),\qquad
E_0=252,\quad R_n=(E_n)!,\quad E_{n+1}=E_nR_n.
\end{equation}
Define the greedy prime chain by
\begin{equation}\label{enc:eq:greedy-chain}
p_{n+1}=\np(p_n^{R_n}),\qquad n\ge0.
\end{equation}

\begin{theorem}[Existence and exact global minimality]\label{enc:thm:greedy}
The set of $x>1$ for which $\lfloor x^c\rfloor$ is prime for every exponent $c$ in \eqref{enc:eq:tail} has a least element $\xi$. It is
\[
\xi=\lim_{n\to\infty}p_n^{1/E_n},
\]
and its first five outputs are $3,5,7,29,p_0$. In particular $\xi>b>29^{1/12}$.
\end{theorem}
\begin{proof}
Bertrand's postulate supplies a prime between each integer $N>1$ and $2N$. Consequently
\[
p_n^{R_n}<p_{n+1}<2p_n^{R_n},\qquad
p_{n+1}+1\le2p_n^{R_n}.
\]
Inductively all root cells remain below $2$, so $p_n<2^{E_n}<(E_n)!=R_n$. Bernoulli's inequality gives
\[
(p_n+1)^{R_n}\ge p_n^{R_n}(1+R_n/p_n)>2p_n^{R_n}.
\]
Lemma~\ref{enc:lem:nesting} and \eqref{enc:eq:seed-cell} therefore construct an admissible limit.

Any admissible $x$ is at least $29^{1/12}$ by Proposition~\ref{enc:prop:finite-minimum}. Its $252$nd power is at least the composite integer $29^{21}$, so its prime floor must be at least $p_0$ and $x\ge p_0^{1/252}$. Inductively, $x\ge p_n^{1/E_n}$ implies $x^{E_{n+1}}\ge p_n^{R_n}$. The latter is composite, so the next prime floor is at least $p_{n+1}$. Hence every admissible $x$ is at least every greedy lower endpoint, and therefore at least their limit. Strict nesting proves $\xi>b$.
\end{proof}

\begin{corollary}[No upper size bound]
The prime outputs $p_n$ increase without bound. More generally, any finite compatible prime-cell construction below $2$ ending at an exponent $E\ge252$ extends indefinitely by the rule $E\mapsto EE!$ and greedy prime selection.
\end{corollary}
\begin{proof}
The same nesting proof applies. Since the limiting base is greater than one and the exponents tend to infinity, its prime floors tend to infinity.
\end{proof}

\section{Algebraic approximation and certified decimal agreement}
\begin{lemma}[A one-sided decimal-boundary estimate]\label{enc:lem:decimal-boundary}
Let $b=m^{1/d}>0$ with integers $m,d\ge1$. For an integer $D\ge1$, set
\[
t=\frac{\lfloor10^Db\rfloor+1}{10^D}.
\]
If $t\le B$, then
\[
t-b\ge\frac{10^{-dD}}{dB^{d-1}}.
\]
Thus $b\le x<b+10^{-dD}/(dB^{d-1})$ implies
$\lfloor10^Dx\rfloor=\lfloor10^Db\rfloor$.
\end{lemma}
\begin{proof}
By definition $t>b$. The positive rational $t^d-m$ has denominator dividing $10^{dD}$, so it is at least $10^{-dD}$. On $[b,t]$ the derivative of $u^d$ is at most $dB^{d-1}$. The mean value theorem proves the estimate. No computation of the $D$ digits is required.
\end{proof}

\begin{proposition}[A Bertrand-only estimate for the factorial minimum]\label{enc:prop:bertrand-error}
With $E_1=252\cdot252!$,
\[
0<\xi-b< b(2^{1/E_1}-1)<\frac1{E_1}<2\cdot10^{-500}.
\]
\end{proposition}
\begin{proof}
The first child interval has upper endpoint at most $b2^{1/E_1}$. For the sharper numerical constant, use $b<4/3$, $\log2<7/10$, and $e^u-1<u/(1-u)$ for $0<u<1$. The desired inequality follows for $E_1>11$. Finally, the integer $E_1$ has $500$ decimal digits and exceeds $5\cdot10^{499}$.
\end{proof}

A short-prime-interval theorem drastically strengthens this estimate. Cully-Hugill proves, in the discussion of consecutive cubes in \cite[Section~4]{CullyHugill2023}, that
\begin{equation}\label{enc:eq:cube-input}
n\ge\exp(\exp(32.892))\quad\Longrightarrow\quad
\mathbb P\cap(n^3,(n+1)^3)\ne\varnothing.
\end{equation}
The cited passage is on page 14 of the arXiv version. This is an unconditional input; no Riemann-hypothesis assumption is used below.

\begin{theorem}[A long common decimal prefix]\label{enc:thm:precision}
For the original least constant $\xi$,
\[
0<\xi-p_0^{1/252}<10^{-10^{498}}.
\]
Moreover, with $D=10^{495}$,
\[
\lfloor10^D\xi\rfloor=\lfloor10^Dp_0^{1/252}\rfloor.
\]
The second assertion concerns exact common digits, not merely an absolute-error convention.
\end{theorem}
\begin{proof}
Write $R=252!$, $E=252R$, and $M=p_0^{R/3}$. Since $3\mid R$, $M$ is an integer. Exact integer comparisons give $R>10^{497}$ and $p_0>10^{30}$, whence
\[
M>10^{10R}>10^{10^{498}}.
\]
This also verifies the threshold in \eqref{enc:eq:cube-input}; for example $e^{32.892}<3^{33}<10^{16}$, whereas $\log M>10^{498}$. Thus
\[
M^3<p_1<(M+1)^3,\qquad p_1+1\le(M+1)^3.
\]
Strict containment of later cells and concavity yield
\[
0<\xi-b<b\big((1+1/M)^{3/E}-1\big)
<\frac{3b}{EM}<\frac1M<10^{-10^{498}}.
\]
For decimal agreement apply Lemma~\ref{enc:lem:decimal-boundary} with $d=252$ and $B=2$. The exact integer inequalities
\[
252\cdot2^{251}<10^{79},\qquad
252\cdot10^{495}+79<10^{498}
\]
show that the next $D$-digit boundary is farther from $b$ than the displayed error. Thus no decimal boundary is crossed.
\end{proof}

\begin{remark}[Precision is not the computation of a huge prime]\label{enc:rem:precision-cost}
Put $N=p_0^{252!}$. Substituting the algebraic approximation $b$ into the next exponent gives $b^{E_1}=N$, a composite integer, not $p_1$. The value $p_1$ has on the order of $10^{498}$ decimal digits. More explicitly, elementary bounds on $252!$ and $\log_{10}p_0$, together with $N<p_1<2N$, place its digit count strictly between $6\cdot10^{498}$ and $7\cdot10^{498}$. Printing such a number is a different task from proving the approximation theorem.

If a $d$-digit prime $q$ is to be identified at exponent $E$ using one whole decimal interval for the base, that interval must fit inside its prime cell. For a cell below $2$,
\[
(q+1)^{1/E}-q^{1/E}<\frac2{Eq}.
\]
For $E=E_1>10^{499}$ this is less than $10^{-d-497}$. Therefore this particular interval-certification method requires at least $d+498$ decimal places. This is not a lower bound for all algorithms for finding or testing primes.
\end{remark}

\begin{proposition}[A certified composite block]\label{enc:prop:composite-block}
For $N=p_0^{252!}$, all of $N+1,\ldots,N+561$ are composite. In particular $p_1\ge N+562$.
\end{proposition}
\begin{proof}
For every prime $\ell<563$, the integer $\ell-1$ divides $252!$. This finite assertion is checked by the exact verifier in Appendix~\ref{enc:app:verification-code}. Since $p_0>\ell$, Fermat's theorem gives $N\equiv1\pmod\ell$. For $1\le h\le561$, a prime divisor $\ell$ of $h+1$ is below $563$, and divides $N+h$. It is a proper divisor because $N+h>\ell$.
\end{proof}

\begin{corollary}[A necessary local twin filter]
If $0\le h\le563^2-4$ and both $N+h$ and $N+h+2$ are prime, then $h+1$ and $h+3$ are primes at least $563$.
\end{corollary}
\begin{proof}
The case $h=0$ is excluded by the compositeness of $N$. A composite positive integer below $563^2$ has a prime divisor below $563$. Applied to $h+1$ or $h+3$, this would give a proper divisor of one of the asserted large primes. Small prime values below $563$ are excluded by the same congruence. This is a necessary sieve condition, not a sufficient construction of twin primes.
\end{proof}


\section{A quantitative criterion for algebraic independence}
\label{enc:sec:quantitative-independence}

The arithmetic argument below separates two inputs: independence of finite
radical approximants, and sufficiently rapid approximation of the limiting
constants.  This separation makes it possible to replace strong prime-gap
estimates by faster growth of the exponent sequence.  Algebraic independence
of prime-representing constants has an established literature; in particular,
see Saito and Takeda~\cite{SaitoTakeda2022}.  No claim of priority for the
general approximation strategy is needed here.

\begin{lemma}[Prime radical basis]
\label{enc:lem:prime-radical-basis}
Let $E\ge2$ be an integer, and let $s_1,\ldots,s_r$ be distinct primes.  Put
$b_i=s_i^{1/E}$, taking positive real roots.  The $E^r$ monomials
\[
 b_1^{j_1}\cdots b_r^{j_r},\qquad 0\le j_i<E,
\]
are linearly independent over $\mathbb Q$.  Consequently,
\[
 [\mathbb Q(b_1,\ldots,b_r):\mathbb Q]=E^r,
\]
and $P(b_1,\ldots,b_r)\ne0$ whenever $P$ is a nonzero rational polynomial
with $\deg_{X_i}P<E$ for every $i$.
\end{lemma}

\begin{proof}
Every displayed monomial has its $E$th power in $\mathbb Q$.  If the quotient
of two such monomials were rational, taking $E$th powers and then valuations
at $s_i$ would imply that each difference of exponents is divisible by $E$.
Those differences lie strictly between $-E$ and $E$, so all are zero.
The monomials are therefore pairwise linearly independent over $\mathbb Q$.
The real-radical independence theorem of Carr and O'Sullivan
\cite[Theorem~1.1]{CarrOSullivan2009} makes them linearly independent.
They also span the generated field, by reduction using $b_i^E=s_i$.
This proves all the assertions.  In particular, $E$ need not be prime or
squarefree.
\end{proof}

For a nonzero polynomial
$P=\sum_{\nu}a_{\nu}X_1^{\nu_1}\cdots X_r^{\nu_r}
\in\mathbb Z[X_1,\ldots,X_r]$, write $|\nu|=\nu_1+\cdots+\nu_r$ and set
\begin{equation}
\label{enc:eq:polynomial-constants}
 B(P)=\max\left\{2,\sum_{\nu}|a_{\nu}|2^{|\nu|}\right\},
 \qquad
 L(P)=\sum_{|\nu|\ge1}|a_{\nu}|\,|\nu|\,2^{|\nu|-1}.
\end{equation}
The second constant is a Lipschitz bound on $[0,2]^r$ for the supremum norm.

\begin{proposition}[Explicit polynomial separation]
\label{enc:prop:polynomial-separation}
Suppose that $s_1,\ldots,s_r$ are distinct primes,
$b_i=s_i^{1/E}<2$, $\alpha_i\in[0,2]$, and
$|\alpha_i-b_i|\le\varepsilon$.  If
\begin{equation}
\label{enc:eq:separation-condition}
 E>\max_i\deg_{X_i}P,
 \qquad
 L(P)\varepsilon\le\tfrac12 B(P)^{1-E^r},
\end{equation}
then
\begin{equation}
\label{enc:eq:separation-conclusion}
 |P(\alpha_1,\ldots,\alpha_r)|
 \ge\tfrac12 B(P)^{1-E^r}>0.
\end{equation}
Moreover, $P(\alpha_1,\ldots,\alpha_r)$ has the same sign as
$P(b_1,\ldots,b_r)$.
\end{proposition}

\begin{proof}
Let $F=\mathbb Q(b_1,\ldots,b_r)$ and
$\beta=P(b_1,\ldots,b_r)$.  Lemma~\ref{enc:lem:prime-radical-basis} gives
$\beta\ne0$ and $[F:\mathbb Q]=E^r$.  Each $b_i$, and hence $\beta$, is an
algebraic integer.  Every embedding of $F$ into $\mathbb C$ sends $b_i$ to
an $E$th root of $s_i$, of absolute value $b_i<2$.  Thus
$|\sigma(\beta)|\le B(P)$ for every embedding.  The nonzero integer norm
then yields
\[
 1\le|N_{F/\mathbb Q}(\beta)|
 \le|\beta|B(P)^{E^r-1},
 \qquad |\beta|\ge B(P)^{1-E^r}.
\]
The mean value inequality gives
\[
 |P(\alpha_1,\ldots,\alpha_r)-\beta|\le L(P)\varepsilon.
\]
Condition~\eqref{enc:eq:separation-condition} proves both the lower bound and
the assertion about the sign.
\end{proof}

\begin{theorem}[Rapid prime-root approximation criterion]
\label{enc:thm:approximation-criterion}
Let $K\subset(1,2)$, let $E_n\to\infty$ be integers, and let
$\varepsilon_n>0$.  Suppose that, for every $\alpha\in K$, there are primes
$s_n(\alpha)$ such that
\[
 \left|\alpha-s_n(\alpha)^{1/E_n}\right|\le\varepsilon_n,
 \qquad s_n(\alpha)^{1/E_n}<2.
\]
Assume also the following two properties.
\begin{enumerate}
 \item For any distinct $\alpha_1,\ldots,\alpha_r\in K$, the primes
 $s_n(\alpha_1),\ldots,s_n(\alpha_r)$ are pairwise distinct for every
 sufficiently large $n$.
 \item For every fixed positive integer $r$,
 \[
 \frac{\log(1/\varepsilon_n)}{E_n^r}\longrightarrow\infty.
 \]
\end{enumerate}
Then $K$ is algebraically independent over $\mathbb Q$.
For any specified nonzero polynomial, every sufficiently late stage also
provides the explicit lower bound in
Proposition~\ref{enc:prop:polynomial-separation}.
\end{theorem}

\begin{proof}
Fix distinct $\alpha_1,\ldots,\alpha_r$ and a nonzero integer polynomial
$P$.  At all sufficiently late stages the prime labels are distinct and
$E_n>\max_i\deg_{X_i}P$.  If $L(P)>0$, the second assumption gives
\[
 \log\bigl(2L(P)\varepsilon_n\bigr)
 +(E_n^r-1)\log B(P)\longrightarrow-\infty,
\]
which is precisely the eventual validity of
$L(P)\varepsilon_n\le\frac12B(P)^{1-E_n^r}$.
For $L(P)=0$ the polynomial is a nonzero constant and the assertion is
immediate.  Proposition~\ref{enc:prop:polynomial-separation} completes the proof.
\end{proof}

\section{Binary families with a common prime prefix}
\label{enc:sec:binary-families}

Write
\[
 p_*=5132842708382182842735812571653,
 \qquad b_*=p_*^{1/252}.
\]
The finite computation and the certificate in Section~\ref{enc:sec:certificates} establish
that $p_*=\operatorname{nextprime}(29^{21})$ and
\begin{equation}
\label{enc:eq:seed-location-theory}
 [b_*,(p_*+1)^{1/252}]
 \subset (29^{1/12},30^{1/12})\subset(1,2).
\end{equation}
All constants constructed in this section therefore have the same five
initial outputs $3,5,7,29,p_*$ at exponents $4,6,7,12,252$.

Let $E_0=252$ and $E_{n+1}=E_nR_n$, where $R_n$ will be specified below.
For each finite binary word $w$, define a prime $s_w$ recursively.  The
empty word has label $s_{\varnothing}=p_*$.  If $|w|=n$, let
\[
 s_{w0}=\operatorname{nextprime}(s_w^{R_n}),
 \qquad
 s_{w1}=\operatorname{nextprime}(s_{w0}).
\]
Associate to $w$ the closed interval
\begin{equation}
\label{enc:eq:binary-root-cell}
 J_w=[s_w^{1/E_n},(s_w+1)^{1/E_n}].
\end{equation}
We use closed intervals only to take compact intersections.  The strict
nesting proved below ensures that no limiting constant occupies an upper
endpoint, where the floor condition would fail.

\begin{lemma}[Nested prime cells]
\label{enc:lem:binary-topology}
Suppose that
\[
 J_{w0}\cup J_{w1}\subset\operatorname{int}(J_w)
\]
for every $w$, that the two child cells are disjoint, and that $E_n\to\infty$.
Then every binary sequence $\omega\in\{0,1\}^{\mathbb N}$ determines exactly
one real number
\[
 \alpha_\omega\in\bigcap_{n\ge0}J_{\omega|n}.
\]
The map $\omega\mapsto\alpha_\omega$ is a homeomorphism onto a compact
perfect set $K$.  Its points satisfy
\[
 \lfloor\alpha_\omega^{E_n}\rfloor=s_{\omega|n}
 \quad(n\ge0).
\]
Different points have different prime labels at every sufficiently late
level.  If $\log E_n/n\to\infty$, then $K$ has upper box dimension zero,
and hence $\dim_{\mathrm H}K=0$.
\end{lemma}

\begin{proof}
The mean value theorem gives
\[
 |J_w|=(s_w+1)^{1/E_n}-s_w^{1/E_n}\le E_n^{-1}.
\]
Nested compactness and this bound give a unique intersection point.
It lies in the interior of every ancestor cell because it belongs to a
child cell, proving the floor identity.  Distinct binary sequences
separate into disjoint cells at their first different bit and never meet
again.  Agreement on the first $n$ bits gives real distance at most
$E_n^{-1}$, so the map is continuous.  A continuous injection from compact
Cantor space into $\mathbb R$ is a homeomorphism onto its image.  Alternatively,
perfectness follows directly by changing a sufficiently late bit.

Let $N_\delta(K)$ be the least number of intervals of length at most
$\delta$ covering $K$.  There are $2^n$ level-$n$ cells, each of diameter
at most $E_n^{-1}$.  For $E_n^{-1}\le\delta<E_{n-1}^{-1}$, therefore,
\[
 \frac{\log N_\delta(K)}{\log(1/\delta)}
 \le\frac{n\log2}{\log E_{n-1}}\longrightarrow0.
\]
The limsup defining the upper box dimension is zero.  The Hausdorff
dimension is bounded by the upper box dimension, so it too is zero;
in particular $K$ has zero Lebesgue measure.
\end{proof}

\begin{proposition}[The all-zero branch is the global minimum]
\label{enc:prop:zero-branch-minimum}
For either of the two constructions below, let
$\xi_0=\alpha_{000\ldots}$.  Then $\xi_0$ is the least $x>1$ for which
every exponent in
\[
 (4,6,7,12,E_0,E_1,\ldots)
\]
has a prime floor output.  Thus minimality holds in the full admissible
set, not only in the binary subset $K$.
\end{proposition}

\begin{proof}
The finite-prefix argument forces every admissible $x$ to satisfy
$x\ge29^{1/12}$.  At exponent $E_0=252$, the value at that lower endpoint
is the composite integer $29^{21}$, so prime output forces
$x\ge p_*^{1/E_0}$.  Let $p_n=s_{0^n}$.  Inductively, if
$x\ge p_n^{1/E_n}$, then
$x^{E_{n+1}}\ge p_n^{R_n}$, a composite integer.  Its floor can be prime
only if it is at least
$p_{n+1}=\operatorname{nextprime}(p_n^{R_n})$.  Therefore
$x\ge p_{n+1}^{1/E_{n+1}}$.  Taking the limit gives $x\ge\xi_0$.
The all-zero branch itself is admissible by
Lemma~\ref{enc:lem:binary-topology}, proving attainment and minimality.
\end{proof}

\subsection{The factorial schedule}

\begin{theorem}[A thin algebraically independent prime family]
\label{enc:thm:factorial-cantor}
Take $R_n=E_n!$, so that $E_{n+1}=E_nE_n!$.  The binary construction is
well-defined and gives a compact perfect algebraically independent set
$K_{\mathrm{fac}}\subset(1,2)$ of upper box dimension zero.  Its minimum is
the original minimum constant $\xi$.  Every member has prime outputs at
all the displayed exponents and the common initial outputs
$3,5,7,29,p_*$.  If $|w|=n$, $s=s_w$, $E=E_n$, $R=E!$, and
$b=s^{1/E}$, then every descendant $\alpha$ of $w$ satisfies
\begin{equation}
\label{enc:eq:factorial-uniform-error}
 0<\alpha-b<\frac{6b}{ERs^{R/3}}
 <\frac{12}{ER}\,2^{-R/3}.
\end{equation}
\end{theorem}

\begin{proof}
We use the following explicit consequence of
Cully-Hugill~\cite[Section~4]{CullyHugill2023}: each interval
$(m^3,(m+1)^3)$ contains a prime whenever the integer $m$ satisfies
$m\ge\exp(\exp(32.892))$.
Here $M=s^{R/3}$ is an integer because $3\mid E!$.  The threshold holds
already for $E=252$ and $s=2$, hence at every node.  For example,
$252!\ge126^{126}$ gives
\[
 \log\log(2^{252!/3})
 \ge126\log126+\log(\log2/3)>32.892.
\]
There is consequently a prime in each of
$(M^3,(M+1)^3)$ and $((M+1)^3,(M+2)^3)$.  The first two primes above
$s^R=M^3$ therefore satisfy
\[
 s^R<s_{w0}<s_{w1}<(M+2)^3,
 \qquad s_{w1}+1\le(M+2)^3.
\]
For the integer $t=R/3\ge2$, the binomial theorem gives
$(s+1)^t-s^t>2$.  Thus
\[
 s^R<s_{wj}<s_{wj}+1<(s+1)^R\qquad(j=0,1),
\]
which proves strict containment of both child cells.  Their prime labels
are odd and different, so $s_{w0}+1<s_{w1}$ and the cells are disjoint.

Every descendant is contained in one of these child cells and hence
\[
 b<\alpha\le b(1+2/M)^{3/(ER)}.
\]
Concavity of $(1+u)^{3/(ER)}$ gives
$(1+2/M)^{3/(ER)}-1<6/(ERM)$, proving
\eqref{enc:eq:factorial-uniform-error}.  In particular, a uniform error bound is
$\varepsilon_n=12\,2^{-E_n!/3}$.  For every fixed $r$,
\[
 \frac{\log(1/\varepsilon_n)}{E_n^r}
 =\frac{(\log2)E_n!/3-\log12}{E_n^r}\longrightarrow\infty.
\]
Lemma~\ref{enc:lem:binary-topology} and
Theorem~\ref{enc:thm:approximation-criterion} now give algebraic independence.
Finally, $E_{n+1}\ge E_n^2$, so $\log E_n\ge2^n\log E_0$ and the Hausdorff
dimension is zero.  Minimality follows from
Proposition~\ref{enc:prop:zero-branch-minimum}.
\end{proof}

\begin{corollary}[A common decimal block of prescribed enormous length]
\label{enc:cor:factorial-common-digits}
Every $\alpha\in K_{\mathrm{fac}}$ satisfies
\[
 0<\alpha-b_*<10^{-10^{498}},
 \qquad
 \left\lfloor10^{\,10^{495}}\alpha\right\rfloor
 =\left\lfloor10^{\,10^{495}}b_*\right\rfloor.
\]
Thus the entire uncountable algebraically independent set has the same
first $10^{495}$ digits after the decimal point.
\end{corollary}

\begin{proof}
At the initial node, put $R=252!$ and $M=p_*^{R/3}$.  The elementary
integer inequalities $p_*>10^{30}$ and $252!>10^{497}$ imply
$M>10^{10^{498}}$.  Since $6b_*<12<252R$, the error in
\eqref{enc:eq:factorial-uniform-error} is less than $M^{-1}$.

To certify the decimal block, apply Lemma~\ref{enc:lem:decimal-boundary} with
$D=10^{495}$, degree $252$, and upper bound $2$.  If $t$ is the next
$D$-place decimal boundary above $b_*$, the lemma and
$252\cdot2^{251}<10^{79}$ give
\[
 t-b_*>10^{-252D-79}>10^{-10^{498}}.
\]
No $\alpha$ reaches the next boundary.  This proves the floor identity.
\end{proof}

\subsection{An alternative requiring only Bertrand's theorem}

The factorial schedule uses a short-prime-interval theorem to obtain its
rapid error decay.  A different schedule obtains the same arithmetic
conclusions from Bertrand's theorem alone.  Its exponents are substantially
larger, so the following is a trade-off in assumptions and exponent growth,
not a faster algorithm for listing prime outputs.

\begin{theorem}[Bertrand-only construction]
\label{enc:thm:bertrand-cantor}
Define
\[
 F_0=252,\qquad R(F)=2^{F^F},\qquad
 F_{n+1}=F_nR(F_n).
\]
Use $F_n$ in place of $E_n$ in the same first-two-primes binary rule, starting
from the same seed $p_*$.  The resulting set $K_{\mathrm B}$ is compact,
perfect, algebraically independent, and of upper box dimension zero.
All its elements have the initial outputs $3,5,7,29,p_*$, and the all-zero
branch is the global minimum for this new exponent sequence.
At a node with label $s$ and exponent $F$, every descendant satisfies
\begin{equation}
\label{enc:eq:bertrand-uniform-error}
 0<\alpha-s^{1/F}<\frac{6}{F\,2^{F^F}}.
\end{equation}
The only theorem on the distribution of primes used in this construction
is Bertrand's theorem.
\end{theorem}

\begin{proof}
Let $R=2^{F^F}$.  Two applications of Bertrand's theorem
\cite{Bertrand} give, for the first two primes $q_0<q_1$ above $s^R$,
\[
 s^R<q_0<2s^R,\qquad q_1<2q_0<4s^R,
 \qquad q_1+1\le4s^R.
\]
Inductively the parent cell is below $2$, hence $s<2^F$.
For $F\ge2$, $R\ge4\cdot2^F>3s$, so Bernoulli's inequality implies
\[
 (s+1)^R=s^R(1+1/s)^R
 \ge s^R(1+R/s)>4s^R.
\]
Both child cells therefore lie strictly inside their parent, and they are
disjoint.  Put $b=s^{1/F}<2$.  Any descendant belongs to one of the two
children, so
\[
 b<\alpha\le b\,4^{1/(FR)},
 \qquad
 \alpha-b<b\frac{3}{FR}<\frac6{FR},
\]
where concavity gives $4^u-1<3u$ for $0<u<1$.
For $\varepsilon_n=6/(F_n2^{F_n^{F_n}})$,
\[
 \frac{\log(1/\varepsilon_n)}{F_n^r}
 =\frac{(\log2)F_n^{F_n}+\log F_n-\log6}{F_n^r}
 \longrightarrow\infty
\]
for each fixed $r$.  The general criterion proves algebraic independence.
The estimates $F_{n+1}\ge F_n^2$ and
Lemma~\ref{enc:lem:binary-topology} prove the topological and dimensional
claims.  The minimality proof is unchanged.
\end{proof}

\begin{corollary}[An elementary shared block of $10^{601}$ digits]
\label{enc:cor:bertrand-common-digits}
For every $\alpha\in K_{\mathrm B}$,
\[
 0<\alpha-b_*<10^{-10^{604}},
 \qquad
 \left\lfloor10^{\,10^{601}}\alpha\right\rfloor
 =\left\lfloor10^{\,10^{601}}b_*\right\rfloor.
\]
\end{corollary}

\begin{proof}
Let $A=252^{252}$.  The exact integer comparison $A>10^{605}$ and
$2^{10}>10^3$ imply
\[
 \frac6{252\,2^A}<2^{-A}<10^{-3A/10}<10^{-10^{604}}.
\]
The comparison $A>10^{605}$ can also be obtained without writing $A$:
$252^5>1.016\cdot10^{12}$ implies
\[
 252^{252}>252^2\,10^{600}(1.016)^{50}
 \ge63504\cdot1.8\cdot10^{600}>10^{605}.
\]
Apply Lemma~\ref{enc:lem:decimal-boundary} as in
Corollary~\ref{enc:cor:factorial-common-digits}, now with $D=10^{601}$.
Since $252D+79<10^{604}$, no member crosses the next boundary.
\end{proof}

\section{Effective coding and consequences}
\label{enc:sec:effective-coding}

\begin{theorem}[Preservation of Turing degrees]
\label{enc:thm:turing-degrees}
For either $K_{\mathrm{fac}}$ or $K_{\mathrm B}$, the homeomorphism
$\omega\mapsto\alpha_\omega$ and its inverse are computable in the usual
representations of binary sequences and real numbers.  In particular,
\[
 \deg_T(\alpha_\omega)=\deg_T(\omega).
\]
Each of these zero-dimensional algebraically independent sets therefore
contains a real of every Turing degree.  Its computable elements form a
dense subset, and its minimum is computable.
\end{theorem}

\begin{proof}
All exponent schedules are computable integer sequences.  Given a finite
binary word, its prime labels can be computed by exhaustive deterministic
primality testing; termination follows from the existence proofs.  The
endpoints of its cell are real algebraic numbers, and can be approximated
with certified rational errors by bisection and integer powering.
Given a requested accuracy $2^{-m}$ and oracle access to $\omega$, compute
enough bits to obtain a cell of diameter less than $2^{-m-2}$, then
approximate its endpoints.  This gives a rational approximation to
$\alpha_\omega$ with error less than $2^{-m}$.

Conversely, after a finite prefix has been recovered, compute the two
child cells.  Their positive separating gap contains an effectively
obtainable rational number $t$.  A Cauchy oracle for $\alpha_\omega$ will
eventually determine whether $\alpha_\omega<t$ or $\alpha_\omega>t$,
because the point belongs to one child cell and cannot equal $t$.
This determines the next bit and gives the inverse reduction.
Since every point is transcendental, the conventional binary expansion
has no dyadic ambiguity and is equivalent to its real-number oracle.
Thus the two Turing reductions give equality of degrees.  Eventually-zero
binary sequences are computable and dense, and include the all-zero
sequence representing the minimum.
\end{proof}

\begin{remark}[Computability versus feasible calculation]
The algorithms in Theorem~\ref{enc:thm:turing-degrees} need not be efficient.
The exponent growth is extreme, and finding a prescribed next prime can
require enormous computation.  Nor does the theorem say that every
constant in the construction is computable: it explicitly includes
noncomputable binary branches of every noncomputable degree.
\end{remark}

\begin{corollary}[Transcendental combinations and uniqueness of finite sums]
\label{enc:cor:transcendental-combinations}
Let $K$ be either of the two constructed sets.
\begin{enumerate}
 \item If $\alpha,\beta\in K$ are distinct, then
 $\alpha+\beta$, $\alpha-\beta$, $\alpha\beta$, and $\alpha/\beta$ are
 transcendental.  More generally, a nonconstant rational function over
 $\mathbb Q$ evaluated at a tuple of distinct elements of $K$ is
 transcendental.
 \item A finite sum of elements of $K$, counted with repetitions,
 determines its multiset of summands uniquely.  The analogous assertion
 holds for finite products.  In particular, equal sums or products of
 distinct multisets are impossible, even when their cardinalities differ.
 \item For every $\delta>0$ there is $\beta\in K$ with
 $0<\beta-\min K<\delta$, and this difference is transcendental.
\end{enumerate}
\end{corollary}

\begin{proof}
For the first assertion, algebraic independence ensures that the
denominator of any nonzero rational function does not vanish.  If a
nonconstant rational function had algebraic value, substituting it into
the minimal polynomial of that value and clearing denominators would
give a nonzero rational polynomial relation among the tuple.
Equal finite sums or products from different multisets likewise give
nonzero polynomial relations.  Finally, change a sufficiently late zero
bit of the minimum branch to one.  The resulting point is larger than
the minimum, remains within the same arbitrarily small ancestor cell,
and has transcendental difference from the minimum by the first part.
\end{proof}

For computable distinct branches, the quantitative certificate
\eqref{enc:eq:separation-condition} can itself be found effectively.  First
compute enough bits that their labels differ, then increase the level
until the rational error bound satisfies that inequality.  The search
terminates by Theorem~\ref{enc:thm:approximation-criterion}; it yields both
the explicit nonzero lower bound and, by certified evaluation of the
radical approximants, the sign of any prescribed nonzero polynomial.
This is an effective statement about the constructed constants, not an
irrationality test for arbitrary named constants.

\begin{corollary}[A thin field containing every Turing degree]
\label{enc:cor:thin-field}
For either constructed set $K$, the real field $F=\mathbb Q(K)$ satisfies
\[
 \dim_{\mathrm H}F=0,\qquad
 \operatorname{trdeg}_{\mathbb Q}F=2^{\aleph_0},\qquad
 F\cap\overline{\mathbb Q}=\mathbb Q.
\]
It is a dense, meager, proper $F_\sigma$ subfield of $\mathbb R$, and it contains a real of
every Turing degree.
\end{corollary}

\begin{proof}
The Cartesian-power cover obtained from covers of $K$ shows that each
$K^r$ also has upper box dimension zero.  Every element of $F$ is a value
of a rational function $P/Q$ over $\mathbb Q$ at a finite tuple from $K$,
where the denominator is nonzero.  On
\[
 \{x\in K^r:|Q(x)|\ge1/m\},\qquad m=1,2,\ldots,
\]
the rational function is Lipschitz: polynomial values and derivatives
are bounded on $[1,2]^r$, and the denominator is bounded away from zero.
Its image therefore has Hausdorff dimension zero.  There are only
countably many rational functions and integers $m$, so their union $F$
has Hausdorff dimension zero. Each displayed domain is compact and its
image is compact. These images have empty interior because their dimension
is zero. Thus $F$ is also a countable union of closed nowhere dense sets,
which proves the $F_\sigma$ and meagerness assertions.  Since $K$ is algebraically independent
and has cardinality $2^{\aleph_0}$, it is a transcendence basis for $F$
of that cardinality.  After repeated entries in a tuple are identified, its rational function
is either constant in $\mathbb Q$ or nonconstant in distinct generators.
The rational-function assertion in
Corollary~\ref{enc:cor:transcendental-combinations} therefore gives
$F\cap\overline{\mathbb Q}=\mathbb Q$.  The field is dense because it
contains $\mathbb Q$, and proper because its Hausdorff dimension is less
than that of $\mathbb R$.  Finally, $K\subset F$ already contains every
Turing degree.
\end{proof}

The dimension assertion here is Hausdorff dimension.  A dense field does
not have upper box dimension zero on bounded intervals; the upper box
statement applies to the compact generating set $K$ and its finite
Cartesian powers.


\section{Certified large primes and finite twin-prime branches}
\label{enc:sec:certificates}

The sparse exponent construction does not by itself make primality testing
unnecessary.  It does, however, accommodate candidates whose neighboring
integers have large, completely factored divisors.  This section records
seven explicit certified primes, including two twin-prime pairs, and proves
that their root cells preserve the preceding outputs.  These are finite
examples, not a theorem that the displayed polynomial family contains
infinitely many primes or twin primes.

\subsection{The seed prime and its exact certificate}

Put
\begin{equation}\label{enc:eq:certificate-seed}
p=5132842708382182842735812571653=29^{21}+24.
\end{equation}
This is the least prime strictly exceeding $29^{21}$.  A complete recursive
order certificate is included in the executable appendix.  There are
$28$ nodes, counting the terminal prime $2$.  Its top-level factorization is
\[
p-1=2^2\cdot157\cdot4457\cdot12541\cdot355452619
       \cdot411378690803,
\]
and the top-level witness is $2$.  Each prime factor is certified in the
same way until the recursion reaches $2$.

For clarity, the elementary criterion used in this recursion is as follows.
If the complete prime factorization of $n-1$ is known and
\[
a^{n-1}\equiv1\pmod n,\qquad
\gcd\bigl(a^{(n-1)/q}-1,n\bigr)=1
\quad(q\mid n-1,\ q\text{ prime}),
\]
then every prime divisor $\ell$ of $n$ has an element of order $n-1$
modulo $\ell$.  Hence $n-1\mid\ell-1$, so $\ell\ge n$ and $n$ is prime.
All the factorizations and witnesses needed for this argument are given
explicitly in the appendix; no probable-prime predicate is invoked.

For the intervening integers $29^{21}+i$, $1\le i\le23$, proper divisors
in increasing order of $i$ are
\[
\begin{gathered}
2,31,2,3,2,5,2,571,2,3,2,13,\\
2,17,2,3,2,72673,2,7,2,3,2.
\end{gathered}
\]
Thus the minimality of the seed is certified as well as its primality.

\subsection{Two complementary primality certificates}

\begin{proposition}[A specialized $n-1$ certificate]
\label{enc:prop:plus-certificate}
Let $s$ be a certified odd prime and let $N>1$ satisfy
$F=2s^2\mid N-1$ and $F>\sqrt N$.  Suppose an integer $a$ satisfies
\begin{equation}\label{enc:eq:plus-certificate}
a^{N-1}\equiv1\pmod N,\qquad
\gcd\bigl(a^{(N-1)/2}-1,N\bigr)
=\gcd\bigl(a^{(N-1)/s}-1,N\bigr)=1.
\end{equation}
Then $N$ is prime.
\end{proposition}
\begin{proof}
For any prime divisor $\ell$ of $N$, the order $h$ of $a$ modulo $\ell$
divides $N-1$.  The two gcd conditions force
$v_2(h)=v_2(N-1)$ and $v_s(h)=v_s(N-1)$.  Consequently $F\mid h$.
Since $h\mid\ell-1$, every prime divisor satisfies
$\ell\ge F+1>\sqrt N$.  This is impossible for a composite $N$.
\end{proof}

For the other side of a twin pair, use the Lucas sequence
\[
U_0=0,\qquad U_1=1,\qquad U_{m+2}=P U_{m+1}-Q U_m,
\]
with discriminant $\Delta=P^2-4Q$.

\begin{proposition}[A specialized $n+1$ certificate]
\label{enc:prop:minus-certificate}
Let $s$ be a certified odd prime and let $N>1$ satisfy
$F=2s^2\mid N+1$ and $F-1>\sqrt N$.
Suppose $\gcd(N,2Q\Delta)=1$ and
\begin{equation}\label{enc:eq:minus-certificate}
U_{N+1}\equiv0\pmod N,\qquad
\gcd\bigl(U_{(N+1)/2},N\bigr)
=\gcd\bigl(U_{(N+1)/s},N\bigr)=1.
\end{equation}
Then $N$ is prime.
\end{proposition}
\begin{proof}
Fix a prime divisor $\ell$ of $N$.  The assumptions ensure that $\ell$ is
odd and does not divide $Q\Delta$.  In $\mathbb F_{\ell^2}$, let
$\alpha,\beta$ be the distinct, nonzero roots of $X^2-PX+Q$.
The identity
\[
U_m=\frac{\alpha^m-\beta^m}{\alpha-\beta}
\]
shows that the least positive $z$ with $U_z\equiv0\pmod\ell$ is the
multiplicative order of $\alpha/\beta$.  Frobenius either fixes both roots
or interchanges them, according as $\Delta$ is a quadratic residue or
nonresidue modulo $\ell$.  Therefore
\[
z\mid \ell-\left(\frac{\Delta}{\ell}\right),
\qquad z\le\ell+1.
\]
The congruence in \eqref{enc:eq:minus-certificate} gives $z\mid N+1$.
The two gcd conditions force the full $2$-adic and $s$-adic valuations of
$N+1$ to occur in $z$, whence $F\mid z$.  It follows that
$\ell\ge F-1>\sqrt N$.  Again this excludes a composite $N$.
\end{proof}

Both criteria are standard order-based primality certificates.  Their
use here is to provide short, reproducible proofs for the selected
branches, not to assert a new general primality test.

\subsection{Quartic lifts and preservation of the coding intervals}

For an odd prime $s\ge3$ and an odd integer $j$ with $1\le j<s$, define
\[
N_\pm(s,j)=s^4+js^2\pm1,\qquad F=2s^2.
\]
The conditions on $s$ and $j$ imply
\begin{equation}\label{enc:eq:quartic-cell}
s^4<N_-(s,j)<N_+(s,j)+1<(s+1)^4.
\end{equation}
They also imply
\[
F\mid N_+(s,j)-1,\quad F^2>N_+(s,j),
\]
and
\[
F\mid N_-(s,j)+1,\quad (F-1)^2>N_-(s,j).
\]
For example, the last strict inequality follows from
\[
(F-1)^2-N_-(s,j)=3s^4-(j+4)s^2+2>0.
\]
All the coefficients used below satisfy $1\le j<s$.

Suppose $s$ is the desired prime output at exponent $E$.  If either
$N=N_\pm(s,j)$ is prime, then \eqref{enc:eq:quartic-cell} gives
\begin{equation}\label{enc:eq:quartic-root-cell}
\left[N^{1/(4E)},(N+1)^{1/(4E)}\right]
\subset\left(s^{1/E},(s+1)^{1/E}\right).
\end{equation}
Thus the quartic lift simultaneously supplies an exactly certified prime
and a child cell lying strictly inside the previous prime cell.
Subsequent nested extensions preserve both outputs.

Starting with \eqref{enc:eq:certificate-seed}, set
\begin{align*}
P_1&=p^4+33p^2+1,\\
P_2&=P_1^4+765P_1^2+1,\\
P_3&=P_2^4+4583P_2^2+1.
\end{align*}
The exact certificates are summarized in Table~\ref{enc:tab:prime-tower}.
The displayed witnesses satisfy every congruence and gcd condition in
Proposition~\ref{enc:prop:plus-certificate}.

\begin{table}[htbp]
\centering
\caption{Certified quartic lifts.}\label{enc:tab:prime-tower}
\begin{tabular}{c r r r}
\hline
Prime & Coefficient $j$ & Decimal digits & Witness $a$\\
\hline
$P_1$ & $33$ & $123$ & $2$\\
$P_2$ & $765$ & $492$ & $5$\\
$P_3$ & $4583$ & $1966$ & $13$\\
\hline
\end{tabular}
\end{table}

Consequently, the outputs
\[
3,5,7,29,p,P_1,P_2,P_3
\]
can be realized at the exponents
\[
4,6,7,12,252,1008,4032,16128,
\]
respectively, and then extended indefinitely by a factorial tail.
These are deliberately selected branches of the coding tree.
No claim is made that they are the globally minimal constants for the
corresponding exponent sequence, or that they lie on the original greedy
factorial branch.

\subsection{Two certified twin-prime pairs}

Define
\begin{align*}
T_\pm&=p^4+20303p^2\pm1,\\
W_\pm&=T_+^4+447479T_+^2\pm1.
\end{align*}
Both $T_\pm$ have $123$ decimal digits, and both $W_\pm$ have $492$ decimal
digits.  Their differences are exactly
\[
T_+-T_-=W_+-W_-=2.
\]
For the plus signs, Proposition~\ref{enc:prop:plus-certificate} applies with
witness $17$ for $T_+$ and witness $7$ for $W_+$.
For the minus signs, Proposition~\ref{enc:prop:minus-certificate} applies with
\[
(P,Q)=(1,-1),\quad\Delta=5\qquad\text{for }T_-,
\]
and
\[
(P,Q)=(1,-3),\quad\Delta=13\qquad\text{for }W_-.
\]
The first Lucas sequence is the Fibonacci sequence; the second satisfies
$U_{m+2}=U_{m+1}+3U_m$.  The exact verifier checks the discriminant gcds,
all rank congruences, all gcd exclusions, and all size inequalities.

Each of the two cells for $T_\pm$ at exponent $1008$ lies strictly inside
the cell for $p$ at exponent $252$.  Each of the two cells for $W_\pm$ at
exponent $4032$ lies strictly inside the cell for $T_+$ at exponent $1008$.
Attaching the same sufficiently sparse factorial tail to the two siblings
produces two prime-generating constants whose outputs at the selected
exponent differ by exactly $2$ and whose earlier outputs agree.

This verifies finite instances of coupled prime-generating constants.
The existence of suitable coefficients $j$ at every stage, or at
infinitely many stages, has not been proved.  In particular, the examples
do not establish twin-prime infinitude and impose no upper bound on the
size of all primes.

\subsection{Reproducibility and scope of the computation}

The executable appendix is a self-contained Python~3 program using only
integer arithmetic and the standard library.  Modular powers use exact
binary exponentiation.  Lucas terms are evaluated by the identities
\[
U_{2m}=U_m(2U_{m+1}-P U_m),\qquad
U_{2m+1}=U_{m+1}^{\,2}-Q U_m^{\,2}.
\]
Thus no enormous Lucas number needs to be constructed: all operations are
performed modulo the candidate.  The program certifies the seed, its
minimality above $29^{21}$, all seven later primes, their stated digit
counts, both gaps of $2$, and the strict inclusion of all seven child
cells.  Its successful execution produces
\begin{quote}\ttfamily\small
Seed: 28 certificate nodes; all 23 intervening integers composite.\\
Prime tower: 123, 492, 1966 digits; all certificates passed.\\
Twin pairs: 123 and 492 digits; all certificates passed.\\
All seven selected child cells are strictly nested.
\end{quote}


\section{Finite prime cells and certified minimum searches}
\label{enc:sp:finite}

For a finite set $A\subset\mathbb N$ of positive exponents, write
\[
 \mathcal W(A)=\{x>1:\lfloor x^c\rfloor\in\mathbb P\text{ for every }c\in A\}.
\]
For a prescribed tuple $(p_c)_{c\in A}$ of primes, its parameter cell is
\begin{equation}\label{enc:sp:cell}
 \bigcap_{c\in A}[p_c^{1/c},(p_c+1)^{1/c})
 =\left[\max_{c\in A}p_c^{1/c},\min_{c\in A}(p_c+1)^{1/c}\right).
\end{equation}
It is nonempty precisely when its left endpoint is strictly smaller than
its right endpoint. All comparisons between endpoints are decidable by
integer powers: $u^{1/a}<v^{1/b}$ if and only if $u^b<v^a$.

\begin{proposition}[Certified finite search]\label{enc:sp:finite-search}
Suppose that $\mathcal W(A)$ is nonempty. Then it has a least element.
Starting at $x=2^{1/\min A}$, repeatedly apply the following rule: if
$n=\lfloor x^c\rfloor$ is not prime for some $c\in A$, replace $x$ by
$q^{1/c}$, where $q$ is the least prime strictly larger than $n$.
This procedure terminates at $\min\mathcal W(A)$. If a feasible upper
bound $B$ is supplied, all its comparisons and its termination can be
certified by a finite calculation with integers.
\end{proposition}
\begin{proof}
The initial value is a lower bound for every feasible parameter. If $x$
is such a lower bound and $\lfloor x^c\rfloor=n$ is not prime, every
feasible $y\ge x$ must satisfy $y^c\ge q$. Thus every replacement preserves
the lower-bound property and strictly increases $x$. Below any feasible
$B$, there are only finitely many breakpoints $m^{1/c}$, with $c\in A$ and
$m\le B^c$. Consequently the procedure terminates. Its terminal point
is both feasible and a lower bound for all feasible points.
\end{proof}

This argument is a verification procedure, not an unconditional assertion
that every finite exponent set is feasible, and not a polynomial-time
complexity assertion. The certificate consists of the forced jumps and
the exact prime checks at the terminal point.

\begin{proposition}[Two exponents cannot force a large first prime]
\label{enc:sp:two-exponents}
If $1\le a<b$ are integers, then $\mathcal W(\{a,b\})$ is nonempty and
\[
 \left\lfloor\bigl(\min\mathcal W(\{a,b\})\bigr)^a\right\rfloor=2.
\]
\end{proposition}
\begin{proof}
Put $r=b/a>1$. If $r<2$, then $x=2^{1/a}$ gives first output $2$ and
second output $\lfloor2^r\rfloor\in\{2,3\}$. If $r\ge2$, put
$n=\lfloor2^r\rfloor$. Bertrand's postulate supplies a prime $q$ with
$n<q<2n$. Hence
\[
 2^r<q<2^{r+1}<3^r,
\]
so $x=q^{1/b}$ is feasible and has first output $2$. The minimum is no
larger than this witness and no smaller than $2^{1/a}$; its first output
must therefore also be $2$.
\end{proof}

There is a sharp version relative to the initial exponent pair $(4,6)$.
\begin{proposition}[Exactly two extra constraints are needed after $(4,6)$]
\label{enc:sp:sharp-four-six}
For every integer $c>6$, there exists an $x$ such that
\[
 (\lfloor x^4\rfloor,\lfloor x^6\rfloor)=(2,3),\qquad
 \lfloor x^c\rfloor\in\mathbb P.
\]
In particular, one additional exponent cannot force the first output of
the minimum above $2$. In contrast, the two additional exponents $7,12$
give
\[
 \min\mathcal W(\{4,6,7,12\})=29^{1/12},\qquad
 (\lfloor x^4\rfloor,\lfloor x^6\rfloor,\lfloor x^7\rfloor,
 \lfloor x^{12}\rfloor)=(3,5,7,29)
\]
at that minimum.
\end{proposition}
\begin{proof}
Every $x\in[3^{1/6},4^{1/6})$ has first two outputs $(2,3)$, since
$2<3^{2/3}<4^{2/3}<3$. For $c\ge15$,
$(4/3)^{c/6}>2$. Apply Bertrand's postulate to
$\lfloor3^{c/6}\rfloor$ to find a prime
$q\in(3^{c/6},4^{c/6})$. Then $x=q^{1/c}$ is a witness. The remaining
cases are certified by
\[
\begin{array}{c|rrrrrrrr}
 c&7&8&9&10&11&12&13&14\\ \hline
 q&5&5&7&7&11&11&11&13
\end{array}
\]
using the exact inequalities $3^c\le q^6<4^c$.
The four-exponent assertion is exactly Proposition~\ref{enc:prop:finite-minimum},
whose forced-jump certificate is given in Table~\ref{enc:tab:jumps}.
\end{proof}

The qualification ``after $(4,6)$'' is essential: this proposition does
not settle whether some other three-exponent prefix can force the first
output above $2$.

\section{Twin-output geometry and an exact moving-window reformulation}
\label{enc:sp:twins}

\begin{proposition}[Exact two-cell criterion]\label{enc:sp:twin-cell}
Let $p\ge2$, $a\ge1$, and $r\ge1$ be integers. There exists an $x>1$ with
\[
 \lfloor x^a\rfloor=p,\qquad \lfloor x^{a+r}\rfloor=p+2
\]
if and only if
\begin{equation}\label{enc:sp:twin-criterion}
 \frac{\log p}{\log(1+3/p)}
 <\frac a r<
 \frac{\log(p+1)}{\log(1+1/(p+1))}.
\end{equation}
Primality of $p,p+2$ is not needed for this geometric statement.
\end{proposition}
\begin{proof}
The two cells are
$[p^{1/a},(p+1)^{1/a})$ and
$[(p+2)^{1/(a+r)},(p+3)^{1/(a+r)})$.
Their intersection is nonempty exactly when
\[
 p^{1/a}<(p+3)^{1/(a+r)},\qquad
 (p+2)^{1/(a+r)}<(p+1)^{1/a}.
\]
Taking logarithms and rearranging gives~\eqref{enc:sp:twin-criterion}.
\end{proof}

The two bounds in~\eqref{enc:sp:twin-criterion} are asymptotic to
$p\log p/3$ and $p\log p$, respectively. Thus, along successful twin
encodings, $p\to\infty$ if and only if $a/r\to\infty$. The geometric
criterion says when an already specified pair can be encoded; it does
not assert primality or the existence of new twin pairs.

The next result chooses its exponent from the scale $N$, without first
choosing a twin pair. It also gives an exact topological counting
interpretation.
\begin{theorem}[A target-free moving window]\label{enc:sp:moving-window}
For an integer $N\ge16$, put
\[
 a_N=\lfloor N\log N\rfloor,\qquad
 I_N=[N^{1/a_N},(2N)^{1/a_N}),
\]
and let
\[
 \mathcal T_N=\{x\in I_N:\lfloor x^{a_N}\rfloor,
                         \lfloor x^{a_N+1}\rfloor\in\mathbb P\}.
\]
Every $x\in\mathcal T_N$ produces a twin-prime pair. Moreover,
$\mathcal T_N$ is nonempty if and only if a twin pair $p,p+2$ exists with
$N\le p<2N$. Its connected components are in one-to-one correspondence
with those twin pairs, and the component for $p$ has length
\begin{equation}\label{enc:sp:twin-component-length}
 \min\{(p+1)^{1/a_N},(p+3)^{1/(a_N+1)}\}
 -\max\{p^{1/a_N},(p+2)^{1/(a_N+1)}\}>0.
\end{equation}
\end{theorem}
\begin{proof}
Write $a=a_N$ and $y=x^a\in[N,2N)$. Then
\[
 \Delta(y)=x^{a+1}-x^a=y(y^{1/a}-1)
\]
is strictly increasing in $y>1$. Since $e^t-1\ge t$,
\[
 \Delta(y)\ge\frac{N\log N}{a}\ge1.
\]
Using $e^t-1<t/(1-t)$ for $0<t<1$ and $a>N\log N-1$ gives
\[
 \Delta(y)<
 \frac{2N\log(2N)}{N\log N-1-\log(2N)}<3.
\]
For the last inequality, the denominator is positive and the difference
between three times the denominator and the numerator is
\[
 h(N)=N\log N-2N\log2-3\log(2N)-3.
\]
Here $h(16)=17\log2-3>0$ and
$h'(N)=\log N+1-2\log2-3/N>0$ for $N\ge16$.
Thus $1\le x^{a+1}-x^a<3$. If both floors are prime, they are odd and
distinct; their difference is at most $3$, and hence is exactly $2$.

Conversely, let $N\le p<2N$ be a twin lower prime and set
$f(y)=y^{1+1/a}$. The preceding bound gives $f(p)<p+3$.
Strict monotonicity of $\Delta$ gives
$f(p+1)>p+2$, including when $p+1=2N$. If $f(p)\ge p+2$, take $y=p$;
otherwise the intermediate value theorem gives a
$y\in(p,p+1)$ with $f(y)=p+2$.
This proves existence, and continuity gives a cell of positive length.
Its endpoints are exactly those in~\eqref{enc:sp:twin-component-length}.
Distinct lower primes give separated cells, since the disjoint intervals
$[p,p+1)$ in the $y$ coordinate have positive gaps between them.
\end{proof}

The theorem is an exact reformulation, including a counting identity;
it is not an unconditional nonemptiness theorem for all or infinitely
many windows. In particular, the twin-prime conjecture is equivalent to
$\mathcal T_{2^m}\ne\varnothing$ for infinitely many integers $m$.
The minimum taken within $I_N$ must not be confused with the unrestricted
minimum of $\mathcal W(\{a_N,a_N+1\})$, whose first output is always $2$
by Proposition~\ref{enc:sp:two-exponents}.

\section{Rigorous obstructions for fixed power spectra}
\label{enc:sp:obstructions}

\begin{proposition}[A fixed-base bound]\label{enc:sp:fixed-base}
Let $x>1$, let $a<b$ be positive integers, and suppose
$\lfloor x^a\rfloor=p$ and $\lfloor x^b\rfloor=p+d$ with $d\ge1$.
Then
\begin{equation}\label{enc:sp:fixed-base-bound}
 p<\frac{d+1}{x-1}.
\end{equation}
Consequently a fixed base has only finitely many twin pairs anywhere in
its entire integer-exponent output spectrum.
\end{proposition}
\begin{proof}
Since $b\ge a+1$,
$xp\le x^{a+1}\le x^b<p+d+1$. Rearranging proves the bound.
Only finitely many integer exponents have outputs below a fixed bound.
\end{proof}

For the constant in the main construction,
$\xi>29^{1/12}>13/10$, where the latter comparison follows from
$29\cdot10^{12}>13^{12}$. Hence $3/(\xi-1)<10$.
Therefore its complete integer-power spectrum can contain only the twin
pairs $(3,5)$ and $(5,7)$: the only possible prime lower members below
$10$ are $2,3,5,7$, and $2+2$ and $7+2$ are composite.
Both pairs occur at exponents $4,6,7$, so the spectrum contains exactly
these two twin pairs. This assertion concerns pairs of values in the
spectrum, not adjacency of their exponents.

\begin{proposition}[Scaling does not enlarge the prime outputs]
\label{enc:sp:scaling}
Suppose a positive integer exponent sequence $C=(C_k)$ has a least
prime-generating parameter $\xi(C)$. For an integer $m\ge1$, put
$mC=(mC_k)$. Then
\[
 \xi(mC)=\xi(C)^{1/m},\qquad
 \lfloor\xi(mC)^{mC_k}\rfloor=\lfloor\xi(C)^{C_k}\rfloor.
\]
Thus a family of bases tending to $1$ need not yield any new prime values
or any larger twin pairs.
\end{proposition}
\begin{proof}
The increasing bijection $x\mapsto x^m$ identifies the two feasible sets
and preserves their order.
\end{proof}

A second obstruction applies even when two different fixed bases and
two independently chosen exponent sequences are allowed.
\begin{theorem}[Exponential approximation forced by bounded output gaps]
\label{enc:sp:two-base-liouville}
Let $x,y>1$ and $D\ge1$ be fixed. If infinitely many distinct pairs
$(a,b)\in\mathbb N^2$ satisfy
\begin{equation}\label{enc:sp:bounded-nonzero-gap}
 1\le|\lfloor x^a\rfloor-\lfloor y^b\rfloor|\le D,
\end{equation}
then $\rho=\log x/\log y$ is a Liouville number. More precisely, along
an unbounded sequence of such exponents,
\begin{equation}\label{enc:sp:exponential-approximation}
 0<\left|\rho-\frac ba\right|
 <\frac{2(D+1)}{a\log y}\,x^{-a}.
\end{equation}
\end{theorem}
\begin{proof}
Infinitely many distinct exponent pairs satisfying the bounded-gap
condition force $a\to\infty$ along a subsequence. Put $A=x^a$ and
$B=y^b$. Then $|A-B|<D+1$. For $A>2(D+1)$, both $A$ and $B$ exceed
$A/2$, so the mean value theorem for the logarithm gives
\[
 |a\log x-b\log y|<2(D+1)x^{-a}.
\]
The two powers cannot be equal because their floors differ. This proves
\eqref{enc:sp:exponential-approximation}.

If $\rho=u/v$ were rational, every nonzero difference
$|\rho-b/a|$ would be at least $1/(va)$, contradicting the displayed
exponential upper bound for sufficiently large $a$. Thus $\rho$ is
irrational. For every positive integer $L$, the upper bound is eventually
smaller than $a^{-L}$. Reducing $b/a$ to a fraction with denominator
$q\le a$ only weakens the bound $a^{-L}\le q^{-L}$; the reduced
denominators must be unbounded because $\rho$ is irrational. This is
the Liouville property.
\end{proof}

The exponential estimate is stronger than the bare Liouville property,
and neither is asserted to be sufficient for bounded nonzero gaps or
for primality. In particular, if $\log x/\log y$ is rational or is an
algebraic irrational number, only finitely many bounded nonzero gaps
can occur. In the algebraic irrational case this follows already from
the elementary Liouville approximation inequality, without needing
Roth's theorem.

\begin{corollary}[A common exponent cannot sustain fixed nonzero gaps]
\label{enc:sp:common-exponent}
Two fixed bases $x,y>1$ cannot satisfy
$1\le|\lfloor x^{E_n}\rfloor-\lfloor y^{E_n}\rfloor|\le D$ for an
unbounded integer sequence $E_n$. If $x\ne y$, the absolute differences
tend to infinity; if $x=y$, they vanish identically.
\end{corollary}
\begin{proof}
For $x>y$, $x^E-y^E=x^E(1-(y/x)^E)\to\infty$, and taking floors changes
the difference by less than $1$. The other cases follow by symmetry.
\end{proof}

Thus finite constructions of two constants with twin outputs at a common
exponent do not automatically iterate to infinitely many twin outputs
for those same two constants.

\section{Adaptive insertion and prescribed residue classes}
\label{enc:sp:adaptive}

The freedom to choose new exponents gives a general extension principle
that needs no estimate for short prime gaps.
\begin{lemma}[Insertion of a prescribed large integer]\label{enc:sp:insertion}
Let $1<u<v$, let $m\ge1$, and let $H>0$. Every sufficiently large integer
$q$ admits an exponent $E\in m\mathbb N$ with $E>H$ such that
\[
 [q^{1/E},(q+1)^{1/E}]\subset(u,v).
\]
\end{lemma}
\begin{proof}
The required strict inequalities are equivalent to
\[
 \frac{\log(q+1)}{\log v}<E<\frac{\log q}{\log u}.
\]
The interval length tends to infinity as $q\to\infty$, and its lower
endpoint also tends to infinity. Eventually its length exceeds $m$ and
it lies above $H$, so it contains a multiple of $m$ strictly between its
endpoints.
\end{proof}

\begin{theorem}[Adaptive coding in arbitrary infinite prime sets]
\label{enc:sp:adaptive-coding}
Let $J$ be a nonempty bounded open interval with $\overline J\subset(1,\infty)$.
Let $S_1,S_2,\ldots$ be arbitrary infinite subsets of the primes, and let
$G_1,G_2,\ldots$ be positive real thresholds. Then there exist an
$\alpha\in J$, strictly increasing positive integer exponents $E_n$, and
primes $p_n\in S_n$, such that
\[
 E_n\mid E_{n+1},\qquad p_n>G_n,\qquad
 \lfloor\alpha^{E_n}\rfloor=p_n\quad(n\ge1).
\]
The exponents can additionally be required at each stage to exceed any
specified finite threshold depending on the preceding choices. If all
points of $J$ satisfy a given finite prime-output prefix, that prefix is
preserved.
\end{theorem}
\begin{proof}
Choose $p_1\in S_1$ large enough to apply Lemma~\ref{enc:sp:insertion} inside
$J$, imposing $p_1>G_1$. This produces a closed prime cell $J_1$ in $J$.
Given $J_n$, apply the lemma in its interior with $m=E_n$, imposing the
next exponent threshold and choosing $p_{n+1}\in S_{n+1}$ sufficiently
large, in particular above $G_{n+1}$. This gives
$J_{n+1}\subset\operatorname{int}(J_n)$ and $E_n\mid E_{n+1}$.
Choose the exponents to tend to infinity.
If $\overline J\subset[u,v]$ with $u>1$, the mean value theorem bounds
\[
 |J_n|\le\frac{1}{E_nu^{E_n-1}}\longrightarrow0.
\]
Nested compact intervals therefore have a unique common point $\alpha$.
For every $n$, this point lies in $J_{n+1}\subset\operatorname{int}(J_n)$,
so the floor at exponent $E_n$ is exactly $p_n$.
\end{proof}

\begin{corollary}[Arbitrarily prescribed reduced residue classes]
\label{enc:sp:residue-coding}
Given positive integers $m_n$ and integers $a_n$ satisfying
$\gcd(a_n,m_n)=1$, one can make the construction in
Theorem~\ref{enc:sp:adaptive-coding} satisfy
\[
 \lfloor\alpha^{E_n}\rfloor\equiv a_n\pmod{m_n}
 \qquad(n\ge1).
\]
\end{corollary}
\begin{proof}
Dirichlet's theorem says that each target set
$S_n=\{p\in\mathbb P:p\equiv a_n\pmod{m_n}\}$ is infinite. Apply the
theorem.
\end{proof}

This extension permits arbitrary prescribed growth thresholds and
residue classes while preserving any already established open prefix
cell. It constructs an admissible constant, not necessarily the global
minimum for the resulting exponent sequence. The exponents are chosen
adaptively and need not be the factorial-tail exponents of the main
construction. Using a target set of twin lower primes here would require
its infinitude as an input; the insertion argument cannot supply that
missing arithmetic assertion.


\section{An exact spectral encoding from the prime-generating constants}
\label{enc:sec:spectral-applications}

The conclusions in this section are elementary consequences of algebraic
independence.  They do not solve an open graph-theoretic problem, and the
use of generic diagonal entries is standard.  Their role here is to show
that the prime-generating construction supplies explicit choices of the
otherwise generic parameters.  In particular, one may use distinct
members of the set $K$ constructed above, even when all the chosen
members have the exceptionally long common decimal prefix established
there.

Fix algebraically independent real numbers $\alpha_1,\ldots,\alpha_n$
over $\mathbb Q$, and write
\[
 D=\operatorname{diag}(\alpha_1,\ldots,\alpha_n).
\]
All matrix entries called rational below belong to $\mathbb Q$.

\begin{theorem}[A quartic trace determines a labeled weighted graph]
\label{enc:thm:quartic-encoding}
For a symmetric rational matrix $A=(a_{ij})$ with zero diagonal, put
\[
 \Phi_\alpha(A)=\operatorname{tr}(A+D)^4.
\]
If $\Phi_\alpha(A)=\Phi_\alpha(B)$ for two such matrices, then
$a_{ij}^2=b_{ij}^2$ for every $i,j$.
Consequently $\Phi_\alpha$ is injective on the class of nonnegative
rational symmetric zero-diagonal matrices, and in particular on the
adjacency matrices of labeled simple graphs.
The exponent $4$ is the smallest positive integer with this universal
injectivity property for labeled simple graphs of arbitrary order.
\end{theorem}

\begin{proof}
Cyclic invariance of the trace gives
\begin{align}
 \operatorname{tr}(A+D)^4
 ={}&\sum_i\alpha_i^4
   +4\sum_i(A^2)_{ii}\alpha_i^2
   +4\sum_{i<j}a_{ij}^2\alpha_i\alpha_j \notag\\
   &+4\sum_i(A^3)_{ii}\alpha_i+\operatorname{tr}A^4.
 \label{enc:eq:quartic-encoding}
\end{align}
Indeed, the six words containing two occurrences each of $A$ and $D$
contribute $4\operatorname{tr}(A^2D^2)+2\operatorname{tr}(ADAD)$.
The terms with three copies of $D$ vanish because $a_{ii}=0$.
Algebraic independence says that equality of the two evaluated
polynomials in~\eqref{enc:eq:quartic-encoding} implies equality of their
coefficients.  Comparing the coefficient of $\alpha_i\alpha_j$ gives
$a_{ij}^2=b_{ij}^2$.

For a simple graph with $m$ edges, labeled degrees $d_i$, and $t$
triangles, the first three traces are
\begin{align*}
 \operatorname{tr}(A+D)&=\sum_i\alpha_i,\\
 \operatorname{tr}(A+D)^2&=\sum_i\alpha_i^2+2m,\\
 \operatorname{tr}(A+D)^3&=\sum_i\alpha_i^3
                  +3\sum_i d_i\alpha_i+6t.
\end{align*}
The Petersen graph and the pentagonal prism both have ten vertices,
are $3$-regular, and contain no triangles.  For any labeling their first
three traces therefore agree.  They are different graphs: the prism
has $4$-cycles, whereas the Petersen graph has none.  Thus no exponent
$1$, $2$, or $3$ is universally injective.  The quartic trace is
injective by the preceding argument.
\end{proof}

\begin{remark}
The nonnegativity qualification is necessary for recovering signed
weights.  If $S$ is a diagonal matrix with entries in $\{1,-1\}$, then
$SDS=D$ and
\[
 SAS+D=S(A+D)S.
\]
Thus $A$ and $SAS$ have the same trace at every power, although their
edge signs can differ.  The theorem recovers squared weights without
this qualification and the weights themselves when they are
nonnegative.
\end{remark}

The invariant is an exact real-number encoding.  Algebraic independence
alone does not give a useful uniform precision bound for decoding it.
For fixed $n$ and fixed computable labels, however, the finitely many
values arising from labeled simple graphs are distinct computable
reals.  Successively refining their certified approximations eventually
produces disjoint intervals.  This is a terminating finite-precision
reconstruction procedure, without a claim of computational efficiency.

\begin{theorem}[Independent eigenvalues]
\label{enc:thm:independent-eigenvalues}
Let $A$ be any rational symmetric $n\times n$ matrix.  The $n$
eigenvalues of $A+D$ are algebraically independent over $\mathbb Q$.
In particular, they are simple and are individually transcendental.
\end{theorem}

\begin{proof}
Replace the $\alpha_i$ temporarily by indeterminates $x_i$, and write
\[
 \det\bigl(tI-A-\operatorname{diag}(x_1,\ldots,x_n)\bigr)
 =t^n+c_1(x)t^{n-1}+\cdots+c_n(x).
\]
The total degree of $c_k$ is $k$, and its homogeneous part of degree
$k$ is $(-1)^k e_k(x)$, where $e_k$ is the $k$th elementary symmetric
polynomial.  The $c_k$ are algebraically independent: otherwise, take
a nonzero relation $P(c_1,\ldots,c_n)=0$, give the $k$th argument of
$P$ weight $k$, and pass to the highest weighted homogeneous part
$P_\mathrm{top}$.  The highest ordinary homogeneous part after
substitution would be
\[
 P_\mathrm{top}(-e_1,e_2,\ldots,(-1)^n e_n)=0,
\]
contrary to the algebraic independence of $e_1,\ldots,e_n$.
Since evaluation at the algebraically independent tuple $\alpha$ is
injective on $\mathbb Q[x_1,\ldots,x_n]$, the numbers
$c_1(\alpha),\ldots,c_n(\alpha)$ are algebraically independent.

Let $\lambda_1,\ldots,\lambda_n$ be the eigenvalues, counted with
multiplicity.  Their elementary symmetric polynomials are the
$c_k(\alpha)$ up to signs.  Hence the field
$\mathbb Q(\lambda_1,\ldots,\lambda_n)$ has transcendence degree at
least $n$.  It is generated by $n$ elements, so these elements must be
algebraically independent.  This also rules out repeated eigenvalues.
\end{proof}

\begin{theorem}[Every vertex is cyclic]
\label{enc:thm:every-vertex-cyclic}
Let $A$ be a rational symmetric matrix whose support graph, with edges
$ij$ whenever $i\ne j$ and $a_{ij}\ne0$, is connected.  Put $H=A+D$.
For every coordinate vector $e_v$, the matrix
\[
 [\,e_v,He_v,\ldots,H^{n-1}e_v\,]
\]
is invertible.  Equivalently, the linear system
$\dot z=Hz+e_v u$ is controllable from any single coordinate.
\end{theorem}

\begin{proof}
Absorb the rational diagonal entries of $A$ into the $\alpha_i$;
this preserves algebraic independence.  We may therefore assume
$a_{ii}=0$.  Consider
\[
 F(x)=\det\bigl(\operatorname{diag}(x_1,\ldots,x_n)-A\bigr).
\]
We first prove that $F$ is irreducible over $\mathbb Q$.
It is multiaffine and its coefficient of $x_1\cdots x_n$ is $1$.
In any factorization $F=GH$, the degree in each variable is additive,
so each variable occurs in at most one factor.  A nontrivial
factorization therefore partitions the variables into two nonempty
sets $S,T$.  Normalize the top monomials of the factors to have
coefficient $1$.  The coefficient of the monomial missing only $x_i$
is $-a_{ii}=0$, so each corresponding coefficient in its factor is
zero.  For $i\in S$ and $j\in T$, the coefficient missing both
$x_i,x_j$ would then be zero.  In $F$, however, this coefficient is
$-a_{ij}^2$.  Thus every edge between $S$ and $T$ would vanish,
contrary to connectedness.

Adjoin an independent variable $t$ and make the invertible substitution
$x_i=t-d_i$.  It follows that
\[
 f(t,d)=\det\bigl(tI-A-\operatorname{diag}(d_1,\ldots,d_n)\bigr)
\]
is irreducible in $\mathbb Q[t,d]$.  Since it is monic in $t$, it is
also irreducible over $\mathbb Q(d)[t]$.  Let $g_v(t,d)$ be the
characteristic polynomial of the principal submatrix with row and
column $v$ deleted.  It is monic of degree $n-1$, so $f$ and $g_v$
are relatively prime over $\mathbb Q(d)$.  Their resultant is therefore
a nonzero polynomial in $\mathbb Q[d]$.  Evaluation at $d=\alpha$
cannot annihilate that polynomial.  Thus $H$ and the principal
submatrix obtained by deleting $v$ have no common eigenvalue.

If an eigenvector of $H$ had zero $v$th coordinate, deleting that
coordinate would give a nonzero eigenvector of the principal
submatrix with the same eigenvalue.  This is impossible.  Every
eigenvector therefore has nonzero $v$th coordinate.  Together with
simplicity of the spectrum, the spectral decomposition of the Krylov
matrix gives a product of an invertible orthogonal matrix, an
invertible diagonal matrix, and a Vandermonde matrix on distinct
eigenvalues.  The Krylov matrix is invertible.
\end{proof}

\begin{corollary}[No exact vertex revival]
\label{enc:cor:no-vertex-revival}
Under the hypotheses of Theorem~\ref{enc:thm:every-vertex-cyclic}, suppose
$n\ge3$.  There are no real $t\ne0$ and complex $\gamma$ with
$|\gamma|=1$ such that
\[
 e^{-itH}e_v=\gamma e_v.
\]
There is also no perfect state transfer between distinct coordinate
vectors.
\end{corollary}

\begin{proof}
Every eigenvector has nonzero overlap with $e_v$.  A revival would
therefore require $e^{-it\lambda_j}=\gamma$ for every $j$.  For three
distinct eigenvalues this gives nonzero integers $k_2,k_3$ such that
\[
 t(\lambda_2-\lambda_1)=2\pi k_2,
 \qquad t(\lambda_3-\lambda_1)=2\pi k_3.
\]
Eliminating $t$ produces the nontrivial rational linear relation
\[
 k_3(\lambda_2-\lambda_1)-k_2(\lambda_3-\lambda_1)=0,
\]
contradicting Theorem~\ref{enc:thm:independent-eigenvalues}.
For a real symmetric $H$, perfect state transfer from $v$ to $w$
would imply revival at $v$ at twice the transfer time: in a real
orthonormal eigenbasis the two coordinate overlaps must agree up to
sign, so the squared phase is common to every eigencomponent.
The already proved absence of revival rules this out.
\end{proof}

These applications are stable under replacing the labels by any other
algebraically independent tuple.  What the prime-generating framework
adds is the simultaneous arithmetic specification of the labels:
each label retains the prescribed prime outputs and, in the set $K$,
the shared long decimal prefix.  No unbounded twin-prime conclusion is
used in, or follows from, these spectral transfers.


\clearpage
\part{Geometry, sparse relations, and exponential independence}\label{part:geometry}
The full spectrum and selected thin subsets are treated separately. Dimension results concern all admissible constants; the later independence results use explicitly refined subsets.

% Consolidated follow-up body. External labels enc:* are intentionally cross-part.
% All theorem/proposition/lemma/corollary/remark environments use the master preamble.
For an increasing positive exponent sequence $E=(E_k)$, write
\[
 \mathcal P(E)=\{x>1:\lfloor x^{E_k}\rfloor\in\mathbb P
                         \text{ for every }k\}.
\]
This part distinguishes the full admissible spectrum from a prescribed
compact prime-coding subset. A dimension statement about the first and
an independence statement about the second do not imply either statement
for the unrestricted least admissible base. We first treat dimension,
then quantitative separation, and finally selection within a fixed schedule.

Throughout this part, $\dim_H$ denotes Hausdorff dimension and
$\overline{\mathbb Q}$ denotes the field of complex algebraic numbers.
A perfect set is nonempty, closed, and has no isolated points.
Algebraic independence of an indexed family means nonvanishing of every
nonzero polynomial in finitely many distinct indexed entries.
Computability refers to ordinary real or complex Cauchy names, with no
complexity bound asserted.

\section{Removing the boundedness hypothesis in the dimension theorem}
\label{fol:sec:interpolation}

For a sequence $\mathbf c=(c_k)_{k\ge1}$ of positive real numbers, put
\[
 C_k=\prod_{i=1}^k c_i,
 \qquad
 \mathcal W(\mathbf c)
 =\{x>1:\lfloor x^{C_k}\rfloor\text{ is prime for every }k\ge1\}.
\]
The distinction between multipliers $c_k$ and cumulative exponents $C_k$
is essential in this section. The intermediate multipliers introduced below
are real numbers, even when all prescribed cumulative exponents are integers.

\begin{theorem}[Saito's bounded-multiplier theorem]\label{fol:thm:saito-bounded}
If $2\le c_k\le R<\infty$ for every $k$, then
$\dim_H\mathcal W(\mathbf c)=1$.
\end{theorem}

This is \cite[Theorem~1.2]{SaitoDimension2021}. The same paper asks in
Question~1.4 of its arXiv version whether the Hausdorff dimension can be
determined when the upper bound on $c_k$ is removed. A bounded refinement
settles that question.

\begin{lemma}[Bounded multiplicative refinement]\label{fol:lem:bounded-refinement}
Let $a>1$ and let $c_k\ge a$ be real numbers. There are real numbers
$d_j\in[a,a^2)$ and strictly increasing integers $N_k$ such that, with
$D_j=\prod_{i=1}^j d_i$, one has
\[
 D_{N_k}=C_k\qquad(k\ge1).
\]
Consequently, $\mathcal W(\mathbf d)\subseteq\mathcal W(\mathbf c)$.
More generally, for any fixed set $T\subseteq\mathbb Z$, satisfying
$\lfloor x^{D_j}\rfloor\in T$ for every $j$ implies
$\lfloor x^{C_k}\rfloor\in T$ for every $k$.
\end{lemma}

\begin{proof}
Set $m_k=\lfloor\log c_k/\log a\rfloor\ge1$ and
$N_0=0$, $N_k=\sum_{i=1}^k m_i$. For $1\le j\le m_k$, define
\[
 d_{N_{k-1}+j}=c_k^{1/m_k}.
\]
Since $a^{m_k}\le c_k<a^{m_k+1}$,
\[
 a\le c_k^{1/m_k}<a^{1+1/m_k}\le a^2.
\]
Each block has product $c_k$, so $D_{N_k}=C_k$ exactly.
If all outputs at the refined exponents $D_j$ are prime, then in particular
the outputs at their subsequence $D_{N_k}=C_k$ are prime.
\end{proof}

\begin{theorem}[Full dimension without an upper multiplier bound]
\label{fol:thm:unbounded-dimension}
For every real sequence $c_k\ge2$, bounded or unbounded,
\[
 \boxed{\dim_H\mathcal W(\mathbf c)=1.}
\]
\end{theorem}

\begin{proof}
Apply Lemma~\ref{fol:lem:bounded-refinement} with $a=2$. The refined multipliers
belong to $[2,4)$, so Theorem~\ref{fol:thm:saito-bounded} gives
$\dim_H\mathcal W(\mathbf d)=1$. Monotonicity of Hausdorff dimension yields
\[
 1=\dim_H\mathcal W(\mathbf d)
 \le\dim_H\mathcal W(\mathbf c)\le\dim_H\mathbb R=1.
\]
\end{proof}

\begin{remark}[Scope and attribution]\label{fol:rem:dimension-scope}
Theorem~\ref{fol:thm:unbounded-dimension} answers Question~1.4 as stated in
\cite[arXiv version]{SaitoDimension2021}. Its proof is an elementary
consequence of Saito's substantial bounded-multiplier theorem, rather than
a new estimate for primes in short intervals. This observation does not
assert historical priority for the refinement argument. This refinement argument alone does not resolve
Question~1.5, which concerns iterated exponentials; the separate
argument in Section~\ref{fol:sec:wright-dimension-extension} does so.
Nor does the refinement imply full dimension for a branch-restricted
Cantor subset or for an arbitrarily prescribed finite prefix cell.
\end{remark}

\begin{corollary}[Arbitrary initial scale]\label{fol:cor:arbitrary-initial-scale}
Let $E_k>0$ satisfy $E_{k+1}/E_k\ge2$ for all $k\ge1$. Then
\[
 \dim_H\{x>1:\lfloor x^{E_k}\rfloor\text{ is prime for all }k\ge1\}=1.
\]
\end{corollary}

\begin{proof}
Set $c_1=2$ and $c_{k+1}=E_{k+1}/E_k$. Then
$C_k=2E_k/E_1$. Under the change of variable $x=y^{2/E_1}$,
one has $x^{E_k}=y^{C_k}$. The change of variable is bi-Lipschitz on
each compact subinterval of $(1,\infty)$, and countable stability of
Hausdorff dimension proves the assertion.
\end{proof}

For example, the conclusion applies to $E_k=k!+k$, beginning with $k=1$:
the first two ratios are $2$ and $9/4$, and for $k\ge2$,
\[
 (k+1)!+(k+1)-2(k!+k)=(k-1)(k!-1)\ge0.
\]
It also applies to $E_k=3^k-k$ and $E_k=3^k-2k$, because respectively
\[
 E_{k+1}-2E_k=3^k+k-1>0,
 \qquad
 E_{k+1}-2E_k=3^k+2k-2>0.
\]
The second sequence starts at $E_1=1$, which is covered by
Corollary~\ref{fol:cor:arbitrary-initial-scale}.
These are statements about the full sets of permissible
bases, not about the arithmetic nature of their least elements.

The endpoint $4$ in the refinement interval $[2,4)$ is best possible for
this universal factorization device. Indeed, a number $c\in(2,4)$ cannot
be a product of two or more factors each at least $2$; if one demanded a
universal interval $[2,R]$ with $R<4$, any $c\in(R,4)$ would be impossible
to refine. This optimality concerns the interpolation lemma alone.


% Local metric consequences. Citation keys: SaitoDimension2021, CullyHugill2023.
% Mathematika 20 (1973), 119--134, Theorem 2.
% DOI: 10.1112/S0025579300004708.

\section{Full dimension for Wright prime orbits}
\label{fol:sec:wright-dimension-extension}

Let $b>1$, put $f_b(x)=b^x$, and write $f_b^{\circ n}$ for its
$n$-fold composition. Define
\[
 \mathcal W_b^{\mathrm{it}}
 =\left\{x>1:
   \left\lfloor f_b^{\circ n}(x)\right\rfloor\text{ is prime
   for every }n\ge1\right\}.
\]
Question~1.5 in the arXiv version of
\cite{SaitoDimension2021} asks for the Hausdorff dimension of
$\mathcal W_2^{\mathrm{it}}$. The argument below answers this question
and allows every fixed real base $b>1$. It uses ordinary prime counting in fixed-ratio intervals and a uniform upper sieve estimate; it does not
use a lower bound for primes in short intervals.
\begin{theorem}[Dimension of prime iterated-exponential orbits]
\label{fol:thm:wright-full-dimension}
For every real $b>1$,
\[
 \boxed{\dim_H\mathcal W_b^{\mathrm{it}}=1.}
\]
More precisely, for every sufficiently large prime $p_0$, there is a compact perfect set $K\subset[p_0+1/4,p_0+3/4]$ of Hausdorff dimension
$1$ such that $\lfloor f_b^{\circ n}(x)\rfloor$ is prime for every
$x\in K$ and every $n\ge0$.
\end{theorem}
We first specify the two external prime-counting inputs. The prime number theorem implies that, for fixed $0<a<d$,
\begin{equation}\label{fol:eq:wright-pnt-input}
 \#\bigl(\mathbb P\cap[aX+O(1),dX+O(1)]\bigr)
 \sim\frac{(d-a)X}{\log X}\qquad(X\longrightarrow\infty).
\end{equation}
The Montgomery--Vaughan form of the Brun--Titchmarsh inequality
\cite[Theorem~2]{MontgomeryVaughan1973}, specialized to modulus $1$,
gives
\[
 \pi(u+h)-\pi(u)<\frac{2h}{\log h}
 \qquad(u\ge1,\ h>1).
\]
In particular, after allowing either convention for interval endpoints and enlarging an absolute constant $C$, every interval $A\subset[1,\infty)$ of length $h\ge0$ satisfies
\begin{equation}\label{fol:eq:wright-sieve-input}
 \#(A\cap\mathbb P)\le C\frac{h+1}{\log(h+2)}.
\end{equation}
For $0\le h\le2$ this last assertion also follows directly by counting integers. The uniformity in the location of $A$ is essential.
\subsection{Closed prime cells and uniform inverse distortion}
Fix $b>1$, write $f=f_b$ and $g(x)=\log x/\log b$, and let
\[
 I_p=[p+1/4,p+3/4]\qquad(p\in\mathbb P).
\]
Choose $P>1$ sufficiently large that
\[
 \rho:=\frac1{P\log b}<1,
 \qquad b^{p+1/4}-1/4\ge p+1\quad(p\ge P),
\]
and that all counting and contraction estimates below hold for $p\ge P$. Starting with a prime $p_0\ge P$, call a finite sequence
$\omega=(p_0,\ldots,p_n)$ admissible if all its entries are prime and
\[
 I_{p_{j+1}}\subset f(I_{p_j})\qquad(0\le j<n).
\]
Its cylinder in the original parameter axis is
\[
 J_\omega=g^{\circ n}(I_{p_n}).
\]
For every $x\in J_\omega$ and $0\le j\le n$ one has
$f^{\circ j}(x)\in I_{p_j}$. In particular, all those integer parts
are prime, including at both endpoints of $J_\omega$.
\begin{lemma}[Distortion and child geometry]
\label{fol:lem:wright-distortion}
There exist positive constants $c_b,C_b$, independent of the admissible word and its length, with the following properties. If
$J=J_\omega$ has final prime label $p$, its children are exactly
$J_{\omega q}=g^{\circ(n+1)}(I_q)$ for primes $q$ with
$I_q\subset f(I_p)$. Their number $N_p$ satisfies
\begin{equation}\label{fol:eq:wright-branch-count}
 c_b\frac{b^p}{p}\le N_p\le C_b\frac{b^p}{p},
 \qquad N_p\ge2.
\end{equation}
Every child satisfies
\begin{equation}\label{fol:eq:wright-child-length}
 c_b b^{-p}|J|\le |J_{\omega q}|\le C_b b^{-p}|J|,
\end{equation}
and distinct children are separated by a gap of length at least
$c_b b^{-p}|J|$. Moreover,
\begin{equation}\label{fol:eq:wright-forward-derivative}
 \sup_{x\in J}(f^{\circ(n+1)})'(x)
 \le C_b\frac{b^p}{|J|}.
\end{equation}
Here and below, constants with the same subscript may be enlarged or decreased when necessary.
\end{lemma}
\begin{proof}
Put $h=g^{\circ n}$. For $u,v\in I_p$, all the intermediate points
$g^{\circ j}(u)$ and $g^{\circ j}(v)$ belong to $[P,\infty)$.
On that half-line $g$ is $\rho$-Lipschitz. The derivative formula
\[
 h'(u)=(\log b)^{-n}
       \prod_{j=0}^{n-1}\bigl(g^{\circ j}(u)\bigr)^{-1}
\]
therefore gives
\begin{align*}
 \left|\log\frac{h'(u)}{h'(v)}\right|
 &\le\sum_{j=0}^{n-1}
       \frac{|g^{\circ j}(u)-g^{\circ j}(v)|}{P}\\
 &\le\frac{|u-v|}{P(1-\rho)}
 \le\frac1{2P(1-\rho)}.
\end{align*}
Thus the distortion of $h$ on $I_p$ is at most
$D=\exp(1/(2P(1-\rho)))$, uniformly in $n$ and $\omega$.
This includes $n=0$, when $h$ is the identity. Since $|I_p|=1/2$,
\begin{equation}\label{fol:eq:wright-inverse-scale}
 \frac{2|J|}{D}\le h'(u)\le2D|J|
 \qquad(u\in I_p).
\end{equation}
Write $a=b^{1/4}$ and $d=b^{3/4}$. The permitted child primes are exactly
\[
 q\in\mathbb P\cap[ab^p-1/4,db^p-3/4].
\]
Equation~\eqref{fol:eq:wright-pnt-input} gives
$N_p\sim(d-a)b^p/(p\log b)$, proving
\eqref{fol:eq:wright-branch-count} after increasing $P$.
On $f(I_p)=[ab^p,db^p]$, the derivative
$g'(t)=1/(t\log b)$ is bounded above and below by positive constants
depending only on $b$ times $b^{-p}$. Combining this fact with
\eqref{fol:eq:wright-inverse-scale} gives
\eqref{fol:eq:wright-child-length}. Distinct integer labels give a gap
of at least $1/2$ between their trimmed cells $I_q$, so the same derivative lower bound proves the asserted gap estimate.
Finally, on $J$ the derivative of $f^{\circ n}=h^{-1}$ is at most
$D/(2|J|)$. Since $f^{\circ n}(J)=I_p$ and
$f'(u)=(\log b)b^u$, the chain rule proves
\eqref{fol:eq:wright-forward-derivative}.
\end{proof}
Increasing $P$ once more, we may suppose that
$C_b b^{-p}\le1/2$ for $p\ge P$ in
\eqref{fol:eq:wright-child-length}. Let $K_n$ be the finite union of all
depth-$n$ cylinders, starting with $K_0=I_{p_0}$, and put
\[
 K=\bigcap_{n\ge0}K_n.
\]
The cylinders are closed, nested, and disjoint at each level, every cylinder has at least two children, and their maximum length tends to
zero. Hence $K$ is a nonempty compact perfect set, and each of its points has one infinite admissible itinerary. The floor condition
above shows $K\subset\mathcal W_b^{\mathrm{it}}$.
\subsection{An interval mass estimate}
Give $I_{p_0}$ mass $1$ and divide each cylinder's mass equally among its $N_p$ children. The consistent cylinder masses define a Borel
probability measure $\mu$ on $K$. Since $N_p\ge2$, the measure has no atoms. By \eqref{fol:eq:wright-branch-count},
\begin{equation}\label{fol:eq:wright-child-mass}
 \mu(J_{\omega q})\le C_b p b^{-p}\mu(J_\omega).
\end{equation}
Fix $0<s<1$. Equations~\eqref{fol:eq:wright-child-length} and
\eqref{fol:eq:wright-child-mass} imply
\begin{equation}\label{fol:eq:wright-cylinder-ratio}
 \frac{\mu(J_{\omega q})}{|J_{\omega q}|^s}
 \le C_b c_b^{-s}p b^{-(1-s)p}
       \frac{\mu(J_\omega)}{|J_\omega|^s}.
\end{equation}
The coefficient on the right is at most $1$ for all sufficiently large $p$, depending on $s$ and $b$. By the choice of $P$, every child
label is at least its parent label plus $1$. Thus the minimum label at depth $n$ tends to infinity. Only finitely many cylinders occur before
all labels reach the required threshold. Taking the maximum of
$\mu(J)/|J|^s$ over those finitely many cylinders and then using
\eqref{fol:eq:wright-cylinder-ratio} gives a finite constant $A_s$ such that
\begin{equation}\label{fol:eq:wright-cylinder-frostman}
 \mu(J)\le A_s|J|^s
 \qquad\text{for every cylinder }J.
\end{equation}
This estimate alone does not yet give a Hausdorff-dimension bound: the number of children is unbounded, so intervals between cylinder
scales must also be controlled. The uniform upper sieve supplies this missing estimate.
Let $B\subset\mathbb R$ be an interval of positive finite length. If $B\cap K$ contains at most one point, then $\mu(B)=0$. Otherwise,
the unique itineraries of the points of $B\cap K$ have a longest common finite prefix. Let $J$ be its cylinder and $p$ its final label. All
points of $B\cap K$ lie in $J$, and at least two different children of
$J$ contain points of $B\cap K$.
Put $t=|B|/|J|$. If $t\ge1$, then
\[
 \mu(B)\le\mu(J)\le A_s|J|^s\le A_s|B|^s.
\]
Suppose henceforth that $0<t<1$. The sibling gap estimate implies
\begin{equation}\label{fol:eq:wright-splitting-scale}
 y:=b^pt\ge c_b>0.
\end{equation}
Even if $B$ is not contained in $J$, the set $B\cap J$ is an interval. By \eqref{fol:eq:wright-forward-derivative}, its image under
$f^{\circ(n+1)}$ has length at most $C_b b^pt=C_b y$, where $n$ is
the depth of $J$. If a child $J_{\omega q}$ meets $B$, then $I_q$ meets this image. Therefore all the corresponding prime labels $q$
belong to one interval of length at most $C_b y+1$. The sieve bound
\eqref{fol:eq:wright-sieve-input}, together with
\eqref{fol:eq:wright-splitting-scale}, yields
\[
 \#\{q:J_{\omega q}\cap B\ne\varnothing\}
 \le C_b\frac{y+1}{\log(y+2)}
 \le C_b\frac{y}{\log(y+2)}.
\]
Using the equal-child masses and
\eqref{fol:eq:wright-branch-count}, we consequently obtain
\begin{align}
 \mu(B)
 &\le C_b\mu(J)\frac{p b^{-p}y}{\log(y+2)}\notag\\
 &=C_b\mu(J)\frac{pt}{\log(b^pt+2)}\notag\\
 &\le C_b A_s |B|^s
       \frac{p t^{1-s}}{\log(b^pt+2)}.
 \label{fol:eq:wright-general-interval}
\end{align}
The final factor is bounded uniformly for $p\ge P$ and $0<t<1$. Indeed, if $t\ge b^{-p/2}$, then
\[
 \frac{p t^{1-s}}{\log(b^pt+2)}\le\frac2{\log b}.
\]
If $t<b^{-p/2}$, then it is at most
\[
 \frac{p b^{-(1-s)p/2}}{\log2},
\]
which is bounded as $p$ varies. Hence there is a finite constant $M_s$ such that
\[
 \mu(B)\le M_s|B|^s
 \qquad\text{for every interval }B\subset\mathbb R.
\]
The mass distribution principle gives $\dim_HK\ge s$. This holds for every $s<1$ for the same compact set $K$ and the same measure $\mu$.
Therefore $\dim_HK=1$, and Theorem~\ref{fol:thm:wright-full-dimension} follows by $K\subset\mathcal W_b^{\mathrm{it}}\subset\mathbb R$.
\begin{remark}[Scope and attribution]
The case $b=2$ answers the exact question stated as Question~1.5 in
\cite[arXiv version]{SaitoDimension2021}. The proof combines classical
prime-counting and sieve inputs with a bounded-distortion Cantor construction; no historical priority assertion is made here. It is a
statement about the full set of permitted starting values and the compact sets explicitly constructed above. It gives no arithmetic
classification of a separately specified least Wright constant.
\end{remark}


\section{Local dimension after a finite prime prefix}
\label{fol:sec:local-dimension}

The global dimension theorem does not by itself localize the admissible set
inside a prescribed prime root cell.  This matters for a finite prefix whose
successive exponent ratios can be smaller than $2$.  A quantitative result
already established by Saito supplies the required bridge.

\begin{proposition}[Saito's quantitative bounded-schedule theorem]
\label{fol:dim:saito-local}
For every $R\ge2$ there is $M_2(R)>0$ with the following property.  For every
real sequence $(c_j)$ satisfying $2\le c_j\le R$, and every $M\ge M_2(R)$,
there is a prime $q\in[M,2M]$ such that, with $D_j=c_1\cdots c_j$,
\[
 \dim_H\left\{y\in[q^{1/c_1},(q+1)^{1/c_1}):
       \lfloor y^{D_j}\rfloor\text{ is prime for every }j\right\}
 \ge \left(1+\frac{R}{q\log q}\right)^{-1}.
\]
\end{proposition}
This is \cite[Theorem~6.1]{SaitoDimension2021}.  Its threshold is uniform
over all bounded schedules with the stated bound $R$.

\begin{lemma}[A local bridge across a large exponent ratio]
\label{fol:dim:bridge}
Suppose a prime-admissibility problem contains an exponent $E>0$, followed
by $E'=rE$, and all successive ratios after $E'$ are at least $2$.  Let $p$
be prime.  Suppose every point of
\[
 J=(p^{1/E},(p+1)^{1/E})
\]
satisfies all the required conditions up to exponent $E$.  If
\begin{equation}\label{fol:dim:bridge-window}
 M\ge M_2(4),\qquad p^r<M,\qquad 2M+1<(p+1)^r,
\end{equation}
then the set satisfying all required conditions and lying in $J$ has
Hausdorff dimension at least
\begin{equation}\label{fol:dim:bridge-bound}
 \left(1+\frac{4}{M\log M}\right)^{-1}.
\end{equation}
\end{lemma}
\begin{proof}
Refine every ratio after $E'$ into a finite block of real factors in
$[2,4)$, and prepend the initial factor $2$.  The products of this bounded
schedule, at its block boundaries, are
\[
 2,\quad 2E''/E',\quad 2E'''/E',\quad\ldots,
\]
where $E'',E''',\ldots$ are the later original exponents.
Proposition~\ref{fol:dim:saito-local} gives a prime $q\in[M,2M]$ and a set of
admissible $y$ in $[q^{1/2},(q+1)^{1/2})$ of dimension at least
$(1+4/(q\log q))^{-1}$.  Apply the map
\[
 x=y^{2/E'}.
\]
On this bounded interval the map and its inverse are Lipschitz, so
Hausdorff dimension is preserved.  Its image lies in
$[q^{1/E'},(q+1)^{1/E'})$, which is strictly inside $J$ by
\eqref{fol:dim:bridge-window}.  Thus all earlier conditions hold, the output
at $E'$ is $q$, and all later original outputs occur at refined block
boundaries and are prime.  Since $q\ge M$, the dimension is at least
\eqref{fol:dim:bridge-bound}.
\end{proof}

The seed condition in this lemma concerns the entire open prime root
cell.  A nonempty intersection of finitely many prime conditions need not
contain such a cell; we do not silently replace one condition by the
other.

\begin{lemma}[Strict continuation of every factorial seed]
\label{fol:dim:factorial-continuation}
Let $E\ge252$ be an integer and $p$ a prime.  Put $r=E!$.  There is a prime
$q$ satisfying
\[
 p^r<q<q+1<(p+1)^r.
\]
Consequently the closed prime root cell
$[q^{1/(Er)},(q+1)^{1/(Er)}]$ is strictly inside
$(p^{1/E},(p+1)^{1/E})$.
\end{lemma}
\begin{proof}
The consecutive-cubes result in \cite[Section~4]{CullyHugill2023} gives a
prime strictly between $m^3$ and $(m+1)^3$ whenever
$m\ge\exp(\exp(32.892))$.  Here set $t=r/3$ and $m=p^t$.
The integer $r$ is divisible by $3$, and already
\[
 \log m\ge\frac{252!}{3}\log2>\exp(32.892),
\]
so the threshold holds.  For example, the elementary bound
$252!>10^{90}$ is more than sufficient for the last inequality.
Since $q<(m+1)^3$ and both sides are integers,
$q+1\le(m+1)^3$.  As $t\ge2$, the binomial theorem gives
$p^t+1<(p+1)^t$.  It follows that
\[
 p^r<q<q+1\le(p^t+1)^3<(p+1)^r.
\]
Taking positive roots proves the last assertion.
\end{proof}

\begin{theorem}[Full local dimension for a factorial continuation]
\label{fol:dim:factorial-full}
Let $0<C_1<\cdots<C_t=E_0$ be a finite list of real exponents, where
$E_0\ge252$ is an integer.  Let $p_0$ be prime, and assume that every
$x$ in
\[
 J_0=(p_0^{1/E_0},(p_0+1)^{1/E_0})
\]
satisfies $\lfloor x^{C_i}\rfloor\in\mathbb P$ for every $i\le t$.
Define
\[
 E_{n+1}=E_n E_n!\qquad(n\ge0).
\]
Then
\begin{equation}\label{fol:dim:factorial-full-eq}
 \dim_H\left\{x\in J_0:
   \lfloor x^{C_i}\rfloor\in\mathbb P\ (1\le i\le t),\quad
   \lfloor x^{E_n}\rfloor\in\mathbb P\ (n\ge1)\right\}=1.
\end{equation}
No lower bound of $2$ is imposed on the ratios inside the finite prefix.
\end{theorem}
\begin{proof}
By Lemma~\ref{fol:dim:factorial-continuation}, construct primes $p_n$ such that
\[
 I_n=[p_n^{1/E_n},(p_n+1)^{1/E_n}],\qquad
 I_{n+1}\subset\operatorname{int}(I_n).
\]
Every point of the interior of $I_n$ satisfies the finite prefix and all
factorial-tail conditions up to $E_n$.  Choose $B>1$ with $I_0\subset(1,B)$.
Then $p_n<B^{E_n}$.  Writing $r_n=E_n!$, we have
\[
 \frac{r_n}{p_n}\ge\frac{E_n!}{B^{E_n}}\longrightarrow\infty.
\]
Indeed, $E!\ge\lfloor E/2\rfloor^{\lceil E/2\rceil}$, whose logarithm
minus $E\log B$ tends to infinity.  Since
$\log(1+1/p)\ge1/(2p)$ for $p\ge1$,
\[
 \frac{(p_n+1)^{r_n}}{p_n^{r_n}}
 =\exp\bigl(r_n\log(1+1/p_n)\bigr)
 \ge\exp\left(\frac{r_n}{2p_n}\right)\longrightarrow\infty.
\]
Thus, for all sufficiently large $n$, the choice
$M_n=\lceil2p_n^{r_n}\rceil$ satisfies
\[
 M_n\ge M_2(4),\qquad p_n^{r_n}<M_n,\qquad
 2M_n+1<(p_n+1)^{r_n}.
\]
Apply Lemma~\ref{fol:dim:bridge} to the seed cell $I_n$ and the jump to
$E_{n+1}$.  The set in \eqref{fol:dim:factorial-full-eq} consequently has
Hausdorff dimension at least
$(1+4/(M_n\log M_n))^{-1}$ for every sufficiently large $n$.
As $M_n\to\infty$, these lower bounds tend to $1$.  The ambient upper
bound is $1$.
\end{proof}

\begin{corollary}[The certified common prime prefix]
\label{fol:dim:original-seed}
Take
\[
 \begin{gathered}
 (C_1,C_2,C_3,C_4,C_5)=(4,6,7,12,252),\\
 p_0=5132842708382182842735812571653=29^{21}+24.
 \end{gathered}
\]
The exact primality certificate and root-cell comparisons in
Section~\ref{enc:sec:certificates}, with the verifier in
Appendix~\ref{enc:app:verification-code}, show that the complete factorial-tail admissible set
inside $(p_0^{1/252},(p_0+1)^{1/252})$ has Hausdorff dimension $1$.
\end{corollary}
\begin{proof}
Those certified comparisons place the seed cell inside the region with first four
prime outputs $(3,5,7,29)$.  Theorem~\ref{fol:dim:factorial-full} applies.
\end{proof}

The full admissible set in this corollary is much larger, in the metric
sense, than the particular compact set obtained by retaining only the
first two admissible primes at each stage.  The latter set's
upper box dimension zero is entirely compatible with
Corollary~\ref{fol:dim:original-seed}.

\section{A metric obstruction to uniform radical approximation}
\label{fol:sec:metric-obstruction}

The very rapid approximation used to prove whole-family algebraic
independence in the explicit prime-tree construction imposes a geometric cost.
The following elementary estimate identifies that cost without making any
assumption about the independence proof itself.

\begin{theorem}[Uniform prime-root approximation forces zero dimension]
\label{fol:dim:approximation-obstruction}
Let $K\subset[1,B]$ be nonempty and bounded, with $B>1$.  Suppose $E_n>0$,
$0<\varepsilon_n\to0$, and for every $x\in K$ and every $n$ there is a prime
$p=p(x,n)$ such that
\[
 |x-p^{1/E_n}|\le\varepsilon_n.
\]
If
\begin{equation}\label{fol:dim:approximation-threshold}
 \frac{E_n}{\log(1/\varepsilon_n)}\longrightarrow0,
\end{equation}
then $\dim_H K=0$ and the lower box dimension of $K$ is zero.  More
strongly, $\dim_H(K^r)=0$ for every positive integer $r$.
\end{theorem}
\begin{proof}
For sufficiently large $n$, every approximating radical lies below
$B+1$, hence $p\le(B+1)^{E_n}$.  It is unnecessary to use a prime-counting
estimate: there are at most $(B+1)^{E_n}$ possible positive integer
radicands.  Thus $K$ is covered by at most $(B+1)^{E_n}$ intervals of
length $2\varepsilon_n$.  For every fixed $s>0$, the sum of the $s$th
powers of their lengths is at most
\[
 (B+1)^{E_n}(2\varepsilon_n)^s,
\]
whose logarithm tends to $-\infty$ by
\eqref{fol:dim:approximation-threshold}.  This proves $\dim_H K=0$.
The same cover gives
\[
 \liminf_{\rho\downarrow0}
 \frac{\log N(K,\rho)}{\log(1/\rho)}=0,
\]
by taking $\rho=2\varepsilon_n$.  Finally, taking products gives a cover
of $K^r$ by at most $(B+1)^{rE_n}$ sets of diameter
$2\sqrt r\,\varepsilon_n$.  The identical estimate proves the last claim.
\end{proof}

The theorem remains valid if ``prime'' is replaced by ``positive
integer.''  It does not assert that packing dimension or upper box
dimension is zero: a small number of covering scales suffices for the
conclusions proved here, whereas those other dimensions require further
control.  In particular, a full-Hausdorff-dimension subset of a prime-admissible
set cannot retain the uniform approximation regime
\eqref{fol:dim:approximation-threshold}. For compact algebraically independent
families, Corollary~\ref{fol:cor:analytic-independent-small} gives the stronger
obstruction that positive Hausdorff dimension is impossible regardless
of the chosen approximation regime.


\section{Algebraic consequences of the dimension contrast}
\label{fol:sec:dimension-algebra}

The dimension conclusion has an algebraic consequence substantially stronger
than uncountability. We use the following result of Edgar and Miller.

\begin{theorem}[Edgar--Miller]\label{fol:thm:edgar-miller}
Let $F\subseteq\mathbb R$ be an analytic subring. Either $F=\mathbb R$, or
$\dim_H(F^m)=0$ for every positive integer $m$.
\end{theorem}

The analytic version follows from \cite[Theorem~1 and the subsequent
remarks]{EdgarMiller2003}. Here ``analytic'' has its descriptive-set-theoretic
meaning, namely a continuous image of a Polish space, not a regularity
condition on functions.

\begin{corollary}[The full spectrum generates the real field as a ring]
\label{fol:cor:prime-ring}
Let $E_1>0$ and $E_{k+1}/E_k\ge2$. Write
\[
 \mathcal P(E)=\{x>1:\lfloor x^{E_k}\rfloor
                 \text{ is prime for every }k\ge1\}.
\]
Then the unital ring generated by $\mathcal P(E)$ is all of $\mathbb R$:
\[
 \boxed{\mathbb Z[\mathcal P(E)]=\mathbb R.}
\]
Equivalently, every real number is an integer polynomial in finitely many
members of one and the same prime-representing spectrum.
\end{corollary}

\begin{proof}
For each $k$, the primality condition is a countable union of half-open
root intervals. Therefore $\mathcal P(E)$ is Borel. Its generated ring is
a countable union of continuous polynomial images of finite powers of
$\mathcal P(E)$, hence is analytic. It has positive Hausdorff dimension
by Corollary~\ref{fol:cor:arbitrary-initial-scale}, because it contains
$\mathcal P(E)$. Theorem~\ref{fol:thm:edgar-miller} forces this ring to be
$\mathbb R$.
\end{proof}

This is a representation theorem, with no asserted effective bound on the
number or size of its parameters. It does not settle the arithmetic nature
of a named real number represented in this way.

\begin{corollary}[A sharp Hausdorff obstruction for independent families]
\label{fol:cor:analytic-independent-small}
If $K\subseteq\mathbb R$ is analytic and algebraically independent over
$\mathbb Q$, then
\[
 \dim_H(K^m)=0\qquad(m\ge1).
\]
\end{corollary}

\begin{proof}
The ring $F=\mathbb Q[K]$ is analytic, by the same polynomial-image
argument. It is proper: if $\sqrt2=P(x_1,\ldots,x_r)$ for a polynomial
$P\in\mathbb Q[X_1,\ldots,X_r]$ at distinct elements of $K$, algebraic
independence would force $P^2-2$ to be the zero polynomial. This is
impossible in the polynomial ring over $\mathbb Q$ (a nonconstant $P$
has positive degree, and a constant rational $P$ cannot square to $2$).
Thus Theorem~\ref{fol:thm:edgar-miller} gives $\dim_H(F^m)=0$, and
$K^m\subseteq F^m$ completes the proof.
\end{proof}

Accordingly, a compact algebraically independent family cannot retain the
positive Hausdorff dimension of its full prime-representing spectrum.
This explains why the zero Hausdorff dimension of the compact independent
sets constructed below is unavoidable. The assertion concerns Hausdorff
dimension and its finite Cartesian powers; it imposes no corresponding
upper box dimension bound by itself.


\section{Separation governed by the exponent support}
\label{fol:sec:support-separation}

The degree of the ambient radical field is not always the relevant
exponent in its norm estimate.  A polynomial whose exponent vectors lie
in a low-dimensional affine space has many conjugates of the same
absolute value.  Retaining that multiplicity gives the following
refinement.

For a nonzero polynomial $P=\sum_{\nu\in S}a_\nu X^\nu$, its
\emph{support rank} is
\[
 \rho(P)=\dim_{\mathbb Q}\operatorname{span}_{\mathbb Q}
       \{\nu-\mu:\nu,\mu\in S\}.
\]
It is the dimension of the Newton polytope, not the number of variables
or the number of terms. Multiplication by a monomial does not change it.
For $T\ge1$ put
\[
 B_T(P)=\max\left\{2,\sum_{\nu}|a_\nu|T^{|\nu|}\right\},\qquad
 L_T(P)=\sum_{|\nu|\ge1}|a_\nu|\,|\nu|T^{|\nu|-1}.
\]
The latter is a Lipschitz constant on $[0,T]^r$ in the supremum norm.

\begin{theorem}[The support-lattice norm bound]
\label{fol:thm:support-norm}
Let $p_1,\ldots,p_r$ be distinct primes, let $E\ge2$ be an integer, and
put $b_i=p_i^{1/E}\le T$, using positive real roots. Let
$0\ne P\in\mathbb Z[X_1,\ldots,X_r]$ satisfy
$\deg_{X_i}P<E$ for every $i$. Fix $\nu_0\in\operatorname{supp}P$, and
let $d_1,\ldots,d_\ell$ be the nonzero Smith invariant factors of the
integer matrix with rows $\nu-\nu_0$, where
$\nu\in\operatorname{supp}P$ and $\ell=\rho(P)$. With an empty product
interpreted as $1$, define
\[
 M_E(P)=\prod_{j=1}^{\ell}\frac{E}{\gcd(E,d_j)}.
\]
Then
\begin{equation}\label{fol:eq:support-norm}
 |P(b_1,\ldots,b_r)|\ge B_T(P)^{1-M_E(P)}
 \ge B_T(P)^{1-E^{\rho(P)}}>0.
\end{equation}
\end{theorem}

\begin{proof}
The prime-radical basis theorem gives
\[
 [\mathbb Q(b_1,\ldots,b_r):\mathbb Q]=E^r
\]
and shows that $P(b_1,\ldots,b_r)\ne0$. For completeness, the monomials
$b_1^{j_1}\cdots b_r^{j_r}$, $0\le j_i<E$, have rational $E$th powers.
Two different monomials have irrational quotient, by valuations at the
distinct primes $p_i$. The real-radical independence theorem of Carr and
O'Sullivan~\cite[Theorem~1.1]{CarrOSullivan2009} makes these $E^r$
monomials linearly independent. They also span the generated field.

Let $F=\mathbb Q(b_1,\ldots,b_r)$, let $\zeta$ be a primitive $E$th
root of unity, and write $\beta=P(b_1,\ldots,b_r)$. Its $E^r$ embeddings
are indexed by $t=(t_1,\ldots,t_r)\in(\mathbb Z/E\mathbb Z)^r$ and send
$b_i$ to $\zeta^{t_i}b_i$. Consider the subgroup
\[
 H_E=\{t:(\nu-\nu_0)\cdot t=0\pmod E
                  \text{ for every }\nu\in\operatorname{supp}P\}.
\]
Smith normal form gives
\[
 h_E:=|H_E|=E^{r-\ell}\prod_{j=1}^{\ell}\gcd(E,d_j),
 \qquad E^r/h_E=M_E(P).
\]
For every $t\in H_E$,
\[
 P(\zeta^{t_1}b_1,\ldots,\zeta^{t_r}b_r)
       =\zeta^{\nu_0\cdot t}\beta.
\]
Thus at least $h_E$ conjugates have absolute value $|\beta|$. All
conjugates have absolute value at most $B_T(P)$. Since $\beta$ is a
nonzero algebraic integer, its norm is a nonzero integer, and hence
\[
 1\le|N_{F/\mathbb Q}(\beta)|
 \le |\beta|^{h_E}B_T(P)^{E^r-h_E}.
\]
Taking $h_E$th roots proves the first inequality. The second follows
from $M_E(P)\le E^\ell$.
\end{proof}

\begin{corollary}[A finite nonvanishing certificate]
\label{fol:cor:support-certificate}
Under the assumptions of Theorem~\ref{fol:thm:support-norm}, suppose
$\alpha_i\in[0,T]$ and $|\alpha_i-b_i|\le\varepsilon$. If
\[
 L_T(P)\varepsilon\le\tfrac12 B_T(P)^{1-M_E(P)},
\]
then
\[
 |P(\alpha_1,\ldots,\alpha_r)|
 \ge\tfrac12 B_T(P)^{1-M_E(P)}>0.
\]
The two real numbers $P(\alpha_1,\ldots,\alpha_r)$ and
$P(b_1,\ldots,b_r)$ have the same sign.
\end{corollary}
\begin{proof}
Their difference has absolute value at most $L_T(P)\varepsilon$.
Apply~\eqref{fol:eq:support-norm}.
\end{proof}

\begin{theorem}[Sparse-support transfer]
\label{fol:thm:sparse-transfer}
Let $\alpha_1,\ldots,\alpha_r>0$ be distinct. Let $E_n$ be positive
integers tending to infinity, and let $0<\varepsilon_n\to0$. Suppose
that, for all sufficiently large $n$, there are pairwise distinct primes
$p_{1,n},\ldots,p_{r,n}$ such that
\[
 \max_i|\alpha_i-p_{i,n}^{1/E_n}|\le\varepsilon_n.
\]
For an integer $\ell\ge1$, assume
\begin{equation}\label{fol:eq:support-limsup}
 \limsup_{n\to\infty}\frac{\log(1/\varepsilon_n)}{E_n^\ell}=\infty.
\end{equation}
Then every nonzero polynomial
$P\in\overline{\mathbb Q}[X_1,\ldots,X_r]$ of support rank at most
$\ell$ satisfies $P(\alpha_1,\ldots,\alpha_r)\ne0$.
The same assertion holds for Laurent polynomials, with their support
viewed as a subset of $\mathbb Z^r$.
\end{theorem}

\begin{proof}
First take an integer polynomial. The approximants and their limits
lie in $[0,T]^r$ for some fixed $T$. The individual degree conditions
hold eventually. Along a subsequence supplied by
\eqref{fol:eq:support-limsup},
\[
 L_T(P)\varepsilon_n
 <\tfrac12 B_T(P)^{1-E_n^\ell}
 \le\tfrac12 B_T(P)^{1-M_{E_n}(P)}.
\]
Corollary~\ref{fol:cor:support-certificate} proves the assertion.

For algebraic coefficients, choose a finite Galois extension $L/\mathbb Q$
containing them and form the nonzero polynomial
\[
 Q(X)=\prod_{\sigma\in\operatorname{Gal}(L/\mathbb Q)}\sigma(P)(X)
       \in\mathbb Q[X].
\]
If $S\subset\nu_0+V$, with $\dim V=\rho(P)$, then the support of $Q$
is contained in $[L:\mathbb Q]\nu_0+V$. Therefore $\rho(Q)\le\rho(P)$.
Clearing its rational coefficient denominators and applying the
integer-polynomial case gives $Q(\alpha)\ne0$, hence $P(\alpha)\ne0$.
Finally, clearing negative exponents by multiplication by a monomial
does not change support rank and does not affect vanishing at positive
coordinates.
\end{proof}

\begin{remark}[What is saved]
A homogeneous polynomial in $r$ variables has support rank at most
$r-1$, so its lower bound uses $E^{r-1}$ in place of $E^r$. A polynomial
such as $X_1^2X_2^3-7X_3^5$ has support rank one regardless of its three
variables. The estimate for it uses $E$, not $E^3$.
No estimate of this form asserts that the growth conditions below are
necessary; they are sufficient conditions produced by the norm method.
\end{remark}

\section{Slower schedules and multiplicative transcendence}
\label{fol:sec:slower-schedules}

Here is a completely specified binary construction to which the sparse
criterion applies. It can use any starting prime cell contained in
$(1,2)$; no large-prime certificate is needed for its existence.
Put $E_0=252$, $s_{\varnothing}=2$, and choose integer multipliers $R_n$
such that
\begin{equation}\label{fol:eq:explicit-R-range}
 3\mid R_n,\qquad R_n\ge3\cdot10^{15},\qquad E_{n+1}=E_nR_n.
\end{equation}
For a binary word $w$ of length $n$, define $s_{w0}$ and $s_{w1}$ to
be the first and second primes strictly above $s_w^{R_n}$, and set
\[
 I_w=[s_w^{1/E_n},(s_w+1)^{1/E_n}].
\]

\begin{proposition}[Uniform root approximation]
\label{fol:prop:slow-tree}
The cells are strictly nested and their siblings are disjoint. Every
binary sequence determines a unique limit, and the set $K$ of these
limits is compact, perfect, and consists of prime-representing constants
at every exponent $E_n$. For
$a=2^{1/252}>1$, every descendant $\alpha$ of a level-$n$ cell satisfies
\begin{equation}\label{fol:eq:slow-error}
 0<\alpha-s_w^{1/E_n}
 <\frac{12}{E_nR_n}a^{-E_nR_n/3}.
\end{equation}
Distinct limits have distinct prime labels at all sufficiently late
levels.
\end{proposition}

\begin{proof}
We use the explicit theorem that $(m^3,(m+1)^3)$ contains a prime for
all integers $m\ge\exp(\exp(32.892))$, proved in
Cully-Hugill~\cite[Section~4]{CullyHugill2023}. At a node with prime
$s$, exponent $E$, and multiplier $R$, the integer $M=s^{R/3}$ satisfies
$M\ge2^{10^{15}}$. In particular
$\log\log M\ge15\log10+\log\log2>32.892$.
The first two primes above $s^R$ are therefore below $(M+2)^3$.
Since $R/3\ge2$,
\[
 (s+1)^{R/3}-s^{R/3}>2.
\]
The upper endpoints of both child cells lie strictly below the upper
endpoint of the parent; their lower endpoints lie strictly above the
parent's lower endpoint. Distinct odd prime labels give disjoint cells.
Their diameters are at most $1/E_n\to0$, proving the nested-limit
assertions and the homeomorphism with binary sequence space.

All cells remain in $[a,3^{1/252}]\subset(1,2)$. If $b=s^{1/E}$, a
descendant belongs to one of the next two child cells, so concavity gives
\[
 0<\alpha-b
 \le b\bigl((1+2/M)^{3/(ER)}-1\bigr)
 <\frac{6b}{ER M}
 <\frac{12}{ER}a^{-ER/3}.
\]
Distinct branches separate at one level and cannot meet thereafter.
\end{proof}

\begin{theorem}[The support-rank growth criterion]
\label{fol:thm:rank-growth}
For the preceding tree, let $\ell\ge1$. If
\begin{equation}\label{fol:eq:rank-growth}
 \limsup_{n\to\infty}\frac{R_n}{E_n^{\ell-1}}=\infty,
\end{equation}
then no nonzero polynomial over $\overline{\mathbb Q}$ of support rank
at most $\ell$ vanishes on any finite tuple of distinct members of $K$.
In particular:
\begin{enumerate}
\item If $\limsup R_n=\infty$, then for distinct
$\alpha_1,\ldots,\alpha_r\in K$ and
$(m_1,\ldots,m_r)\in\mathbb Z^r\setminus\{0\}$,
\begin{equation}\label{fol:eq:monomial-transcendence}
 \alpha_1^{m_1}\cdots\alpha_r^{m_r}
       \notin\overline{\mathbb Q}.
\end{equation}
The same holds for nonzero rational exponent vectors.
\item If \eqref{fol:eq:rank-growth} holds with $\ell=r$, every $r$ distinct
members of $K$ are algebraically independent.
\item If \eqref{fol:eq:rank-growth} holds with $\ell=r-1$ and $r\ge2$,
then for distinct $\alpha_1,\ldots,\alpha_r\in K$ the $r-1$ ratios
$\alpha_2/\alpha_1,\ldots,\alpha_r/\alpha_1$ are algebraically independent.
\end{enumerate}
\end{theorem}

\begin{proof}
Taking logarithms in~\eqref{fol:eq:slow-error} gives
\[
 \frac{\log(1/\varepsilon_n)}{E_n^\ell}
 \ge\frac{\log a}{3}\frac{R_n}{E_n^{\ell-1}}
       +\frac{\log(E_nR_n/12)}{E_n^\ell}.
\]
Apply Theorem~\ref{fol:thm:sparse-transfer} to any fixed finite tuple.
If the monomial in (i) equaled an algebraic number $c$, the nonzero
Laurent polynomial $X^m-c$, of support rank one, would vanish. Rational
exponents reduce to integer exponents by raising to a common positive
denominator. Part (ii) uses the bound $\rho(P)\le r$.
For (iii), a nonzero polynomial relation among the ratios would
homogenize to a nonzero homogeneous polynomial in $r$ variables, whose
support rank is at most $r-1$.
\end{proof}

\begin{corollary}[A factorial-scale example]
\label{fol:cor:factorial-scale}
Take $N=10^{15}$ and
\[
 R_n=3(N+n),\qquad E_n=252\,3^n\frac{(N+n-1)!}{(N-1)!}.
\]
Then $K$ is a compact perfect set for which every nontrivial finite
monomial in distinct members is transcendental. In particular every
member, every product of finitely many distinct members, and every
quotient of two distinct members is transcendental. The set has upper
box dimension zero.
\end{corollary}
\begin{proof}
Here $R_n\to\infty$, so Theorem~\ref{fol:thm:rank-growth}(i) applies. There
are $2^n$ level-$n$ cells of diameter at most $E_n^{-1}$, and
$\log E_n/n\to\infty$ by the factorial formula. If
$E_n^{-1}\le\delta<E_{n-1}^{-1}$, the covering number satisfies
\[
 \frac{\log N_\delta(K)}{\log(1/\delta)}
 \le\frac{n\log2}{\log E_{n-1}}\longrightarrow0.
\]
\end{proof}

\begin{corollary}[A smaller full-independence schedule]
\label{fol:cor:quasipolynomial-schedule}
With $E_0=252$, define
\[
 R_n=3\left\lceil 10^{15}+
                   \exp\bigl((\log E_n)^2\bigr)\right\rceil,
 \qquad E_{n+1}=E_nR_n.
\]
The entire binary limit set is algebraically independent over
$\overline{\mathbb Q}$ and has upper box dimension zero.
\end{corollary}
\begin{proof}
For every fixed $\ell$, $R_n/E_n^{\ell-1}\to\infty$. Thus
Theorem~\ref{fol:thm:rank-growth} excludes every nonzero polynomial relation.
Also $E_{n+1}\ge E_n^2$ eventually, which gives the dimensional
conclusion as above.
\end{proof}

\begin{remark}[Comparison with factorial continuation]
The last multiplier is asymptotic to $3\exp((\log E)^2)$, substantially
smaller than $E!$. The general $r$-variable growth condition
$R_n/E_n^{r-1}\to\infty$ is already present for the countable
sub-boundary in Saito and Takeda~\cite[Theorem~1.3]{SaitoTakeda2022}.
Here it applies to the entire specified binary tree; the support-rank
refinement gives additional conclusions at weaker growth rates.
Neither the factorial-scale example nor the sparse-support criterion
alone proves full algebraic independence when $R_n$ grows only linearly
in $n$. For instance, a relation involving $X_1+X_2-1$ has support rank
two and is not excluded by the rank-one argument.
\end{remark}

\section{A consequence for a pre-existing prime-constant boundary}
\label{fol:sec:boundary-transfer}

We record a direct consequence for the family studied by Saito and
Takeda, rather than only for the new binary subfamilies. Let $c_1$ be a
positive integer and let $c_k\ge2$ be integers for $k\ge2$. Put
\[
 C_k=c_1c_2\cdots c_k,\qquad
 W(c_k)=\{A>1:\lfloor A^{C_k}\rfloor\in\mathbb P
                    \text{ for every }k\ge1\}.
\]
The left sub-boundary is
\[
 \widetilde\partial_L W(c_k)
 =\{\sup U\in\mathbb R:U\text{ is a connected component of }
                          \mathbb R\setminus W(c_k)\}.
\]
The right sub-boundary is defined by replacing $\sup U$ with $\inf U$;
write it as $\widetilde\partial_R W(c_k)$ and set
$\widetilde\partial W(c_k)=\widetilde\partial_L W(c_k)
\cup\widetilde\partial_R W(c_k)$. Only finite endpoints are included.
Under the hypotheses below,
$W(c_k)$ is closed and perfect by
Saito--Takeda~\cite[Theorem~1.2]{SaitoTakeda2022}, and these endpoints
belong to $W(c_k)$.

\begin{theorem}[Multiplicative transcendence on the left sub-boundary]
\label{fol:thm:left-boundary-monomials}
If $\limsup_{k\to\infty}c_{k+1}=\infty$, then, for every finite tuple
of distinct
$A_1,\ldots,A_r\in\widetilde\partial_L W(c_k)$ and every nonzero
$m\in\mathbb Z^r$,
\[
 A_1^{m_1}\cdots A_r^{m_r}\notin\overline{\mathbb Q}.
\]
More generally, if
\[
 \limsup_{k\to\infty}\frac{c_{k+1}}{C_k^{\ell-1}}=\infty,
\]
no nonzero algebraic-coefficient polynomial of support rank at most
$\ell$ vanishes at such a tuple.
\end{theorem}

\begin{proof}
Set $\theta=21/40$ and $I=\{k:c_{k+1}\ge3\}$, which is infinite.
We use the following precise published input,
Saito--Takeda~\cite[Lemma~5.2(1)]{SaitoTakeda2022}: for each
$A\in\widetilde\partial_L W(c_k)$, writing
$p_k(A)=\lfloor A^{C_k}\rfloor$, one has for all sufficiently large
$k\in I$
\[
 0\le A-p_k(A)^{1/C_k}
 \le A\,p_1(A)^{(\theta-1)C_{k+1}/c_1}.
\]
Its hypotheses are $c_1>0$, integral $c_{k+1}\ge2$, and infinitely many
$c_{k+1}\ge3$. Our additional integrality of $c_1$ ensures that every
radical exponent $C_k$ is an integer.

For a fixed finite tuple the prime labels become pairwise distinct,
since $A_i\ne A_j$ implies
$|A_i^{C_k}-A_j^{C_k}|\to\infty$. With
$T=1+\max_i A_i$ and
\[
 \lambda=\frac{1-\theta}{c_1}\min_i\log p_1(A_i)>0,
\]
all sufficiently large $k\in I$ satisfy
\[
 \max_i|A_i-p_k(A_i)^{1/C_k}|
 \le T\exp(-\lambda C_{k+1}).
\]
Along the subsequence where $c_{k+1}/C_k^{\ell-1}$ is unbounded,
\[
 \frac{\log(1/\varepsilon_k)}{C_k^\ell}
 \ge\lambda\frac{c_{k+1}}{C_k^{\ell-1}}
       -\frac{\log T}{C_k^\ell}\longrightarrow\infty.
\]
This subsequence lies in $I$ eventually.
Theorem~\ref{fol:thm:sparse-transfer} proves the more general assertion.
Its rank-one case applied to $X^m-a$, for algebraic $a$, proves the
monomial statement.
\end{proof}

\begin{corollary}[The factorial exponent family]
\label{fol:cor:boundary-factorial}
Let
\[
 W_{!}=\{A>1:\lfloor A^{k!}\rfloor\in\mathbb P
                   \text{ for all }k\ge1\}.
\]
Every nontrivial finite monomial in distinct elements of
$\widetilde\partial_L W_{!}$ is transcendental. In particular, for
distinct $A,B\in\widetilde\partial_L W_{!}$, the quotient $A/B$ is
transcendental.
\end{corollary}
\begin{proof}
Take $c_k=k$, including $c_1=1$, in
Theorem~\ref{fol:thm:left-boundary-monomials}.
\end{proof}

\begin{remark}[The precise comparison]
Saito--Takeda~\cite[Proposition~8.2 and Corollary~8.3]{SaitoTakeda2022}
prove rational linear independence of two appropriate boundary
constants under $c_{k+1}\to\infty$. On the all-left-sub-boundary case,
Corollary~\ref{fol:cor:boundary-factorial} strengthens irrationality of the
quotient to its transcendence, and the theorem treats arbitrary finite
monomials. Their proposition also allows one right-sub-boundary point.
The general support-rank argument does not cover that extra case:
its natural approximant is a root of $p_k+1$, which need not be prime.
The quotient case can nevertheless be recovered by the separate
argument below. These are proved refinements of a published theorem,
not a claim that the cited paper posed the refinements as open
conjectures. We have not established priority for the general
support-lattice norm observation.
\end{remark}

\begin{theorem}[Transcendental quotients with one arbitrary boundary point]
\label{fol:thm:boundary-quotient}
Assume $\limsup_{k\to\infty}c_{k+1}=\infty$. Let
\[
 A_1\in\widetilde\partial W(c_k),\qquad
 A_2\in\widetilde\partial_L W(c_k),\qquad A_1<A_2,
\]
where $\widetilde\partial W(c_k)$ is the union of the left and right
sub-boundaries. Then $A_2/A_1$ is transcendental.
\end{theorem}

\begin{proof}
Let $E=C_k$, and write $p_i(k)=\lfloor A_i^{C_k}\rfloor$.
If $A_1$ is a left sub-boundary point, put $q_1(k)=p_1(k)$; if it is a
right sub-boundary point, put $q_1(k)=p_1(k)+1$. Put $q_2(k)=p_2(k)$.
The two parts of Saito--Takeda~\cite[Lemma~5.2]{SaitoTakeda2022} show
that there are constants $K,\lambda>0$, depending on the fixed pair,
such that for all sufficiently large $k\in I=\{k:c_{k+1}\ge3\}$,
\[
 |q_i(k)^{1/E}-A_i|\le K\exp(-\lambda C_{k+1})\qquad(i=1,2).
\]
For the right endpoint, the precise bound used is
\[
 0\le(p_1(k)+1)^{1/C_k}-A_1
 \le(p_1(1)+1)^{1/c_1}A_1^{(\theta-1)C_{k+1}},
 \qquad\theta=21/40.
\]
The left-endpoint bound was stated in the preceding proof.

For all sufficiently large $k$, the integer $q_2(k)$ is a prime greater
than $q_1(k)$, because $A_2>A_1$. The polynomial
\[
 F_k(X)=q_1(k)X^E-q_2(k)
\]
is Eisenstein at $q_2(k)$ and is therefore irreducible over $\mathbb Q$.
In particular, its positive root
$z_k=(q_2(k)/q_1(k))^{1/E}$ has degree $E$.
Both root approximants stay in a fixed compact subinterval of
$(0,\infty)$, so the quotient map is Lipschitz there. Consequently
\begin{equation}\label{fol:eq:quotient-approximation}
 |z_k-A_2/A_1|\le K'\exp(-\lambda C_{k+1})
\end{equation}
for a fixed constant $K'>0$.

Suppose $z=A_2/A_1$ were algebraic, and choose a nonzero
$Q\in\mathbb Z[X]$ with $Q(z)=0$, of degree $d$. For $E>d$,
irreducibility of $F_k$ makes $\operatorname{Res}(F_k,Q)$ a nonzero
integer. Choose a fixed $T>1$ with $q_i(k)^{1/E}\le T$ and $z_k\le T$
eventually, and set $B=\max\{2,\sum_j|[X^j]Q|T^j\}$.
All complex roots of $F_k$ have modulus $z_k$, whence
\[
 1\le|\operatorname{Res}(F_k,Q)|
 =q_1(k)^d\prod_{\zeta^E=1}|Q(\zeta z_k)|
 \le q_1(k)^d|Q(z_k)|B^{E-1}.
\]
Since $q_1(k)\le T^E$, this yields
\[
 |Q(z_k)|\ge T^{-dE}B^{1-E}\ge\exp(-H E)
\]
for a fixed $H>0$. On the other hand,
\eqref{fol:eq:quotient-approximation} and a derivative bound for $Q$ give
$|Q(z_k)|\le K''\exp(-\lambda E c_{k+1})$.
Choose $k\to\infty$ along a subsequence with $c_{k+1}\to\infty$.
The last two estimates are incompatible. Therefore $z$ is
transcendental.
\end{proof}

\begin{remark}
For $c_k=k$, this gives transcendence of $A_2/A_1$ for the full pair
configuration of Saito--Takeda's Proposition~8.2 and Corollary~8.3,
including their allowed right sub-boundary point. The proof uses the
published approximation estimates and a classical integer resultant;
it does not assume algebraic independence of the pair.
\end{remark}


\section{Relative independence of complete exponential towers}\label{fol:sec:relative-towers}
The perfect prime set can be thinned without changing its exponent
sequence. This permits relative independence over prescribed constants
and, more strongly, independence of the logarithm and every finite exponential iterate of
every retained point. The selection argument is a specialization of the
classical perfect-set method of Mycielski \cite{Mycielski1964}; we include
the proof and its effective refinement because the retained prime
identities and the coefficient-field issue are essential here.

For $x>1$, define
\[
 T_{-1}(x)=\log x,\qquad T_0(x)=x,\qquad
 T_{j+1}(x)=\exp(T_j(x))\quad(j\ge0).
\]
An indexed family is algebraically independent if every polynomial in
finitely many distinct indexed entries, with coefficients in the base
field, is nonzero at those entries. Thus the assertion below also
excludes coincidences between entries with different indices.

\begin{lemma}[Functional independence of exponential towers]
\label{fol:lem:tower-functions}
For each $m\ge0$, the real-analytic functions
$T_{-1},T_0,\ldots,T_m$ are algebraically independent over $\mathbb C$.
More generally, for independent real variables $x_1,\ldots,x_r$, the
functions $T_j(x_i)$, $1\le i\le r$, $-1\le j\le m$, are algebraically
independent over $\mathbb C$.
\end{lemma}
\begin{proof}
For a multi-index $a=(a_{-1},a_0,\ldots,a_m)$ of nonnegative integers
and $x>e$, write $M_a(x)=\prod_{j=-1}^mT_j(x)^{a_j}$. Then
\[
 \log M_a(x)=a_{-1}\log\log x+a_0\log x
 +\sum_{j=1}^m a_jT_{j-1}(x).
\]
If $a\ne b$, the largest index at which they differ determines the sign
of $\log(M_a/M_b)$ for sufficiently large $x$, and its absolute value
tends to infinity. Consequently, among the finitely many monomials of a
nonzero polynomial, the lexicographically largest, reading indices from
$m$ down to $-1$, dominates all the others as $x\to\infty$. Its nonzero
coefficient prevents cancellation. This proves the one-variable claim.
For several variables, group a purported identity by the monomials in
$T_j(x_1)$. Holding the other variables fixed, the one-variable result
forces every coefficient function to vanish. Induction on $r$ finishes
the proof.
\end{proof}

\begin{lemma}[Analytic avoidance on a perfect set]
\label{fol:lem:analytic-avoidance}
Let $J$ be a connected open real interval and let $P\subset J$ be a
nonempty compact perfect set. If $f:J^r\to\mathbb C$ is real analytic
and is not identically zero, its zero set in $P^r$ is closed and nowhere
dense relative to $P^r$.
\end{lemma}
\begin{proof}
Closedness follows from continuity. If the zero set had nonempty
relative interior, $f$ would vanish on a product $U_1\times\cdots\times
U_r$ of nonempty relatively open subsets of $P$. Each $U_i$ has an
accumulation point in $J$. Apply the one-variable analytic identity
theorem successively to the variables. The resulting identity on
$J^r$ contradicts the hypothesis.
\end{proof}

\begin{proposition}[Cantor fusion with retained constraints]
\label{fol:prop:analytic-fusion}
Under the assumptions on $J,P$ in Lemma~\ref{fol:lem:analytic-avoidance},
let $(f_\nu)_{\nu\ge1}$ be a countable family of nonzero real-analytic
functions, where $f_\nu$ has some finite arity $r_\nu$. There is a compact
perfect set $K\subset P$, homeomorphic to $2^{\mathbb N}$, such that
\[
 f_\nu(x_1,\ldots,x_{r_\nu})\ne0
\]
whenever the $x_i$ are pairwise distinct points of $K$.
Every property holding pointwise throughout $P$ is retained on $K$.
\end{proposition}
\begin{proof}
We construct a binary family of nonempty relatively clopen subsets of
$P$, with disjoint siblings, nested closures, and diameters at level $n$
at most $2^{-n}$. Such sets exist because a compact perfect subset of
the line contains a Cantor subset; we may first replace $P$ by that
subset. Suppose that the level $n-1$ sets have been selected. Choose two
prospective points in each parent, giving $2^n$ point variables in a
product of relatively open perfect sets. Exclude all coincidences, and,
for each $\nu\le n$, exclude the zeros of $f_\nu$ on every ordered choice
of $r_\nu$ distinct point variables. There are only finitely many
exclusions. Each is nowhere dense by
Lemma~\ref{fol:lem:analytic-avoidance}, applied to the corresponding
coordinate variables; the diagonal exclusions are nowhere dense because
$P$ has no isolated points. A surviving configuration therefore exists.
Continuity permits pairwise disjoint, sufficiently small clopen
neighborhoods of its points on whose corresponding products every
required value remains nonzero. These are the new level sets.
The intersection of the unions of the level sets is a compact Cantor
set $K$. A finite tuple of distinct points eventually lies in distinct
level sets. At any later stage that also handles $f_\nu$, the relevant
product has been protected from its zeros. This proves the claim.
\end{proof}

\begin{theorem}[Prime constants independent over any countable field]
\label{fol:thm:relative-towers}
Let $P\subset(1,\infty)$ be a nonempty compact perfect set and let
$(C_n)_{n\ge0}$ be positive integers such that
$\lfloor x^{C_n}\rfloor$ is prime for every $x\in P$ and every $n$.
For any countable subfield $F\subset\mathbb C$, there exists a compact
perfect $K\subset P$ such that the entire indexed family
\begin{equation}\label{fol:eq:tower-family}
 \mathcal T(K)=\bigl(T_j(x)\bigr)_{(x,j)\in K\times\{-1,0,1,2,\ldots\}}
\end{equation}
is algebraically independent over $F$ and over every algebraic
extension of $F$ inside $\mathbb C$.
The original exponent sequence works unchanged for every base point
$x\in K$. Any upper box dimension or common-prefix bound already proved
for $P$ is inherited by $K$.
\end{theorem}
\begin{proof}
For every $r,m$ and every nonzero polynomial
$Q\in F[Y_{i,j}:1\le i\le r,-1\le j\le m]$, include the function
\[
 f_Q(x_1,\ldots,x_r)=Q\bigl(T_j(x_i)\bigr)
\]
in the avoidance list. This is a countable family, and
Lemma~\ref{fol:lem:tower-functions} shows that no listed function is
identically zero. Apply Proposition~\ref{fol:prop:analytic-fusion}.
Every finite selection of distinct pairs $(x,j)$ embeds in such a
rectangular array, so the resulting assertion is precisely algebraic
independence of \eqref{fol:eq:tower-family}. Algebraic independence persists
under algebraic extension of the coefficient field: a finite tuple
transcendental over $F$ has the same transcendence degree after an
algebraic base extension. The remaining statements follow from
$K\subset P$.
\end{proof}

\subsection{Effectiveness without deciding relations among constants}
A sequence $(c_i)_{i\ge0}$ of complex numbers is \emph{uniformly
computable} if one algorithm, on input $(i,N)$, returns a rational
complex number within $2^{-N}$ of $c_i$. Proposition~\ref{fol:prop:slow-tree}
with the explicit schedule in Corollary~\ref{fol:cor:factorial-scale} supplies
a prime Cantor host of upper box dimension zero. It has a computable
binary presentation: cylinder endpoints can be
computed, cylinders at each level are separated, their diameters have a
computable bound tending to zero, and the eventually-zero branches give
a uniformly computable dense set.

\begin{theorem}[Effective relative fusion]\label{fol:thm:effective-towers}
Suppose $P$ has the computable presentation just described and
$F=\mathbb Q(c_0,c_1,\ldots)$ for a uniformly computable sequence.
Then the set $K$ in Theorem~\ref{fol:thm:relative-towers} can be chosen with a
computable homeomorphism
\[
 \Phi:2^{\mathbb N}\longrightarrow K
\]
whose inverse is computable on $K$. Consequently
\[
 \deg_T\Phi(\omega)=\deg_T\omega
 \quad\text{and}\quad
 \deg_T T_j(\Phi(\omega))=\deg_T\omega
 \qquad(j\ge-1).
\]
In particular, $K$ has points of every Turing degree, has a dense set of
computable points, and has a computable least point. These properties
hold while \eqref{fol:eq:tower-family} is algebraically independent over the
algebraic closure of $F$. No algorithm deciding equality in $F$ is
assumed.
\end{theorem}
\begin{proof}
It suffices to enumerate the genuinely nonzero polynomials with
coefficients in the ring $R=\mathbb Z[c_0,c_1,\ldots]$. Indeed, clearing
denominators turns any relation with coefficients in $F$ into a nonzero
polynomial relation with coefficients in $R$. Enumerate finite arrays of distinct variable monomials with formal
expressions for their coefficients; combine equal variable monomials
before performing any test. The resulting coefficient values are
uniformly computable. Nonzeroness of a computable complex number is semidecidable:
a rational approximation with an error bound excluding zero is a
certificate. Dovetail these tests for every coefficient of every formal
polynomial, and emit a polynomial as soon as at least one coefficient is
certified nonzero. Every nonzero polynomial eventually appears; a
polynomial all of whose coefficients vanish never enters the list.
Thus no decision concerning algebraic relations among the $c_i$ is
needed.

At a fusion stage, candidate configurations from the uniformly
computable dense set of $P$ can be enumerated. Nonzeroness of the finitely
many required function values is again semidecidable, using rational
interval evaluation of polynomials, logarithms, and finite exponential iterates.
Dovetail the searches over all candidates. The successful configurations
form a nonempty relatively open set by the proof of
Proposition~\ref{fol:prop:analytic-fusion}; density therefore guarantees
termination. After finding one, effective continuity gives rational
neighborhoods on which every required value has a certified positive
lower bound in modulus. Search sufficiently deep original cylinders
contained in these neighborhoods, and impose the desired diameter
bound. This completes a computable fusion construction.

The successive selected cylinders compute $\Phi(\omega)$ from
$\omega$. Conversely, the positive separation between the two selected
children lets a Cauchy name for the point determine each successive
branch bit. Hence $\Phi$ and its inverse are computable. The maps $T_j$
are computable; their inverse on the relevant range is the computable
$j$-fold logarithm for $j\ge0$, or the exponential for $j=-1$. The degree identities follow. Eventually-zero codes
yield a dense computable subset. Ordering the two children from left to
right makes the all-zero code the least point of $K$.
\end{proof}

\begin{corollary}[Independence relative to named constants]
\label{fol:cor:named-constants}
The effective construction applies to
\[
 F=\mathbb Q(e,\pi,\gamma,\delta),\qquad
 \delta=\int_0^\infty e^{-t}\log(1+t)\,dt,
\]
and to any field generated by a uniformly computable sequence of named
constants. For example, it produces a computable prime-representing
constant $\alpha$ for which
\[
 \log\alpha,\ \alpha,\ e^\alpha,\ e^{e^\alpha},\ e^{e^{e^\alpha}},\ldots
\]
are algebraically independent over
$\overline{\mathbb Q(e,\pi,\gamma,\delta)}$.
\end{corollary}
\begin{proof}
The constants displayed are computable. For $\gamma$, the sequence
$H_n-\log n$ has an explicit error bounded by $1/n$. For $\delta$, use
computable quadrature on a finite interval and the tail estimate
\[
 \int_A^\infty e^{-t}\log(1+t)\,dt
 \le\int_A^\infty e^{-t}(1+t)\,dt=(A+2)e^{-A}.
\]
Apply Theorem~\ref{fol:thm:effective-towers} to the all-zero code.
\end{proof}

\begin{remark}[Scope of the relative assertion]
The last corollary makes no assertion that $e,\pi,\gamma,\delta$ are
algebraically independent among themselves, and proves neither the
irrationality of $\gamma$ nor Schanuel's conjecture. The new constants
are selected to avoid relations over the field of the old ones. The
prime-output condition is asserted for $x\in K$; it is not asserted for
its logarithm or exponential iterates with the same exponents. Nor is the least
point of $K$ claimed to equal the unrestricted least prime-representing
constant for the schedule.
\end{remark}

\begin{remark}[An unavoidable effectiveness boundary]
For an arbitrary countable field, the existence conclusion remains
valid, but the effective conclusion need not. If the field contains
all computable real numbers, any set algebraically independent over it
contains no computable real. This does not contradict
Theorem~\ref{fol:thm:effective-towers}: the computable reals do not admit a
uniformly computable enumeration. In particular, when a prescribed
field already contains the old computable minimum, a relatively
independent subtree necessarily omits that point.
\end{remark}

\paragraph{Finite nonvanishing certificates.}
The effective proof gives more than a nonzero limit statement. Once a
listed polynomial and a tuple of separated branch cylinders have been
handled, it records a positive rational $\varepsilon$ such that
$|Q(T_j(x_i))|\ge\varepsilon$ throughout that cylinder product. Given
oracles for a finite tuple of distinct branches and a nonzero polynomial
with the represented coefficients, one can reach such a certificate in
finite time. This is an algorithmic termination guarantee, not a
polynomial-time bound or a uniform closed-form estimate in degree and
height.

\paragraph{Effective construction in pseudocode.}
The following specification uses certified rational interval arithmetic;
all searches are dovetailed, so a nonterminating failed candidate cannot
block the construction.
\begin{enumerate}
\item Enumerate formal coefficient expressions in
$\mathbb Z[c_0,c_1,\ldots]$ and finite polynomial arrays. Emit an array
only after one evaluated coefficient is certified nonzero.
\item At stage $n$, obtain the first $n$ emitted arrays. In each selected
parent cylinder, enumerate pairs of eventually-zero branch points.
For each simultaneous candidate configuration, first certify pairwise
distinctness, then test the required arrays on all injective choices
of its $2^n$ point variables.
\item Retain the first configuration for which all tests certify a
nonzero value. Compute positive rational lower bounds on the finitely
many moduli and sufficiently small rational neighborhoods preserving
half these bounds.
\item Replace each selected point by a descendant prime cylinder
contained in its neighborhood, with diameter at most $2^{-n}$. Order
siblings from left to right and record all certified lower bounds.
\item Given a binary input sequence, follow its selected cylinders.
Their unique limit is the output. Given a Cauchy name for an output,
recover its bits from the separated sibling intervals.
\end{enumerate}

\subsection{Combining further analytic constraints}
\begin{proposition}[Fusion for closed nowhere-dense relations]
\label{fol:prop:closed-relation-fusion}
Let $P$ be a compact Cantor set and, for each $\nu\ge1$, let
$Z_\nu\subset P^{r_\nu}$ be closed and nowhere dense. There is a compact
Cantor set $K\subset P$ such that no ordered tuple of pairwise distinct
points of $K$ belongs to any $Z_\nu$. If $P$ has the computable
presentation used in Theorem~\ref{fol:thm:effective-towers}, and membership
in $P^{r_\nu}\setminus Z_\nu$ is uniformly semidecidable from Cauchy
names, with a successful test providing a rational open neighborhood
contained in that complement, the fusion is effective and preserves
Turing degrees as in that theorem.
\end{proposition}
\begin{proof}
At stage $n$, apply every relation $Z_\nu$, $\nu\le n$, to every
injective coordinate selection among the $2^n$ prospective children.
The coordinate projections are open maps, so inverse images of closed
nowhere-dense sets under these projections are nowhere dense. Together
with the diagonals, finitely many such inverse images leave a nonempty
open set of configurations in the prescribed parents. Choose a
configuration there and sufficiently small clopen neighborhoods of its
coordinates. The nested-set argument of
Proposition~\ref{fol:prop:analytic-fusion} applies unchanged. In the
effective case, enumerate configurations from the dense computable set
and dovetail the complement tests. One terminates, and its certified
neighborhoods allow computable selection of descendant cylinders.
\end{proof}

\begin{corollary}[Local analytic tests]\label{fol:cor:local-analytic-tests}
In Proposition~\ref{fol:prop:closed-relation-fusion}, one may include
relations of the form
\[
 Z=P^r\cap B\cap\{x:f(x)=0\},
\]
where $B$ is a closed rational box contained in an open box $U$, and
$f$ is analytic and not identically zero on $U$. Such a relation is
closed and nowhere dense. If $f$ is computable with effective local
continuity bounds, its complement has the certificates required in the
effective proposition. Countably many global tower conditions and
countably many local conditions can therefore be enforced in one and
the same subtree.
\end{corollary}
\begin{proof}
The zero set intersected with $B$ is closed because $B$ is compactly
contained in the domain of continuity. Were its intersection with
$P^r$ to have relative interior, the successive identity-theorem
argument in Lemma~\ref{fol:lem:analytic-avoidance}, now inside $U$, would
force $f$ to vanish throughout $U$. To certify that a point is outside
$Z$, either certify that one coordinate is strictly outside the
corresponding closed interval of $B$, or certify that the point lies in
$U$ and that $f$ is nonzero there. Every point in the complement admits
one of these certificates. Effective continuity gives a suitable open
neighborhood.
\end{proof}


\section{Exponential freeness of a prime-generated field}
\label{fol:sec:relative-exponential-field}

Independence of the individual towers can be supplemented by a field
property involving exponentials of arbitrary rational functions of the
prime constants.  We keep two coefficient fields distinct: $F_0$ is the
prescribed countable field over which the towers are independent, whereas
$F=\overline{\mathbb Q}(K)$ is the rational-function field generated by
the constructed prime constants and the algebraic numbers.
The latter is not the algebraic closure of that field.

\begin{theorem}[Relative Lindemann--Weierstrass property]
\label{fol:thm:relative-exponential-field}
Let $P\subset(1,\infty)$ be a nonempty compact perfect prime-coding set
with a fixed exponent schedule, and let $F_0\subset\mathbb C$ be any
countable field.  A compact perfect subset $K\subset P$ can be chosen
so that both of the following hold:
\begin{enumerate}
\item The family
$\{T_j(\alpha):\alpha\in K,\ j\ge-1\}$ is algebraically independent
over $F_0$, with the notation of Theorem~\ref{fol:thm:relative-towers}.
\item Put $F=\overline{\mathbb Q}(K)$.  For every finite family of
pairwise distinct $u_1,\ldots,u_m\in F$, the numbers
$e^{u_1},\ldots,e^{u_m}$ are linearly independent over $F$.
\end{enumerate}
Every member retains the fixed prime-output schedule of $P$.
Under the hypotheses of Theorem~\ref{fol:thm:effective-towers}, $K$ can be
chosen computably homeomorphic to Cantor space, preserving all Turing
degrees as in that theorem.
\end{theorem}

\begin{proof}
Add the following tests to the tower conditions in the fusion argument:
\begin{equation}
\label{fol:eq:field-exponential-tests}
 H(\mathbf x)=\sum_{j=1}^m A_j(\mathbf x)
                          \exp(R_j(\mathbf x)),
\end{equation}
where $A_j,R_j\in\overline{\mathbb Q}(X_1,\ldots,X_r)$, the $A_j$
are not all zero, and the $R_j$ are pairwise distinct as formal rational
functions.  This is a countable effectively enumerable family, because
algebraic numbers admit exact effective representations and equality
of rational functions over algebraic numbers is decidable.
Each expression is analytic off its poles.

We first prove that $H$ is not identically zero on any open real box
avoiding its poles.  A nonzero polynomial has no open zero set.
Consequently, the box contains a rational point $\mathbf a$ at which
at least one $A_j$ is nonzero and every difference $R_j-R_k$ is nonzero.
The exponents $R_j(\mathbf a)$ are then distinct algebraic numbers,
and the algebraic coefficients $A_j(\mathbf a)$ are not all zero.
The classical Lindemann--Weierstrass theorem gives $H(\mathbf a)\ne0$;
see, for example, the proof in B\'ezivin--Robba~\cite{BezivinRobba1989}.

Here is a closed formulation of the bad relation.  Let $D$ be the product
of the nonzero denominator polynomials in the test.  In $(1,\infty)^r$
put
\[
 Z_H=\{D=0\}\ \cup\ \{\mathbf x:D(\mathbf x)\ne0,
                                      H(\mathbf x)=0\}.
\]
This set is closed: a limit point with $D\ne0$ is handled by continuity,
and every other limit point is included in $D=0$.
Its restriction to $P^r$ is nowhere dense.  First remove the polynomial
zero set of $D$, whose restriction is nowhere dense by
Lemma~\ref{fol:lem:analytic-avoidance}.  On any remaining small open box,
$H$ is analytic and nonzero as a function, so the same iterated identity
argument applies.  For complex-valued $H$, use the real-analytic
function $|H|^2$ with the same zero set.
Also, the complement of $Z_H$ is effectively open: at a point outside
it, rational approximation eventually certifies both $D\ne0$ and
$H\ne0$. Computable derivative bounds on a sufficiently small box
then give a rational neighborhood certificate for these strict inequalities.

Apply Proposition~\ref{fol:prop:closed-relation-fusion} simultaneously to
these relations and to all the tower relations over $F_0$.
This gives (1) and the topological conclusions.  In the effective case,
the newly added tests have computable coefficients and semidecidable
complements, so the same effective fusion applies.

To deduce (2), take pairwise distinct $u_j\in F$ and coefficients
$a_j\in F$ which are not all zero.  Only finitely many members
$\boldsymbol\alpha$ of $K$ occur in their rational-function expressions;
if none occur, adjoin one unused member of $K$.
Algebraic independence of $K$ over $\mathbb Q$, hence over
$\overline{\mathbb Q}$, makes evaluation at this tuple injective.
Thus the $u_j$ have pairwise distinct formal representatives
$R_j(\mathbf X)$ and the $a_j$ have representatives $A_j(\mathbf X)$
which are not all zero.  A relation $\sum_j a_je^{u_j}=0$ would make
one of the excluded tests~\eqref{fol:eq:field-exponential-tests} vanish
at $\boldsymbol\alpha$, a contradiction.
\end{proof}

\begin{corollary}[Exponentials of rational functions]
\label{fol:cor:field-exponential-independence}
If $u_1,\ldots,u_m\in F$ are linearly independent over $\mathbb Q$,
then $e^{u_1},\ldots,e^{u_m}$ are algebraically independent over $F$.
In particular, for every $\alpha\in K$ and $d\ge1$,
\[
 e,\ e^\alpha,\ e^{\alpha^2},\ldots,e^{\alpha^d},\ e^{1/\alpha}
\]
are algebraically independent over $F$.
\end{corollary}

\begin{proof}
For distinct integer multi-indices $\nu$, the exponents
$\sum_j\nu_j u_j$ are pairwise distinct.  Expanding a nonzero polynomial
in the $e^{u_j}$ would give a prohibited $F$-linear relation among
exponentials of distinct members of $F$.
For the displayed example, the exponents
$1,\alpha,\alpha^2,\ldots,\alpha^d,1/\alpha$ are linearly independent
over $\mathbb Q$ because $\alpha$ is transcendental.
\end{proof}




\clearpage
\part{Modular return and the arithmetic of least constants}\label{part:arithmetic}
The return lemma first supplies the congruence mechanism. Local minimality then links prime intervals to Diophantine approximation. The ordinary Mills conclusions are explicitly conditional on algebraicity.

\section{Arbitrary linear perturbations of polynomial exponential clocks}
\label{sur:sec:general-clocks}

The following return principle handles every nonzero integer linear
coefficient. It supplies a universal modular-return property, stronger
than the divisibility return sufficient for several transcendence results.
The threshold on the seed is essential. Throughout this section,
$v_p(0)=\infty$.

\begin{lemma}[Unit-slope residue lemma]\label{sur:lem:clock-unit}
Let $A\ge2$ be an integer, $Q\in\mathbb Z[X]$, and $D\in\mathbb Z$.
If $q\ge1$ and $\gcd(D,q)=1$, then $Q(A^j)+Dj$ attains every residue
class modulo $q$ at arbitrarily large nonnegative integer indices $j$.
\end{lemma}
\begin{proof}
Use strong induction on $q$, the case $q=1$ being immediate. For $q>1$,
write $q=uv$, where every prime divisor of $u$ divides $A$ and
$\gcd(A,v)=1$. The sequence $A^j$ is eventually periodic modulo $q$
with a period $t<q$: take $t=\operatorname{ord}_v(A)$ if $v>1$, and
$t=1$ otherwise. The same period works for $Q(A^j)$. Set
$g=\gcd(t,q)<q$. By induction, for any prescribed residue $c$ there
are arbitrarily large $j$ for which
\[
 Q(A^j)+Dj\equiv c\pmod g.
\]
Choose such $j$ beyond the preperiod. Translation by $th$ changes the
value modulo $q$ by $Dth$. Since $\gcd(Dt,q)=g$, a nonnegative $h$
can be chosen so that $Q(A^{j+th})+D(j+th)\equiv c\pmod q$.
The resulting index is at least $j$, proving arbitrarily late occurrence.
\end{proof}

\begin{theorem}[Eventual universal return]\label{sur:thm:clock-return}
Let $a\ge2$, $d\ne0$, and $e$ be integers, and let
$P\in\mathbb Z[X]$. Define
\[
 R(n)=P(a^n)+dn+e.
\]
There is an integer $m_0$ such that for every $m\ge m_0$ and every
$q\ge1$, the congruence
\begin{equation}\label{sur:eq:clock-return}
 R(n)\equiv R(m)\pmod q
\end{equation}
has arbitrarily large solutions $n$. Its solution set is eventually
periodic and has positive natural density.
\end{theorem}
\begin{proof}
If $|d|=1$, Lemma~\ref{sur:lem:clock-unit} proves the assertion directly.
Otherwise let $S$ consist of all primes not exceeding the largest prime
divisor of $|d|$. In particular, every prime divisor of $d$ belongs to
$S$, and every prime divisor of $p-1$ belongs to $S$ when $p\in S$.
For each $\ell\in S$, put
\[
 B_\ell=\max\bigl(\{v_\ell(p-1):p\in S,\ p\nmid a\}\cup\{0\}\bigr).
\]
Choose $m_0$ so that, whenever $m\ge m_0$ and $\ell\in S$ divides $a$,
\begin{equation}\label{sur:eq:clock-seed-valuation}
 v_\ell\bigl(P(a^m)-P(0)\bigr)\ge v_\ell(d)+B_\ell.
\end{equation}
This is possible because $a^m$ divides $P(a^m)-P(0)$.
Indeed, the uniform choice
\[
 m_0=\max\{1,\lceil 2\log_2|d|\rceil\}
\]
is sufficient: for $\ell\mid a$ the left side of
\eqref{sur:eq:clock-seed-valuation} is at least $m$, while
$v_\ell(d)+B_\ell\le2\log_2|d|$.

Fix $m\ge m_0$ and factor $q=q_Sq_{\rm out}$, with $q_S$ supported
on $S$ and $q_{\rm out}$ coprime to every prime in $S$. For
$\ell\in S$ let $t_\ell=v_\ell(q_S)$, and put
\[
 K_\ell=
 \begin{cases}
 \max(t_\ell,B_\ell),&\ell\nmid a,\\
 \max(t_\ell-v_\ell(d),B_\ell,0),&\ell\mid a.
 \end{cases}
 \qquad M=\prod_{\ell\in S}\ell^{K_\ell}.
\]
For $\ell\nmid a$, impose $n_0\equiv m\pmod{\ell^{K_\ell}}$.
For $\ell\mid a$, impose
\begin{equation}\label{sur:eq:clock-local-system}
 d(n_0-m)\equiv P(a^m)-P(0)\pmod{\ell^{t_\ell}},
 \qquad n_0\equiv m\pmod{\ell^{B_\ell}}.
\end{equation}
These two congruences are compatible by
\eqref{sur:eq:clock-seed-valuation}: after removing the $\ell$-power
from $d$, its remaining factor is a unit, and the required difference
$n_0-m$ is divisible by $\ell^{B_\ell}$. Thus each local system has a
solution modulo $\ell^{K_\ell}$. The Chinese remainder theorem gives
a positive integer $n_0$ satisfying all of them.

For every sufficiently large $j$, set $n=n_0+Mj$. We claim that
\begin{equation}\label{sur:eq:clock-smooth-return}
 R(n)\equiv R(m)\pmod{q_S}.
\end{equation}
At a prime $\ell\mid a$, the exponential part satisfies
$P(a^n)\equiv P(0)\pmod{\ell^{t_\ell}}$ for large $n$, and
\eqref{sur:eq:clock-local-system} gives the claim. At a prime $p\nmid a$
with $t_p>0$, the order of $a$ modulo $p^{t_p}$ divides
$(p-1)p^{t_p-1}$. Every prime factor of this order lies in $S$.
For a factor $\ell\mid a$, its valuation in that order is at most
$B_\ell$, so the required divisibility of $n-m$ follows from the
second congruence in \eqref{sur:eq:clock-local-system}. For a factor
$\ell\nmid a$, it follows from $K_\ell\ge\max(t_\ell,B_\ell)$.
Consequently $a^n\equiv a^m\pmod{p^{t_p}}$. Also
$n\equiv m\pmod{p^{t_p}}$, proving
\eqref{sur:eq:clock-smooth-return}.

On this arithmetic progression we have
\[
 R(n_0+Mj)=Q(A^j)+Dj+(dn_0+e),\qquad
 Q(X)=P(a^{n_0}X),\quad A=a^M,\quad D=dM.
\]
All prime factors of $D$ lie in $S$, so $\gcd(D,q_{\rm out})=1$.
Lemma~\ref{sur:lem:clock-unit} supplies arbitrarily large $j$ for which
the desired congruence holds modulo $q_{\rm out}$ as well. Together
with \eqref{sur:eq:clock-smooth-return}, this proves
\eqref{sur:eq:clock-return}. Finally, $R$ is eventually periodic modulo
$q$, since both $a^n$ and $dn$ are. An eventually periodic solution
set with arbitrarily late members has positive natural density.
\end{proof}

\begin{remark}[The seed threshold cannot simply be omitted]
For $R(n)=2^n+6n$ we have $R(1)=8$. For $n\ge3$, equality to $8$
modulo $8$ forces $n\equiv0\pmod4$, whereas equality modulo $3$
forces $n$ to be odd. Hence no sufficiently late index returns to the
first value modulo $24$.
\end{remark}

\begin{corollary}[Transcendental minima for tails]\label{sur:cor:clock-minimum}
Suppose additionally that $P$ has degree $r\ge1$ and positive leading
coefficient, and that $a^r>40/19$. For every sufficiently large integer
$N$, the minimum
\[
 \xi_N=\min\{x>1:\lfloor x^{R(n)}\rfloor\text{ is prime for every }n\ge N\}
\]
exists and is transcendental.
\end{corollary}
\begin{proof}
Choose $N\ge m_0$ so large that $R(n)$ is positive and
$R(n+1)/R(n)\ge2$ for $n\ge N$. Write
$C_k=R(N+k-1)$, $c_1=C_1$, and $c_{k+1}=C_{k+1}/C_k$.
Then $c_1\ge1$, $c_{k+1}\ge2$, $C_k\in\mathbb N$, and
\[
 \lim_{k\to\infty}c_{k+1}=a^r>40/19.
\]
These are conditions (B1)--(B4) of Saito's Type-B theorem,
\cite[Theorem~2.3]{SaitoVariants}. Choose a sufficiently small
$\varepsilon>0$. Its index set
$I=\{k:c_{k+1}\ge40/19+\varepsilon\}$ is cofinite. For every $m,L\ge1$,
Theorem~\ref{sur:thm:clock-return}, with modulus $LC_m$, gives arbitrarily
large $k$ such that
\[
 C_k\equiv C_m\pmod{LC_m}.
\]
One may therefore take $k>m$ with $k\in I$, proving (B6). Taking
$L=1$ also proves (B5), namely $C_m\mid C_k$ at a later index in $I$.
The cited theorem gives existence and transcendence of $\xi_N$.
\end{proof}

\begin{remark}[Scope of the corollary]
This includes every sufficiently late tail of $ba^n+dn+e$ with
$a\ge3$, $b\ge1$, $d\ne0$, and arbitrary integer $e$. It also includes
$4^n-10\cdot2^n+37n+5$ and
$27^n+2\cdot9^n-8\cdot3^n-1000n+7$.
Deleting a finite prefix changes the defining minimum; no identification
with the original full-prefix minimum is claimed. The case $d=0$ is
outside the return theorem, so ordinary Mills' constant is not resolved.
\end{remark}

\begin{corollary}[Composite windows for Pisot and Salem powers]
\label{sur:cor:clock-windows}
Let $P$ have positive leading coefficient and degree at least one.
For every Pisot or Salem number $\alpha$ and every fixed integer
$H\ge0$, there is an arithmetic progression of indices on which,
after finitely many initial indices are removed, all the integers
\[
 \lfloor\alpha^{R(n)}\rfloor+h,\qquad |h|\le H,
\]
are composite. No condition $a^{\deg P}>40/19$ is needed here.
\end{corollary}
\begin{proof}
Let $D$ be the degree of $\alpha$ and
$T(t)=\operatorname{Tr}_{\mathbb Q(\alpha)/\mathbb Q}(\alpha^t)$.
This is an integer linear recurrence with nonzero trailing coefficient.
There is an index $s_0$ such that $T$ is periodic modulo every prime
$p$ from index $s_0$ onward. Indeed, the companion matrix is invertible
modulo primes not dividing the trailing coefficient, giving pure
periodicity. Only finitely many remaining primes exist, and one may
maximize their finite preperiods. Equivalently, use
\cite[Lemma~2.1(b)]{SaitoIntervals}.

Choose $m\ge m_0$ such that $R(m)\ge s_0$ and
$T(R(m))>H+D+1$. For each integer $j$ with $|j|\le H+D$, choose a
prime divisor $p_j$ of $T(R(m))+j$. Let $L$ be a common multiple of the
periods of $T$ modulo these finitely many primes. By
Theorem~\ref{sur:thm:clock-return}, an eventually periodic positive-density
set of indices satisfies $R(n)\equiv R(m)\pmod L$. For all sufficiently
large such indices,
\[
 T(R(n))+j\equiv T(R(m))+j\equiv0\pmod{p_j}.
\]
Since $R(n)\to\infty$ and $T(t)\to\infty$, these integers eventually
exceed their fixed divisors $p_j$, and hence are composite.

Every nonprincipal conjugate of a Pisot or Salem number has modulus
at most one. Consequently
\[
 |T(t)-\alpha^t|\le D-1,
 \qquad |\lfloor\alpha^t\rfloor-T(t)|\le D.
\]
The required floor window is therefore contained in the trace window
already treated. Finally, the eventually periodic return set contains
all sufficiently late members of some arithmetic progression.
\end{proof}

\begin{proposition}[Intrinsic non-reversibility]\label{sur:prop:clock-nonrev}
Under the hypotheses of Corollary~\ref{sur:cor:clock-windows}, $R$ is an
inhomogeneous integer linear recurrence sequence, but no integer
inhomogeneous recurrence representation for it has trailing coefficient
$1$ or $-1$.
\end{proposition}
\begin{proof}
An integer polynomial in $a^n$ plus a linear polynomial in $n$ is a
linear recurrence sequence. Suppose a monic integer polynomial $Q$
satisfies $Q(E)R=\text{constant}$, where $E$ is the shift operator.
For every nonzero coefficient of $X^j$ in $P$ with $j\ge1$, linear
independence of distinct exponential sequences gives $Q(a^j)=0$.
The coefficient of $n$ also gives $dQ(1)=0$. In particular, some
$a^j$ divides $Q(0)$, so $a\mid Q(0)$. The trailing coefficient of the
recurrence is $-Q(0)$ and cannot equal $\pm1$.
\end{proof}

The last two results enlarge the family answering Saito's
Problem~1.7 in \cite{SaitoIntervals}; they are not an additional
numbered-problem resolution. The window radius $H$ is fixed, and its
arithmetic progression can depend on $H$ and $\alpha$.


\section{Local Type B and factorial perturbations}
\label{mil:sec:factorial-local}

Put $\theta=21/40$. We use the short-interval prime theorem and the prime-chain extension lemma in Saito, arXiv:2508.16068v3, Theorem~1.6 and Lemma~4.2. We also use his arithmetic approximation lemma, Lemma~6.1: sufficiently strong exponential approximation of integer powers of an algebraic number by integers forces a Pisot power, with degree bounded by one plus the reciprocal approximation exponent. The arguments below specify the new localization step and its consequences; they do not replace those analytic and Diophantine inputs.

Write
\[
 W(C)=\{A>1:\lfloor A^{C_k}\rfloor\text{ is prime for every }k\geq1\},
 \qquad c_{k+1}=C_{k+1}/C_k.
\]
Call $A\in W(C)$ a \emph{local left minimum} if there is $\eta>0$ such that $(A-\eta,A)\cap W(C)=\varnothing$.

\begin{theorem}[Local Type B]
\label{mil:thm:local-type-b}
Let $(C_k)$ be positive integers satisfying $c_{k+1}\geq2$ and
\[
 r=\limsup_{k\to\infty}c_{k+1}>40/19.
\]
Choose $\varepsilon>0$ so that
\[
 I=\{k:c_{k+1}\geq40/19+\varepsilon\}
\]
is infinite. Suppose that
\begin{equation}
 \text{for each }m\text{ there is }k>m,\ k\in I,
 \text{ with }C_m\mid C_k.\label{mil:eq:local-b5}
\end{equation}
If a local left minimum $A\in W(C)$ is algebraic, then $r<\infty$ and there is an integer $g\geq1$ such that $\beta=A^g$ is a Pisot number of degree $\ell$ satisfying
\[
 3\leq\ell\leq1+\left(\frac{19r}{40}-1\right)^{-1},
\]
and $g\mid C_k$ for all sufficiently large $k\in I$. In particular, if $r>60/19$, every local left minimum is transcendental.

If, in addition,
\begin{equation}
 \text{for each }m,L\geq1\text{ there is }k>m,\ k\in I,
 \text{ with }C_k\equiv C_m\pmod{LC_m},\label{mil:eq:local-b6}
\end{equation}
then every local left minimum is transcendental for all $r>40/19$.
\end{theorem}

\begin{proof}
Set $p_k=\lfloor A^{C_k}\rfloor$. First, (B5) excludes $A^{C_m}\in\mathbb N$ for every $m$: that integer would be the prime $p_m$, whereas a later proper multiple $C_k$ would make $A^{C_k}=p_m^{C_k/C_m}$ composite.

We show explicitly why the global-minimum argument localizes. Consider the exceptional set
\[
 J=\left\{k:\|A^{C_k}\|\leq
 \frac1{c_{k+1}p_k^{c_{k+1}-1}}\right\}.
\]
If $A$ is algebraic, $J$ is finite. Indeed, an infinite $J$ gives
\[
 \|A^{C_k}\|\leq\frac1{2p_k}\leq A^{-C_k}
 \quad(k\in J\text{ sufficiently large}).
\]
The arithmetic approximation lemma gives a Pisot power $A^g$ of degree at most two, with $g\mid C_k$ for late $k\in J$. Degree one gives a composite integer power at late indices. In degree two the equality case of that lemma makes $A^g$ a quadratic Pisot unit, so
\[
 \|A^{C_k}\|=A^{-C_k}>\frac1{2p_k}
\]
for sufficiently large $k\in J$, again a contradiction. This is the argument behind Saito's Lemma~6.2 and uses no minimality.

Choose $k_0$ beyond $J$. Saito's elementary Lemma~5.1 then gives the nested inequalities
\begin{equation}
 p_k^{c_{k+1}}\leq p_{k+1}<(p_k+1)^{c_{k+1}}-1
 \qquad(k\geq k_0).\label{mil:eq:local-nesting}
\end{equation}
Suppose infinitely many $m\in I$ violate
\begin{equation}
 p_{m+1}\leq p_m^{c_{m+1}}+p_m^{\theta c_{m+1}}.
 \label{mil:eq:local-short-child}
\end{equation}
At each sufficiently large such $m$, the short-interval theorem and prime-chain extension lemma give a replacement prime sequence $(q_k)$ with $q_k=p_k$ for $k\leq m$,
\[
 p_m^{c_{m+1}}\leq q_{m+1}
 \leq p_m^{c_{m+1}}+p_m^{\theta c_{m+1}}<p_{m+1},
\]
and nested prime intervals thereafter. Because $m\in I$, the replacement at level $m+1$ also lies strictly below $(p_m+1)^{c_{m+1}}-1$ when $m$ is large. The nested intervals define a limit $\zeta_m$ whose prime outputs agree with $p_k$ for $k_0\leq k\leq m$ and with $q_k$ afterwards. Moreover,
\[
 p_m^{1/C_m}\leq\zeta_m<A,
 \qquad
 \zeta_m\geq 2^{-1/C_m}A\longrightarrow A.
\]
For each of the finitely many indices $k<k_0$, the number $A^{C_k}$ lies strictly inside the interval $[p_k,p_k+1)$, since integer powers have been excluded. Continuity therefore ensures the same outputs at $\zeta_m$ for all sufficiently large $m$. Thus $\zeta_m\in W(C)$, and eventually $A-\eta<\zeta_m<A$, contradicting local left minimality. This proves \eqref{mil:eq:local-short-child} for all late $m\in I$.

Taking $c_{k+1}$-th roots and using concavity yields
\begin{equation}
 0< A^{C_k}-p_k\leq
 \frac{2}{c_{k+1}p_k^{(1-\theta)c_{k+1}-1}}
 \qquad(k\in I\text{ sufficiently large}).
 \label{mil:eq:local-fractional-bound}
\end{equation}
If $r<\infty$, the arithmetic approximation lemma applied to this estimate gives a Pisot power $\beta=A^g$ with
\[
 \ell\leq1+((1-\theta)r-1)^{-1},
 \qquad g\mid C_k\quad(k\in I\text{ sufficiently large}).
\]
Here replacing $p_k$ by $A^{C_k}$ affects only a bounded factor, since the exponents $(1-\theta)c_{k+1}-1$ are bounded. If $r=\infty$, fix $T>1$ and restrict to infinitely many indices for which this exponent is at least $T$. Then \eqref{mil:eq:local-fractional-bound} is at most $p_k^{-T}\leq 2^T A^{-TC_k}$, so the approximation lemma applies with the fixed constant $K=2^T$. It forces degree at most $1+1/T<2$, hence degree one, which is impossible. Hence algebraicity requires finite $r$ and degree at least two.

The small fractional part in \eqref{mil:eq:local-fractional-bound} and the Pisot property imply
\begin{equation}
 p_k=\operatorname{Tr}(\beta^{C_k/g})
 \qquad(k\in I\text{ sufficiently large}).
 \label{mil:eq:local-trace}
\end{equation}
For completeness, the other conjugate powers have sum tending to zero; since the fractional part tends to zero, the trace is the floor rather than the next integer.

If $\ell=2$, write the other conjugate as $\gamma$. Equation~\eqref{mil:eq:local-trace} forces $\gamma<0$ and $C_k/g$ odd at all late $k\in I$. Choose such an $m$ and then $k>m$ as in (B5). The quotient $q=C_k/C_m$ is an odd integer at least three. With $u=\beta^{C_m/g}$ and $v=\gamma^{C_m/g}<0$,
\[
 p_k=u^q+v^q=(u+v)
 \sum_{j=0}^{q-1}(-1)^j u^{q-1-j}v^j.
\]
The second factor is a rational algebraic integer, hence an integer, and all its terms are positive. It exceeds one, while $u+v=p_m$ is prime. This contradiction proves $\ell\geq3$ and the stated degree bound.

Finally assume (B6). Choose a sufficiently large $m\in I$, put $K=\mathbb Q(\beta)$, and let $\alpha=\beta^{C_m/g}$. The trace sequence $T(n)=\operatorname{Tr}_{K/\mathbb Q}(\alpha^n)$ is an integer linear recurrence with characteristic polynomial $N_{K/\mathbb Q}(X-\alpha)$, whose constant coefficient, up to sign, is $N_{K/\mathbb Q}(\alpha)$. Since $\alpha$ is Pisot,
\[
 0<|N_{K/\mathbb Q}(\alpha)|<T(1)=p_m
\]
for large $m$. Its recurrence matrix is therefore invertible modulo the prime $p_m$, and $T(n)$ is purely periodic modulo $p_m$, with some period $L$. By (B6), choose $k>m$ in $I$ with $C_k/C_m\equiv1\pmod L$. Equation~\eqref{mil:eq:local-trace} gives
\[
 p_k=T(C_k/C_m)\equiv T(1)=p_m\equiv0\pmod{p_m}.
\]
As $p_k>p_m$, this is impossible.
\end{proof}

\begin{corollary}[Least completions of prime prefixes]
Under the hypotheses of Theorem~\ref{mil:thm:local-type-b}, fix a finite sequence of prime outputs that has at least one completion in $W(C)$. Its least completion exists and is a local left minimum. Whenever either transcendence conclusion of that theorem applies, all these least completions are transcendental. They form a countable dense subset of $W(C)$.
\end{corollary}

\begin{proof}
The finite constraints define an interval $[a,b)$ containing a point of $W(C)$. The infimum of $W(C)\cap[a,b)$ belongs to that set: take a sequence tending down to the infimum and use right continuity of each floor map. Call it $A$. Condition (B5) excludes $A=a$, because $a$ is a lower endpoint $p_j^{1/C_j}$ for some prescribed output. Hence $(a,A)\cap W(C)=\varnothing$ and $A$ is a local left minimum. Countability follows from the countable collection of finite prime prefixes. For any $B\in W(C)$, the minima of its successive prefixes converge to $B$, since the width of its level-$k$ prime interval is at most $1/C_k\to0$.
\end{proof}

\subsection{The original perturbed cubic families}

The following examples were among the initial families considered in this
project. They are recorded as consequences of the return criterion, rather
than counted again as new first-time transcendence results.

\begin{corollary}[Three full cubic perturbations]\label{mil:cor:original-cubic-perturbations}
For each of the exponent sequences
\[
 C_k=3^k-k,\qquad C_k=3^k-2k,\qquad C_k=3^k-2k+1
 \qquad(k\ge1),
\]
every local left minimum of $W(C)$ is transcendental. In particular,
all three full minimum constants are transcendental.
\end{corollary}
\begin{proof}
Each sequence is positive, has consecutive ratios at least two,
and has limiting ratio three. Indeed, for $C_k=3^k-dk+e$,
\[
 C_{k+1}-2C_k=3^k+dk-d-e>0
\]
for $(d,e)=(1,0),(2,0),(2,1)$ and every $k\ge1$.
The existence theorem therefore gives a nonempty prime-representing set.
It remains to verify \eqref{mil:eq:local-b6} from every seed $m$.

For $d=1$, Lemma~\ref{sur:lem:clock-unit} says that $3^n-n$ attains
every residue modulo any positive integer $q$ at arbitrarily large
indices. Apply this with $q=LC_m$ and the target residue $C_m$.

For $d=2$, write $q=2^u v$, where $u\ge0$ and $v$ is odd, and
restrict to indices $n=m+2^u j$. At all such indices,
\[
 3^n-2n+e\equiv3^m-2m+e\pmod{2^u}.
\]
For $u\ge1$ the exponential congruence follows because the order of
$3$ modulo $2^u$ divides $2^u$; for $u=0$ there is no constraint.
On this progression the value equals
\[
 3^m\bigl(3^{2^u}\bigr)^j-2^{u+1}j-2m+e.
\]
The coefficient $-2^{u+1}$ is a unit modulo $v$, so
Lemma~\ref{sur:lem:clock-unit} supplies arbitrarily large $j$ giving the
required residue modulo $v$ as well. The Chinese remainder theorem
then gives $C_n\equiv C_m\pmod q$. Taking $q=LC_m$ proves
\eqref{mil:eq:local-b6}. Since the index set $I$ is cofinite,
Theorem~\ref{mil:thm:local-type-b} applies to all three families.
\end{proof}

\subsection{Factorials plus polynomial drift}

An integer polynomial is called \emph{intersective} here if it has a root modulo every positive integer.

\begin{theorem}[Factorial perturbations]
\label{mil:thm:factorial-drift}
Let $P\in\mathbb Z[X]$ be intersective. For all sufficiently large integers $N$, put
\[
 W_N=\{A>1:\lfloor A^{n!+P(n)}\rfloor
 \text{ is prime for every }n\geq N\}.
\]
Then $W_N$ is nonempty, and every local left minimum of $W_N$, including its least element and the least completion of any feasible finite prime prefix, is transcendental.

In particular, without deleting any initial constraint,
\[
 \boxed{\xi(k!+k)\text{ is transcendental}.}
\]
The same conclusion holds for $\xi(k!+k^2)$.
\end{theorem}

\begin{proof}
For $C_n=n!+P(n)$, factorial growth gives $C_n>0$ eventually and
\[
 C_{n+1}/C_n\sim n+1\longrightarrow\infty.
\]
Choose $N$ so that all these ratios are at least two. This ensures nonemptiness by Saito's existence theorem. Given $m\geq N$, choose a root $t$ of $P$ modulo $q=C_m$ and take $n>m$ sufficiently large with $n\equiv t\pmod q$. Then $q\mid n!$ and $q\mid P(n)$, so $C_m\mid C_n$. Condition (B5) holds, and the unbounded-ratio conclusion of Theorem~\ref{mil:thm:local-type-b} proves the claim.

For $P(X)=X$, the polynomial has root zero modulo every modulus, and
\[
 C_{k+1}-2C_k=(k-1)(k!-1)\geq0
 \qquad(k\geq1).
\]
Thus $N=1$ is allowed. For $P(X)=X^2$,
\[
 C_{k+1}-2C_k=(k-1)k!-k^2+2k+1>0.
\]
The assertion is immediate for $k=1,2$; for $k\geq3$ use $k!\geq k(k-1)$, which makes the right side at least $k(k-1)(k-2)+k+1>0$.
\end{proof}

The least-constant assertions in this theorem are direct applications of Saito's existing Type~B theorem. The additional local-minimum conclusion uses the localization above. No new claim about the usual Mills constant is implied.

\begin{proposition}[Exact factorial criterion for (B5)]
Fix any integer polynomial $P$ and choose $N$ so that $C_n=n!+P(n)>0$ and $C_{n+1}/C_n\geq2$ for $n\geq N$. Condition (B5) for this tail is equivalent to the statement that $P$ has a root modulo $C_m$ for every $m\geq N$.
\end{proposition}

\begin{proof}
The forward construction in the proof of Theorem~\ref{mil:thm:factorial-drift} proves sufficiency. Conversely, iterate (B5) starting at any fixed $m$ to obtain an unbounded chain of indices $m<n_1<n_2<\cdots$ with $C_m\mid C_{n_j}$. At a level $n_j\geq C_m$, the factorial is divisible by $C_m$, so $P(n_j)\equiv0\pmod{C_m}$. This proves necessity.
\end{proof}

This criterion separates (B5) from (B6). For $C_k=k!+k^2$, (B5) holds, but (B6) fails at $m=1$, $L=4$. Indeed, $C_1=2$, whereas the residues of $C_2,\ldots,C_7$ modulo eight are $6,7,0,1,4,1$. For $k\geq8$, one has $C_k\equiv k^2\pmod8\in\{0,1,4\}$. No later exponent is congruent to $C_1$ modulo $LC_1=8$.

Every polynomial having an integer root is intersective, but the class is strictly larger. For example,
\[
 P(X)=(X^2-13)(X^2-17)(X^2-221)
\]
has no integer root and is intersective. For an odd prime not dividing $13\cdot17$, either $13$ or $17$ is a quadratic residue, or both are nonresidues and $221$ is a residue. A nonzero root of the corresponding quadratic lifts to every prime power. At $13$ use $17\equiv4$; at $17$ use $13\equiv8^2$; at $2$ use that $17\equiv1\pmod8$ is a square in $\mathbb Z_2$. The Chinese remainder theorem combines these local roots. Thus Theorem~\ref{mil:thm:factorial-drift} includes genuinely non-linear drifts with no rational root.

More generally, let $H_n\in\mathbb Z$ satisfy: for every $q\geq1$, $q\mid H_n$ for every sufficiently large $n$. If $C_n=H_n+P(n)$ is eventually positive with successive ratios at least two and upper limiting ratio greater than $60/19$, the same proof applies to every intersective $P$. For instance, $H_n=(n!)^s$, with any positive integer $s$, is admissible after a sufficiently late start.

The criterion does not cover every polynomial drift. For $P=1$, (B5) would yield an unbounded divisibility chain from a modulus $C_m>1$, forcing $C_m\mid1$ at a sufficiently late level. Hence (B5) fails. This is an obstruction to the present method, not a disproof of transcendence for $\xi(n!+1)$. Likewise, arbitrary early constraints added to a good tail can create prime-radical boundary candidates $A=p^{1/E}$; they require a separate argument unless some tail exponent is a proper multiple of $E$.

\subsection{A modest quotient restriction at bounded integer ratios}

The following statement does not imply algebraic independence. It records exactly what the elementary integer comparison still yields when the ratios remain bounded.

\begin{proposition}[Rational quotient height]
Suppose all $c_{k+1}$ are integers at least two and $r=\limsup c_{k+1}<\infty$, with $r\geq3$. Set $\kappa=19r/40-1$. If $1<A<B$ are local left minima of $W(C)$ and $B/A=u/v$ is a reduced positive rational, then
\[
 u\geq A^{\kappa}.
\]
For the cubic clock, this gives $u\geq A^{17/40}$.
\end{proposition}

\begin{proof}
For integer ratios, the nesting inequalities \eqref{mil:eq:local-nesting} hold for every element of $W(C)$, without assuming algebraicity. The lower inequality follows because $p_k^{c_{k+1}}$ is an integer. The only possible obstruction to the strict upper inequality would be the prime $(p_k+1)^{c_{k+1}}-1$, which factors nontrivially. Consequently the downward replacement argument establishing \eqref{mil:eq:local-short-child} applies directly to every local left minimum. It gives \eqref{mil:eq:local-fractional-bound} for both $A$ and $B$ at all late indices with $c_{k+1}=r$. There are infinitely many such indices, since the ratios are integers with finite upper limit $r$.

If $u<A^{\kappa}$, choose $0<\kappa'<\kappa$ with $u<A^{\kappa'}$. At those indices, for sufficiently large $k$,
\[
 0<A^{C_k}-p_k\leq A^{-\kappa'C_k},\qquad
 0<B^{C_k}-q_k\leq B^{-\kappa'C_k}.
\]
Because $v^{C_k}B^{C_k}=u^{C_k}A^{C_k}$, the integer
\[
 v^{C_k}q_k-u^{C_k}p_k
\]
has absolute value at most
\[
 (vB^{-\kappa'})^{C_k}+(uA^{-\kappa'})^{C_k}\longrightarrow0.
\]
Indeed, $vB^{-\kappa'}<uA^{-\kappa'}<1$. Yet this integer is nonzero: for large $k$, $p_k,q_k$ are distinct primes larger than $u,v$, so equality would force $p_k\mid q_k$. This contradiction proves the bound.
\end{proof}

Whenever $A<2^{40/17}$, the bound $A^{17/40}<2$ excludes no positive integer numerator greater than one. Thus this proposition is a quantitative restriction for larger branch minima, not a resolution of the principal quotient or transcendence problem.

\section{Sharper Type A bounds and late-tail Type C transcendence}

In this section $\xi(C_k)$ denotes the least $A>1$ for which
$\lfloor A^{C_k}\rfloor$ is prime for every indicated index. We use
Theorems~2.1 and~2.3 and Lemma~9.1 of Saito,
\emph{Transcendency of variants of Mills' constant},
arXiv:2508.16068v3. These results extend the Pisot reduction in
\emph{Mills' constant is irrational}, arXiv:2404.19461v2.
All theorem-number comparisons in this section refer specifically
to those arXiv versions. The version of record of the first paper,
published in the \emph{Ramanujan Journal} on 10 August 2026, has
different numbering; its publicly accessible preview does not
establish whether the Type A bound was changed. No priority claim
relative to that version is made.
Let $\kappa=1.324717957\ldots$ be the smallest Pisot number, namely
the positive real root of $X^3-X-1$.

\subsection{An exact floor inequality improves Type A}

\begin{theorem}\label{mil:thm:ac-floor}
Let $b>40/19$, $h\ge0$, and $\alpha b-h\ge1$. Put
$C_k=\lfloor\alpha b^k-h\rfloor$ for $k\ge1$. Then $\xi(C_k)$
exists and is irrational. More precisely, it is either transcendental,
or some $\xi(C_k)^{C_m}$ is a Pisot number of degree
\[
 \ell\le 1+\left(\frac{19b}{40}-1\right)^{-1}.
\]
In particular, $\xi(\lfloor b^k\rfloor)$ is irrational for every
$b>40/19$.
\end{theorem}
\begin{proof}
Set $x_k=\alpha b^k-h\ge1$. Since $b>2$ and $h\ge0$,
\[
 C_{k+1}=\lfloor b x_k+(b-1)h\rfloor
 \ge\lfloor2x_k\rfloor\ge2\lfloor x_k\rfloor=2C_k.
\]
Also $C_1\ge1$, the exponents are integers, and
$C_{k+1}/C_k\to b>40/19$. All hypotheses of Saito's
Theorem~2.1 are therefore satisfied.
\end{proof}

This improves the sufficient lower bound $1+\sqrt2$ displayed in
Saito's arXiv v3 Theorem~1.9(A). With $h=0$, it also reduces the amplitude
condition in his Corollary~2.2 to $\alpha\ge1/b$: the additional
condition $\alpha\ge1/[b(b-2)]$ is unnecessary. This is an elementary
sharpening of an existing theorem, not a new irrationality mechanism.
Neither the endpoint $b=40/19$ nor the remaining range
$2<b\le40/19$ is settled by this argument.

For completeness, if an integer sequence satisfies
$C_k=\alpha b^k+o(b^k)$ with $\alpha>0$ and $b>40/19$, then
$C_{k+1}/C_k\to b$, so every sufficiently late tail has an irrational
minimum by the same theorem. This statement does not assert anything
about an earlier prefix whose consecutive ratios fail to be at least two.

\subsection{A finite asymptotic gcd controls every sufficiently late tail}

For a positive integer sequence write
\[
 \operatorname{agcd}(C_n)=
 \lim_{m\to\infty}\gcd(C_m,C_{m+1},\ldots).
\]

\begin{theorem}\label{mil:thm:ac-tail}
Let $(C_n)$ be positive integers. Suppose that, after deletion of a
finite prefix, the following hold:
\begin{enumerate}
\item $C_{n+1}/C_n\ge c$ for a fixed $c>40/19$;
\item for every $m$ there exists $k>m$ with $C_m\mid C_k$;
\item $D=\operatorname{agcd}(C_n)<\infty$.
\end{enumerate}
Then the minimum $\xi_K=\xi((C_n)_{n\ge K})$ is transcendental for
every sufficiently large $K$.

More explicitly, choose $x_0\ge1$ such that every $x\ge x_0$ has a
prime in $[x,x+cx^{1-1/c}]$, and put $M=2x_0^{1/c}+1$.
Once the hypotheses hold on the tail, the sufficient condition
\[
 C_K>D\frac{\log M}{\log\kappa}
\]
implies that $\xi_K$ is transcendental.
\end{theorem}
\begin{proof}
Such an $x_0$ exists by the Baker--Harman--Pintz short-interval theorem,
since $1-1/c>21/40$. Bertrand's postulate supplies a fixed prime
$p$ with $x_0^{1/c}<p\le2x_0^{1/c}$. This prime is independent of
$K$. Saito's Lemma~9.1 applies to every tail: its first consecutive
ratio is at least $c$, and consequently
$p>x_0^{1/(C_{K+1}/C_K)}$. It gives a constant $A_K$ representing
primes on that tail with first output $p$. Hence
\begin{equation}\label{mil:eq:ac-uniform-small}
 \xi_K\le A_K<(p+1)^{1/C_K}\le M^{1/C_K}.
\end{equation}

Choose $\varepsilon$ with $0<\varepsilon<c-40/19$. The index set
$I$ in Saito's Theorem~2.3 is then the entire set of tail indices.
The recurring-divisibility hypothesis is exactly its condition
(B5). If $\xi_K$ were algebraic, that theorem would give a positive
integer $g$ such that $\beta=\xi_K^g$ is a Pisot number of degree
at least three and $g\mid C_n$ for all sufficiently large $n$.
Thus $g\mid D$. By \eqref{mil:eq:ac-uniform-small},
\[
 \kappa\le\beta=\xi_K^g\le\xi_K^D
 \le M^{D/C_K}<\kappa,
\]
a contradiction. Finally $C_K\to\infty$ geometrically, so the
displayed sufficient condition holds for all sufficiently large $K$.
\end{proof}

The theorem gives a late-tail conclusion under a lower ratio
$c>40/19$, whereas Saito's Type C theorem assumes all ratios are
at least three. It does not prove transcendence of the minimum
with a fixed earlier prefix still attached.

\begin{corollary}\label{mil:cor:ac-shifted}
Let $a\ge3$, $b\ge1$ and $e\ne0$ be integers. For every sufficiently
large $K$, the minimum associated with
$(ba^n+e)_{n\ge K}$ is transcendental.
\end{corollary}
\begin{proof}
Write $C_n=ba^n+e$. Positivity and a lower consecutive-ratio bound
$c>40/19$ hold eventually, since $C_{n+1}/C_n\to a$.
Increase the starting index so that
\[
 v_p(ba^n)>v_p(e)\qquad(p\mid a)
\]
for every subsequent $n$. It follows that $v_p(C_n)=v_p(e)$
for such primes.

Fix an index $m$ after this point and split $C_m=q_aq_u$, where
all prime factors of $q_a$ divide $a$ and $\gcd(q_u,a)=1$.
The stabilized valuations show that $q_a$ divides every subsequent
$C_n$. If $q_u>1$, choose $T=\operatorname{ord}_{q_u}(a)$;
otherwise take $T=1$. Then
\[
 C_{m+T}=ba^{m+T}+e\equiv ba^m+e=C_m\equiv0\pmod{q_u}.
\]
Thus $C_m\mid C_{m+T}$, establishing recurring divisibility.
Finally,
\[
 C_{n+1}-aC_n=(1-a)e,
\]
so $\operatorname{agcd}(C_n)$ divides $|(a-1)e|$ and is finite.
Theorem~\ref{mil:thm:ac-tail} applies.
\end{proof}

The restriction $e\ne0$ matters: for a pure geometric clock the
asymptotic gcd is infinite. This corollary covers constant
perturbations, rather than relying on a nonzero linear perturbation
of the index.

\begin{corollary}\label{mil:cor:ac-large-base}
Let $a\ge4$, $b\ge1$ and $e\ne0$ be integers. If $K$ satisfies
\[
 ba^K+e\ge1,\qquad (a-2)ba^K\ge e,\qquad
 v_p(ba^K)>v_p(e)\quad(p\mid a),
\]
then $\xi((ba^n+e)_{n\ge K})$ is transcendental.
In particular, if
\[
 \gcd(a,e)=1,\qquad ba+e\ge1,\qquad e\le ba(a-2),
\]
then the full minimum $\xi((ba^k+e)_{k\ge1})$ is transcendental.
Consequently $\xi(ra^k-1)$ and $\xi(ra^k+1)$ are transcendental
for every integer $a\ge4$ and $r\ge1$.
\end{corollary}
\begin{proof}
The displayed conditions ensure positivity, consecutive ratios
at least two, and recurring divisibility by the preceding proof.
If the minimum were algebraic, Theorem~2.3 of Saito would give
\[
 3\le\ell\le1+\left(\frac{19a}{40}-1\right)^{-1}
 \le1+\frac{10}{9}<3,
\]
a contradiction. For the full-clock assertion, $\gcd(a,e)=1$
implies $\gcd(a,C_m)=1$ for every $m$, so periodicity of powers
of $a$ modulo $C_m$ gives recurring divisibility directly.
\end{proof}

\begin{corollary}\label{mil:cor:ac-two-bases}
Let $2\le a<b$ be coprime integers, and put $C_n=a^n+b^n$.
Every sufficiently late tail minimum is transcendental. If
$b\ge4$, then the full minimum $\xi((a^n+b^n)_{n\ge1})$ is
transcendental.
\end{corollary}
\begin{proof}
The consecutive ratio is a weighted average of $a$ and $b$, so
it is greater than two and tends to $b$. For every $m$,
$\gcd(C_m,ab)=1$. The recurrence
\[
 C_{n+2}=(a+b)C_{n+1}-abC_n
\]
is therefore invertible modulo $C_m$. Its finite state sequence
is purely periodic modulo $C_m$, giving $C_m\mid C_{m+T}$ for
some $T\ge1$.

Let $D_N=\gcd(C_N,C_{N+1},\ldots)$. Since $D_N$ is coprime to
$ab$, the same recurrence propagates divisibility backwards and
gives
\[
 D_N\mid\gcd(C_1,C_2)
 =\gcd(a+b,a^2+b^2)\mid2.
\]
The last divisibility uses $\gcd(a+b,ab)=1$. Hence the asymptotic
gcd is finite, and Theorem~\ref{mil:thm:ac-tail} applies. When
$b\ge4$, the degree contradiction
$3\le\ell\le1+(19b/40-1)^{-1}<3$ already applies to the full
clock by Saito's Theorem~2.3.
\end{proof}

For example, every sufficiently late tail of $2^n+3^n$ has a
transcendental minimum. This is a sum of independent geometric
components and does not rely on a polynomial perturbation of one
geometric component.

\subsection{An explicit parity refinement in the cubic-base example}
The following is an earlier elementary refinement. The stronger complete
finite-state analysis in Theorem~\ref{f13:thm:type-c-cutoff-new} lowers
the full-schedule cutoff to $1.624\times10^{14}$ in general and
$8.12\times10^{13}$ for even $r$.


Saito computes
\[
 \operatorname{agcd}(r3^k-1)=
 \begin{cases}1,&r\text{ even},\\2,&r\text{ odd},\end{cases}
\]
but uses the uniform upper bound two in the final estimate of
Theorem~1.9(C). Retaining the first case gives the following
numerical refinement.

\begin{corollary}\label{mil:cor:ac-even}
Let $x_0=\exp(3\exp(32.76))$. For every even integer $r$ satisfying
\[
 r>\frac{\log(2x_0^{1/3}+1)}{3\log\kappa}+\frac13,
\]
the number $\xi(r3^k-1)$ is transcendental. In particular, every
even integer $r\ge2.002\times10^{14}$ satisfies this condition.
\end{corollary}
\begin{proof}
The Type C conclusion gives a cubic Pisot number $\xi^g$ in
the algebraic case, with $g\mid\operatorname{agcd}(r3^k-1)=1$.
Thus $g=1$, and the same upper bound used in the source becomes
\[
 \xi\le(2x_0^{1/3}+1)^{1/(3r-1)}<\kappa,
\]
contradicting minimality of $\kappa$ among Pisot numbers.
The displayed real threshold is approximately
$2.0014802709094633\times10^{14}$, below the safe rounded bound
$2.002\times10^{14}$.
\end{proof}

This last result is a numerical corollary, not an unconditional
resolution for all $r\ge1$. In particular, this section does not
prove transcendence of $\xi(3^k-1)$. It does not remove the
global prime-between-cubes hypothesis from Saito's Corollary~2.7.

\section{Finite prefixes and nonlinear polynomial perturbations}
\label{mil:sec:return-extensions}

\subsection{Finite-prefix stability and its exact endpoint obstruction}
A \emph{left-gap point} of a set $W\subset(1,\infty)$ means a point
$A\in W$ for which $W\cap(A-\delta,A)=\varnothing$ for some $\delta>0$.
This is the same notion as a local left minimum. We use the localized
Type B theorem proved in Theorem~\ref{mil:thm:local-type-b}; its proof includes
the replacement-point argument and the exclusion of integral powers at
schedule coordinates. Thus no separate localization hypothesis is needed
in this section.

\begin{theorem}[Finite-prefix dichotomy]\label{mil:thm:finite-prefix}
Let $W$ be a prime-representing set whose every left-gap point is
transcendental. Given positive integers $E_1,\ldots,E_s$, put
\[
 W^*=W\cap\{x>1:\lfloor x^{E_j}\rfloor\text{ is prime for }1\le j\le s\}.
\]
Every algebraic left-gap point of $W^*$ belongs to
\[
 \bigcup_{j=1}^s\{p^{1/E_j}:p\text{ prime}\}.
\]
If, for each $j$, some exponent $C_k$ defining $W$ is an integer
multiple of $E_j$ with $C_k/E_j\ge2$, then every left-gap point of
$W^*$ is transcendental. In particular, whenever $W^*$ is nonempty,
its minimum is transcendental under this divisibility condition.
\end{theorem}
\begin{proof}
Suppose $A\in W^*$ and $A^{E_j}$ is nonintegral for every $j$.
All the finitely many added floor conditions are then constant in a
neighborhood of $A$. A left gap for $W^*$ is therefore a left gap for
$W$, proving transcendence. Otherwise $A^{E_j}=p$ for a prime $p$,
because the added condition is satisfied, giving the stated exceptional
form. If $C_k/E_j=t\ge2$, this would give $A^{C_k}=p^t$, contradicting
the prime condition in $W$. The minimum, when the intersection is
nonempty, exists by the same right-continuity argument as
\cite[Lemma~4.1]{SaitoVariants}, and is a left-gap point.
\end{proof}

Apply this theorem to a sufficiently late tail of
$R(n)=P(a^n)+dn+e$, with $d\ne0$, positive leading coefficient of $P$,
and $a^{\deg P}>40/19$. Theorem~\ref{sur:thm:clock-return} gives eventual universal return, which
verifies all Type B hypotheses for the tail, so all its left-gap points are
transcendental. The theorem retains an arbitrary finite prime prefix;
the only possible algebraic obstruction is a prime radical arising
exactly at a newly imposed floor boundary. This is stronger than a
tail-minimum assertion, but is not an unconditional removal of every
possible prefix obstruction.

A useful sufficient check for a prefix exponent $C_m=R(m)>0$ is
\begin{equation}\label{mil:eq:prefix-bad-prime-check}
 v_p(C_m)\le m v_p(a)\qquad(p\mid a).
\end{equation}
Then $a^m$ is already zero modulo the part of $C_m$ supported on primes
dividing $a$. On its coprime part, choose a multiplicative period of
$a$, and choose $n\equiv m$ modulo both that period and $C_m$.
Consequently $C_m\mid R(n)$ for arbitrarily large $n$, supplying the
required tail multiple. Only finitely many early checks are needed,
as the next theorem shows.

The check is not automatic. For $R(n)=3^n+9n+6$, $R(1)=18$, whereas
$R(n)\equiv6\pmod9$ for every $n\ge2$. Thus $18$ divides no later
exponent. Without an additional argument excluding the relevant prime
radicals from the full set, the preceding theorem does not settle that
full-prefix minimum.

\subsection{A general theorem for nonlinear drift}
\begin{theorem}[Polynomial drift and tail transcendence]
\label{mil:thm:nonlinear-drift}
Let $a\ge2$ be an integer and $P,Q\in\mathbb Z[X]$. Suppose $P$ has
positive leading coefficient and degree $r\ge1$, and put
\[
 C_n=P(a^n)+Q(n),\qquad D(X)=P(0)+Q(X).
\]
Assume $D$ is not the zero polynomial and $L=a^r>40/19$.
For every sufficiently large $N$, the minimum
\[
 \xi_N=\min\{x>1:\lfloor x^{C_n}\rfloor\text{ is prime for all }n\ge N\}
\]
exists and is transcendental. This includes every nonconstant integer
polynomial $Q$.

If $L\ge4$, every left-gap point of every sufficiently late tail is
transcendental. If $L\ge4$ and the full sequence additionally has
positive integer exponents, all successive ratios at least two, and the
earlier recurring-divisibility conditions, then every left-gap point of
the full prime-representing set is transcendental, without deleting a prefix.
\end{theorem}
\begin{proof}
First we prove eventual (B5). Since $D$ is a nonzero polynomial,
$D(m)\ne0$ for every sufficiently large $m$, and
$\log|D(m)|=O(\log m)$. For each prime $p\mid a$, eventually
\[
 v_p(D(m))<m v_p(a)\le v_p\bigl(P(a^m)-P(0)\bigr).
\]
Hence
\begin{equation}\label{mil:eq:drift-valuation}
 v_p(C_m)=v_p(D(m))<m v_p(a).
\end{equation}
Write $C_m=uv$, where $u$ is supported on the primes dividing $a$
and $\gcd(a,v)=1$. By \eqref{mil:eq:drift-valuation}, $u\mid a^m$.
Choose a positive period $t$ of $a$ modulo $v$ (take $t=1$ if $v=1$).
For every positive integer $j$, put $n=m+j\operatorname{lcm}(C_m,t)$.
Then $a^n\equiv a^m\pmod{C_m}$ and $Q(n)\equiv Q(m)\pmod{C_m}$.
Therefore
\[
 C_n\equiv C_m\equiv0\pmod{C_m}.
\]
This gives arbitrarily late multiples of every sufficiently late
exponent, proving (B5), including membership in the cofinite set $I$
of Saito's theorem.

Next we bound the asymptotic gcd. Let
$J=\{j\ge1:[X^j]P\ne0\}$, let $s=\deg D$, and let $b$ be the
leading coefficient of $D$. With $E$ denoting the forward shift and
$\Delta=E-1$, we have the exact integer identity
\begin{equation}\label{mil:eq:drift-constant}
 \Delta^s\prod_{j\in J}(E-a^j)C_n
   =K,\qquad K=s!\,b\prod_{j\in J}(1-a^j)\ne0.
\end{equation}
Indeed, the product annihilates every nonconstant exponential term,
and leaves a degree-$s$ polynomial with leading coefficient
$b\prod_{j\in J}(1-a^j)$; the finite difference produces the stated
constant. Every tail gcd divides $K$. In particular, any integer $g$
dividing every sufficiently late $C_n$ satisfies $g\mid K$.

For a sufficiently late tail, positivity, $c_{k+1}\ge2$, integrality,
and $c_{k+1}\to L>40/19$ verify (B1)--(B4). By (B5) and
\cite[Theorem~2.3]{SaitoVariants}, algebraicity of its minimum would
produce a Pisot number $\beta=\xi_N^g$ of degree $\ell$ satisfying
\begin{equation}\label{mil:eq:drift-degree}
 3\le\ell\le1+\left(\frac{19}{40}L-1\right)^{-1},
 \qquad g\mid K.
\end{equation}
When $L\ge4$, the upper bound is less than $3$, a contradiction.
The localized theorem gives exactly the same contradiction for every
left-gap point, proving the stronger assertion in this case.

It remains to treat $L=3$. We show uniformly that $\xi_N\to1$.
Choose $40/19<c<3$, so all sufficiently late ratios are at least $c$.
Let $\theta=21/40$. The Baker--Harman--Pintz input used in
\cite[Theorem~1.6]{SaitoVariants} supplies a prime in
$[x,x+x^\theta]$ for every sufficiently large $x$.
Choose a fixed sufficiently large prime $p$. Uniformly for every
$r\ge c$ and every $q\ge p$,
\[
 q^r+q^{\theta r}+1<(q+1)^r,
\]
since $(q+1)^r-q^r\ge r q^{r-1}$ and
$(1-\theta)c-1>0$. Starting with the same prime $p$ at the first
coordinate of any late tail, select the subsequent primes in these
intervals. Nested prime-root intervals give a member of that tail set
less than $(p+1)^{1/C_N}$. Consequently
\begin{equation}\label{mil:eq:drift-uniform-upper}
 1<\xi_N<(p+1)^{1/C_N}\longrightarrow1.
\end{equation}
For $L=3$, \eqref{mil:eq:drift-degree} forces $\ell=3$. Every cubic Pisot
number $\beta$ satisfies $\beta>5/4$: if $f$ is its monic minimal
polynomial and $\beta_2,\beta_3$ its other conjugates, then
\[
 1\le|f(1)|=(\beta-1)|1-\beta_2||1-\beta_3|<4(\beta-1).
\]
But \eqref{mil:eq:drift-uniform-upper} and $g\mid K$ imply
\[
 \beta=\xi_N^g<(p+1)^{|K|/C_N}<5/4
\]
for every sufficiently large $N$, a contradiction.
\end{proof}

\begin{remark}[What was extended]
The proof does not assert universal modular return for nonlinear
$Q$. It proves precisely the divisibility condition (B5) that is
needed, bounds the possible Pisot-power exponent by a fixed integer,
and then excludes the remaining cubic possibility for late minima.
The condition $D\ne0$ excludes a pure exponential polynomial with zero
constant term. In that case the constant $K$ in
\eqref{mil:eq:drift-constant} vanishes, and this argument gives no bound
for $g$. The case $a=2$, $\deg P=1$ has limiting ratio $2$ and also
lies outside the present theorem.
\end{remark}

\subsection{A full nonlinear example where (B6) fails}
\begin{proposition}\label{mil:prop:nonlinear-b6-counterexample}
For the full schedule
\[
 C_n=8^n+9n^2\qquad(n\ge1),
\]
the least prime-representing constant and every left-gap point are
transcendental. Nevertheless Saito's condition (B6) fails at $m=2$.
\end{proposition}
\begin{proof}
The ratios tend to $8$. Moreover, $8^n\ge(3/2)n^2$ for $n\ge1$, so
\[
 C_{n+1}-2C_n=6\cdot8^n+9(-n^2+2n+1)>0.
\]
For every $n\ge1$,
\[
 v_2(C_n)=2v_2(n)<3n,
\]
because the two summands have different $2$-adic valuations.
Thus \eqref{mil:eq:prefix-bad-prime-check} holds at every index, and the
preceding congruence construction proves (B5) for the full schedule.
The Type-B degree bound is incompatible with degree at least three,
so the localized theorem proves the first assertion.

For the second assertion, $C_2=100$. Suppose $n\ge3$ and
$C_n\equiv100\pmod{128\cdot65537}$. Put $j=3n$. Since $j\ge9$,
the congruence modulo $128$ yields $j^2\equiv100\pmod{128}$, and hence
\[
 j\equiv10\text{ or }22\pmod{32}.
\]
The prime $65537$ satisfies $\operatorname{ord}_{65537}(2)=32$.
It follows that $2^j$ is congruent to $1024$ or $-64$, respectively.
The other congruence would therefore require
\[
 j^2\equiv-924\text{ or }164\pmod{65537}.
\]
Both right-hand sides are quadratic nonresidues. Indeed, quadratic
reciprocity gives
\[
 \left(\frac{-924}{65537}\right)
 =\left(\frac3{65537}\right)
  \left(\frac7{65537}\right)
  \left(\frac{11}{65537}\right)=(-1)^3=-1,
\]
while
\[
 \left(\frac{164}{65537}\right)
 =\left(\frac{41}{65537}\right)
 =\left(\frac{19}{41}\right)
 =\left(\frac3{19}\right)=-1.
\]
Thus no such later $n$ exists. Taking
$L_0=32\cdot65537=2097184$ makes $L_0C_2$ a multiple of
$128\cdot65537$, so (B6) fails at $(m,L)=(2,L_0)$.
\end{proof}

This example disproves an unrestricted all-seed modular-return
extension to nonlinear drift. It does not disprove eventual universal
return for a fixed polynomial $Q$; that stronger assertion is not
needed for Theorem~\ref{mil:thm:nonlinear-drift} and remains unproved here.

\section{Recalled cubic identities and the exact gap exponent}
\label{mil:sec:recalled-cubic}
This section collects the cubic identities and gap asymptotics used in the arithmetic restrictions that follow.
Saito's reduction implies that, if $A$ is algebraic, some $\beta=A^{3^m}$ has conjugates
\[
 \beta,\ -\rho,\ \gamma,\qquad
 0<\sigma=|\gamma|<\rho<1,
 \qquad \rho\le\beta^{-17/40}.
\]
The number $\beta$ is a totally real cubic Pisot number. Set $M=3^r$, $p_r=\lfloor\beta^M\rfloor$, $g_r=p_{r+1}-p_r^3$, and $N=|\Norm(\beta)|$.

\begin{lemma}[Eventual least-prime choice]
For the actual least Mills constant, each sufficiently late next output is the least prime strictly above the cube of the preceding output.
\end{lemma}
\begin{proof}
The prime-gap bound gives a prime between $q^3$ and $q^3+O(q^{63/40})<(q+1)^3-1$ for every sufficiently large prime $q$. Thus every sufficiently large prime starts an indefinitely extendible nested cubic prime chain. If a smaller prime than the chosen next output existed above a late tower cube, replace that output by the smaller prime and extend it. The nested intervals preserve every earlier output and give a smaller prime-representing constant, contradicting minimality.
\end{proof}

\begin{proposition}[Previously established gap exponent]
Under algebraicity,
\[
 g_r\sim3(\beta^2\rho)^{3^r},\qquad
 \frac{\log g_r}{\log(p_r^3)}\longrightarrow
 \kappa=\frac{2+\log\rho/\log\beta}{3},
 \qquad \frac12<\kappa<\frac{21}{40}.
\]
Moreover $\kappa$ is transcendental.
\end{proposition}
\begin{proof}
The small-conjugate sum $-\rho^M+\gamma^M$ lies in $(-1,0)$ eventually, so $p_r=\Tr(\beta^M)$. For $X=\beta^M$, $Y=-\rho^M$, $Z=\gamma^M$, Newton's identity gives
\[
 g_r=-3(X+Y)(X+Z)(Y+Z)\sim3\beta^{2M}\rho^M.
\]
Since $N=\beta\rho\sigma$ is a positive integer and $\sigma<\rho$, one has $\beta\rho^2>N\ge1$, giving $\kappa>1/2$. The conjugate upper bound gives $\kappa\le21/40$.

We show that $t=\log\rho/\log\beta$ is irrational. Otherwise $\beta^u\rho^v=1$ for some positive integers $u,v$. In an $S_3$ splitting field, interchange the first two conjugates in $\beta^u(-\rho)^v=(-1)^v$ and divide the identities. Distinct absolute values force $u=v$, hence $\beta\rho=1$, contradicting $N=\beta\rho\sigma<1$. In a cyclic cubic splitting field, put $l_i=\log|\beta_i|$ and apply the three cyclic embeddings. The resulting equations $ul_1+vl_2=ul_2+vl_3=ul_3+vl_1=0$ have determinant $u^3+v^3>0$, contradicting $l_1>0$. Thus $t$ is irrational. If $t$ were algebraic, Gelfond--Schneider would make $\beta^t$ transcendental, contrary to $\beta^t=\rho$. Therefore $t$, and hence $\kappa$, is transcendental; equality at the rational upper endpoint is impossible.
\end{proof}
\section{Frobenius restrictions and quantitative prime divisors in the cubic shadow}

Throughout this section we assume that the least Mills constant $A$ is
algebraic.  Theorem~1.7 of Saito, \emph{Transcendency of variants of Mills'
constant}, arXiv:2508.16068v3,
\url{https://arxiv.org/html/2508.16068v3}, supplies a totally real cubic
Pisot power $\beta=A^{3^m}$ with conjugates
\[
 \alpha_1=\beta,\qquad \alpha_2=-\rho,\qquad \alpha_3=\gamma,
 \qquad 0<\sigma=|\gamma|<\rho<1.
\]
Let $f(X)=X^3-aX^2+bX-C$ be its minimal polynomial, let
$\Delta=\operatorname{disc}(f)$, and put $N=|C|$.
For $M_r=3^r$, write
\[
 p_r=\operatorname{Tr}(\beta^{M_r}),\qquad
 g_r=p_{r+1}-p_r^3,\qquad h_r=g_r/3.
\]
The sum $(-\rho)^{3^r}+\gamma^{3^r}$ lies in $(-1,0)$ for all sufficiently
large $r$, so $p_r=\lfloor\beta^{3^r}\rfloor$ is the Mills prime.
Also $h_r>0$: the prime $p_{r+1}$ is at least $p_r^3$ and cannot equal that
composite cube.  The classical cubic identities give
\[
 F_r:=\frac{h_{r+1}}{h_r}
 =\prod_{i<j}\bigl(\alpha_i^{2M_r}
       -\alpha_i^{M_r}\alpha_j^{M_r}+\alpha_j^{2M_r}\bigr)
 \in\mathbb N.
\]
Indeed, with $X_i=\alpha_i^{M_r}$, expansion gives
\[
 h_r=-(X_1+X_2)(X_2+X_3)(X_3+X_1).
\]
The right side is a symmetric algebraic integer, hence an integer.
Factoring each $X_i^3+X_j^3$ proves the formula for $F_r$, which is again
a symmetric algebraic integer and is positive because the roots are real.
The dominant terms give $h_r\sim B^{M_r}$ and $F_r\sim B^{2M_r}$, where
$B=\beta^2\rho$.  These identities supply the arithmetic input for the transition-factor
restrictions below.

\begin{theorem}[Frobenius and congruence sieve]\label{mil:thm:cubic-frobenius-sieve}
Let $\ell$ be a prime with $\ell\nmid6C\Delta$ and $\ell\mid F_r$.
Then $f$ is reducible modulo $\ell$.  More precisely,
\[
 \begin{array}{c|c}
 \text{factor degrees of }f\pmod\ell&\text{necessary congruence}\\ \hline
 (1,1,1)&\ell\equiv1\pmod{2\cdot3^{r+1}}\\
 (1,2)&\ell\equiv-1\pmod{2\cdot3^{r+1}}.
 \end{array}
\]
In particular, $\ell\geq2\cdot3^{r+1}-1$.
The assertion concerning $F_r$ is valid for any irreducible cubic with
nonzero constant term, without a Mills hypothesis.
\end{theorem}

\begin{proof}
Reduce the three distinct nonzero roots in a finite splitting field over
$\mathbb F_\ell$.  A vanishing factor in $F_r$ means that, for some pair,
$R^{3^r}$ is a primitive sixth root of unity, where $R$ is the ratio of
the corresponding original roots.  Consequently
\[
 \operatorname{ord}(R)=2\cdot3^{r+1}.
\]
If Frobenius permutes the three roots in a $3$-cycle, orient the chosen
pair cyclically.  Then
$R^{1+\ell+\ell^2}=1$.  The exponent is odd, which is impossible for an
element of the displayed even order.  If all roots lie in $\mathbb F_\ell$,
the order divides $\ell-1$, proving the first congruence.

In the remaining factorization type, let $a$ be the fixed root and $b,c$
the two roots exchanged by Frobenius.  If the vanishing factor uses the
ratio $b/c$, then $(b/c)^\ell=c/b$, so its order divides $\ell+1$.
If it uses $R=b/a$, its order divides $\ell^2-1$.  Suppose that
$\ell\equiv1\pmod3$.  Then the whole factor $3^{r+1}$ in its order divides
$\ell-1$, which is also even.  Thus $R^{\ell-1}=1$, so $R$ is fixed by
Frobenius and $b=c$, a contradiction.
Therefore $\ell\equiv-1\pmod3$.  All of the factor $3^{r+1}$ in the order
must divide $\ell+1$, and $\ell+1$ is even.  This proves the second
congruence.
\end{proof}

This restriction is stronger than the previously obtained fixed quadratic
norm condition: a cubic-inert prime is completely excluded from $F_r$,
rather than merely subjected to an even-valuation condition.

\begin{lemma}[Local multiplier valuations]\label{mil:lem:cubic-local-valuations}
For each fixed prime $\ell$, extend the normalized valuation $v_\ell$ to a
splitting field and order the root valuations as
\[
 0\leq u_\ell\leq v_\ell',
\]
where the three entries are $0,u_\ell,v_\ell'$.
For all sufficiently large $r$,
\[
 v_\ell(F_r)=2\cdot3^r u_\ell\quad(\ell\ne3),\qquad
 v_3(F_r)=2\cdot3^r u_3+c_3,
\]
where $c_3\in\{0,1,2\}$ is fixed.  Moreover
$2u_\ell\leq v_\ell(C)$.
\end{lemma}

\begin{proof}
Integrality makes all root valuations nonnegative.  Their minimum must
be zero: otherwise their power sum $p_r$ is divisible by $\ell$ for every
sufficiently large $r$, whereas $p_r$ is a prime tending to infinity.
The sum of the three valuations is $v_\ell(C)$, giving
$2u_\ell\leq u_\ell+v_\ell'=v_\ell(C)$.

For a pair of unequal root valuations, the corresponding quadratic factor
has valuation exactly $2\cdot3^r$ times the smaller valuation.
For a pair of equal valuations $t$, divide by the $2\cdot3^r$-th power of
one root.  The remaining factor is
$\Phi_6(z^{3^r})$, where $z$ is a local unit and
$\Phi_6(X)=X^2-X+1$.

When $\ell\ne3$, sufficiently large $r$ kills the $3$-primary part of the
order of the residue of $z$.  A root of $\Phi_6$ has order $6$ in
characteristic at least $5$, and order $3$ in characteristic $2$.
Thus $\Phi_6(z^{3^r})$ is eventually a unit.
When $\ell=3$, its reduction is $(X+1)^2$.
If the residue of $z$ differs from $-1$, its $3^r$-th power also differs
from $-1$, so the factor is a unit.  Otherwise $z^{3^r}\to-1$ in the local
field, and $\Phi_6(z^{3^r})\to3$; its normalized valuation is eventually
exactly $1$.
At most two of the three pair ratios can reduce to $-1$, since three
cyclically oriented such ratios would have product $-1$ instead of $1$.

Finally, the sum of the smaller valuations over the three pairs is
$u_\ell$.  Summing the preceding formulas proves the lemma.
\end{proof}

\begin{theorem}[A large allowed prime at every late transition]
\label{mil:thm:cubic-large-primes}
Let $S$ be the primes dividing $6C\Delta$, and let $(n)_S$ denote the part
of a positive integer $n$ supported on $S$.  For all sufficiently large $r$,
\[
 (F_r)_S\leq9N^{3^r}.
\]
Every such $F_r$ has a prime divisor outside $S$.  Consequently,
\[
 P^+(F_r)\geq2\cdot3^{r+1}-1,
 \qquad P^+(g_r)\geq2\cdot3^r-1
\]
for all sufficiently large $r$, where $P^+$ denotes the largest prime
divisor.
\end{theorem}

\begin{proof}
Apply Lemma~\ref{mil:lem:cubic-local-valuations} simultaneously to the finite
set $S$.  Its valuation inequalities give
\[
 (F_r)_S\leq3^2\prod_{\ell\in S}
       \ell^{3^r v_\ell(C)}=9N^{3^r}.
\]
On the other hand,
\[
 \frac{B^2}{N}
 =\frac{\beta^4\rho^2}{\beta\rho\sigma}
 =\frac{\beta^3\rho}{\sigma}>\beta^3>1.
\]
Since $F_r\sim B^{2\cdot3^r}$, it follows that
$F_r/(F_r)_S\to\infty$.  Thus a prime divisor outside $S$ exists at each
sufficiently late stage, and Theorem~\ref{mil:thm:cubic-frobenius-sieve} gives
the first largest-prime bound.  The second follows from
$F_{r-1}\mid h_r\mid g_r$.
\end{proof}

\begin{theorem}[Number and radical of gap prime divisors]
\label{mil:thm:cubic-radical-growth}
Under the same hypothetical algebraic Mills shadow,
\[
 \omega(g_r)\geq\frac r2-O(1),\qquad
 \log\operatorname{rad}(g_r)
     \geq\frac{\log3}{4}r^2-O(r).
\]
Equivalently,
\[
 \omega(g_r)\geq\frac{\log\log g_r}{2\log3}-O(1),
 \qquad
 \log\operatorname{rad}(g_r)
 \geq\frac{(\log\log g_r)^2}{4\log3}-O(\log\log g_r).
\]
\end{theorem}

\begin{proof}
First we give a direct finite-field proof that a prime $\ell\notin S$
can divide at most two distinct multipliers $F_j$.
For each of the three original pair ratios, its order in a finite splitting
field is fixed.  If its factor vanishes in $F_j$, that order must be exactly
$2\cdot3^{j+1}$, so each pair can account for at most one index.
Moreover, among the $3$-primary orders of the three ratios the largest
occurs at least twice.  Indeed, write the ratios as $R,Q,RQ$; in a cyclic
$3$-group the product of elements of unequal orders has the larger order.
Consequently their $3$-primary orders have at most two distinct values,
and the prime can occur at at most two indices $j$.

Choose an index $j_0$ after which Theorem~\ref{mil:thm:cubic-large-primes}
applies, and choose one prime $q_j\notin S$ dividing $F_j$ for each
$j_0\leq j<r$.  Each chosen prime occurs at most twice and every one divides
$h_r$, giving $\omega(g_r)\geq(r-j_0)/2$.
Furthermore,
\[
 \prod_{j=j_0}^{r-1}q_j\leq\operatorname{rad}(g_r)^2,
 \qquad q_j\geq2\cdot3^{j+1}-1.
\]
Taking logarithms and summing yields
$\log\operatorname{rad}(g_r)\geq(\log3)r^2/4-O(r)$.
The equivalent forms use $g_r\sim3B^{3^r}$, hence
$\log\log g_r=r\log3+O(1)$.
\end{proof}

\begin{proposition}[A $3$-adic embedding and cyclic cubic shadows]
\label{mil:prop:cubic-split-three}
The polynomial $f$ has a root in $\mathbb Q_3$, so the cubic field admits
an embedding into $\mathbb Q_3$ and has a prime above $3$ with ramification
index and residue degree both equal to one.
If $\mathbb Q(\beta)/\mathbb Q$ is Galois, then $3$ splits completely in
$\mathbb Q(\beta)$.  In the Galois case every prime divisor of $F_r$ outside $S$
splits completely in the cubic field and satisfies
$\ell\equiv1\pmod{2\cdot3^{r+1}}$.
\end{proposition}

\begin{proof}
We first establish the elementary prime-tail rigidity needed here, rather
than leave it as a reference to an earlier draft.
Write $T_r=\operatorname{Tr}(\beta^{3^r})$.
The integral companion matrix of $f$ gives the standard trace congruence
\[
 T_{r+1}-T_r\in3^{r+1}\mathbb Z.
\]
One proof groups the terms in the matrix trace into cyclic rotation orbits
of closed index words; the orbits of length $3^j$ contribute multiples of
$3^j$, and the shorter repeated words satisfy
$u^{3^t}\equiv u^{3^{t-1}}\pmod{3^t}$ for integers $u$.
Thus $T_r$ converges to a limit $L\in\mathbb Z_3$ with
$L-T_r\in3^{r+1}\mathbb Z_3$.

For a sufficiently late prime $q=T_r$, exclude the fixed divisors of
$3C\Delta$, and let $d\in\{1,2,3\}$ be the degree of the finite splitting
field of $f\pmod q$.
If $v_3(q^d-1)\leq r$, cubing is periodic on all the powered roots
$\alpha_i^{3^r}$: writing $q^d-1=3^t u$ with $3\nmid u$, any positive
$K$ with $3^K\equiv1\pmod u$ makes $q\mid T_{r+jK}$ for every $j\geq0$.
For large $j$ this contradicts primality.  Hence $v_3(q^d-1)>r$.
For $d=1,2$ this forces $q\equiv\pm1\pmod{3^{r+1}}$; for $d=3$,
lifting the exponent gives $q\equiv1\pmod{3^r}$.
In all cases $T_r\equiv\varepsilon\pmod{3^r}$, with a fixed sign
$\varepsilon\in\{1,-1\}$ because $T_r\equiv a\pmod3$.
Therefore $L=\varepsilon$.
Using $T_0=a$ and $T_1=a^3-3ab+3C$, the congruences modulo $3$ and $9$
give $a\equiv\varepsilon\pmod3$ and $C\equiv ab\pmod3$.  Consequently
\[
 a\equiv\varepsilon\pmod3,
 \qquad f(X)\equiv(X-\varepsilon)(X^2+b)\pmod3.
\]
If $b\not\equiv2\pmod3$, the root $\varepsilon$ of the displayed residual
polynomial is simple.  If $b\equiv2\pmod3$, the root $-\varepsilon$ is
simple.  Hensel's lemma therefore supplies a root of $f$ in $\mathbb Z_3$.
This gives the asserted embedding and degree-one completion.  In a Galois
extension all primes above $3$ have the same local degree, so a degree-one
completion implies complete splitting.

For a prime outside $S$, the Frobenius permutation in a cyclic cubic
field is either the identity or a $3$-cycle.  The latter is excluded by
Theorem~\ref{mil:thm:cubic-frobenius-sieve}; the remaining conclusion follows
from the same theorem.
\end{proof}

These results do not prove that a new rational prime divisor appears at
every stage: a prime may first divide one pair-factor and later a different
pair-factor.  They also do not contradict the known distribution of primes
in residue classes or splitting classes.  They are additional necessary
conditions on a hypothetical algebraic Mills constant, not a proof of
unconditional transcendence or a resolution of the full totally real cubic
Pisot problem.  A contradiction still requires new information about the
least-prime choice or the entire trace sequence.

\section{Local gap criteria and effective bounded-degree certificates}

Saito's Theorem~1.8 in arXiv:2508.16068v3 already proves that the
Density Hypothesis, and hence the Riemann Hypothesis, implies the
transcendence of Mills' constant. The results below separate a much
more localized sufficient condition from an effective finite test for
algebraic candidates of bounded degree.

Write $A$ for the least Mills constant and
\[
 P_k=\lfloor A^{3^k}\rfloor,\qquad X_k=P_k^3,
 \qquad G_k=P_{k+1}-P_k^3.
\]
We use two previously established facts: eventually
$P_{k+1}=\operatorname{nextprime}(P_k^3)$; and if $A$ is algebraic, then
\[
 \frac{\log G_k}{\log X_k}\longrightarrow\kappa,
 \qquad \frac12<\kappa<\frac{21}{40},
 \qquad \kappa\text{ is transcendental}.
\]
These facts were proved in Section~\ref{mil:sec:recalled-cubic}.

\begin{corollary}[A localized subsequence criterion]
Mills' constant is transcendental if
\[
 \liminf_{k\to\infty}
 \frac{\log(\operatorname{nextprime}(X_k)-X_k)}{\log X_k}
 \leq\frac12.
\]
Equivalently, it suffices that for every $\varepsilon>0$ there are
infinitely many $k$ for which $(X_k,X_k+X_k^{1/2+\varepsilon}]$
contains a prime. It also suffices that the sequence
$\log G_k/\log X_k$ does not converge, or that it has an algebraic
real accumulation point.
\end{corollary}

\begin{proof}
Under algebraicity the displayed sequence converges to the
transcendental number $\kappa>1/2$. Each of the stated alternatives
contradicts this. The equivalence of the first two formulations is the
definition of the lower limit and eventual least-prime choice.
\end{proof}

Only a single prime at each of infinitely many prescribed tower
cubes is required. No density of primes in these intervals, and no
assertion for all real starting points, is needed. Conversely, an
almost-everywhere prime-gap theorem does not automatically control
this prescribed sparse sequence. This criterion is a consequence of
the exact gap exponent, not an unconditional proof of transcendence.

\begin{theorem}[Degree, signature, and polynomial structure]\label{mil:thm:mills-degree-structure}
If the least Mills constant $A$ is algebraic of degree $d$, then
\[
 d=3^{t+1}\quad\text{for some }t\geq0,
 \qquad \beta=A^{d/3}
\]
is a totally real cubic Pisot number with a negative dominant small
conjugate. If $f(X)=X^3-aX^2+bX-c$ is its minimal polynomial, then
\[
 f(X^{d/3})=X^d-aX^{2d/3}+bX^{d/3}-c
\]
is the minimal polynomial of $A$. In particular, the signature of
$\mathbb Q(A)$ is $(3,(d-3)/2)$.
\end{theorem}

\begin{proof}
Saito's Theorem~1.7 supplies $A^{3^m}$ with the stated cubic Pisot
configuration. Let $K=\mathbb Q(A^{3^m})$, a totally real cubic field.
As long as $A^{3^u}$ is a cube in $K$ and $u>0$, its unique real
positive cube root $A^{3^{u-1}}$ belongs to $K$. It remains a totally
real cubic Pisot number: every embedding cubes to an embedding of
the preceding Pisot number, and the degree is three since its cube
already generates $K$. The negative dominant small conjugate is
preserved under taking odd real roots.

This procedure stops either at $A\in K$, when $d=3$, or at
$\beta=A^{3^t}\in K$ with $t\geq1$ and $\beta$ not a cube in $K$.
We claim that $X^{3^t}-\beta$ is irreducible over $K$. More generally,
if $b$ is not a cube in a characteristic-zero field, then
$X^{3^v}-b$ is irreducible for every $v\geq1$. For $v=1$ this is the
criterion for a cubic to have a root. Inductively, if $\theta$ is a
root of the irreducible polynomial $X^{3^{v-1}}-b$, then
\[
 \operatorname{Norm}_{K(\theta)/K}(\theta)=b,
\]
because the degree is odd. If $\theta$ were a cube in $K(\theta)$,
its norm would make $b$ a cube in $K$, a contradiction. Thus
$X^3-\theta$ is irreducible over $K(\theta)$, proving the induction
by the tower law.

Applying this to $\beta$ gives $[\mathbb Q(A):K]=3^t$, hence
$d=3^{t+1}$ and $A^{d/3}=\beta$. Integrality follows from the
integrality of the Pisot power. The monic polynomial $f(X^{d/3})$
annihilates $A$ and has degree $d$, so it is its minimal polynomial.
Each of the three nonzero real roots of $f$ has exactly one real
$(d/3)$rd root, because $d/3$ is odd. This proves the signature.
\end{proof}

\begin{theorem}[An explicit bounded-degree onset]\label{mil:thm:bounded-degree-onset}
Fix an integer $D\geq3$ and a rational number $U>1$. Define
\[
 G=3^j=\max\{3^v:v\in\mathbb Z_{\geq0},\ 3^v\leq D/3\},
 \qquad T=U^G,
\]
and
\[
 K(D,U)=j+\min\{r\in\mathbb Z_{\geq0}:3^r\geq4(T+1)\}.
\]
Suppose that $A$ is algebraic, $[\mathbb Q(A):\mathbb Q]\leq D$,
and $A\leq U$. Then for every $k\geq K(D,U)$ the quantities
\[
 h_k=\frac{P_{k+1}-P_k^3}{3},\qquad
 e_k=\operatorname{cres}_{P_k}(h_k)
\]
are well defined integers and satisfy
\[
 \boxed{e_{k+1}=e_k^3.}
\]
Here $\operatorname{cres}_{p}(h)$ is the unique representative of $h$
modulo the odd integer $p$ in $(-p/2,p/2)$.
\end{theorem}

\begin{proof}
By Theorem~\ref{mil:thm:mills-degree-structure}, the degree is $d=3^{t+1}$
and $A^{3^t}$ has the stated totally real cubic Pisot configuration.
Since $3^t=d/3\leq D/3$, we have $t\leq j$. Thus $\beta=A^G$ is an
odd power of that cubic Pisot number. Its degree remains three:
conjugate fibers under powering have equal size, while the dominant
Pisot conjugate is the unique member of its fiber, forcing this size
to be one. The real-root configuration and signs are preserved.

We may therefore label the conjugates of $\beta$ as
\[
 \beta,\ -\rho,\ \gamma,
 \qquad 0<|\gamma|<\rho<1,\qquad \beta\leq T.
\]
If $f\in\mathbb Z[X]$ is its monic irreducible cubic polynomial, then
$f(-1)$ is a nonzero integer, so
\[
 1\leq|f(-1)|
   =(\beta+1)(1-\rho)|1+\gamma|
   \leq2(T+1)(1-\rho).
\]
It follows that $\rho\leq1-1/[2(T+1)]$.

For $k\geq K(D,U)$ put $n=3^{k-j}\geq4(T+1)$. Then
$\rho^n\leq e^{-2}<1/4$. Since $n$ is odd and $|\gamma|<\rho$,
\[
 -\frac12<-\rho^n+\gamma^n<0.
\]
The trace is an integer, and hence
\[
 P_k=\lfloor\beta^n\rfloor
     =\beta^n-\rho^n+\gamma^n=\operatorname{Tr}(\beta^n).
\]
Let $b_k$ be the sum of pairwise products of these three conjugate
powers, and put $c_k=\operatorname{Norm}(\beta)^n$. Newton's identity
now gives
\[
 P_{k+1}=P_k^3-3P_kb_k+3c_k,
 \qquad h_k=c_k-P_kb_k.
\]
Furthermore $P_k\geq\beta^n/2$, whence
\[
 \frac{|c_k|}{P_k}
 \leq2(\rho|\gamma|)^n
 \leq2\rho^{2n}<\frac12.
\]
Thus centered reduction recovers $c_k$ exactly, namely $e_k=c_k$.
Since $c_{k+1}=c_k^3$, the assertion follows. The primes $P_k$ are odd
at all these stages, since $P_{k+1}>P_k^3$ and $K(D,U)\geq2$.
\end{proof}

\begin{corollary}[Finite exclusion certificates]
Suppose $U$ is a certified upper bound for the actual least Mills
constant. If three consecutive actual Mills floors
$P_k,P_{k+1},P_{k+2}$ are certified at some $k\geq K(D,U)$ and either
one of the required divisibilities by $3$ fails or
$e_{k+1}\ne e_k^3$, then $A$ is not algebraic of degree at most $D$.

Alternatively, if one seeks to exclude algebraic numbers of degree
at most $D$ and primitive minimal-polynomial coefficient height at
most $H$, it suffices to use $U=H+1$, by Cauchy's root bound.
\end{corollary}

This is an effective necessary-condition theorem. It does not assert
that a required failure has been computed. A prime chain obtained by
repeated greedy choices is not sufficient input unless its entries
have also been certified as floors of the actual least Mills
constant. No decimal estimate depending on RH is used.

\begin{theorem}[An effective gap margin at bounded degree]\label{mil:thm:bounded-degree-gap}
Retain the algebraicity, degree, and upper-bound hypotheses and the
notation of Theorem~\ref{mil:thm:bounded-degree-onset}. Put
\[
 \Delta=(T+3)^{-2},\qquad
 \eta=\frac{\log(1+\Delta)}{6\log T}>0,
 \qquad R_{\mathrm{gap}}=
 \max\left\{4(T+1),\frac{\log2}{\log(1+\Delta)}\right\}.
\]
For every $k\geq j$ with $3^{k-j}\geq R_{\mathrm{gap}}$, one has
\[
 \boxed{G_k\geq\frac{27}{32}X_k^{\,1/2+\eta}.}
\]
In particular, one certified actual gap violating this inequality
excludes every algebraic candidate of degree at most $D$ bounded by
$U$.
\end{theorem}

\begin{proof}
Write $\sigma=|\gamma|$ for the smaller conjugate modulus of $\beta$.
First we show $\rho-\sigma\geq\Delta$. If $\gamma=-\sigma$, the
positive integer discriminant of the irreducible cubic gives
\[
 1\leq\sqrt{\operatorname{disc}(f)}
   =(\beta+\rho)(\beta+\sigma)(\rho-\sigma)
   \leq(T+1)^2(\rho-\sigma).
\]
If $\gamma=\sigma$, set $a=\beta-\rho+\sigma=\operatorname{Tr}(\beta)
\in\mathbb Z$. Since $f(a)$ is a nonzero integer,
\[
 1\leq|f(a)|
   =(\rho-\sigma)|a+\rho|\,|a-\sigma|
   \leq(T+3)^2(\rho-\sigma).
\]
Both cases imply the desired separation. Since $\sigma<1$, it follows
that $\rho/\sigma\geq1+\Delta$. The absolute norm
$N=\beta\rho\sigma$ is a positive integer, so
\[
 \beta\rho^2=N\frac{\rho}{\sigma}\geq1+\Delta.
\]
Consequently the algebraic gap exponent obeys the uniform bound
\[
 \kappa-\frac12
 =\frac{\log(\beta\rho^2)}{6\log\beta}
 \geq\frac{\log(1+\Delta)}{6\log T}=\eta.
\]

Now put $n=3^{k-j}$ and
$X=\beta^n$, $Y=-\rho^n$, $Z=\gamma^n$. The threshold for $n$ gives
$\rho^n<1/4$ and $(\sigma/\rho)^n\leq1/2$. The preceding theorem
identifies $P_k=X+Y+Z<X$ and the next floor with the cubic trace.
Thus the exact identity
\[
 G_k=X^3+Y^3+Z^3-(X+Y+Z)^3
     =-3(X+Y)(X+Z)(Y+Z)
\]
applies. We have $X+Y\geq3X/4$, $X+Z\geq3X/4$, and
$-(Y+Z)\geq\rho^n/2$, regardless of the sign of $\gamma$.
Therefore
\[
 G_k\geq\frac{27}{32}X^2\rho^n
       =\frac{27}{32}X^{3\kappa}
       \geq\frac{27}{32}(P_k^3)^{1/2+\eta},
\]
as claimed.
\end{proof}

The explicit gap bound is conditional on bounded-degree algebraicity.
It supplies a finite exclusion test, but does not assert that the
actual Mills chain has been shown to violate that test.

\section{Square clocks: a block criterion in every algebraic degree}

Write
\[
 \mathcal W_2=\{A>1:\lfloor A^{2^k}\rfloor\text{ is prime for every }k\ge1\},
 \qquad \xi_2=\min\mathcal W_2.
\]
Existence of this minimum is known; its irrationality is the explicit
Question~1.6 of Saito, \emph{Mills' constant is irrational},
arXiv:2404.19461v2. The results below do not answer that question
unconditionally. They give an all-degree criterion based on long consecutive blocks of
square gaps.

We use the following established inputs from Saito,
\emph{Transcendency of variants of Mills' constant},
arXiv:2508.16068v3: Lemma~6.1 (exponentially close algebraic powers
yield a Pisot power), Lemma~5.8 (identification of the nearest integer
with a trace), Lemma~5.10 (irrational arguments of nonreal Pisot
conjugates), Lemma~5.12, attributed there to Smyth (equal moduli occur
only for a complex conjugate pair), and Lemma~5.13 (a finite
irrationality exponent for such an argument, from linear forms in
logarithms). All references to numbering in this section concern
those specific arXiv versions.

\begin{lemma}[A doubling survivor calculation]\label{mil:sq-survivor}
Fix $1/5<a<1/4$ and put $I=[a,1-a]$. If
\[
 \{2^j x\}\in I\qquad(0\le j\le L),
\]
where $L\ge1$ is an integer, then
\[
 \operatorname{dist}\bigl(\{x\},\{1/3,2/3\}\bigr)\le2^{-L}.
\]
Consequently, the only points whose entire forward doubling orbit
stays in $I$ are $1/3$ and $2/3$.
\end{lemma}
\begin{proof}
A point of $I$ whose next iterate also lies in $I$ must lie in
\[
 I_0=[a,(1-a)/2],\qquad I_1=[(1+a)/2,1-a].
\]
Since $2a>(1-a)/2$, doubling modulo one cannot carry a point of
$I_0$ back into $I_0$. By reflection about $1/2$, it cannot carry
a point of $I_1$ back into $I_1$ either. Thus the first $L$ binary
digits of $\{x\}$ alternate. They agree with the first $L$ digits
of $1/3=0.010101\ldots{}_2$ or $2/3=0.101010\ldots{}_2$.
The corresponding binary cylinder has length $2^{-L}$, proving
the bound. Letting $L$ tend to infinity proves the last assertion.
\end{proof}

\begin{theorem}[Consecutive square-gap blocks]\label{mil:sq-block}
Let $A\in\mathcal W_2$, and put
\[
 p_k=\lfloor A^{2^k}\rfloor,\qquad d_k=p_{k+1}-p_k^2.
\]
Suppose that, for some fixed $\delta>0$, there are integers
$m_j\to\infty$ and $L_j\ge1$ such that
\begin{equation}\label{mil:sq-good-block}
 d_k\le p_k^{1-\delta}
 \quad(m_j\le k\le m_j+L_j),
 \qquad \frac{L_j}{m_j}\longrightarrow\infty.
\end{equation}
Then $A$ is transcendental. In particular, the conclusion holds
if $d_k\le p_k^{1-\delta}$ for every sufficiently large $k$.
\end{theorem}
\begin{proof}
Call an index satisfying the gap inequality good. For a good index,
write $A^{2^k}=p_k+u_k$ with $0<u_k<1$. The strict lower bound
holds because an integral value at one index would make the next
output a composite square. Since
\[
 (p_k+u_k)^2<p_{k+1}+1=p_k^2+d_k+1,
\]
we obtain
\begin{equation}\label{mil:sq-small-frac}
 0<u_k<\frac{d_k+1}{2p_k}\ll p_k^{-\delta}
 \ll A^{-\delta 2^k}.
\end{equation}
Here an infinite set of good indices necessarily has $\delta<1$;
alternatively one can replace $\delta$ by any smaller positive
number below one. Assume for contradiction that $A$ is algebraic.
Apply Saito's Lemma~6.1 to the increasing sequence $2^k$ with $k$
good. It gives a positive integer $g$, dividing every sufficiently
late such $2^k$, for which $\beta=A^g$ is a Pisot number. Hence
$g=2^h$ for an integer $h\ge0$.

If $\beta$ has degree one, it is an integer at least two, and
$p_k=\beta^{2^{k-h}}$ is composite for all sufficiently large $k$.
Thus its degree $\ell$ is at least two. Write its conjugates as
$\beta=\beta_1,\beta_2,\ldots,\beta_\ell$ and put
\[
 S(n)=\sum_{i=2}^{\ell}\beta_i^n.
\]
By \eqref{mil:sq-small-frac} and Saito's Lemma~5.8, at every sufficiently
large good index, with $n=2^{k-h}$, one has
\begin{equation}\label{mil:sq-trace-sign}
 \operatorname{Tr}(\beta^n)=p_k,
 \qquad S(n)=-u_k<0.
\end{equation}

Smyth's conjugate-modulus lemma implies that the small conjugates
of maximum modulus consist of either a single real conjugate or
one nonreal conjugate pair. In the first case, the even powers of
that real conjugate are positive and dominate the remaining
terms, so $S(2^r)>0$ for all sufficiently large $r$. This contradicts
\eqref{mil:sq-trace-sign}.

In the second case write the maximal pair as
$z=\rho e^{2\pi i\phi}$ and $\overline z$, with $0<\rho<1$.
All other small conjugates have modulus at most $\sigma<\rho$.
If there are no others, the remainder below is zero. We have
\[
 S(n)=2\rho^n\cos(2\pi n\phi)+O(\sigma^n).
\]
It follows from \eqref{mil:sq-trace-sign} that, at good indices,
\[
 \cos(2\pi 2^{k-h}\phi)\le O((\sigma/\rho)^{2^{k-h}})=o(1).
\]
Choose any fixed $1/5<a<1/4$. Outside $[a,1-a]$ the cosine
is bounded below by the positive number $\cos(2\pi a)$.
Therefore every sufficiently late good index satisfies
\begin{equation}\label{mil:sq-angle-strip}
 \{2^{k-h}\phi\}\in[a,1-a].
\end{equation}
Saito's Lemma~5.10 says that $\phi$ is irrational.

Applying Lemma~\ref{mil:sq-survivor} along a block in
\eqref{mil:sq-good-block}, and multiplying its conclusion by three,
gives
\[
 \|3\cdot2^{m_j-h}\phi\|\le3\cdot2^{-L_j}.
\]
Saito's Lemma~5.13 supplies constants $c>0$ and $\lambda>0$,
depending on $z$, such that $\|q\phi\|\ge c q^{-\lambda}$
for every integer $q\ge2$. Explicitly, choose $\phi\in(0,1)$,
put $w=z/|z|$, and take the fixed logarithms
$\log w=2\pi i\phi$ and $\log(-1)=\pi i$. For a nearest integer
$p$ to $q\phi$, the nonzero linear form
\[
 q\log w-2p\log(-1)=2\pi i(q\phi-p)
\]
has integer coefficients of size $O(q)$. Baker's
\emph{Transcendental number theory} (1975), Theorem~3.1,
as stated in Saito's Theorem~5.5, gives the asserted bound.
Combining the two bounds yields
$L_j\le\lambda m_j+O(1)$, a contradiction.

For the special case of all sufficiently large indices being
good, the logarithmic-forms bound is unnecessary: the entire
tail orbit in \eqref{mil:sq-angle-strip} survives, so
Lemma~\ref{mil:sq-survivor} makes $2^{k-h}\phi$ congruent to $1/3$
or $2/3$, contradicting irrationality directly.
\end{proof}

The use of the unique maximal conjugate or conjugate pair is the
reason that Theorem~\ref{mil:sq-block} has no restriction such as
$\delta>1/3$ or $\delta>1/2$. An isolated-index criterion and this block criterion have different
hypotheses: an arbitrary infinite set of good indices need not contain
the blocks required here. This distinction is essential.

\begin{corollary}[A conditional answer for the square constant]
\label{mil:sq-gap-corollary}
Suppose that, for some fixed $0<\theta<1/2$ and $C>0$, the least
prime strictly exceeding every sufficiently large real $x$ is at
most $x+Cx^\theta$. Then $\xi_2$ is transcendental.
\end{corollary}
\begin{proof}
For every sufficiently large prime $p$, put
$q=\operatorname{nextprime}(p^2)$. The hypothesis gives
\[
 p^2<q\le p^2+Cp^{2\theta}<(p+1)^2-1.
\]
Consequently, from any sufficiently large prime at any square-clock
index one can continue indefinitely by this least-prime rule.
The intervals
\[
 [p_k^{1/2^k},(p_k+1)^{1/2^k}]
\]
are nested; moreover, each child's upper endpoint is strictly
below its parent's upper endpoint. Their intersection therefore
gives a prime-representing constant with the prescribed initial
prime, with no ambiguity at the half-open upper endpoints.

For the least square constant, its prime sequence must therefore
obey
\[
 p_{k+1}=\operatorname{nextprime}(p_k^2)
\]
at every sufficiently large index. Indeed, any smaller admissible next prime could be
continued by the preceding construction, inside all earlier
cylinders, to obtain a smaller element of $\mathcal W_2$.
Choose $0<\delta<1-2\theta$. Then
$d_k\le Cp_k^{2\theta}\le p_k^{1-\delta}$ eventually, and apply
Theorem~\ref{mil:sq-block}.
\end{proof}

This is a conditional result. The required prime-interval bound is
not supplied by the usual RH estimate of order $\sqrt{x}\log x$.
No unconditional irrationality or transcendence of $\xi_2$ is
asserted here.

\subsection{Arbitrarily sparse multiplier-four approximants}

\begin{proposition}\label{mil:sq-sparse-approx}
Let $T\subset\{2,3,\ldots\}$ be any infinite set with no two
consecutive elements. For $N\ge1$, retain the exponent indices
\[
 S_N=\mathbb N\setminus\{t\in T:t>N\},
\]
list them in increasing order as $n_{N,1},n_{N,2},\ldots$, and put
\[
 \eta_N=\min\{A>1:\lfloor A^{2^{n_{N,j}}}\rfloor
                     \text{ is prime for every }j\ge1\}.
\]
Then each $\eta_N$ is transcendental and
\[
 \eta_N\le\eta_{N+1}\le\xi_2,
 \qquad \lim_{N\to\infty}\eta_N=\xi_2.
\]
The set $T$ may be chosen with arbitrarily large successive gaps.
\end{proposition}
\begin{proof}
The retained exponent clock has integer successive ratios two
or four. Since the omitted indices are isolated and infinitely
many remain after each finite cutoff, ratio four occurs infinitely
often. Saito's arXiv:2404.19461v2, Theorem~1.5, therefore supplies
existence and transcendence of each $\eta_N$. This is an existing
transcendence theorem, not a new proof of that input.

The inclusions $S_N\subset S_{N+1}\subset\mathbb N$ imply the
displayed monotonicity. Let $\zeta=\lim\eta_N\le\xi_2$.
If the sequence is eventually equal to $\zeta$, the union of
the retained constraints gives $\zeta\in\mathcal W_2$. Otherwise
$\eta_N<\zeta$ for every $N$. Fix $k$. If $\zeta^{2^k}=t$ is
an integer, then $t\ge2$ and, for every sufficiently large $N$,
\[
 \lfloor\eta_N^{2^{k+2}}\rfloor=t^4-1
 =(t^2-1)(t^2+1),
\]
which is composite, contradicting the retained constraint
$k+2\in S_N$. Hence $\zeta^{2^k}$ is nonintegral. Continuity
of the floor away from integers now implies
$\lfloor\zeta^{2^k}\rfloor$ is prime. This holds for every $k$,
so $\zeta\in\mathcal W_2$ and $\zeta=\xi_2$ by minimality.
\end{proof}

A limit of transcendental numbers need not be irrational, so this
approximation statement does not settle Question~1.6. A multiplier
three behaves differently: once it is inserted into a pure square
clock, later exponents have an odd factor and are not a subsequence
of $2,4,8,\ldots$. The set inclusions in this proof then do not
apply without additional hypotheses.

\subsection{What the Type D sign argument can and cannot give}

For an eventual $2/3$ multiplier word, write the tail as
$3^{\nu_1}2\,3^{\nu_2}2\cdots$, with infinitely many twos.
Under the additional integer-ratio hypotheses of Saito's
arXiv:2508.16068v3, Theorem~2.8, including ratios at least two
throughout, that theorem assumes
\begin{equation}\label{mil:sq-type-d}
 \limsup_{N\to\infty}
 \frac{\nu_{N+1}}{\sum_{j=1}^{N}\nu_j+(\log2/\log3)N}=\infty
\end{equation}
and concludes transcendence of the minimum. Item~2.9 is a remark
about independence of the chosen initial index, not a further
theorem. If $t_N$ is the position of the $N$th two, condition
\eqref{mil:sq-type-d} is equivalent to
\[
 \limsup_{N\to\infty}t_{N+1}/t_N=\infty.
\]
Indeed, $\nu_{N+1}=t_{N+1}-t_N-1$, while the denominator in
\eqref{mil:sq-type-d} is comparable to $t_N$, with fixed positive
comparison constants. This is a convenient reformulation of an
existing theorem, not a weakening of its hypothesis.

There is a version-specific discrepancy worth recording accurately:
the introductory Theorem~1.9(D) of arXiv v3 displays
$\lfloor\log k\rfloor$, whereas its abstract and detailed proof
use $\lfloor\sqrt{\log k}\rfloor$. The latter gives two-positions
$\lceil e^{j^2}\rceil$ (up to initial indexing), with unbounded
successive ratios; the former
gives $t_N\sim e^N$ and does not satisfy \eqref{mil:sq-type-d}.
Accordingly the argument in that version establishes the
square-root-logarithm example, and does not by itself establish
the logarithm example.

One cannot extend that proof to every bounded pattern using only
the sign of the leading conjugate pair. For example, an alternating
two/three word makes the exponents at successive relevant
three-steps differ by a factor of six. The sign requirement then
has the form
\[
 \{6^j x\}\in J=[1/4,3/4]\qquad(j\ge0).
\]
The two maps
\[
 f_2(x)=(x+2)/6,\qquad f_3(x)=(x+3)/6
\]
send $J$ into the disjoint subintervals
$[3/8,11/24]$ and $[13/24,5/8]$ of $J$. Their attractor is an
uncountable invariant set, of Hausdorff dimension $\log2/\log6$,
all of whose forward multiplication-by-six orbits remain in $J$.
In particular it contains irrational angles. Thus the sign
condition alone permits many nonperiodic orbits. This does not
claim that those angles are arguments of algebraic conjugates;
it identifies precisely the extra arithmetic information missing
from this attempted extension. No new unconditional resolution
for the alternating Type D clock is claimed.


\clearpage
\part{Polynomial patterns and digital structure}\label{part:patterns}
These results arise from related questions in the source literature. They share floor-function and recurrence techniques, but each theorem states its own hypotheses and uses no unproved assertion about Mills constants.

\section{Palindromic progressions and the base-two obstruction}

Let $P_b$ denote the positive integers whose standard base-$b$ expansion is
palindromic.  Question~27 of Pongsriiam and Saito, \emph{Palindromes and
Antipalindromes in Short Intervals}, Journal of Integer Sequences \textbf{26}
(2023), Article 23.9.5, p.~26, asks whether every finite coloring of $P_b$ has
a monochromatic nonconstant three-term arithmetic progression, and whether
the stronger assertion holds for every subset of positive upper relative
density.  The source is
\cite{PongsriiamSaito2023}.
Here an arithmetic progression is always required to have positive common
difference.

\begin{theorem}\label{sur:thm:palindrome-classification}
Both assertions in Question~27 hold if and only if $b\geq3$.
More precisely, for $b\geq3$, every $A\subseteq P_b$ with
\[
 \limsup_{N\to\infty}
 \frac{|A\cap[1,N]|}{|P_b\cap[1,N]|}>0
\]
contains a $b$-term arithmetic progression.  For $b=2$, there is an explicit
four-coloring of $P_2$ with no monochromatic three-term progression.
There is even a three-term-progression-free subset of $P_2$ having lower
relative density $1/7$.
\end{theorem}

\begin{proof}
First suppose $b\geq3$.  Write $P_{b,L}$ for the palindromes with exactly
$L$ digits.  Their number is
\[
 a_L=(b-1)b^{\lceil L/2\rceil-1}.
\]
Put $Q_L=\sum_{j\leq L}a_j$.  For $L\geq2$ we have
$a_L\leq b a_{L-1}\leq bQ_{L-1}$, and consequently
$Q_L\leq(b+1)Q_{L-1}$.
Let the upper relative density of $A$ be $\delta>0$.  There are arbitrarily
large $N$ with $|A\cap[1,N]|\geq(\delta/2)|P_b\cap[1,N]|$.
If $b^{L-1}\leq N<b^L$, extending to the end of the digit-length block gives
\[
 \frac{|A\cap[1,b^L-1]|}{Q_L}\geq\frac{\delta}{2(b+1)}.
\]
It follows, by summing exact-length block counts and discarding any fixed
finite initial segment, that arbitrarily large $L$ satisfy
\[
 \frac{|A\cap P_{b,L}|}{a_L}\geq\eta,
 \qquad \eta=\frac{\delta}{4(b+1)}>0.
\]
For each such $L$, fix a leading digit $a\in\{1,\ldots,b-1\}$ for which the
corresponding slice of $A$ has density at least $\eta$ among palindromes
with that leading digit.  For $L\geq2$ this slice is parametrized by the cube
$\{0,\ldots,b-1\}^{d}$, where $d=\lceil L/2\rceil-1$, through an injective
affine map
\[
 \Phi(v_1,\ldots,v_d)
   =a(b^{L-1}+1)+\sum_{j=1}^{d}w_jv_j,
 \qquad w_j>0.
\]
The paired digit weights are $b^i+b^{L-1-i}$; the central weight, when
present, is $b^{(L-1)/2}$.

The density Hales--Jewett theorem implies that, once $d$ is large enough,
the preimage of this slice of $A$ contains a combinatorial line.  Thus there
is a nonempty set of free coordinates $I$ such that replacing all coordinates
in $I$ by the same digit $t$, for $t=0,\ldots,b-1$, gives points of the
preimage.  Their images are
\[
 X,\ X+D,\ldots,X+(b-1)D,
 \qquad D=\sum_{j\in I}w_j>0.
\]
This proves the density assertion.  In any finite coloring, at least one
color has positive upper relative density, proving the coloring assertion.

For the binary case we first observe that palindromes of one fixed length
$L$ contain no three-term progression.  Indeed, suppose $x<y<z$ have that
length and $x+z=2y$.  Let $h$ be the highest binary position where $x$ and
$z$ differ.  Symmetry makes $j=L-1-h$ the lowest position where they differ,
so $v_2(z-x)=j$.  Since $y$ lies between $x$ and $z$, it shares their common
prefix above position $h$.  Since $y$ is palindromic, it also shares their
common suffix below position $j$.  Therefore $v_2(y-x)\geq j$, contrary to
$z-x=2(y-x)$.

Color each $n\in P_2$ by the pair $(L\bmod2,s)$, where $L$ is its digit
length and $s$ its second leading digit; for $n=1$ set $s=0$.
Suppose $x<y<z$ of one color satisfy $x+z=2y$, and let $L$ be the largest
digit length.  If only $z$ has length $L$, the parity condition gives
$y<2^{L-2}$, whence $z=2y-x<2^{L-1}$, a contradiction.
If $y,z$ have length $L$ but $x$ does not, then $x<2^{L-2}$.
Put $B=2^L$.  When $s=0$, both $y,z$ lie in $[B/2,3B/4)$, so
$x=2y-z>B/4$.  When $s=1$, they lie in $[3B/4,B)$, so $x>B/2$.
Both conclusions are impossible.  Finally, if all three lengths are equal,
the preceding observation applies.  Thus this four-coloring has no
monochromatic three-term progression.

For an explicit density counterexample, take the color class $C$ consisting
of even-length binary palindromes with leading digits $11$.
At the end of length $2m$,
\[
 |C\cap[1,2^{2m}-1]|=2^{m-1},\qquad
 |P_2\cap[1,2^{2m}-1]|=2^{m+1}-2.
\]
Before the next $C$-block begins, the denominator has increased by $2^m$
from length $2m+1$, and by $2^{m-1}$ from the leading-$10$ part of length
$2m+2$, while the numerator is unchanged.  Thus the successive minima and
maxima of the relative density tend respectively to
\[
 \frac{1}{7}\quad\hbox{and}\quad\frac{1}{4}.
\]
Within a $C$-block the ratio increases, and between such blocks it decreases.
This proves the claimed lower relative density and completes the proof.
\end{proof}

The combinatorial input is the density Hales--Jewett theorem of Furstenberg
and Katznelson; an elementary proof is D.~H.~J. Polymath, \emph{A new proof
of the density Hales--Jewett theorem}, Annals of Mathematics
\textbf{175} (2012), 1283--1327,
\cite{Polymath2012}.
This answers the literal Question~27 in every base. The distinction between
binary and larger bases is essential.


\section{Positive-density polynomial patterns in Piatetski--Shapiro sequences}

Saito and Yoshida, \emph{Distributions of finite sequences represented by
polynomials in Piatetski--Shapiro sequences}, arXiv:2006.13930v2~\cite{SaitoYoshidaPatterns2021},
Question~6.1, ask for the order of magnitude of polynomial patterns whose
indices are restricted to a positive-density subset of the integers.
The following argument gives an affirmative answer. Throughout, an
arithmetic progression has positive common difference, and a polynomial
pattern of degree at most $d$ means that all its successive differences
of order $d+1$ vanish.

\begin{lemma}[Local polynomial approximation and counting]\label{sur:lem:local-density}
Fix integers $d\geq0$, $k\geq d+2$, and real numbers $K\geq0$ and
$0<\eta\leq1$. There are $c>0$ and $L_0$ such that, for every integer
$L\geq L_0$, the following holds. Let $B\subseteq\{1,\ldots,L\}$ have
$|B|\geq\eta L$, let $a(1),\ldots,a(L)$ be integers, and suppose that a
real polynomial $P$ of degree at most $d$ satisfies
\[
 |a(t)-P(t)|\leq K\qquad(1\leq t\leq L).
\]
Then at least $cL^2$ arithmetic progressions of length $k$ contained in
$B$ have $a$-values forming a polynomial pattern of degree at most $d$.
The constants are independent of the coefficients of $P$.
\end{lemma}

\begin{proof}
Partition $[-K,K]$ into $Q=O_{d,K}(1)$ disjoint sets, each contained in
an interval of length $\varepsilon=2^{-d-2}$. Color $t\in B$ by the set
containing $e(t)=a(t)-P(t)$. A color class $B'$ has at least
$(\eta/Q)L$ elements. The usual supersaturation consequence of
Szemer\'edi's theorem, obtained by Varnavides averaging, gives at least
$c(k,\eta/Q)L^2$ arithmetic progressions of length $k$ in $B'$.

For each such progression $t_j=u+jv$, the numbers $e(t_j)$ belong to
one interval $[b,b+\varepsilon]$. Since differences of a constant
vanish, for $0\leq j\leq k-d-2$,
\[
 \left|\Delta^{d+1}a(t_j)\right|
 =\left|\Delta^{d+1}e(t_j)\right|
 \leq 2^{d+1}\varepsilon=\frac12.
\]
Here $\Delta$ denotes difference in the progression index $j$; the
polynomial contribution vanishes because $P(u+jv)$ has degree at most
$d$. The left-hand difference is an integer, hence it equals zero.
Newton interpolation now gives a polynomial of degree at most $d$
representing all the $a(t_j)$.
\end{proof}

\begin{theorem}[Affirmative answer to Question 6.1]\label{sur:thm:ps-density}
Let $d\geq1$, $d<\alpha<d+1$, and $k\geq d+2$. If
$A\subseteq\mathbb N$ has positive lower natural density, define
\[
 T_A(N)=\#\left\{(u,v)\in\mathbb N^2:
 \begin{array}{l}
 u,u+v,\ldots,u+(k-1)v\in A\cap[1,N],\\
 \bigl(\lfloor(u+jv)^\alpha\rfloor\bigr)_{j=0}^{k-1}
 \text{ has degree at most }d
 \end{array}\right\}.
\]
Then
\[
 T_A(N)\asymp_{\alpha,d,k,A}N^{\,2-\alpha/(d+1)}.
\]
The upper bound holds without a density assumption. If $A$ has only
positive upper natural density, the corresponding lower bound holds
along an unbounded sequence of values of $N$.
\end{theorem}

\begin{proof}
The upper bound is Saito--Yoshida's Theorem~1.3, since
$T_A(N)\leq T_{\mathbb N}(N)$. For the lower bound choose $\delta>0$
such that $|A\cap[1,N]|\geq\delta N$ for every sufficiently large $N$.
Discarding integers below $\delta N/2$ leaves at least $\delta N/2$
elements of $A$. Put
\[
 \beta=1-\frac{\alpha}{d+1}>0,
 \qquad L=\lfloor N^\beta\rfloor.
\]
Partition the remaining interval $[\delta N/2,N]$ into disjoint full
integer blocks of length $L$, discarding the last incomplete block.
Since $L=o(N)$, at least $\delta N/3$ elements of $A$ remain for all
large $N$. Call a block good if it contains at least $(\delta/6)L$
elements of $A$. Bad blocks contain in total at most $\delta N/6$
elements of $A$, so there are at least $\delta N/(6L)$ good blocks.

On a block beginning at $x\geq\delta N/2$, expand $(x+t)^\alpha$ to
degree $d$ at $t=0$, writing the resulting polynomial as $P_x(t)$.
For $0\leq t<L$, Taylor's theorem gives
\[
 |(x+t)^\alpha-P_x(t)|
 \leq C_{\alpha,d,\delta}
 N^{\alpha-d-1}L^{d+1}
 \leq C_{\alpha,d,\delta}.
\]
Thus the integer sequence $\lfloor(x+t)^\alpha\rfloor$ is uniformly
within $K=C_{\alpha,d,\delta}+1$ of $P_x$. By
Lemma~\ref{sur:lem:local-density}, each good block contributes at least
$cL^2$ desired progressions. Progressions counted in distinct blocks
are distinct. Consequently,
\[
 T_A(N)\geq\frac{\delta N}{6L}\,cL^2
 \gg NL\asymp N^{\,2-\alpha/(d+1)}.
\]
The same proof applies on any subsequence with
$|A\cap[1,N]|\geq\delta N$, proving the upper-density assertion.
\end{proof}

\begin{corollary}[Simultaneous finite families]\label{sur:cor:ps-simultaneous}
Fix finitely many real exponents $d_i<\alpha_i<d_i+1$, with integers
$d_i\geq1$, and let $k\geq2+\max_i d_i$. For a set $A$ of positive
lower natural density, there are at least
\[
 cN^{\,1+\min_i(1-\alpha_i/(d_i+1))}
\]
index progressions of length $k$ in $A\cap[1,N]$ on which every sequence
$\lfloor n^{\alpha_i}\rfloor$ forms a polynomial pattern of degree at
most $d_i$.
\end{corollary}

\begin{proof}
Use the common block length
$L=\lfloor N^{\min_i(1-\alpha_i/(d_i+1))}\rfloor$.
Each Taylor remainder is bounded independently of $N$. Color a point
by the tuple of residual bins, using bin width $2^{-d_i-2}$ in the
$i$th coordinate, and repeat the proof of
Lemma~\ref{sur:lem:local-density} and Theorem~\ref{sur:thm:ps-density}.
This gives a lower bound; no matching upper bound for the simultaneous
family is asserted.
\end{proof}

\subsection{An older theorem answers the prime-index condition}

In Saito--Yoshida, \emph{Arithmetic progressions in the graphs of
slightly curved sequences}, J. Integer Sequences \textbf{22} (2019),
Article~19.2.1, Question~12 (arXiv:1807.06971v3~\cite{SaitoYoshidaCurved2019}, Question~6.2), the
authors ask whether the primes satisfy their condition~(C).
This follows directly from an earlier result of Tao and Ziegler,
\emph{Narrow progressions in the primes}, arXiv:1409.1327v2~\cite{TaoZiegler2014},
Theorem~5. It should therefore be recorded as a literature implication,
not as a new theorem about prime progressions.

Here are the details. Fix an integer $k\geq3$ and a real $\gamma>0$.
Tao--Ziegler gives infinitely many progressions
$a,a+r,\ldots,a+(k-1)r$ consisting of primes, with
$r\ll(\log a)^{L_k}$. Set
\[
 b=1+\min(\gamma/4,1),\qquad h_\gamma(x)=x^b.
\]
The functions $h_\gamma$ and $h_\gamma'$ strictly increase, and
\[
 \frac{h_\gamma'(x+1)^2}{h_\gamma(x)^\gamma}
 \asymp x^{\,2b-2-b\gamma}\longrightarrow0.
\]
For a sequence of such prime progressions with $a\to\infty$, put
$x=(a-1)^{1/b}$. Then
\[
 h_\gamma(x+1)-h_\gamma(x)
 \sim b(a-1)^{1-1/b}
 \gg (\log a)^{L_k},
\]
so the whole progression lies in
$(h_\gamma(x),h_\gamma(x+1)]$. Passing to increasing starting points
makes $x\to\infty$, and all requirements of condition~(C) hold.
Saito--Yoshida's Proposition~9 (arXiv Proposition~5.1) then gives the
affirmative answer concerning the graphs at prime indices.

\section{Optimal three-term waiting times}

For $1<\alpha<2$ and an integer $r\geq1$, put $b_n=\lfloor n^\alpha\rfloor$ and define
\[
 L_{\alpha,3,r}(x)=\min\left\{y\geq0:
 \begin{array}{l}
 \text{there exists }n\in[x,x+y]\cap\mathbb Z\\
 \text{with }b_n+b_{n+2r}=2b_{n+r}
 \end{array}\right\}.
\]
In Appendix~A of arXiv:2006.13930v2, immediately before
Proposition~A.2, Saito and Yoshida conjecture the following upper bound.
Their Proposition~A.2 supplies a matching lower bound in the limsup
sense, so the exponent obtained here is sharp.

\begin{theorem}\label{sur:thm:ps-waiting}
For every $1<\alpha<2$ and every fixed integer $r\geq1$,
\[
 L_{\alpha,3,r}(x)=O_{\alpha,r}(x^{\,1-\alpha/2}).
\]
\end{theorem}

\begin{proof}
Put $f(t)=t^\alpha$, $N=\lceil x\rceil$, and
\[
 \lambda=f''(N)=\alpha(\alpha-1)N^{\alpha-2},\qquad
 M=\left\lceil\frac{32}{r\sqrt\lambda}\right\rceil.
\]
Let $n_j=N+rj$ and $b_j=\lfloor f(n_j)\rfloor$ for $0\leq j\leq M$.
Since $\lambda\to0$ and $rM=O(N^{1-\alpha/2})=o(N)$, throughout
$[N,N+rM]$ we have, for all sufficiently large $N$,
\[
 \frac{\lambda}{2}\leq f''(t)\leq\lambda.
\]
The real increments $A_j=f(n_{j+1})-f(n_j)$ increase strictly and
satisfy
\[
 0<A_{M-1}-A_0\leq(M-1)r^2\lambda<1
\]
for large $N$. Indeed, $Mr^2\lambda=O(r\sqrt\lambda+r^2\lambda)\to0$.
It follows that the nondecreasing sequence $\lfloor A_j\rfloor$ has
at most two constant runs. One run, say $j_0\leq j<j_0+J$, has
$J\geq M/2$.

Suppose that no desired three-term progression begins at any $n_j$
with $0\leq j\leq M-2$. The integer increments
$d_j=b_{j+1}-b_j$ then satisfy $d_j\ne d_{j+1}$. On the selected run,
$\lfloor A_j\rfloor=q$ is constant and
\[
 d_j=\lfloor A_j+\{f(n_j)\}\rfloor\in\{q,q+1\}.
\]
Thus these increments alternate, and their consecutive pair sums
are all $2q+1$. In particular,
\[
 b_{j_0+2v}=b_{j_0}+v(2q+1)\qquad(0\leq2v\leq J).
\]
Take $h=\lfloor J/4\rfloor$. For sufficiently large $N$, $J\geq8$
and $h\geq J/8$. Therefore the three values
$b_{j_0},b_{j_0+2h},b_{j_0+4h}$ form an arithmetic progression.
Set $t=n_{j_0}$ and $D=2rh$. Cancellation of their integer parts gives
\[
 |f(t+2D)-2f(t+D)+f(t)|<2.
\]
On the other hand, $D\geq rM/8$, and the integral formula for the
second difference yields
\begin{align*}
 f(t+2D)-2f(t+D)+f(t)
 &=\int_0^D\int_0^D f''(t+u+v)\,du\,dv\\
 &\geq\frac{\lambda}{2}D^2
 \geq\frac{\lambda r^2M^2}{128}\geq8,
\end{align*}
a contradiction. Hence a three-term progression with the original
fixed step $r$ starts at some $n_j$ with $j\leq M-2$, proving
\[
 L_{\alpha,3,r}(x)\leq1+rM
 \leq1+r+\frac{32}{\sqrt\lambda}
 =O_{\alpha,r}(x^{1-\alpha/2}).\qedhere
\]
\end{proof}


\section{Balanced equations in Piatetski--Shapiro sequences}

Write $\mathrm{PS}(\alpha)=\{\lfloor n^\alpha\rfloor:n\in\mathbb N\}$.
Saito's Theorem~2.2 in \emph{Finiteness of solutions to linear
Diophantine equations on Piatetski--Shapiro sequences}
(arXiv:2306.17813v2~\cite{SaitoBalanced}) gives a general exceptional-set dimension bound
$3r/(s-4)$ for equations $y=\sum_{i=1}^r a_i x_i$ with positive
coefficients. The following argument gives the stronger bound $(r+1)/s$
when the coefficients sum to one. In particular, the three-term
progression conclusion of his Theorem~1.3 extends from almost every
$\alpha>10$ to almost every $\alpha>3$.

\begin{theorem}\label{sur:thm:ps-balanced}
Let $r\ge2$ be an integer, and fix positive real numbers
$a_1,\ldots,a_r$ with $\sum_i a_i=1$. For $r+1<s<t$, let
$E_{\boldsymbol a}(s,t)$ be the set of $\alpha\in[s,t]$ for which
\[
 y=\sum_{i=1}^r a_i x_i
\]
has infinitely many nontrivial solutions in
$\mathrm{PS}(\alpha)^{r+1}$. A solution is nontrivial if not all its
coordinates coincide. Then
\[
 \dim_{\mathrm H}E_{\boldsymbol a}(s,t)\le\frac{r+1}{s}.
\]
Consequently, for almost every $\alpha>r+1$, the equation has only
finitely many nontrivial solutions in $\mathrm{PS}(\alpha)$.
\end{theorem}

\begin{proof}
All implied constants may depend on $s,t,r$ and the fixed coefficients.
Writing $y=\lfloor n^\alpha\rfloor$ and
$x_i=\lfloor n_i^\alpha\rfloor$, a solution necessarily satisfies
\begin{equation}\label{sur:eq:ps-balanced-error}
 \left|\sum_i a_i n_i^\alpha-n^\alpha\right|\le2.
\end{equation}
Indeed, the actual error is less than one, since the weights sum to
one. Strict increase of the sequence of power floors implies that
nontriviality is equivalent to $(n_1,\ldots,n_r)\ne(n,\ldots,n)$.
Moreover, \eqref{sur:eq:ps-balanced-error} implies $n_i\le Cn$ uniformly:
$a_i n_i^\alpha\le n^\alpha+2$. Thus there are only finitely many
solutions with bounded $n$.

Put $D=\max_i|n_i-n|$. For each fixed $n$, group index vectors into
dyadic shells $H\le D<2H$, where $H=1,2,4,\ldots$ and $H\le Cn$.
We first show that the number of vectors in a shell for which
\eqref{sur:eq:ps-balanced-error} holds at some $\alpha\in[s,t]$ is at most
\begin{equation}\label{sur:eq:ps-balanced-count}
 C H^{r-1}(H^2/n+1).
\end{equation}
For $H\ge\varepsilon n$, with fixed small $\varepsilon>0$, this follows
from the trivial $O(n^r)$ bound. If $H<\varepsilon n$, choose
$\varepsilon$ small enough that all $n_i$ belong to $[n/2,3n/2]$.
Set $d_i=n_i-n$. Uniform Taylor expansion gives
\[
 \sum_i a_i n_i^\alpha-n^\alpha
 =\alpha n^{\alpha-1}\sum_i a_i d_i
   +O(n^{\alpha-2}H^2).
\]
Consequently every admissible vector satisfies
\begin{equation}\label{sur:eq:ps-balanced-slab}
 \left|\sum_i a_i d_i\right|\le C H^2/n+C n^{1-s}.
\end{equation}
The right side is independent of $\alpha$. After the first $r-1$
deviations have been fixed in $O(H^{r-1})$ ways, the last deviation
lies in an interval containing $O(H^2/n+1)$ integers, since $a_r>0$
is fixed. This proves \eqref{sur:eq:ps-balanced-count}. No rationality
assumption on the coefficients is needed.

We next prove that the exponent set associated with one admissible
vector lies in an interval of length
\begin{equation}\label{sur:eq:ps-balanced-length}
 C n^{2-s}/H^2.
\end{equation}
Put $q_i=n_i/n$ and $F(z)=\sum_i a_i q_i^z-1$. Since the coefficients
sum to one, $F(0)=0$, and nontriviality gives $F''(z)>0$. Hence
\begin{equation}\label{sur:eq:ps-balanced-identity}
 zF'(z)-F(z)=\int_0^z uF''(u)\,du.
\end{equation}
For $0<q\le C$, define
\[
 J_s(q)=\int_0^s u(\log q)^2q^u\,du
       =1+(s\log q-1)q^s.
\]
The quotient $J_s(q)/(q-1)^2$ extends continuously and positively at
$q=0$ and $q=1$, with limits $1$ and $s^2/2$, respectively. It is
positive elsewhere by its integral representation. Compactness
therefore gives $J_s(q)\ge c(q-1)^2$ uniformly for $0<q\le C$.
Choosing an index for which $|q_i-1|=D/n$, we obtain from
\eqref{sur:eq:ps-balanced-identity}, for $s\le z\le t$,
\[
 zF'(z)-F(z)\ge cH^2/n^2.
\]
Let $\delta=2n^{-s}$ and
$B=\{z\in[s,t]:|F(z)|\le\delta\}$. This set contains the exponent
set in \eqref{sur:eq:ps-balanced-error}. On $B$,
\[
 F'(z)\ge (c/t)H^2/n^2-\delta/s
          \ge c' H^2/n^2
\]
for all sufficiently large $n$, since $H\ge1$ and $s>2$.
The set $B$ is an interval, possibly empty. Indeed, if $x<y$ lie in
$B$, convexity gives $F'(z)\ge F'(x)>0$ on $[x,y]$, so all
intermediate values remain in $[-\delta,\delta]$. The mean value
theorem now proves \eqref{sur:eq:ps-balanced-length}.

Fix $(r+1)/s<\sigma<1$. The sum of the $\sigma$-powers of the covering
interval lengths for a fixed $n$ is, by
\eqref{sur:eq:ps-balanced-count} and \eqref{sur:eq:ps-balanced-length}, at most
\begin{equation}\label{sur:eq:ps-balanced-sum}
 Cn^{(2-s)\sigma}
 \left(n^{-1}\sum_H H^{r+1-2\sigma}
                   +\sum_H H^{r-1-2\sigma}\right),
\end{equation}
where both sums are over dyadic $H\le Cn$. The first term is
$O(n^{r-s\sigma})$, whose sum over $n$ converges because
$s\sigma>r+1$.

For $r\ge3$, the second term is $O(n^{r-1-s\sigma})$, also summable.
For $r=2$ and $\sigma<1/2$, the same estimate applies. At
$\sigma=1/2$, it is $O(n^{(2-s)\sigma}\log n)$, and
$\sigma>3/s$ forces $s>6$, ensuring convergence. For $r=2$ and
$\sigma>1/2$, it is $O(n^{(2-s)\sigma})$, which is summable because
\[
 (s-2)\sigma>(s-2)\frac3s=3-\frac6s>1.
\]
Thus the full Hausdorff sum converges. The interval lengths tend
uniformly to zero, and their tails cover every exponent admitting
infinitely many nontrivial solutions. It follows that
$\dim_{\mathrm H}E_{\boldsymbol a}(s,t)\le\sigma$. Letting
$\sigma\downarrow(r+1)/s$ proves the dimension bound. A countable
compact-interval cover of $(r+1,\infty)$ gives the final assertion.
\end{proof}

\begin{corollary}\label{sur:cor:ps-ap-three}
For almost every $\alpha>3$, the set $\mathrm{PS}(\alpha)$ contains
only finitely many nonconstant three-term arithmetic progressions.
For $3<s<t$, the set of exceptions in $[s,t]$ has Hausdorff dimension
at most $3/s$.
\end{corollary}
\begin{proof}
Apply Theorem~\ref{sur:thm:ps-balanced} with $r=2$ and
$a_1=a_2=1/2$.
\end{proof}

The theorem establishes an upper bound, not the proposed exact
dimension. It does not prove infinitude for almost every
$2<\alpha<3$, and it does not settle whether any exponent
$\alpha>2$ admits arbitrarily long arithmetic progressions.
The identity $F(0)=0$ distinguishes the balanced case from the
remaining case of mixed positive coefficients with sum greater
than one. The argument supplies the stated upper bound without any rationality
assumption on the positive coefficients.


\section{Exact rational-power orders for algebraic tower limits}\label{sur:exp-section}
Let $h(x)$ denote the limit of the tower $x,x^x,x^{x^x},\ldots$ on
$[e^{-e},e^{1/e}]$, and put
\[
 \operatorname{ord}(A)=\min\{m\ge1:A^m\in\mathbb Q\},
\]
with value $\infty$ if the set is empty.
Kobayashi, Saito, and Takeda, \emph{Journal of Integer Sequences}
\textbf{26} (2023), Article~23.3.3~\cite{KobayashiSaitoTakeda2023}, Theorem~4, define
\[
 \mathcal E(k)=\left\{
 \left(\frac{kt}{s}\right)^{s/(kt)}:
 \begin{array}{l}
 1\le t\le\vartheta\log k,\quad kt/e\le s\le kte,\quad s,t\in\mathbb N,\\
 (kt)^{1/t},s^{1/t}\in\mathbb N,\quad\gcd(kt,s)=1
 \end{array}\right\},
 \qquad \vartheta=(\log2-e^{-1})^{-1}.
\]
They prove that, for an algebraic base $A$ of order $k\ge2$,
algebraicity of $h(A)$ implies $A\in\mathcal E(k)$.
Their Remark~6 asks whether $\mathcal E(k)$ itself is the exact-order-$k$
exceptional set. The literal set equality fails because some candidates
have smaller order. This observation does not contradict Theorem~4.

\begin{proposition}[A counterexample to the literal equality]\label{sur:exp-counterexample}
For $k=4$, the number
\[
 A=\left(\frac49\right)^{9/4}
  =\left(\frac23\right)^{9/2}
\]
belongs to $\mathcal E(4)$, satisfies $h(A)=4/9$, and has
$\operatorname{ord}(A)=2$.
\end{proposition}
\begin{proof}
Take $t=1$ and $s=9$ in the definition of $\mathcal E(4)$.
The inequalities $4/e<9<4e$ and $\gcd(4,9)=1$ verify every condition.
The square $A^2=(2/3)^9$ is rational, whereas $A$ is irrational:
the $2$-adic valuation of $A^2$ is the odd integer $9$, which a square
of a rational number cannot have. Thus $\operatorname{ord}(A)=2$.
The map $f(y)=y^{1/y}$ is strictly increasing on $[1/e,e]$,
and the tower identity implies $h(f(q))=q$ for every
$q\in[1/e,e]$; this is Lemma~14 of the cited paper.
Applying it with $q=4/9$ gives $h(A)=4/9$.
\end{proof}

The same base is already identified in the paper's order-two discussion.
Consequently the proposition is a modest correction to the literal
formulation of Remark~6, not a new transcendence mechanism.
The following valuation condition gives a finite explicit classification
for every order, including orders that are not square-free.

For relatively prime positive integers $u,v$ with $u\ne v$, let
\[
 G(u,v)=\gcd\{v_p(u),v_p(v):p\text{ prime, nonzero entries only}\}.
\]
Thus $G(u,v)$ is the largest integer $g\ge1$ such that both $u$ and $v$
are $g$th powers of positive integers. The convention discards zero
valuations; it allows $u=1$ or $v=1$.

\begin{lemma}[Exact order from valuations]\label{sur:exp-order-lemma}
If $a,b$ are relatively prime positive integers with $a\ne b$, then
\[
 \operatorname{ord}\!\left((a/b)^{b/a}\right)
  =\frac{a}{\gcd(a,G(a,b))}.
\]
\end{lemma}
\begin{proof}
The $m$th power is rational if and only if
$a\mid mb\,v_p(a/b)$ for every prime $p$.
Since $\gcd(a,b)=1$, these conditions are equivalent to
$a/\gcd(a,G(a,b))\mid m$.
Taking the least positive $m$ proves the formula.
\end{proof}

\begin{theorem}[All-order exceptional-set classification]\label{sur:exp-classification}
Fix an integer $k\ge2$. The set of algebraic numbers
$A\in[e^{-e},e^{1/e}]$ with $\operatorname{ord}(A)=k$ and algebraic
$h(A)$ is exactly
\[
 \left\{\left(\frac uv\right)^{v^t/k}:
 \begin{array}{l}
 t,u,v\in\mathbb N,\quad 1\le t\le\vartheta\log k,\quad u^t=kt,\\
 ue^{-1/t}\le v\le ue^{1/t},\quad\gcd(u,v)=1,\\
 \gcd(k,G(u,v))=1
 \end{array}\right\}.
\]
For each listed base, $h(A)=(u/v)^t$.
Equivalently, one obtains the exact exceptional set by retaining from
$\mathcal E(k)$ precisely the parameters satisfying
\[
 \gcd\bigl(kt,G(kt,s)\bigr)=t.
\]
\end{theorem}
\begin{proof}
First suppose that $A$ has order $k$ and algebraic tower limit.
The cited Theorem~4 gives parameters $t,s$ in $\mathcal E(k)$.
Write $u=(kt)^{1/t}$ and $v=s^{1/t}$, so that
$u^t=kt$, $\gcd(u,v)=1$, and
\[
 A=\left(\frac uv\right)^{v^t/k}.
\]
Because $k\mid u^t$, one has $\gcd(v,k)=1$. Prime valuations now give
\begin{equation}\label{sur:exp-order-candidate}
 \operatorname{ord}(A)=\frac{k}{\gcd(k,G(u,v))}.
\end{equation}
Indeed $A^m$ is rational exactly when
$k\mid mv^t v_p(u/v)$ for every prime $p$; the factor $v^t$ is
invertible modulo $k$. The asserted coprimality condition follows.
The bounds on $s=v^t$ give the stated bounds on $v$.

Conversely, suppose the displayed conditions hold and set
$q=(u/v)^t$. Then $q\in[1/e,e]$, and
\[
 q^{1/q}
  =\left(\frac uv\right)^{t v^t/u^t}
  =\left(\frac uv\right)^{v^t/k}=A.
\]
The identity $h(q^{1/q})=q$ gives an algebraic tower value.
Equation~\eqref{sur:exp-order-candidate} gives exact order $k$.
Finally $G(kt,s)=tG(u,v)$, so
$\gcd(kt,G(kt,s))=t\gcd(k,G(u,v))$.
This proves the equivalent formulation.
\end{proof}

The classification is an elementary refinement of the paper's finite
candidate set, using its Theorem~4 and Lemma~14. It does not require new
bounds for transcendental numbers or prime gaps. The refinement concerns exact rational-power order; it does not claim
a new transcendence criterion for tower values.


\section{A sparse zeta criterion for arbitrary real normal numbers}

Kanado and Saito, \emph{Normality of algebraic numbers and the Riemann zeta
function}, arXiv:2412.02337v2, Corollary~1.2, give a normality criterion for
positive algebraic irrational numbers using every integer $k\geq0$.
Their more general Theorem~2.2 permits the following strengthening.  This
is a corollary of their analytic theorem with an additional invariant-measure
argument, rather than a resolution of a numbered open problem or of the
normality of any specified algebraic irrational number.
The primary source is \cite{KanadoSaitoNormality}.

For $b\geq2$, $\alpha>0$, and an integer $k\geq1$, put
\[
 Z_{k,b,\alpha}(N)=\frac1{\log N}
 \sum_{1\leq|n|\leq N}
 \zeta\!\left(-k+\frac{2\pi i n}{\log b}\right)
 \frac{\exp(2\pi i n\log\alpha/\log b)}{n^{k+1}}.
\]

\begin{theorem}\label{sur:thm:sparse-normality}
Let $b\geq2$ be an integer, let $\alpha>0$ be any real number, and let
$R\subseteq\mathbb N$ be unbounded.  Then $\alpha$ is normal in base $b$
if and only if
\[
 \lim_{N\to\infty}Z_{2r-1,b,\alpha}(N)=0
 \qquad\text{for every }r\in R.
\]
In particular, algebraicity is unnecessary, the line $k=0$ may be omitted,
and any unbounded selection of the positive odd values of $k$ suffices.
\end{theorem}

\begin{proof}
Let $B_m$ be the Bernoulli polynomials.  For $m\geq2$ their periodic
extensions $B_m(\{x\})$ are continuous on the circle $\mathbb T=\mathbb R/\mathbb Z$.
Theorem~2.2 of Kanado--Saito, with
$d=1/\log b$ and $\theta=\log\alpha/\log b$, states that for every fixed
integer $k\geq1$,
\[
 \begin{aligned}
 &\frac{(k+1)!}{(2\pi i)^{k+1}}
 \sum_{1\leq|n|\leq N}
 \zeta\!\left(-k+\frac{2\pi i n}{\log b}\right)
 \frac{\exp(2\pi i n\log\alpha/\log b)}{n^{k+1}}
 \\
 &\qquad=-\frac1{(\log b)^k}
 \sum_{0<h<\log_bN}B_{k+1}(\{\alpha b^h\})
 +O_{k,b,\alpha}(\log\log N).
 \end{aligned}
\]
There is no algebraicity hypothesis in that theorem.
After division by $\log N$, changes of the endpoint of the $h$-sum cost
$o(1)$.  Conversely, every large integer cutoff $H$ can be reached by
taking $N$ near $b^{H+1/2}$.  Hence the asserted zeta limits are equivalent
to
\[
 \frac1H\sum_{h=1}^{H}B_{2r}(\{\alpha b^h\})\longrightarrow0
 \qquad(r\in R).
\]
Necessity follows immediately from normality, or equivalently from uniform
distribution of $\{\alpha b^h\}$.

For sufficiency, let $\mu$ be any weak subsequential limit of the empirical
probability measures
\[
 \mu_H=\frac1H\sum_{h=1}^{H}\delta_{\{\alpha b^h\}}
 \quad\text{on }\mathbb T.
\]
Such limits exist by compactness.  They are invariant under the continuous
circle map $T_b(x)=bx\pmod1$, because the difference between the averages
of $f\circ T_b$ and $f$ telescopes to two endpoint terms for every continuous
$f$.  Continuity of the even Bernoulli polynomials gives
\[
 \int B_{2r}(x)\,d\mu(x)=0\qquad(r\in R).
\]
The absolutely and uniformly convergent Fourier expansion
\[
 B_{2r}(\{x\})=
 \frac{(-1)^{r+1}2(2r)!}{(2\pi)^{2r}}
 \sum_{n=1}^{\infty}\frac{\cos(2\pi nx)}{n^{2r}}
\]
therefore yields
\[
 \sum_{n=1}^{\infty}\frac{a_n}{n^{2r}}=0,
 \qquad a_n=\int\cos(2\pi nx)\,d\mu(x),\quad |a_n|\leq1.
\]
If some $a_n$ were nonzero, choose the least such index $m$.
Multiplying by $m^{2r}$ and letting $r\to\infty$ through $R$ would give
$a_m=0$: the remaining terms tend to zero by dominated convergence,
being dominated for $r\geq1$ by the summable sequence $(m/n)^2$, $n>m$.
Thus every $a_n$ vanishes.

Let $\check\mu$ be the image of $\mu$ under reflection $x\mapsto-x$,
and let $\lambda$ be Lebesgue probability measure on $\mathbb T$.
The vanishing cosine coefficients and Fourier uniqueness imply
\[
 \mu+\check\mu=2\lambda.
\]
Both measures on the left are positive, so $\mu\leq2\lambda$; in particular
$\mu$ has an $L^1$ density.  The $T_b$-invariance gives, for every nonzero
integer $n$,
\[
 \widehat\mu(n)=\widehat\mu(bn)=\widehat\mu(b^j n)
   \longrightarrow0
\]
by the Riemann--Lebesgue lemma.  Thus $\mu=\lambda$.
Every subsequential weak limit is therefore $\lambda$, proving uniform
distribution and hence normality.
\end{proof}

\begin{corollary}\label{sur:cor:general-exponential-normality}
For any real $a>1$ and $\alpha>0$, the sequence $\{\alpha a^h\}$ is uniformly
distributed modulo one if and only if
$Z_{k,a,\alpha}(N)\to0$ for every integer $k\geq1$, with the same definition
using $a$ in place of $b$.  No Diophantine hypothesis on $a$ or $\alpha$
is necessary.
\end{corollary}

\begin{proof}
Use Kanado--Saito Theorem~2.2 as above.  The polynomials $1,B_m$ for $m\geq2$
span exactly the polynomials $p$ with $p(0)=p(1)$.  These are uniformly dense
in continuous functions on the circle: approximate a continuous periodic
function on $[0,1]$ by a polynomial $p$, then subtract
$x(p(1)-p(0))$ to match its endpoint values.  Thus vanishing of all the
Bernoulli averages with degree at least two is equivalent to uniform
distribution.
\end{proof}

The sparse-odd reduction uses the integer-base circle invariance and is
not asserted for noninteger $a$.  These criteria do not establish the
required zeta cancellation for $\sqrt2$, $(3/2)^h$, Mills' constant, or a
chosen prime-representing constant.  A full Hausdorff dimension statement
for a set of such constants likewise does not by itself imply that it
contains a normal number.


\section{Realization of asymptotic and zeta dimension profiles}
\label{sur:sec:fractal-profiles}

Question~1.4 of Saito's \emph{New fractal dimensions and some applications to arithmetic patches}, arXiv:1802.08839v4~\cite{SaitoFractalProfiles}, p.~3, asks whether there are relations between asymptotic dimension and zeta dimensions. Appendix~B of that paper already disproves either universal ordering between these invariants. We establish the following precise strengthening. It characterizes every possible finite numerical profile on integer lattices; it does not preclude relations under additional geometric or algebraic hypotheses.

\begin{theorem}[Realization of all lattice dimension profiles]
For every integer $d\geq1$, every $m\in\{0,\ldots,d\}$ and every pair $0\leq a\leq b\leq d$, there is an infinite set $X\subseteq\mathbb Z^d$ such that
\[
 \operatorname{asdim}X=m,\qquad
 \underline{\dim}_{\zeta}X=a,\qquad
 \overline{\dim}_{\zeta}X=b.
\]
Consequently, these are exactly the possible triples on infinite subsets of $\mathbb Z^d$.
\end{theorem}

We use the Euclidean metric, or equivalently the sup metric, throughout this section. For $Y\subseteq\mathbb Z^d$, the zeta dimensions are its lower and upper counting exponents:
\[
 \underline{\dim}_{\zeta}Y
 =\liminf_{R\to\infty}\frac{\log\#(Y\cap[-R,R]^d)}{\log R},
 \qquad
 \overline{\dim}_{\zeta}Y
 =\limsup_{R\to\infty}\frac{\log\#(Y\cap[-R,R]^d)}{\log R}.
\]
Indeed, at every fixed covering scale a set of bounded diameter contains a bounded number of lattice points. The resulting covering and counting functions differ by a multiplicative constant depending only on that scale and $d$.

\begin{lemma}[Sparse sequences with prescribed counting exponents]
For $0\leq\alpha\leq\beta\leq1$, there is a strictly increasing sequence of positive integers $(x_n)$ such that
\[
 x_{n+1}-x_n\longrightarrow\infty,\qquad
 \liminf_{n\to\infty}\frac{\log n}{\log x_n}=\alpha,
 \qquad
 \limsup_{n\to\infty}\frac{\log n}{\log x_n}=\beta,
 \qquad
 \frac{\log x_{n+1}}{\log x_n}\longrightarrow1.
\]
Its range $S$ has asymptotic dimension zero and lower and upper zeta dimensions $\alpha,\beta$.
\end{lemma}

\begin{proof}
If $\alpha=\beta=t>0$, take
\[
 x_n=\left\lceil n^{1/t}(\log n+1)^2\right\rceil.
\]
If $\alpha=\beta=0$, take
\[
 x_n=\left\lceil\exp\bigl((\log n)^2+\log n\bigr)\right\rceil.
\]
Both constructions have the stated properties after an irrelevant finite initial adjustment.

Suppose now that $\alpha<\beta$. For sufficiently large $s$, construct a nonnegative locally absolutely continuous function $g$ satisfying $0\leq g'(s)\leq s$ almost everywhere. Put
\[
 \phi(s)=s+2\log(s+1)+g(s),\qquad q(s)=\frac{s}{\phi(s)}.
\]
We alternate phases on which $g'=s$ and $g'=0$. Globally $g(s)\leq C+s^2/2$. During a phase of the first type,
\[
 \phi(s)-s\phi'(s)
 =2\log(s+1)-\frac{2s}{s+1}+g(s)-s^2<0
\]
for all sufficiently large $s$. Thus $q$ decreases, and tends to zero if that phase is continued indefinitely. During a phase of the second type the same numerator, with the term $-s^2$ omitted, is positive. Hence $q$ increases and tends to one if that phase is continued indefinitely.

Choose an integer $j_0$ large enough that
\[
 \alpha_j=\max(\alpha,1/j)
 <\beta_j=\min(\beta,1-1/j)
 \quad(j\geq j_0).
\]
Start with $g=0$ at a sufficiently large value of $s$, so that $q>\beta_{j_0}$. Run a decreasing phase until $q=\alpha_j$, then an increasing phase until $q=\beta_j$, and repeat. Each stopping point is finite. The switching times tend to infinity: otherwise $q$, which is Lipschitz on each fixed compact interval, would undergo infinitely many oscillations of amplitudes bounded below by a positive constant in a finite interval. It follows that
\[
 \liminf_{s\to\infty}q(s)=\alpha,
 \qquad \limsup_{s\to\infty}q(s)=\beta.
\]

Set $x_n=\lceil\exp(\phi(\log n))\rceil$ for all sufficiently large $n$. At almost every sufficiently large real $u$,
\[
 \frac{d}{du}\exp(\phi(\log u))
 =\frac{\exp(\phi(\log u))\phi'(\log u)}{u}
 \geq (\log u+1)^2.
\]
Integration on $[n,n+1]$ proves that the integer gaps tend to infinity. Moreover, $\phi'(s)\leq s+3$ almost everywhere, whence
\[
 \phi(\log(n+1))-\phi(\log n)
 =O\left(\frac{\log n}{n}\right)=o(1).
\]
Rounding changes the logarithms by $o(1)$. These estimates prove the adjacent-log ratio assertion. They also show that sampling $q(s)$ at $s=\log n$ preserves its lower and upper limits.

For $C_S(R)=\#(S\cap[1,R])$ and $x_n\leq R<x_{n+1}$, the count equals $n+O(1)$. Bounding its logarithmic exponent between the values with denominators $\log x_n$ and $\log x_{n+1}$ proves the asserted counting exponents. Finally, at each fixed scale $r$, the $r$-chain components of $S$ are singletons outside a finite initial subset. Their diameters are therefore uniformly bounded. Those components give a uniformly bounded, $r$-disjoint cover, proving $\operatorname{asdim}S=0$.
\end{proof}

\begin{proof}[Proof of the profile theorem]
Apply the lemma with $\alpha=a/d$ and $\beta=b/d$, and let $B=S^d\subseteq\mathbb Z^d$. In the sup metric,
\[
 \#\bigl(B\cap[0,R]^d\bigr)=C_S(R)^d,
\]
so the two zeta dimensions of $B$ are $a,b$. Products of the uniformly bounded $r$-chain-component partitions of $S$ form a uniformly bounded $r$-disjoint cover of $B$. Thus $\operatorname{asdim}B=0$.

For $m\geq1$, define
\[
 A_m=\bigcup_{n\geq1}
 \left(2^{2^n}e_1+\{0,\ldots,n\}^m\times\{0\}^{d-m}\right).
\]
Both zeta dimensions of $A_m$ are zero, since its counting function is
\[
 O\bigl((1+\log\log R)^{m+1}\bigr).
\]
By inclusion in an $m$-coordinate lattice, $\operatorname{asdim}A_m\leq m$. Conversely, $A_m$ contains isometric copies of arbitrarily large unit-mesh $m$-dimensional cubes. If its asymptotic dimension were at most $m-1$, then, for every fixed $r>0$, these cubes would admit covers by $m$ families of $r$-disjoint sets of a common diameter bound $D(r)$.

Translate and restrict to centered subcubes, and pass to a diagonal limit. To make this compactness step explicit, assign each lattice point a color in $\{1,\ldots,m\}$ and an equivalence class indicating its covering set. On each fixed finite cube there are finitely many such color-and-equivalence patterns. In the diagonal limit, equivalence, diameter at most $D(r)$, and separation at least $r$ of distinct classes of the same color all persist, because each is a finite constraint. The limit provides $m$ families of uniformly bounded $r$-disjoint sets covering $\mathbb Z^m$. This contradicts $\operatorname{asdim}\mathbb Z^m=m$. Consequently $\operatorname{asdim}A_m=m$. For $m=0$, use $A_0=\{0\}$.

Finally let $X=A_m\cup B$. Monotonicity and the finite-union theorem for asymptotic dimension give $\operatorname{asdim}X=m$. Writing $C_Y(R)=\#(Y\cap[-R,R]^d)$, we have
\[
 C_B(R)\leq C_X(R)\leq C_B(R)+C_{A_m}(R).
\]
Since $C_{A_m}(R)=R^{o(1)}$, adding it does not change either nonnegative logarithmic counting exponent of $B$. Hence the zeta dimensions of $X$ are $a,b$.

Every nonempty subset of $\mathbb Z^d$ has integer asymptotic dimension between zero and $d$, and its lower and upper zeta dimensions satisfy $0\leq a\leq b\leq d$. The construction attains every triple allowed by these intrinsic constraints.
\end{proof}

The asymptotic-dimension facts used above are standard; see Bell--Dranishnikov, \emph{Asymptotic dimension in B\k{e}dlewo}, arXiv:math/0507570~\cite{BellDranishnikov}, Proposition~5, Proposition~6 and Corollary~9. The realization theorem supplies a complete answer to the numerical-profile interpretation of Saito's question on lattice subsets. The formulation concerns arbitrary lattice subsets and does not rule out relations on more restricted classes such as homogeneous spaces or finitely generated groups.

\subsection{An existing-theorem consequence for avoidance dimensions}

Question~6.1 of Saito's \emph{New bounds for dimensions of a set uniformly avoiding multi-dimensional arithmetic progressions}, arXiv:1910.13071v1~\cite{SaitoAvoidance}, p.~14, asks for a specified upper bound on $D_S(k,\varepsilon)$. Its $k=3$ case follows from existing results, as follows. Write $T=\log(1/\varepsilon)$ and $M_1=\lceil1/\varepsilon\rceil$. Saito's equation~(2.1) gives
\[
 D_S(3,\varepsilon)\leq D_A(3,\varepsilon)
 \leq1+\frac{\log(r_3(M_1)/M_1)}{2\log M_1}.
\]
Kelley--Meka, \emph{Strong Bounds for 3-Progressions}, arXiv:2302.05537v6~\cite{KelleyMeka}, Theorems~1.1--1.2, imply $r_3(N)/N\leq\exp(-c(\log N)^{1/12})$ for a constant $c>0$ and large $N$. Therefore
\[
 D_A(3,\varepsilon)\leq1-c'T^{-11/12}.
\]
The deficit requested in Saito's question, with $M=\lceil1/(8\varepsilon)\rceil$, is
\[
 \frac{\log(\log M(\log\log M)^2)}{\log(1/\varepsilon)}
 =O\left(\frac{\log T}{T}\right).
\]
Since $T^{1/12}/\log T\to\infty$, the displayed bound is stronger for all sufficiently small $\varepsilon$. This is a deduction from established additive-combinatorial results, not a new proof of Roth's theorem or a resolution for every $k$.

\clearpage
\clearpage
\part{Growth constants, endpoint geometry, and finite prime covers}
\label{part:classifications}
\section{Polynomial recursions: radical rigidity and a coefficient obstruction}
\label{tw:sec:growth-new}

This section concerns an integer orbit
\begin{equation}\label{tw:eq:growth-orbit}
 x_{n+1}=P(x_n),\qquad x_n\in\mathbb Z,\qquad x_n\longrightarrow+\infty,
\end{equation}
where
\[
 P(X)=aX^d+bX^{d-1}+\cdots\in\mathbb Q[X],\qquad a>0,\quad d\ge2.
\]
Put
\[
 r=a^{1/(d-1)},\qquad A=r^{-1},\qquad s=\frac{b}{da},\qquad
 \alpha=\lim_{n\to\infty}x_n^{d^{-n}}.
\]
The known polynomial-iteration asymptotic is
\begin{equation}\label{tw:eq:growth-asymptotic}
 \alpha>1,\qquad
 x_n=A\alpha^{d^n}-s+O(\alpha^{-d^n}).
\end{equation}
See \cite[Section~2]{WagnerZiegler2021} and
\cite[equation~(1.2)]{Kumar2023}.
It follows by centering and scaling the recurrence to
$y_{n+1}=y_n^d+O(y_n^{d-2})$, then summing the logarithmic errors.

\paragraph{Status of the contribution.}
Dubickas already completely classified the algebraic possibilities under
$r\in\mathbb Q$ in \cite{Dubickas}.
We do not count that classification as new.  The results below are
explicit consequences and refinements of the established approximation
theorems, formulated for arbitrary positive rational leading coefficient.
Their proofs are given here; no claim of bibliographic priority is made.

\subsection{The approximation input}

We use the following two consequences of the theorem of
Kulkarni--Mavraki--Nguyen, \emph{Trans. Amer. Math. Soc.} 371 (2019),
3787--3804.  See \cite{KMN2019,CorvajaZannier2004}; both consequences follow from
\cite[Theorem~1.2]{Kumar2023};
the second is the fixed-coefficient argument underlying his Theorem 2.2.

\begin{lemma}[Two established approximation consequences]\label{tw:lem:growth-KMN}
\begin{enumerate}
\item Let $B>1$ be rational and $c\ne0$ algebraic.  If
$\|cB^m\|<\theta^m$ for infinitely many positive integers $m$, where
$0<\theta<1$ is fixed, then $B\in\mathbb Z$ and $c\in\mathbb Q$.
\item Let $\lambda>1$ be algebraic, let $q\ne0$ be an algebraic
integer, and let $\beta\in\mathbb Q\setminus\mathbb Z$.
For every $0<\theta<1$, only finitely many positive integers $m$ satisfy
$\|q\lambda^m+\beta\|<\theta^m$.
\end{enumerate}
Here $\|t\|$ denotes distance to the nearest integer.
\end{lemma}

\begin{proof}[Derivation from the stated approximation theorem]
For (1), apply Kumar's Theorem 1.2 to its singleton tuple
$\delta_1=B$, with fixed coefficient $c$.  Part (i) makes $B$ an
algebraic integer, hence an integer.  Part (iv), because every Galois
automorphism fixes $B$, gives $\sigma(cB^m)=cB^m$ for every $\sigma$
at all sufficiently large successful indices.  Consequently $c$ is
rational.  Constant coefficient height satisfies the required sublinear
height condition, for example with $f(m)=\sqrt m+1+h(c)$.

For (2), we spell out the argument from Theorem 1.2, so that
only the stated algebraic-integrality hypothesis on $q$ is needed.  To recall why
integrality of $q$ is decisive, apply Theorem 1.2 to the nondegenerate
pair $(\lambda,1)$ with coefficients $(q,\beta)$.  Part (i) makes
$\lambda$ an algebraic integer.  Part (iii) says that the sum of the
distinct conjugates of $q\lambda^m$, together with $\beta$, is an
integer at all sufficiently large successful indices.  That conjugate
sum is itself an integer because $q\lambda^m$ is an algebraic integer.
Thus $\beta$ would be integral, a contradiction.  There is no collision
between $\beta$ and a conjugate of $q\lambda^m$ at large indices:
an algebraic number with a rational conjugate is rational, and would
then equal $\beta$, whereas $q\lambda^m$ grows in absolute value.
\end{proof}

\subsection{Exact classification of rational-power growth constants}

\begin{theorem}[Radical rigidity]\label{tw:thm:growth-radical}
For an orbit \eqref{tw:eq:growth-orbit}, the following are equivalent:
\begin{enumerate}
\item $\alpha^h\in\mathbb Q$ for some positive integer $h$;
\item
\begin{equation}\label{tw:eq:growth-monomial}
 P(X)=a(X+s)^d-s,\qquad s=\frac{b}{da}.
\end{equation}
\end{enumerate}
In this case
\begin{equation}\label{tw:eq:growth-radical-value}
 \alpha=a^{1/(d-1)}|x_0+s|,\qquad
 \alpha^{d-1}=a|x_0+s|^{d-1}\in\mathbb Q.
\end{equation}
Furthermore, $\alpha$ is rational if and only if both
\eqref{tw:eq:growth-monomial} and $a^{1/(d-1)}\in\mathbb Q$ hold.
Whenever $\alpha$ is rational it is an integer.
\end{theorem}

\begin{proof}
Suppose $B=\alpha^h\in\mathbb Q$.  The sequence $d^n$ modulo $h$ is
eventually periodic.  Therefore some integers $n_0\ge0$, $\ell\ge1$
and $0\le e<h$ satisfy
\[
 d^{n_0+\ell j}=e+hm_j\qquad(j\ge0),
\]
where $m_j$ are eventually positive and tend to infinity.
Choose a positive integer $D_0$ such that $D_0s\in\mathbb Z$.
Equation \eqref{tw:eq:growth-asymptotic} gives
\[
 \left\|D_0A\alpha^e B^{m_j}\right\|=O(B^{-m_j}).
\]
Choose $B^{-1}<\theta<1$ and discard finitely many indices.
Lemma~\ref{tw:lem:growth-KMN}(1) shows that $B$ is an integer and
$A\alpha^e$ is rational.  Hence the numbers
$A\alpha^eB^{m_j}-s$ have bounded denominators.  Their distances
from the integers $x_{n_0+\ell j}$ tend to zero, so discreteness gives
\begin{equation}\label{tw:eq:growth-exact-subsequence}
 x_{n_0+\ell j}=A\alpha^{d^{n_0+\ell j}}-s
 \qquad\text{for all sufficiently large }j.
\end{equation}

Set $D=d^\ell$ and $R=P^{\circ\ell}$.  Comparing two consecutive
terms of \eqref{tw:eq:growth-exact-subsequence} gives
\[
 R(x_{n_0+\ell j})=A^{1-D}(x_{n_0+\ell j}+s)^D-s.
\]
These are infinitely many distinct arguments, since the orbit escapes.
Polynomial identity therefore gives
\[
 R(X)=A^{1-D}(X+s)^D-s.
\]
Normalize the original polynomial by the affine map $L(X)=r(X+s)$:
$Q=L\circ P\circ L^{-1}$.  Then $Q$ is monic of degree $d$ and
$Q^{\circ\ell}(Y)=Y^D$.  A polynomial commutes with its iterates,
so
\begin{equation}\label{tw:eq:growth-commute}
 Q(Y^D)=Q(Y)^D.
\end{equation}
If $Q\ne Y^d$, let $cY^t$ be its highest nonzero term below $Y^d$.
The highest nonleading term on the right of
\eqref{tw:eq:growth-commute} has degree $dD-d+t$ and coefficient $Dc$,
whereas on the left it has degree $tD$.
Since $dD-d+t>tD$, this is impossible.  Thus $Q=Y^d$, which is
exactly \eqref{tw:eq:growth-monomial}.

Conversely, under \eqref{tw:eq:growth-monomial}, put $y_n=r(x_n+s)$.
Then $y_{n+1}=y_n^d$ and $y_n=y_0^{d^n}$.
Positive escape forces $|y_0|>1$, and forces $y_0>0$ if $d$ is odd.
Taking the defining limit yields \eqref{tw:eq:growth-radical-value}.
Since $x_0+s\in\mathbb Q\setminus\{0\}$, this also proves the
criterion for rationality.

Finally, if $\alpha\in\mathbb Q$, apply
Lemma~\ref{tw:lem:growth-KMN}(1) directly to
$\|D_0A\alpha^{d^n}\|=O(\alpha^{-d^n})$.
It follows that $\alpha$ is an integer.
\end{proof}

\begin{corollary}[An irrationality test without root normalization]
If $a^{1/(d-1)}\notin\mathbb Q$, the growth constant of every
escaping integer orbit of $P$ is irrational.
If $P$ is not of the form \eqref{tw:eq:growth-monomial}, then no positive
integer power of its growth constant is rational.
\end{corollary}

The second statement excludes all rational radicals, but does not by
itself prove transcendence: a quadratic Pisot unit associated with a
Chebyshev recurrence is a familiar obstruction to such an inference.

\subsection{An explicit coefficient test for transcendence}

Write the positive leading coefficient as $a=u/v$ in lowest terms,
where $u,v\in\mathbb Z_{>0}$, and define
\begin{equation}\label{tw:eq:growth-M}
 M_d(a)=\prod_{p\mid u}p^{\lceil v_p(u)/(d-1)\rceil}.
\end{equation}
The empty product is $1$.

\begin{theorem}[Denominator obstruction]\label{tw:thm:growth-denominator}
If the growth constant $\alpha$ of \eqref{tw:eq:growth-orbit} is
algebraic, then
\begin{equation}\label{tw:eq:growth-shift-integral}
 M_d(a)\frac{b}{da}\in\mathbb Z.
\end{equation}
Consequently, failure of \eqref{tw:eq:growth-shift-integral} proves
that $\alpha$ is transcendental.
For $P\in\mathbb Z[X]$, the condition is
\begin{equation}\label{tw:eq:growth-divisibility}
 da\mid M_d(a)b.
\end{equation}
In particular, $d\nmid b$ is sufficient for transcendence in this
integer-coefficient case, although \eqref{tw:eq:growth-divisibility}
can be strictly stronger.
\end{theorem}

\begin{proof}
Put $M=M_d(a)$.  The number $q=MA$ is an algebraic integer, because
\[
 q^{d-1}=\frac{M^{d-1}v}{u}\in\mathbb Z.
\]
Multiplying \eqref{tw:eq:growth-asymptotic} by $M$ gives
\[
 Mx_n=q\alpha^{d^n}-Ms+O(\alpha^{-d^n}).
\]
If $\alpha$ is algebraic and $Ms$ is a nonintegral rational number,
Lemma~\ref{tw:lem:growth-KMN}(2), with $\beta=-Ms$ and exponents
$d^n$, contradicts this approximation.  This proves
\eqref{tw:eq:growth-shift-integral}.

If $a\in\mathbb Z_{>0}$, then $M_d(a)\mid a$.
Condition \eqref{tw:eq:growth-divisibility} therefore implies
$d(a/M_d(a))\mid b$, in particular $d\mid b$.
\end{proof}

\paragraph{Optimality within this scalar-integrality argument.}
For a positive integer $T$, the coefficient $TA$ is an algebraic
integer if and only if
\[
 (TA)^{d-1}=T^{d-1}v/u\in\mathbb Z,
\]
equivalently $M_d(a)\mid T$.  Necessity uses that a rational algebraic
integer is integral, and sufficiency uses the displayed monic equation.
Thus \eqref{tw:eq:growth-shift-integral} is the strongest condition
obtainable by this method of multiplying the asymptotic by an integer
and applying Lemma~\ref{tw:lem:growth-KMN}(2).  This is an optimality
statement about this particular argument, not about all possible
transcendence criteria.

\begin{corollary}[A family outside rational leading normalization]
Let $q$ be a positive integer that is not a square and let
\[
 x_0\ge1,\qquad x_{n+1}=q^3x_n^3+3x_n^2+1.
\]
Then $\lim x_n^{3^{-n}}$ is transcendental.
More generally, if $a\ge1$ is an integer, $a/M_3(a)>1$, and
$b\in\mathbb Z$ is not divisible by $3a/M_3(a)$, every escaping
integer orbit of $aX^3+bX^2+cX+e$ has transcendental growth constant.
\end{corollary}

\begin{proof}
For the displayed family, choose a prime $p$ with $v_p(q)$ odd.
Then
\[
 v_p(q^3/M_3(q^3))
 =3v_p(q)-\lceil3v_p(q)/2\rceil
 =\lfloor3v_p(q)/2\rfloor\ge1.
\]
In fact $q^3/M_3(q^3)>1$, so $3q^3/M_3(q^3)$ does not divide
the coefficient $3$.  The denominator obstruction applies.
The integer orbit escapes because its recurrence is strictly larger
than $x_n$ for every $x_n\ge1$.  Its leading normalization
$(q^3)^{1/2}$ is irrational, so Dubickas's rational-normalization
hypothesis does not cover this example.  The more general assertion
is \eqref{tw:eq:growth-divisibility}.
\end{proof}

For instance, the recurrence
\[
 x_0=1,\qquad x_{n+1}=8x_n^3+3x_n^2+1
\]
starts with $1,12,14257,\ldots$ and has a transcendental growth
constant.  Here $M_3(8)=4$ and $M_3(8)b/(3\cdot8)=1/2$;
the weaker condition $3\nmid b$ would miss it.

\subsection{Periodic polynomial schedules}

\begin{corollary}[A periodic transfer]\label{tw:cor:growth-periodic}
Suppose $P_0,\ldots,P_{t-1}\in\mathbb Q[X]$ have degrees
$d_i\ge2$ and positive leading coefficients.  Let an integer sequence
satisfy
\[
 x_{n+1}=P_{n\bmod t}(x_n),\qquad x_n\longrightarrow+\infty.
\]
Put $F=P_{t-1}\circ\cdots\circ P_0$, $D=\prod_i d_i$, and
$\gamma=\lim_{n\to\infty}x_{tn}^{D^{-n}}$.
Then Theorems~\ref{tw:thm:growth-radical} and
\ref{tw:thm:growth-denominator} apply to $\gamma$ with polynomial $F$.
In particular, if the top two coefficients $F(X)=a_FX^D+b_FX^{D-1}
+\cdots$ fail
$M_D(a_F)b_F/(Da_F)\in\mathbb Z$, then $\gamma$ is transcendental.
\end{corollary}

\begin{proof}
The subsequence $z_n=x_{tn}$ is an escaping integer orbit of the
single polynomial $F$.  Apply the preceding results.
\end{proof}

For an explicit alternating example, take $P_0(X)=2X^2+X$ and
$P_1(X)=X^2+1$, with $x_0\ge1$.  Their two-step polynomial is
\[
 F(X)=4X^4+4X^3+X^2+1.
\]
Here $D=4$, $a_F=4$, $b_F=4$, $M_4(4)=2$, and
$M_4(4)b_F/(4a_F)=1/2$.  Hence the two-step growth constant
$\gamma=\lim x_{2n}^{4^{-n}}$ is transcendental.  This transfer
is a direct composition corollary, not a separate deep theorem.

\paragraph{Scope.}
These results classify rational radicals exactly and certify
transcendence in explicit coefficient families.  They neither classify
all algebraic exceptions for arbitrary leading coefficient nor settle
the transcendence of the ordinary Mills constant.  A polynomial orbit
is not automatically an orbit of consecutive least primes.


\section{Exact modular return and persistent failure for nonlinear factorial clocks}
\label{tw:sec:factorial-nonlinear-return}

\begin{lemma}[The Type-B input and its local form]
\label{tw:lem:factorial-local-input}
Let $(C_n)$ be positive integers with $C_{n+1}/C_n\ge2$ and
$C_{n+1}/C_n\to\infty$.  Suppose that every $C_m$ divides
arbitrarily late terms $C_n$.  Then
\[
 W(C)=\{A>1:\lfloor A^{C_n}\rfloor\text{ is prime for every }n\ge1\}
\]
has a minimum, and every left-gap point of $W(C)$ is transcendental.
Here a left-gap point is an $A\in W(C)$ for which some interval
$(A-\eta,A)$ is disjoint from $W(C)$.
\end{lemma}
\begin{proof}
Existence follows from Saito's Theorem~1.3 \cite{SaitoVariants}.
For the arithmetic assertion, apply the already proved local Type-B
Theorem~\ref{mil:thm:local-type-b}. The ratio limit is infinite, its
index set $I$ is cofinite, and the assumed arbitrarily late divisibility
return supplies its condition (B5). A left-gap point is precisely a local
left minimum in that theorem's terminology. Its unbounded-ratio conclusion
therefore proves transcendence.
\end{proof}

The polynomial-exponential argument in the preceding work established
eventual divisibility return, but did not establish universal return for
nonlinear polynomial perturbations.  Here we settle the corresponding
return question exactly for a different class: a modularly vanishing
sequence plus an integer polynomial.  This yields full, unshifted
transcendence examples for which universal return fails at arbitrarily
late seeds.  The distinction between the two return conditions is therefore
structural, rather than merely an early-index obstruction.

We say that an integer sequence $H(n)$ is \emph{modularly vanishing} if,
for each positive integer $q$, one has $q\mid H(n)$ for all sufficiently
large $n$.  We call $F\in\mathbb Z[X]$ \emph{intersective} when it has an
integer root modulo every positive integer.  Put
\[
 C_n=H(n)+Q(n),\qquad Q\in\mathbb Z[X],
\]
and assume that the $C_n$ are positive and strictly increasing.  At a seed
$m$, write $\mathcal R_m$ for the universal-return condition
\[
 \text{for every }L\ge1\text{ there is }n>m
 \text{ with }C_n\equiv C_m\pmod{LC_m}.
\]
For schedules with ratios tending to infinity, the index set $I$ in
Saito's Type-B condition is cofinite; failures of $\mathcal R_m$ at
arbitrarily late seeds therefore also refute eventual Type-B condition
(B6).

\begin{theorem}[An exact polynomial-image criterion]
\label{tw:thm:exact-factorial-return}
Under the preceding assumptions, the following conditions are equivalent
for each fixed seed $m$:
\begin{enumerate}
\item $\mathcal R_m$ holds.
\item For every $q\ge1$, there are arbitrarily large $n$ such that
      $C_n\equiv C_m\pmod q$.
\item The polynomial $Q(X)-C_m$ is intersective.
\end{enumerate}
For every modulus $q$, the set of return indices has natural density
\[
 \frac{\rho_Q(C_m;q)}q,\qquad
 \rho_Q(c;q)=\#\{x\in\mathbb Z/q\mathbb Z:Q(x)\equiv c\pmod q\}.
\]
\end{theorem}
\begin{proof}
Suppose (1) holds, and fix $q$ and $N>m$.  Choose a multiple $L$ of $q$
so large that
\[
 LC_m>\max_{m<n\le N}(C_n-C_m).
\]
Strict increase ensures that none of the finitely many indices
$m<n\le N$ can return modulo $LC_m$.  Condition (1) supplies an index
$n>N$ returning modulo $LC_m$, hence also modulo $q$.  This proves (2).
Condition (2) immediately implies (1), by taking $q=LC_m$.

For a fixed $q$, modular vanishing gives
$C_n\equiv Q(n)\pmod q$ for every sufficiently large $n$.
The polynomial residue sequence $Q(n)\bmod q$ is periodic with period
$q$.  Consequently a return occurs at arbitrarily large indices exactly
when $Q(X)-C_m$ has a root modulo $q$.  This proves the equivalence with
(3), and counting residue classes gives the asserted density.
\end{proof}

\begin{corollary}[Power perturbations: exact local criterion]
\label{tw:cor:monomial-return-criterion}
If $Q(X)=X^r$ with $r\ge1$, then $\mathcal R_m$ holds if and only if
$C_m$ is a perfect $r$th power of a positive integer.
\end{corollary}
\begin{proof}
If $C_m=b^r$, the residue class $n\equiv b\pmod q$ gives a root for
every $q$.  Conversely, if $C_m$ is not an $r$th power, some prime $p$
satisfies $r\nmid v_p(C_m)$.  Writing $e=v_p(C_m)$, the congruence
$x^r\equiv C_m\pmod{p^{e+1}}$ would force
$rv_p(x)=e$, which is impossible.  Apply
Theorem~\ref{tw:thm:exact-factorial-return}.
\end{proof}

For completeness we establish an all-index growth inequality; this
avoids silently deleting an initial segment of the following family.

\begin{lemma}[A uniform factorial--power doubling inequality]
\label{tw:lem:factorial-power-doubling}
For every pair of integers $n,r\ge1$,
\[
 (n+1)!+(n+1)^r\ge 2(n!+n^r).
\]
\end{lemma}
\begin{proof}
It suffices to prove
$2n^r-(n+1)^r\le(n-1)n!$.  At $n=1$ this is immediate.
For $n=2,3,4$ the left side is nonpositive once $r\ge2,3,4$,
respectively; for the remaining exponents its respective maxima are
$1,2,7$, which are smaller than $(n-1)n!$.

Suppose $n\ge5$ and the left side is positive.  Set
$u=\log n$, $t=\log(1+1/n)$, and $r_0=\log2/t$, so that
$x=r_0-r>0$.  The inequalities
$1-e^{-tx}\le tx$, $xe^{-ux}\le1/(eu)$,
$t<1/n$, and $r_0<(n+1)\log2<\tfrac7{10}(n+1)$ give
\[
 2n^r-(n+1)^r
 =2n^{r_0}e^{-ux}(1-e^{-tx})
 <\frac{2n^{7(n+1)/10}}{en\log n}.
\]
Also $n!\ge(n/e)^n$.  The logarithm of the ratio between
$(n-1)(n/e)^n$ and the last upper bound is
\[
 F(n)=\tfrac3{10}(n+1)\log n-n+\log(n-1)
      +\log\log n+1-\log2.
\]
This is positive for every real $n\ge5$.  Indeed,
$F(5)=\tfrac95\log5-4+\log2+\log\log5>0.04$, and
\[
 F'(x)=\tfrac3{10}\log x-\tfrac7{10}
       +\frac3{10x}+\frac1{x-1}+\frac1{x\log x}>0\quad(x\ge5).
\]
For the last inequality one may split into $[5,8]$, $[8,11]$,
and $[11,\infty)$: on the first interval use
$\log5>1.6$ and $8\log8<17$; on the second use
$\tfrac3{10}\log8>0.62$ and $1/(x-1)\ge0.1$;
on the third the first two terms are already positive.
Thus $(n-1)n!$ exceeds the required upper bound.
\end{proof}

\begin{theorem}[Full transcendence despite infinitely many return failures]
\label{tw:thm:factorial-power-full}
Let $s,r\ge1$ be integers and put
\[
 C_n=(n!)^s+n^r\qquad(n\ge1).
\]
The least prime-representing constant for the full schedule exists and
is transcendental.  More generally, every left-gap point of its
prime-representing set is transcendental.

If $r>s$, universal return fails at every prime seed $m=p$.
More precisely,
\[
 C_n\not\equiv C_p\pmod{p^{s+1}}\qquad(n>p),
\]
so Saito's condition (B6) fails with $m=p$ and $L=p$.  These failures
occur arbitrarily late, while the divisibility condition (B5) holds
at every seed.
\end{theorem}
\begin{proof}
The preceding lemma gives $C_{n+1}\ge2C_n$ also for general $s$,
because
\[
 \big((n+1)^s-2\big)(n!)^s\ge(n-1)n!
\]
for $n\ge2$ and $s\ge1$; the case $n=1$ is direct.
Moreover $C_{n+1}/C_n\sim(n+1)^s\to\infty$.
For any seed $m$, take an arbitrarily large multiple $n$ of $C_m$.
Then $C_m\mid n^r$ and, once $n\ge C_m$, also $C_m\mid n!$.
Thus $C_m\mid C_n$, proving (B5).  The other hypotheses of
Lemma~\ref{tw:lem:factorial-local-input} follow from positivity,
integrality, the doubling inequality, and the unbounded ratios.
That lemma proves existence and both transcendence assertions.

Now suppose $r>s$ and let $p$ be prime.  Since $v_p(p!)=1$,
\[
 v_p(C_p)=v_p\big((p!)^s+p^r\big)=s.
\]
For $n>p$ with $p\nmid n$, the factorial term is divisible by $p$
whereas $n^r$ is not, so $v_p(C_n)=0$.
If instead $p\mid n$, then $n\ge2p$ and $v_p(n!)\ge2$.
Both terms defining $C_n$ are then divisible by
$p^{\min(2s,r)}$, and $\min(2s,r)>s$.  Neither case can give
$C_n\equiv C_p\pmod{p^{s+1}}$.
Finally $p^{s+1}\mid pC_p$, so return modulo $pC_p$ is impossible
as well.  Infinitely many prime seeds prove the last assertion.
\end{proof}

\begin{remark}[Scope and novelty]
The transcendence conclusion is an explicit full-schedule consequence
of Saito's existing Type-B criterion; it is not a new general
transcendence method.  The additional results are the exact
polynomial-image return criterion and the persistent separation
between (B5) and (B6), including an infinite family with completely
explicit obstructing seeds and moduli.  In particular
$C_n=n!+n^2$ has a transcendental minimum although (B6) fails at every
prime seed.  These arguments do not settle eventual universal return
for $P(a^n)+Q(n)$ in general, nor do they prove full-schedule
transcendence for $3^n+n^2$.
\end{remark}


\section{Ceiling prime constants: endpoints and right extrema}
\label{tw:sec:ceiling-new}

For an integer $c\geq3$, put
\[
 \mathcal V_c=\{B>1:\lceil B^{c^k}\rceil\text{ is prime for every }k\geq1\}.
\]
T\'oth proves that this set is nonempty and constructs constants by
descending from a prime $p$ to a prime below $p^c$; see
\emph{A Variation on Mills-Like Prime-Representing Functions},
J. Integer Sequences \textbf{20} (2017), Article 17.9.8 \cite{Toth2017}.
The following results address an endpoint distinction that prevents
automatic transfer of results about the smallest floor constant.
The arithmetic inputs used below are the prime-gap theorem with exponent
$\theta=21/40$, the Corvaja--Zannier--Dubickas approximation criterion,
and the Pisot conjugate-sum bound, in the precise forms recalled in
Saito, \emph{Mills' constant is irrational}, arXiv:2404.19461v2,
Theorems 2.1 and 2.6 and Lemma 2.7 \cite{SaitoMills}.
These are deductions from existing machinery; no literature priority
claim is made for the ceiling formulation.

\subsection{The infimum need not be a prime-representing constant}

\begin{proposition}[Endpoint alternatives]\label{tw:prop:ceil-endpoint}
Let $L=\inf\mathcal V_c$. Then the following statements hold.
\begin{enumerate}
\item If $L>1$ and $L^{c^k}$ is nonintegral for every $k\geq1$,
then $L\in\mathcal V_c$.
\item If $L>1$ and $L\notin\mathcal V_c$, there are $k\geq1$ and
an even integer $t\geq2$ such that $L^{c^k}=t$ and
\[
 t^{c^j}+1\text{ is prime for every }j\geq0.
\]
In particular, this alternative is impossible unless $c$ is a power of $2$.
\item If $\mathcal V_c\cap(1,2]\ne\varnothing$, then
$L=1$ and $\mathcal V_c$ has no smallest element.
\end{enumerate}
\end{proposition}

\begin{proof}
Choose $B_i\in\mathcal V_c$ tending to $L$ from above (or equal to $L$
when it already belongs to the set). At a noninteger value $L^{c^k}$,
the ceiling is locally constant, proving the first assertion.
An element $B\in\mathcal V_c$ cannot have an integer power $B^{c^k}$:
that integer would be prime, whereas its $c$th power at the next index
would be composite. Consequently nonattainment at $L>1$ requires an
integer boundary $t=L^{c^k}$. Since $B_i>L$, right limits of ceilings
give $t^{c^j}+1$ prime for every fixed $j\geq0$. Already $t+1$ prime
forces $t$ even. Write $c=2^a u$ with $u$ odd. If $u>1$, then
$(t^{2^a})^u+1$ has the proper factor $t^{2^a}+1$, contradicting the
condition at $j=1$. This proves the second assertion.

For the last assertion, take $1<B\leq2$ in $\mathcal V_c$.
Then $B^{1/c}\in\mathcal V_c$: its first output is $\lceil B\rceil=2$,
and all subsequent outputs are those of $B$ shifted by one index.
Iteration produces $B^{c^{-r}}\downarrow1$. The number $1$ itself is
excluded and does not have prime outputs.
\end{proof}

\begin{corollary}[Unconditional absence of a minimum for infinitely many bases]
For every sufficiently large integer $m$, the nonempty set
$\mathcal V_{3^m}$ has infimum $1$ and no smallest element.
\end{corollary}

\begin{proof}
Choose any $B\in\mathcal V_3$, whose existence is T\'oth's theorem.
For each $m\geq1$, put $A_m=B^{3^{1-m}}$. Then
\[
 A_m^{(3^m)^k}=B^{3^{m(k-1)+1}}
\]
has prime ceiling for every $k\geq1$. For sufficiently large $m$,
$1<A_m\leq2$, so Proposition~\ref{tw:prop:ceil-endpoint} applies.
\end{proof}

\begin{remark}
The criterion involving $t^{c^j}+1$ is necessary, not sufficient, for a
missing endpoint; no assertion about an infinite family of generalized
Fermat primes is used. Similarly, the displayed numerical branch beginning
with $2,7,337$ in T\'oth's paper is not being used here to certify
unconditionally that $\mathcal V_3\cap(1,2]$ is nonempty. That would
require an infinite extension argument for that particular finite prefix.
\end{remark}

\subsection{Rightmost completions and their arithmetic}

A point $B\in\mathcal V_c$ is a \emph{local right maximum} if
$(B,B+\varepsilon)\cap\mathcal V_c=\varnothing$ for some $\varepsilon>0$.

\begin{lemma}[Attainment and local extension]\label{tw:lem:ceil-right}
Every nonempty bounded subset of $\mathcal V_c$ has its supremum in
$\mathcal V_c$. Every feasible finite prime prefix has a greatest
completion, which is a local right maximum. For every sufficiently large
prime $p$, the one-term prefix $\lceil B^c\rceil=p$ is feasible.
\end{lemma}

\begin{proof}
If $S\subset\mathcal V_c$ is nonempty and bounded, choose elements
increasing to $b=\sup S>1$. Left continuity of the ceiling gives
$\lceil b^{c^k}\rceil$ prime for each fixed $k$. Thus $b\in\mathcal V_c$.
The same proof preserves any prescribed finite prefix. Its completion
set is bounded, and its greatest member has no integer power by the
argument in Proposition~\ref{tw:prop:ceil-endpoint}. All finitely many
prefix conditions therefore remain valid in a neighborhood of that
member, proving local right maximality.

The prime-gap theorem implies, with $\theta=21/40$, that the largest
prime $q<p^c$ satisfies
\[
 p^c-Kp^{c\theta}\leq q<p^c
\]
for all sufficiently large primes $p$, with a fixed constant $K$.
Since $c\theta<c-1$ for $c\geq3$, this gives
$(p-1)^c+1<q<p^c$, and also $q>p$. Iteration yields primes $p_n$ with
\[
 (p_n-1)^c+1<p_{n+1}<p_n^c.
\]
The intervals $((p_n-1)^{1/c^n},p_n^{1/c^n}]$ are nested, their
left endpoints strictly increase, and their lengths tend to zero.
Their unique limiting point lies strictly above every left endpoint
and at or below every right endpoint, and therefore has the required
prime ceilings.
\end{proof}

\begin{corollary}[No global minimum for all sufficiently large bases]
There is an integer $c_0$ such that $\inf\mathcal V_c=1$ and
$\mathcal V_c$ has no smallest element for every integer $c\geq c_0$.
\end{corollary}

\begin{proof}
The large-prime threshold in the extension part of
Lemma~\ref{tw:lem:ceil-right} can be chosen uniformly for all integer
$c\geq3$. Indeed the backward prime-gap bound has constants independent
of $c$, and
\[
 p^c-(p-1)^c\geq p^{c-1},\qquad
 \frac{p^{c-1}}{p^{c\theta}}=p^{c(1-\theta)-1}
 \geq p^{17/40}.
\]
Thus one fixed sufficiently large prime $P$ can serve as the first seed
for every $c\geq3$. The resulting constant satisfies $1<B\leq P^{1/c}$.
For $c\geq\max\{3,\lceil\log_2 P\rceil\}$ this lies in $(1,2]$,
and Proposition~\ref{tw:prop:ceil-endpoint} applies.
\end{proof}

\begin{theorem}[Arithmetic of ceiling right maxima]\label{tw:thm:ceil-arith}
Let $c\geq3$ be an integer and let $B$ be a local right maximum of
$\mathcal V_c$. Then $B$ is irrational. More precisely:
\begin{enumerate}
\item If $c\geq5$, then $B$ is transcendental.
\item If $c=4$, then either $B$ is transcendental or $B^{4^m}$ is a
quadratic Pisot number for some $m\geq1$.
\item If $c=3$, then either $B$ is transcendental or $B^{3^m}$ is a
cubic Pisot number for some $m\geq1$.
\end{enumerate}
In particular, these conclusions apply to the greatest completion of
every feasible finite prefix, and such completions exist unconditionally
for every sufficiently large first prime.
\end{theorem}

\begin{proof}
Put $p_k=\lceil B^{c^k}\rceil$. We first prove that eventually
\begin{equation}\label{tw:eq:ceil-greedy}
 p_{k+1}=\max\{q<p_k^c:q\text{ prime}\}.
\end{equation}
Indeed, let $q$ denote the maximum on the right. Necessarily
$p_{k+1}\leq q$, since $p_k^c$ is composite. If the inequality is strict,
start the extension in Lemma~\ref{tw:lem:ceil-right} with $q$ at index
$k+1$. Its resulting constant $D$ satisfies
\[
 D>(q-1)^{1/c^{k+1}}\geq p_{k+1}^{1/c^{k+1}}>B,
\]
while its level-$k$ power lies in $(p_k-1,p_k)$. This preserves all
earlier prime outputs: the boundaries of a level-$j$ ceiling cell,
raised to the integer power $c^{k-j}$, are integers, so no earlier
boundary lies in the interior of the unit interval $(p_k-1,p_k)$.
Both $D$ and $B$ lie in the level-$k$ cell, whose width tends to zero.
For arbitrarily large $k$ such a $D>B$ would contradict local right
maximality. This proves~\eqref{tw:eq:ceil-greedy}.

Set
\[
 \delta=c(1-\theta)-1=\frac{19c}{40}-1>0.
\]
The prime-gap estimate and $B^{c^{k+1}}>p_{k+1}-1$ imply
\begin{equation}\label{tw:eq:ceil-approx}
 0<p_k-B^{c^k}\ll p_k^{-\delta}\ll B^{-\delta c^k}.
\end{equation}
The estimate follows directly by applying the mean value theorem to
$(p_k^c-Kp_k^{c\theta}-1)^{1/c}$; all quantities are positive for
sufficiently large $k$.

Suppose that $B$ is algebraic. The approximation criterion of
Corvaja--Zannier and Dubickas, in Saito's Theorem 2.6 with $q=1$ and
exponents $c^k$, implies that $\beta=B^{c^m}$ is a Pisot number for
some $m\geq1$. Its degree $\ell$ is at least two, since an integer
power of $B$ is impossible. Write its other conjugates as
$\beta_2,\ldots,\beta_\ell$ and put
$\rho=\max_{i\geq2}|\beta_i|<1$. For $n=c^{k-m}$, the integer trace
is eventually the unique nearest integer to $\beta^n$, hence equals
$p_k$. Therefore
\[
 0<\sum_{i=2}^{\ell}\beta_i^n=p_k-\beta^n\ll\beta^{-\delta n}.
\]
The Pisot conjugate-sum lower bound gives, for some $\lambda>0$,
$\rho^n n^{-\lambda}\ll\beta^{-\delta n}$ along these exponents.
Taking $n$th roots yields $\rho\leq\beta^{-\delta}$. Since the
nonzero integer norm has absolute value at least one,
\[
 1\leq|\operatorname{Norm}(\beta)|
 \leq\beta\rho^{\ell-1}
 \leq\beta^{1-\delta(\ell-1)},
 \qquad \ell\leq1+\frac1\delta.
\]
For $c\geq5$ this upper bound is strictly less than two, a
contradiction. For $c=4$ it allows only $\ell=2$, and for $c=3$ it
allows only $\ell=2$ or $3$.

It remains to eliminate $\ell=2$ when $c=3$. If the conjugate of
$\beta$ is $\gamma$ and $N=\beta\gamma\in\mathbb Z$, then for
$n=3^{k-m}$ and all sufficiently large $k$,
\[
 p_{k+1}=\beta^{3n}+\gamma^{3n}
 =p_k\bigl(p_k^2-3N^n\bigr).
\]
This contradicts primality because $p_{k+1}>p_k>1$.
Finally, a rational $B$ would have rational Pisot power $\beta$,
which must be an integer. Thus all cases are irrational.
\end{proof}

\begin{corollary}[Dense canonical transcendental completions]
For every integer $c\geq5$, the greatest completions of feasible finite
prime prefixes form a countable dense subset of $\mathcal V_c$ consisting
entirely of transcendental numbers.
\end{corollary}

\begin{proof}
There are countably many finite prime prefixes. Given $B\in\mathcal V_c$,
let $D_n$ be the greatest completion of its first $n$ outputs. Both $B$
and $D_n$ lie in the same level-$n$ cell, whose width is at most $c^{-n}$,
so $D_n\to B$. Theorem~\ref{tw:thm:ceil-arith} supplies transcendence.
\end{proof}

\begin{remark}[The sign cannot be ignored]
For a ceiling right maximum the Pisot conjugate sum is positive.
For a floor left minimum it is negative. This is why the quadratic
possibility at $c=4$ is not eliminated by the usual even-power sign
argument. The theorem does not assert that every ceiling constant is
transcendental, does not assert the existence of a global least ceiling
constant, and does not prove the transcendence of T\'oth's displayed
cubic numerical branch.
\end{remark}


\section{Sharp finite prime covers for algebraic power floors}
\label{tw:sec:sharp-prime-covers}

We classify the fixed finite sets of primes that can cover all power
floors of a suitable power of a Salem number. This sharpens existence
questions about composite power floors studied in
\cite{Dubickas2006,SaitoIntervals,HIKK2026}, and we contrast the result
with Pisot units.  These are classifications of \emph{finite prime covers};
they do not classify every possible mechanism of eventual compositeness.
No claim of bibliographic priority is made.

For a finite nonempty set of primes $S$, put $M_S=\prod_{p\in S}p$.
A sequence of positive integers has an $S$-cover if every term has a
proper prime divisor belonging to $S$.

\subsection{The distribution input}

\begin{lemma}[Independent Salem angles]\label{tw:lem:salem-independent-angles}
Let $\tau$ be a Salem number of degree $d=2s+2\ge4$, with conjugates
$\tau,\tau^{-1},z_1,z_1^{-1},\ldots,z_s,z_s^{-1}$, where
$z_j=e^{2\pi i\theta_j}$ and $|z_j|=1$.  Then
$1,\theta_1,\ldots,\theta_s$ are linearly independent over $\mathbb Q$.
Consequently, for all integers $a\ge1$ and $b$, the vectors
$((an+b)\theta_1,\ldots,(an+b)\theta_s)$ are equidistributed modulo one.
\end{lemma}
\begin{proof}
More strongly, suppose $\prod_jz_j^{m_j}=\omega$ for integers $m_j$
and a root of unity $\omega$.  A rational linear relation among
$1,\theta_1,\ldots,\theta_s$ is a special case of such a relation.
For each fixed $j$, choose an automorphism of the splitting field that
sends $z_j$ to $\tau$.  No other chosen $z_k$ can map to either $\tau$ or
$\tau^{-1}$: the latter possibility would imply $z_k=z_j^{-1}$, which
is not among the chosen representatives.  Thus all the other images
have modulus one, and the image of $\omega$ also has modulus one.
Taking absolute values gives $\tau^{m_j}=1$, hence
$m_j=0$.  The asserted equidistribution now follows from the Weyl
criterion: every nontrivial integral character has irrational frequency,
and its exponential sum is a bounded geometric sum divided by its length.
\end{proof}

Write
\[
 T_n=\operatorname{Tr}(\tau^n),\qquad
 E_n=\tau^{-n}+2\sum_{j=1}^s\cos(2\pi n\theta_j).
\]
The exact floor identity is
\begin{equation}\label{tw:eq:sharp-floor-identity}
 \lfloor\tau^n\rfloor=T_n-\lceil E_n\rceil.
\end{equation}
We record the strict correction window, including its proof, because its
endpoint is essential to all-index constructions.

\begin{lemma}[Strict Salem window]\label{tw:lem:sharp-salem-window}
For every $n\ge1$,
\[
 -(d-2)<E_n<d-2,\qquad
 d-\lceil E_n\rceil\in I_d:=\{2,3,\ldots,2d-3\}.
\]
For every $c\in\{-d+3,\ldots,d-2\}$, the set of $n$ on any fixed
positive-step arithmetic progression for which $\lceil E_n\rceil=c$
has a positive relative natural density $w_c$, independent of the
progression.
\end{lemma}
\begin{proof}
Put $u=\tau^{-n}\in(0,1)$ and $x_j=|1-z_j^n|^2>0$.  The nonzero
integer norm of $\tau^n-1$ satisfies
\[
 1\le\left|\operatorname{N}(\tau^n-1)\right|
   =\frac{(1-u)^2}{u}\prod_{j=1}^sx_j,
\]
so $\prod_jx_j>u$.  If $\sum_jx_j\le u$, the arithmetic--geometric
mean inequality gives $\prod_jx_j\le(u/s)^s\le u$, a contradiction.
Hence $E_n=u+2s-\sum_jx_j<2s$; the lower bound follows from
$E_n\ge u-2s>-2s$.

For the density assertion, let $Y(t)=2\sum_{j=1}^s\cos(2\pi t_j)$ on
$[0,1)^s$ with Lebesgue probability measure.  Its range is $[-2s,2s]$.
Every open interval $(c-1,c)$ under consideration meets the interior of
that range, and its inverse image contains a nonempty open set.  Its
measure $w_c$ is therefore positive.  Each level set of $Y$ has measure
zero: after fixing $t_2,\ldots,t_s$, the equation in $t_1$ has at most
two solutions, and Fubini's theorem applies.  Equidistribution from
Lemma~\ref{tw:lem:salem-independent-angles}, together with $\tau^{-n}\to0$,
therefore gives the claimed limiting frequencies.  The vanishing
perturbation causes no problem because the two boundary level sets have
measure zero.
\end{proof}

\subsection{An exact characterization of all finite covers}

\begin{theorem}[Salem cover criterion]\label{tw:thm:salem-cover-exact}
Let $\tau$ be a Salem number of degree $d\ge4$, let $J\subset\mathbb Z$
be finite and nonempty, and let $S$ be a finite nonempty set of primes.
The following conditions are equivalent.
\begin{enumerate}
\item There is an integer $h\ge1$ such that, for every $n\ge1$ and
$j\in J$, the integer $\lfloor\tau^{hn}\rfloor+j$ has a proper prime
divisor belonging to $S$.
\item The same assertion holds for all sufficiently large $n$, for some
$h\ge1$.
\item For every $r\in I_d$ and $j\in J$,
\begin{equation}\label{tw:eq:salem-cover-gcd}
 \gcd(M_S,r+j)>1.
\end{equation}
\end{enumerate}
If condition~\textup{(3)} fails, then for every $h\ge1$ there is a
$j\in J$ for which the indices $n$ satisfying
\[
 \gcd(M_S,\lfloor\tau^{hn}\rfloor+j)=1
\]
have positive lower natural density.  If condition~\textup{(3)} holds,
a suitable $h$ is effectively computable.
\end{theorem}
\begin{proof}
Let $C$ be an integral companion matrix for the minimal polynomial of
$\tau$.  Since a Salem polynomial is reciprocal with constant term one,
$C$ is invertible modulo every positive integer, and
$T_n=\operatorname{tr}(C^n)$.

Suppose that~\eqref{tw:eq:salem-cover-gcd} fails at $(r,j)$, and fix an
arbitrary $h$.  Choose $L\ge1$ with $(C^h)^L\equiv I\pmod{M_S}$.
For all $n=Lm$, therefore, $T_{hn}\equiv d\pmod{M_S}$.  On a positive
relative density $w_{d-r}$ of these indices,
$\lceil E_{hn}\rceil=d-r$.  Formula~\eqref{tw:eq:sharp-floor-identity}
then yields
\[
 \lfloor\tau^{hn}\rfloor+j\equiv r+j\pmod{M_S},
\]
which is coprime to $M_S$.  The lower density in the full index sequence
is at least $w_{d-r}/L>0$.  Thus condition~(2) fails.

Conversely, suppose condition~(3) holds.  Choose $h_0$ such that
$C^{h_0}\equiv I\pmod{M_S}$, and take a sufficiently large multiple
$h$ of $h_0$ that
\[
 \lfloor\tau^h\rfloor+\min J>\max S.
\]
For every $n\ge1$, the trace is congruent to $d$, while the strict
window gives $r=d-\lceil E_{hn}\rceil\in I_d$.  Condition~(3)
provides a prime $p\in S$ dividing $r+j$, hence dividing the actual
integer $\lfloor\tau^{hn}\rfloor+j$.  The size inequality makes that
divisor proper.  Finite matrix orders can be found by exact modular
arithmetic.  A rational lower bound $1<u<\tau$, obtained by algebraic
root isolation, makes the size step effective as well.  The implication
(1)$\Rightarrow$(2) is immediate.
\end{proof}

\begin{corollary}[Sharp prime budget for consecutive blocks]
\label{tw:cor:salem-block-budget}
Let $H\ge0$.  Among finite prime sets that cover all integers
\[
 \lfloor\tau^{hn}\rfloor,\ \lfloor\tau^{hn}\rfloor+1,\ldots,
 \lfloor\tau^{hn}\rfloor+H\qquad(n\ge1)
\]
for some positive integer $h$, the unique smallest set under inclusion is
\[
 S_{d,H}=\{p\text{ prime}:p\le2d-3+H\}.
\]
In particular, the minimum number of covering primes is exactly
$\pi(2d-3+H)$, and some effectively computable power of every Salem
number starts an all-composite block of length $H+1$ at every index.
\end{corollary}
\begin{proof}
The integers $r+j$, for $r\in I_d$ and $0\le j\le H$, exhaust
$\{2,\ldots,2d-3+H\}$.  Condition~\eqref{tw:eq:salem-cover-gcd} forces
$S$ to contain each prime in that interval by taking $r+j$ equal to
that prime.  Conversely, $S_{d,H}$ contains a prime divisor of every
integer in the interval.  Apply Theorem~\ref{tw:thm:salem-cover-exact}.
\end{proof}

\begin{corollary}[Exact shift obstruction]\label{tw:cor:salem-shift-obstruction}
For a finite nonempty set $J\subset\mathbb Z$, an all-index cover by
some finite prime set exists after taking a power of $\tau$ if and only if
\[
 J\cap\{-1,-2,\ldots,-(2d-2)\}=\varnothing.
\]
If this intersection is nonempty, no finite prime set covers the
corresponding shifted floors, even eventually, for any positive power
of $\tau$.
\end{corollary}
\begin{proof}
For a fixed $j$, the consecutive integer interval $I_d+j$ contains
$1$ or $-1$ exactly when $-(2d-2)\le j\le-1$.  Such an integer violates
\eqref{tw:eq:salem-cover-gcd} for every $S$.  Outside that range all
integers in $I_d+j$ have absolute value at least two.  Taking the union
of their prime divisors over the finite set $J$ supplies a finite
$S$ satisfying the criterion.
\end{proof}

The negative assertion concerns \emph{finite prime covers}, not the
possibility that every term is composite with unbounded smallest prime
factor.  In particular, it does not assert infinitely many primes among
$\lfloor\tau^{hn}\rfloor-1$.

\begin{corollary}[Simultaneous floor and ceiling]\label{tw:cor:salem-floor-ceiling}
Some effectively computable power $\alpha$ of every Salem number has
both $\lfloor\alpha^n\rfloor$ and $\lceil\alpha^n\rceil$ composite for
every $n\ge1$.  The smallest possible finite prime cover for both
sequences is $\{p:p\le2d-3\}$.
\end{corollary}
\begin{proof}
No positive power of a Salem number is an integer: conjugating it to a
unit-circle conjugate would contradict its modulus.  Thus ceiling is
floor plus one.  Apply Corollary~\ref{tw:cor:salem-block-budget} with $H=1$;
the extra endpoint $2d-2$ is even and greater than two, so it introduces
no additional prime.
\end{proof}

\subsection{Exact frequencies in degree four}

\begin{proposition}[Quartic divisor frequencies]\label{tw:prop:quartic-frequencies}
For a quartic Salem number, the correction values $c=-1,0,1,2$ have
respective frequencies
\[
 (w_{-1},w_0,w_1,w_2)=(1/3,1/6,1/6,1/3)
\]
on every fixed positive-step arithmetic progression.  More generally,
for every Salem number of any degree and every $h\ge1$, the even values
of $\lfloor\tau^{hn}\rfloor$ have natural density exactly $1/2$.
\end{proposition}
\begin{proof}
In degree four, calculate the measures of the four inverse images of
\[(-2,-1),\quad(-1,0),\quad(0,1),\quad(1,2)\] under $2\cos(2\pi t)$; the respective
lengths are $1/3,1/6,1/6,1/3$.  In arbitrary degree the distribution of
$Y=2\sum\cos(2\pi t_j)$ is symmetric under $Y\mapsto-Y$, by translating
every $t_j$ by $1/2$.  Away from the measure-zero integral levels,
$\lceil-Y\rceil=1-\lceil Y\rceil$, so the parity of the correction
is split equally.  The trace modulo two is periodic.  On each of its
finitely many residue progressions the correction has that same limiting
distribution, proving the parity statement.
\end{proof}

For a concrete example
\[
 P(X)=X^4-X^3-X^2-X+1,\qquad \alpha=\tau^{260},
\]
the substitution $Y=X+X^{-1}$ gives $Y^2-Y-3$.
Its roots are $(1\pm\sqrt{13})/2$, one above $2$ and one in $(-2,2)$.
The quartic is irreducible: over $\mathbb Q(\sqrt{13})$ the discriminant
$Y^2-4=(\sqrt{13}-1)/2$ for the large root has norm $-3$ and is not a
square. Hence the large real root $\tau$ is Salem. A companion matrix
and its exact power modulo $30$ are
\[
 C=\begin{pmatrix}0&0&0&-1\\1&0&0&1\\0&1&0&1\\0&0&1&1\end{pmatrix},
 \qquad
 D=C^{260}\equiv
 \begin{pmatrix}21&10&20&10\\20&11&20&10\\20&10&21&10\\20&10&20&11\end{pmatrix}
 \pmod{30}.
\]
Direct multiplication gives $D^3=I$ and
\[
\operatorname{tr}(D)=\operatorname{tr}(D^2)=\operatorname{tr}(D^3)=4\pmod{30}.
\]
Thus
$\operatorname{Tr}(\alpha^n)\equiv4\pmod{30}$ for all $n\ge1$.
Hence the floors are congruent to $5,4,3,2$ modulo $30$ with frequencies
$1/3,1/6,1/6,1/3$, respectively.  The disjoint divisor classes are
therefore
\[
 \begin{array}{c|ccc}
 \text{covering prime}&2&3&5\\ \hline
 \text{natural density}&1/2&1/6&1/3.
 \end{array}
\]
The former three-prime certificate is now sharp: every finite prime
cover for powers of any quartic Salem number must contain $2,3,5$.
The same example also makes every ceiling composite. The divisors
are proper because $\tau>3/2$ and $\tau^{260}>6$.

\subsection{A sharp contrast with Pisot units}

\begin{theorem}[Finite covers for two Pisot classes]
\label{tw:thm:pisot-cover-classification}
Let $\beta>1$ be a Pisot unit, and let $S$ be a finite nonempty set of
primes, with product $M_S$.
\begin{enumerate}
\item If $\beta$ is totally real of degree $d\ge2$, an $S$-cover of
all $\lfloor\beta^{hn}\rfloor$ exists for some $h\ge1$ if and only if
$\gcd(M_S,d-1)>1$.  Thus degree two admits no finite prime cover after
any power; in every higher totally real degree one prime suffices.
\item In the same totally real situation, one $S$-cover of both floor
and ceiling for some power exists if and only if
\[
 \gcd(M_S,d-1)>1\quad\text{and}\quad\gcd(M_S,d)>1.
\]
For $d\ge3$ the minimum number of primes is exactly two.
\item If $\beta$ has degree three and its two other conjugates are
non-real, a finite set $S$ covers all floors of some power if and only
if $\{2,3\}\subseteq S$.  This same condition is necessary and sufficient
to cover both floors and ceilings.
\end{enumerate}
All the sufficient constructions are effective.  Replacing all indices
by all sufficiently large indices leaves the characterizations unchanged.
\end{theorem}
\begin{proof}
The companion matrix $C$ is invertible modulo every $M_S$ because
$\beta$ is a unit.  In the totally real case, take an even multiple $h$
of its order modulo $M_S$, sufficiently large that
\[
 0<E_{hn}:=\sum_{i=2}^d\beta_i^{hn}<1
 \qquad(n\ge1).
\]
This is possible because every $|\beta_i|<1$.  The trace is congruent
to $d$, the floor is congruent to $d-1$, and the ceiling is congruent to
$d$.  A further size enlargement makes the indicated covering primes
proper divisors, proving sufficiency in (1) and (2).

Conversely, for any proposed power $h$, restrict $n$ to even multiples
of an order of $C^h$ modulo $M_S$.  The same two residue statements
hold for all sufficiently large such $n$.  If either required gcd is
one, the corresponding integers have no divisor in $S$, proving
necessity.  Since $d$ and $d-1$ are coprime, two distinct primes are
necessary in (2), and they suffice when $d\ge3$.

For (3), write the conjugates as $\beta,z,\bar z$, where
$z=\rho e^{2\pi i\theta}$ and $0<\rho<1$.  The angle $\theta$ is
irrational.  Indeed, otherwise some positive power satisfies
$z^m=\bar z^m$.  The degree of $\beta^m$ divides the prime degree three,
and the conjugate collision rules out degree three, so $\beta^m$ is
rational.  It is then a rational algebraic-integer unit, hence $1$ or
$-1$, contradicting $\beta>1$.

Now $E_n=2\rho^n\cos(2\pi n\theta)$ has absolute value less than one
for all large $n$.  On every arithmetic progression its positive and
negative signs each have relative density $1/2$.  On a matrix-return
progression the floors consequently take residues $2$ and $3$ modulo
$M_S$, each with positive relative density.  This forces both primes
$2$ and $3$ into $S$.  Conversely, choose a sufficiently large
matrix-return power.  Its floors have residue $2$ or $3$ modulo six,
and its ceilings have residue $3$ or $4$.  Both are covered by $2,3$;
the size condition again ensures proper divisibility.
\end{proof}

This classification answers why the different constructions need
one, two, or three small primes.  It does not contradict the existence
of infinitely many composite quadratic Pisot floors: having infinitely
many composite terms is much weaker than covering every term of a power
subsequence by one fixed finite set of primes.


% Primary inputs: SaitoDimension2021, Matomaki2010, SaitoTakeda2022,
% BakerHarmanPintz2001.  New draft: full local dimension and computability.
\section{Local dimension and effective spectra with arbitrarily long square stretches}
\label{tw:newloc:section}

Put $E_1>0$ and $E_{n+1}=c_{n+1}E_n$, and write
\[
 W_E=\{A>1:\lfloor A^{E_n}\rfloor\in\mathbb P\text{ for every }n\ge1\}.
\]
In this section the multipliers $c_{n+1}$ are integers at least $2$, and
$c_{n+1}\ge3$ for infinitely many $n$. No upper bound on the multipliers,
and no bound on the lengths of runs of $2$'s, is imposed.
Saito--Takeda's theorem \cite{SaitoTakeda2022} proves that this set
is nonempty, closed, and perfect. The results below strengthen this topological conclusion locally
and identify an ordinary computability consequence. They do not apply to
an eventually all-square schedule.

\subsection{Two prime-interval inputs and a uniform seed lemma}

We use the following two published estimates. There are positive constants
$b,d,D$ and sufficiently large thresholds such that
\begin{align}
 \#\bigl(\mathbb P\cap[x,x+x^{21/40}]\bigr)
   &\ge b\frac{x^{21/40}}{\log x},\label{tw:newloc:bhp}\\
 \#\mathcal B(x,\gamma)&\le D x^{2/3-\gamma}
       \qquad(1/2\le\gamma\le1),\label{tw:newloc:exceptional}
\end{align}
where $\mathcal B(x,\gamma)$ is any disjoint family of intervals
$[u,u+u^\gamma]\subset[x,2x]$ that contain at most
$d u^\gamma/\log u$ primes. The first is the Baker--Harman--Pintz estimate,
as stated in \cite[Lemma~3.1]{SaitoDimension2021}; see also \cite{BHP}.
The second is Matom\"aki's Lemma~9 \cite{Matomaki2010}; its constants and threshold are uniform in
$\gamma\in[1/2,1]$.

\begin{lemma}[A bounded-tail good seed]\label{tw:newloc:goodseed}
There are $Q_0>0$ and $a>0$ with the following property. Let
$2\le t_j\le4$ for every $j\ge1$, and let $q\ge Q_0$ be prime. If
\[
 \#\bigl(\mathbb P\cap[q^{t_1},q^{t_1}+q^{t_1-1}]\bigr)
 \ge \frac{d}{2}\frac{q^{t_1-1}}{t_1\log q},
\]
then the set of $y\in[q^{1/2},(q+1)^{1/2})$ satisfying
\[
 \lfloor y^{2}\rfloor=q,\qquad
 \lfloor y^{2t_1\cdots t_j}\rfloor\in\mathbb P\quad(j\ge1)
\]
has Hausdorff dimension at least
$\bigl(1+4/(q\log q)\bigr)^{-1}$.
\end{lemma}

\begin{proof}
Call a prime $v$ good for a next multiplier $t$ if its corresponding
interval has at least $(d/2)v^{t-1}/(t\log v)$ primes. Suppose $v$ is good
for $t\in[2,4]$, and consider its prime children
$w\in[v^t,v^t+v^{t-1}]$. For a following multiplier $s\in[2,4]$, the
intervals $[w^s,w^s+w^{s-1}]$ are disjoint, because distinct prime labels
differ by at least $2$, and lie in $[v^{ts},2v^{ts}]$ once $v$ is large.
An interval belonging to a child that is not good is exceptional in
\eqref{tw:newloc:exceptional}, with $\gamma=1-1/s$. Thus the number of such
children is at most
\[
 Dv^{ts(2/3-(1-1/s))}=Dv^{t(1-s/3)}\le Dv^{t/3}.
\]
Uniformly for $t,s\in[2,4]$, this is smaller than
$(d/4)v^{t-1}/(t\log v)$ for sufficiently large $v$, since
$t-1-t/3\ge1/3$. Discarding at most a fixed number of endpoint children
still leaves at least $a v^{t-1}/(t\log v)$ good children, for a fixed
$a>0$. Their closed root cells can be taken strictly inside the parent's
cell: $w>v^t$ and $w+1<(v+1)^t$ after the endpoint discards.

Iterate the selection. The resulting tree satisfies conditions
(C1)--(C4) of Saito's Proposition~5.1, with $R=4$, $\theta=1$, $L=1$,
and $Q=a$. That proposition gives the stated dimension bound. This is
precisely the part of Saito's construction after the initial good prime
has been found; it does not require that the given seed be selected from
an interval $[M,2M]$.
\end{proof}

\begin{lemma}[A local reset at a multiplier at least three]
\label{tw:newloc:reset}
There is $P_0$ such that the following holds. Start with a prime root cell
$J=(p^{1/E},(p+1)^{1/E})$, where $p\ge P_0$. Suppose the next exponent is
$rE$ with $r\ge3$, and all following exponent ratios are at least $2$.
There is a subset of $J$ satisfying all these future prime conditions,
whose dimension is at least
\[
 \left(1+\frac{4}{q\log q}\right)^{-1}
 \quad\text{for some prime }q>p^3.
\]
The threshold $P_0$ is independent of the future schedule and of $E,r$.
\end{lemma}

\begin{proof}
Write $r=s2^m$, where $m\ge0$ is an integer and $3\le s<6$; for example,
take $m=\lfloor\log_2(r/3)\rfloor$. Replace the jump $r$ by the factors
$s,2,\ldots,2$. Refine every subsequent ratio $u\ge2$ into equal factors
$u^{1/\lfloor\log_2u\rfloor}\in[2,4)$. The refined schedule after the
first factor $s$ has all factors in $[2,4]$.

By subdividing $[p^s,p^s+p^{s-1}]$ into intervals of the length in
\eqref{tw:newloc:bhp}, that interval contains at least
$b' p^{s-1}/(s\log p)$ primes, uniformly for $s\in[3,6]$ and large $p$.
If $t\in[2,4]$ is the next refined factor, the intervals
$[q^t,q^t+q^{t-1}]$ associated with these prime candidates are disjoint
and lie in $[p^{st},2p^{st}]$. Formula~\eqref{tw:newloc:exceptional} excludes
at most $Dp^{s(1-t/3)}\le Dp^{s/3}$ bad candidates.
Since $s-1-s/3\ge1$, there is a good candidate $q>p^s\ge p^3$.
Also $q+1<(p+1)^s$, so its root cell at exponent $sE$ is strictly inside
$J$. Apply Lemma~\ref{tw:newloc:goodseed}, followed by the bi-Lipschitz
change of variable $A=y^{2/(sE)}$. All original exponents are among the
refined exponents, which proves the assertion.
\end{proof}

\subsection{Full dimension in every relative open set}

\begin{theorem}[Full local dimension]\label{tw:newloc:fulllocal}
Assume $c_{n+1}\in\mathbb Z$, $c_{n+1}\ge2$ for every $n$, and
$c_{n+1}\ge3$ for infinitely many $n$. For every open interval $U$ with
$U\cap W_E\ne\varnothing$,
\[
 \boxed{\dim_H(U\cap W_E)=1.}
\]
In particular, every feasible finite prime-output prefix that occurs in
$W_E$ has a full-dimensional set of completions.
\end{theorem}

\begin{proof}
Choose $A\in U\cap W_E$ and let $p_n=\lfloor A^{E_n}\rfloor$.
For every $j<n$, the quotient $M=E_n/E_j$ is an integer. On the half-open
cell
\[
 I_n=[p_n^{1/E_n},(p_n+1)^{1/E_n}),
\]
the integer part at exponent $E_j$ is constant: after taking the
$E_n$-th power, its possible discontinuities are integers of the form
$h^M$, and none lies strictly between the consecutive integers
$p_n$ and $p_n+1$. Since $A\in I_n$, the entire cell satisfies all
conditions through index $n$.

Its length tends to zero, and $p_n\to\infty$. Consequently, for every
sufficiently large $n$ such that $c_{n+1}\ge3$, the cell $I_n$ lies in
$U$ and $p_n\ge P_0$. Lemma~\ref{tw:newloc:reset} supplies a subset of
$U\cap W_E$ with dimension at least
$(1+4/(q_n\log q_n))^{-1}$, where $q_n>p_n^3$. Letting $n$ run through
these infinitely many indices gives dimension at least $1$. The ambient
upper bound is $1$.

For the last assertion, note that every $A\in W_E$ is an interior point
of each of its finite prime cells. Indeed, $A^{E_j}$ cannot be an integer,
since a later integer exponent ratio would make a prime output into a
proper prime power. Its prescribed finite prefix therefore holds in an
open neighborhood of $A$, to which the theorem applies.
\end{proof}

\begin{remark}[What the integer assumption does]
The nesting argument uses the integrality of $E_n/E_j$ essentially.
The local reset lemma itself allows real ratios, but it applies only
when the entire seed cell already respects the earlier conditions.
Global full dimension for arbitrary real ratios at least $2$ does not,
by itself, prove full dimension in every prescribed prefix. No such
unrestricted local assertion is made here.
\end{remark}

\subsection{Computability of the minimum and of a dense spectrum}

\begin{lemma}[An integer-arithmetic continuation rule]
\label{tw:newloc:computablebranch}
There are fixed integers $K,P\ge1$ such that, for every integer multiplier
schedule at least $2$, every prime $p\ge P$ immediately before a
multiplier at least $3$ has a computable continuation. The continuation
is uniform in the computable schedule and the seed. At a multiplier at
least $3$, at least two child primes can be chosen, each with a
computable continuation.
\end{lemma}

\begin{proof}
Set
\[
 H(v)=\left\lceil\frac{v}{K\lceil\log_2 v\rceil}\right\rceil,
\]
and say that a prime $v$ is square-good if $[v^2,v^2+v]$ contains at least
$H(v)$ primes. Choose $K$ so large that, for sufficiently large $v$,
\[
 H(v)\le\frac{d v}{2\log v}.
\]
Failure of square-goodness then makes $[v^2,v^2+v]$ exceptional in
\eqref{tw:newloc:exceptional} with $\gamma=1/2$.

For an integer $c\ge3$, \eqref{tw:newloc:bhp} implies that
$[p^c,p^c+p^{c-1}]$ contains at least
$b''p^{c-1}/(c\log p)$ primes, uniformly in $c\ge3$ for large $p$.
Indeed $p^{c-1}/(p^c)^{21/40}=p^{19c/40-1}$ tends to infinity uniformly,
and subdivision into short intervals proves the bound. If the following
multiplier is $2$, remove those candidate primes $q$ that are not
square-good. Their disjoint square intervals lie in $[p^{2c},2p^{2c}]$,
so at most $Dp^{c/3}$ are removed. The difference is at least two,
uniformly for $c\ge3$, after increasing $P$.

At a multiplier $c=2$, require that the current seed $p$ is square-good.
There are at least $H(p)$ prime children in $[p^2,p^2+p]$; if the next
multiplier is also $2$, discard at most $Dp^{2/3}$ that are not
square-good. Since $H(p)\gg p/\log p$, at least two remain for large $p$.
If the next multiplier is at least $3$, no discard is needed.

Thus at each step select the least available child, imposing
square-goodness precisely when the following multiplier is $2$.
All powers and interval endpoints are integers; primality and the finite
prime counts are decidable, and $\lceil\log_2v\rceil$ is computable by
integer comparison. The initial multiplier is at least $3$, so no
square-good assumption on the initial seed is necessary.
The child root cells are strictly nested, and their lengths tend to zero
with a computable modulus. Their unique limit gives the required
computable real. Choosing the first or second available child at a step
with multiplier at least $3$ gives the final assertion.
\end{proof}

\begin{theorem}[Computable least constants, despite arbitrarily long square stretches]
\label{tw:newloc:computableminimum}
Suppose $E_1>0$ is computable, $(c_{n+1})$ is a computable integer sequence
with $c_{n+1}\ge2$, and $c_{n+1}\ge3$ for infinitely many $n$. Then the
least element of $W_E$ is a computable real. Moreover, computable elements
are dense in $W_E$.
\end{theorem}

\begin{proof}
Write $M_n=E_n/E_1\in\mathbb N$. For an index $n$ immediately before a
multiplier at least $3$, enumerate primes $q$ in increasing order and
test the finite conditions
\[
 \left\lfloor q^{1/(E_n/E_j)}\right\rfloor\in\mathbb P
       \qquad(1\le j<n).
\]
These tests are exact integer-root and primality computations. The cell
$[q^{1/E_n},(q+1)^{1/E_n})$ satisfies the first $n$ prime conditions
exactly when they hold. At any feasible cell, $q\ge2^{M_n}$, because the
first prime output is at least $2$.

Choose $n$ with $2^{M_n}\ge P$. By
Lemma~\ref{tw:newloc:computablebranch}, every feasible cell at this index
has an infinite continuation. Such cells exist by the nonemptiness
of $W_E$, so the enumeration stops. Let $q_n$ be its first successful
prime. The actual least constant $\xi_E$ therefore lies in
\[
 q_n^{1/E_n}\le\xi_E<(q_n+1)^{1/E_n}.
\]
The existence of the minimum follows also directly from closedness,
nonemptiness, and the lower bound $A\ge2^{1/E_1}>1$.
For $E_n\ge1$, the mean value theorem gives
\[
 (q_n+1)^{1/E_n}-q_n^{1/E_n}\le E_n^{-1}.
\]
Given a requested precision, search for a later index before a multiplier
at least $3$ with $E_n$ sufficiently large, perform this finite search,
and approximate the two computable endpoints. This produces a Cauchy
name for $\xi_E$.

For density, choose $A\in W_E$ and an arbitrary neighborhood $U$ of $A$.
A sufficiently late cell immediately before a multiplier at least $3$
is contained in $U$, respects every earlier condition, and has prime
label at least $P$. Applying the computable continuation rule produces
a computable member of $U\cap W_E$.
\end{proof}

\begin{remark}[Computability is not a numerical evaluation]
The integers $K,P$ in the preceding proof exist by the published analytic
estimates; they have not been numerically evaluated in this draft.
Fixing such integers yields an ordinary Turing algorithm, and the theorem
asserts computability in this standard existence sense. It supplies
neither a practical running-time bound nor a new certified decimal digit
of the ordinary least Mills constant. In particular, the potentially
enormous finite search preceding the first usable branching level is
not replaced by assuming that the familiar small greedy prime chain
extends forever. This distinction is the reason that Elsholtz's explicit,
practically evaluable variants \cite{Elsholtz2020} remain substantially stronger computational
information for their specified schedules.
\end{remark}

\begin{corollary}[Every local piece contains every Turing degree]
\label{tw:newloc:degrees}
Under the hypotheses of Theorem~\ref{tw:newloc:computableminimum}, every
open interval $U$ meeting $W_E$ contains a computably embedded binary
Cantor set of elements of $W_E$. The embedding and its inverse on its
image are computable. Hence $U\cap W_E$ contains a real of every Turing
degree.
\end{corollary}

\begin{proof}
Begin with a sufficiently late feasible seed cell inside $U$. At each
subsequent multiplier at least $3$, use a bit to select the first or the
second permitted child from Lemma~\ref{tw:newloc:computablebranch}. At all
intermediate steps take the least permitted child. There are infinitely
many bit choices, all searches terminate, and the computable upper bounds
on root-cell lengths show that the map from binary sequences to limits is
computable. Distinct bits select separated prime cells, so the map is
injective. Its inverse is computable by comparing a Cauchy name of the
limit with the finitely many separated child cells at each bit level.
The limit never belongs to an integer-power endpoint, because a later
integer multiplier would give a composite prime power. Thus every such
comparison terminates. The two computable reductions preserve Turing
degrees. No positive Hausdorff dimension is asserted for this particular
two-child Cantor subset.
\end{proof}


\section{A complete bounded exclusion for cubic Mills shadows}\label{tw:sec:cubic-new}

The results of this section address a finite, explicitly bounded part of
Saito's cubic Pisot obstruction. They do not exclude all cubic Pisot
numbers and do not prove that Mills' constant is transcendental. The
finite verification uses only small-prime modular arithmetic; it does
not assume values of the actual Mills prime sequence.

\begin{lemma}[An elementary lower bound]\label{tw:lem:cubic-below-two}
Every totally real cubic Pisot number exceeds $2$.
\end{lemma}
\begin{proof}
Let its conjugates be $\beta,u,v$, with $\beta>1$ and
$0<|u|,|v|<1$, and let $f$ be its monic irreducible polynomial.
Since $|\beta uv|$ is a nonzero integer, $|uv|\ge1/\beta$.
The nonzero integer $f(1)f(-1)$ gives
\[
 1\le |f(1)f(-1)|
   =(\beta^2-1)(1-u^2)(1-v^2)
   \le (\beta^2-1)(1-1/\beta)^2.
\]
Indeed, $(1-u^2)(1-v^2)\le(1-|uv|)^2\le(1-1/\beta)^2$.
For $1<\beta\le2$ the last expression is at most $3/4$, a
contradiction.
\end{proof}

\begin{lemma}[Exact coefficient catalogue]\label{tw:lem:cubic-catalogue}
Fix an integer $B\ge3$. The monic minimal polynomials of totally real
cubic Pisot numbers $\beta<B$ whose small conjugate of larger modulus
is negative are exactly the following two disjoint finite families.
\begin{enumerate}
\item Opposite-sign small conjugates:
\[
 X^3-aX^2-dX+c,\qquad
 2\le a<B,\quad d,c\ge1,\quad d+c\le a.
\]
\item Two negative small conjugates:
\[
 X^3-aX^2-dX-c,
\]
where $1\le a<B$, $1\le d<2B$, $1\le c<B$, and
\[
 \begin{gathered}
 d\le2a+2,\qquad d\le a+c,\qquad B^3-aB^2-dB-c>0,\\
 a^2d^2+4d^3-4a^3c-27c^2-18adc>0.
 \end{gathered}
\]
\end{enumerate}
\end{lemma}
\begin{proof}
In the first case write the roots as $\beta,-\rho,\sigma$, with
$0<\sigma<\rho<1$. Then $a=\beta-\rho+\sigma$,
$d=\beta(\rho-\sigma)+\rho\sigma>0$, and $c=\beta\rho\sigma>0$.
The inequality $f(-1)<0$ is exactly $d+c<a+1$, or $d+c\le a$.
Conversely these integer conditions give one root in each of
$(-1,0),(0,1),(a,a+1)$ by direct sign evaluation. Since
$f(a)=-ad+c<0$, the positive large root exceeds the trace $a$;
therefore the negative small root has larger modulus. Thus $\beta<B$.

In the second case write the roots as $\beta,-\rho,-\sigma$, with
$0<\sigma<\rho<1$. The preceding lemma gives
$a=\beta-\rho-\sigma>0$. Furthermore
\[
 d=a(\rho+\sigma)+\rho^2+\rho\sigma+\sigma^2<2a+3.
\]
The other coefficient bounds, the positive discriminant, and $f(B)>0$
are necessary. Conversely, positive discriminant gives three distinct
real roots. There is exactly one positive root by Descartes' rule,
and the other two are negative. The negative critical point is
$(a-\sqrt{a^2+3d})/3>-1$ because $d\le2a+2$.
Together with $f(-1)<0$, this puts both negative roots in $(-1,0)$.
The large root exceeds $1$ because the product of the three root
moduli is the positive integer $c$; $f(B)>0$ puts it below $B$.
Finally, each polynomial in either family is irreducible: a rational
root would be an integer, necessarily the large root, leaving a
monic integral quadratic with a nonzero integral constant term of
absolute value less than one.
\end{proof}

\begin{theorem}[A finite cubic exclusion with a density bound]\label{tw:thm:cubic-bounded}
Let $\beta<100$ be a totally real cubic Pisot number whose dominant
small conjugate is negative. Then there are a prime $\ell\le47$,
an integer $j\ge0$, and an integer $1\le L\le1148$ such that
\[
 \ell\mid\lfloor\beta^{3^{j+tL}}\rfloor
 \qquad\text{for every sufficiently large }t.
\]
Consequently, the set of composite indices in
$(\lfloor\beta^{3^r}\rfloor)_{r\ge0}$ has lower density at least
$1/1148$.
\end{theorem}
\begin{proof}[Proof with a finite arithmetic certificate]
Write the polynomial as $X^3-aX^2+bX-c$, and set
\[
 (T_r,B_r,E_r)=
 \left(\sum_i\alpha_i^{3^r},
       \sum_{i<j}(\alpha_i\alpha_j)^{3^r},c^{3^r}\right).
\]
Newton's identities give the integral map
\[
 \Phi(T,B,E)=
 (T^3-3TB+3E,\ B^3-3TBE+3E^2,\ E^3).
\]
For a fixed prime $\ell$, its orbit modulo $\ell$ is eventually
periodic. A state with trace coordinate zero on the cycle yields
an infinite arithmetic progression of traces divisible by $\ell$.
Since the dominant small conjugate is negative, these traces agree
with the desired floors for every sufficiently large $r$.

The exact coefficient catalogue with $B=100$ contains $213569$
polynomials: $161700$ from the first family and $51869$ from the
second. Exhaustive iteration modulo the following primes gives the
listed counts of polynomials certified for the first time:
\[
\begin{array}{c|rrrrrrrr}
\ell&2&3&5&7&11&13&17&19\\\hline
\text{count}&134584&27162&27780&14024&4856&2501&1163&544
\end{array}
\]
\[
\begin{array}{c|rrrrrrr}
\ell&23&29&31&37&41&43&47\\\hline
\text{count}&492&182&142&86&35&14&4.
\end{array}
\]
There are no survivors. The largest cycle period used is $1148$,
and the largest preperiod is $3$. A complete standard-library Python verifier appears in
Appendix~\ref{tw:app:cubic-verification}; all coefficients,
discriminants, and modular orbit operations are exact integers.
An independent companion-matrix audit checks the characteristic
polynomial at the two cycle endpoints and the zero trace at the
chosen intermediate index. The resulting progression has density
$1/L\ge1/1148$. Its entries eventually exceed $\ell$, so they are
composite rather than merely divisible by $\ell$.
\end{proof}

The degree and minimal-polynomial structure needed below were proved in
Theorem~\ref{mil:thm:mills-degree-structure}. We now combine that
structure with the bounded modular exclusion.

\begin{corollary}[Height and degree consequences for Mills' constant]
Suppose that the least Mills constant $A$ is algebraic. In its
minimal-degree cubic representation
\[
 d=3^{t+1},\qquad\beta=A^{3^t},\qquad
 \operatorname{MinPoly}_A(X)
   =X^d-aX^{2d/3}+bX^{d/3}-c,
\]
one necessarily has $\beta>100$ and $a\ge99$. In particular, the
coefficient height of $A$ is at least $99$, independently of its
degree. If $U>1$ is any independently certified upper bound on $A$,
then its degree must additionally satisfy
\[
 U^{d/3}>100.
\]
\end{corollary}
\begin{proof}
Theorem~\ref{mil:thm:mills-degree-structure} supplies the degree and
cubic-shadow structure. The bounded exclusion
rules out $\beta<100$, and equality is impossible for a cubic
irrational. Its trace is $a>\beta-2>98$, hence $a\ge99$.
The final inequality follows from $\beta=A^{d/3}\le U^{d/3}$.
\end{proof}

The height statement is a bounded exclusion valid at all degrees;
it is not an assertion that a numerical value of the actual Mills
constant has been certified. No such bound $U$ is assumed here.

\begin{example}[The hand-checkable range below four]
There are precisely four polynomials in the catalogue below $4$:
\[
 X^3-2X^2-X+1,\quad X^3-3X^2-X+1,\quad
 X^3-3X^2-X+2,\quad X^3-3X^2-2X+1.
\]
For their trace sequences, respectively, the indices
$r\equiv0\pmod2$, all $r$, all $r\ge1$, and
$r\equiv1\pmod2$ give divisibility by $2,3,2,2$.
These statements follow immediately by applying $\Phi$ modulo $2$
or $3$. This is an independently hand-checkable subcase of the
finite theorem.
\end{example}

\section{Norm divisors that force an entire composite tail}

\begin{proposition}[An exact norm-prime obstruction]
Let $f(X)=X^3-aX^2+bX-c$ be the minimal polynomial of a cubic
Pisot number whose dominant small conjugate is negative, and put
$T_r=\operatorname{Tr}(\beta^{3^r})$.
If a prime $\ell\mid c$ divides any $T_j$, then
\[
 \ell\mid T_r\quad\text{for every }r\ge j,
\]
and the floor sequence is eventually composite. In particular,
$\gcd(a,c)>1$ is sufficient for eventual compositeness.

More exactly, suppose $\ell\ge5$, $\ell\mid c$, and
$\ell\nmid b(a^2-4b)$. Let $u,v$ be the nonzero distinct roots of
$X^2-aX+b$ in $\mathbb F_{\ell^2}$, and put $R=u/v$.
There is an index $j$ with $\ell\mid T_j$ if and only if
\[
 \operatorname{ord}(R)=2\cdot3^s
 \quad\text{for some integer }s\ge0.
\]
In that case the first such index is $s$ and every subsequent
trace is divisible by $\ell$.
\end{proposition}
\begin{proof}
Modulo a prime dividing $c$, Newton's recurrence reads
$T_{r+1}=T_r(T_r^2-3B_r)$, proving propagation and the gcd criterion.
Under the additional hypotheses the three residual roots are
$0,u,v$, so $T_r=u^{3^r}+v^{3^r}$.
Its vanishing is equivalent to $R^{3^r}=-1$, and in a finite field
of odd characteristic this is equivalent to the stated order
condition with $s\le r$.
\end{proof}

\begin{corollary}[An explicit infinite family]
For every integer $a\ge8$ with $a\equiv\pm1\pmod5$, let $\beta_a$
be the large root of
\[
 X^3-aX^2-3X+5.
\]
Then $\beta_a$ is cubic Pisot and
\[
 5\mid\lfloor\beta_a^{3^r}\rfloor\qquad(r\ge1).
\]
Thus every one of these terms is composite. This family is not
obtained from the initial gcd obstruction: $\gcd(a,5)=1$.
\end{corollary}
\begin{proof}
The opposite-sign coefficient catalogue applies since $3+5\le a$;
in this case floor and trace agree for every odd power. The norm is
$-5$ and $T_1=a^3+9a-15\equiv a(a^2-1)\equiv0\pmod5$.
The propagation proposition applies, and all the relevant floors
exceed $5$.
\end{proof}

The norm-prime criterion is an elementary application of finite-field
orders and Newton identities. The bounded catalogue and its complete
modular certification are the additional concrete computations here;
no claim of priority over all published computations is made.
\clearpage
\part{Boundary independence and further arithmetic extensions}
\label{part:extensions}
% Extracted and consolidated from Saito_Questions_Further_Results.tex.
% This fragment has no preamble, document environment, bibliography,
% or external input/include. Appendix delimiter is retained.
% All labels use f13:.

\section{Relations among all left boundary constants under sparse growth bursts}
\label{f13:bdsup:section}

Theorem~\ref{fol:thm:left-boundary-monomials} already establishes
support-rank independence for the entire left sub-boundary under an
unbounded-limsup hypothesis. Here we record explicit continuation and
local approximation bounds and deduce homogeneous-ratio consequences
and a sparse-burst example. The recalled boundary formulation is the
same theorem, not a second extension from a selected binary family.

Let $C_1$ be a positive integer, let
$C_{n+1}=c_{n+1}C_n$ with integers $c_{n+1}\ge2$, and put
\[
 W(C)=\{A>1:\lfloor A^{C_n}\rfloor\in\mathbb P\text{ for all }n\}.
\]
Write $L(C)$ for its local left minima: $A\in L(C)$ if some $\eta>0$
satisfies $(A-\eta,A)\cap W(C)=\varnothing$. Equivalently, these are the right endpoints of complementary intervals
lying immediately to their left, including the least element of
$W(C)$. They form Saito--Takeda's left sub-boundary
$\widetilde\partial_L W(C)$; the terminology refers to the side on
which the gap lies.
For a nonzero polynomial $F=\sum_\nu a_\nu X^\nu$, define
\[
 \rho(F)=\dim_{\mathbb Q}\operatorname{span}_{\mathbb Q}
    \{\nu-\mu:a_\nu a_\mu\ne0\}.
\]
This is the affine dimension of its exponent support.

\begin{lemma}[Prime-root separation recalled]\label{f13:bdsup:separation}
Fix $T>1$ and a nonzero integer polynomial $F$ of support rank $\rho$.
For all sufficiently large integers $E$ and all distinct primes
$p_1,\ldots,p_r$ with $p_i^{1/E}\le T$,
\[
 0<|F(p_1^{1/E},\ldots,p_r^{1/E})|
 \ge B^{1-E^\rho},\qquad
 B=\max\left(2,\sum_\nu |a_\nu|T^{|\nu|}\right).
\]
Consequently, if distinct positive numbers $A_i$ admit these prime-root
approximations with errors $\varepsilon_n$ at integers $E_n\to\infty$,
and
\[
 \limsup_n \frac{\log(1/\varepsilon_n)}{E_n^\ell}=\infty,
\]
then no nonzero polynomial over $\overline{\mathbb Q}$ of support rank
at most $\ell$ vanishes at the tuple $(A_1,\ldots,A_r)$.
\end{lemma}
\begin{proof}
The separation bound is Theorem~\ref{fol:thm:support-norm}; the
algebraic-coefficient transfer, including Laurent polynomials, is
Theorem~\ref{fol:thm:sparse-transfer} applied along the subsequence
on which the displayed quotient tends to infinity.
\end{proof}

\begin{lemma}[Continuation through arbitrary integer square stretches]
\label{f13:bdsup:continuation}
There are fixed integers $K,P\ge1$ with the following property. For a
prime $p\ge P$, call $p$ square-good when $[p^2,p^2+p]$ contains at least
\[
 H(p)=\left\lceil\frac{p}{K\lceil\log_2p\rceil}\right\rceil
\]
primes. Given any sequence of integer multipliers at least $2$, a seed
$p\ge P$ has an infinite prime continuation if either it is square-good
or its next multiplier is at least $3$. At a next multiplier $c\ge3$,
the first child can be chosen in
$[p^c,p^c+p^{21c/40}]$. All root cells in the continuation can be chosen
strictly nested.
\end{lemma}
\begin{proof}
Choose $K$ so large that $H(q)\le dq/(2\log q)$ for all sufficiently
large $q$, where $d$ is Matom\"aki's exceptional-interval constant.
The estimate gives only $O(X^{1/3})$ primes $q\in[X,2X]$ that fail to
be square-good: apply it to their disjoint intervals
$[q^2,q^2+q]$ at squared scale $X^2$, splitting into a fixed number of
dyadic intervals and discarding at most one boundary-crosser per split.

At a step $c\ge3$, Baker--Harman--Pintz \cite{BHP} supplies
$\gg p^{21c/40}/(c\log p)$ prime candidates in
$[p^c,p^c+p^{21c/40}]$. If the following multiplier is $2$, discard
candidates that are not square-good; at most $O(p^{c/3})$ are removed.
Since $21/40>1/3$, at least two candidates remain uniformly for $c\ge3$
once $p\ge P$. Uniformity follows, for example, because
$p^{c(21/40-1/3)}/c$ is increasing for $c\ge3$ when $p$ is sufficiently
large. If the following multiplier is at least $3$, no goodness
restriction is needed.

At a square step, the current seed is required to be square-good.
Its interval $[p^2,p^2+p]$ supplies $H(p)\gg p/\log p$ candidates;
when the following multiplier is also $2$, discard at most
$O(p^{2/3})$ poor candidates. Again at least two remain for $p\ge P$.
This maintains exactly the goodness condition needed for the next step.
Every retained prime satisfies $p^c<q$ and $q+1<(p+1)^c$.
The assertion for $c\ge3$ uses $21c/40<c-1$, uniformly with positive
margin; for $c=2$ it follows directly from $q\le p^2+p$.
Iteration gives strictly nested prime root cells of lengths tending
to zero and hence the claimed continuation.
\end{proof}

\begin{lemma}[Uniform approximation of every left-gap endpoint]
\label{f13:bdsup:localapprox}
Let $A\in L(C)$, and put $p_n=\lfloor A^{C_n}\rfloor$.
At every sufficiently late index with $c=c_{n+1}\ge3$,
\[
 p_{n+1}\le p_n^c+p_n^{21c/40},\qquad
 0<A-p_n^{1/C_n}
 \le \frac{2}{cC_n}p_n^{1/C_n-19c/40}.
\]
The index threshold may depend on $A$, but the estimates hold for all
larger indices with $c\ge3$, whether the multipliers are bounded or not.
\end{lemma}
\begin{proof}
Integer multipliers ensure that the half-open prime root cell
$[p_n^{1/C_n},(p_n+1)^{1/C_n})$ respects all previous outputs. They also exclude $A^{C_n}\in\mathbb N$, since a later
integer multiplier would make a proper prime power composite.

Suppose that the first displayed inequality fails at arbitrarily large
indices with $c\ge3$. Put $p=p_n$. By
Lemma~\ref{f13:bdsup:continuation}, a prime
$q\in[p^c,p^c+p^{21c/40}]$ has an infinite continuation respecting
the entire remaining integer multiplier schedule.
Each replacement prime lies below $p_{n+1}$, and its root cell lies
strictly inside the parent cell, because $21c/40<c-1$ for $c\ge3$.
Its continuation gives $A_n'<A$ in $W(C)$ with
$|A_n'-A|\le1/C_n$. Infinitely many such replacements contradict the
left-gap condition. The first inequality follows.

Consequently
\[
 A^{cC_n}<p_{n+1}+1\le p_n^c+2p_n^{21c/40}.
\]
Concavity of the $c$th root gives
$A^{C_n}-p_n\le(2/c)p_n^{1-19c/40}$.
Applying the mean value theorem to the $C_n$th root yields the stated
bound. The finitely many cases with $C_n=1$ can be discarded.
\end{proof}

\begin{theorem}[Left-boundary support theorem, recalled]
\label{f13:bdsup:main}
Let $\ell\ge1$ and assume
\[
 \boxed{\limsup_{n\to\infty}
       \frac{c_{n+1}}{C_n^{\ell-1}}=\infty.}
\]
For every finite tuple of distinct members $A_1,\ldots,A_r$ of $L(C)$,
and every nonzero polynomial
$F\in\overline{\mathbb Q}[X_1,\ldots,X_r]$ with $\rho(F)\le\ell$,
\[
 F(A_1,\ldots,A_r)\ne0.
\]
The assertion also holds for Laurent polynomials.
\end{theorem}
\begin{proof}
Apply Theorem~\ref{fol:thm:left-boundary-monomials}, using
$L(C)=\widetilde\partial_L W(C)$; multiply by a monomial for
Laurent polynomials.
\end{proof}

\begin{corollary}[Transcendental multiplicative combinations]
\label{f13:bdsup:monomials}
If $\limsup_n c_{n+1}=\infty$, then for distinct
$A_1,\ldots,A_r\in L(C)$ and nonzero
$(u_1,\ldots,u_r)\in\mathbb Q^r$,
\[
 \boxed{A_1^{u_1}\cdots A_r^{u_r}\notin\overline{\mathbb Q}.}
\]
In particular every quotient of two distinct members of $L(C)$ is
transcendental. This applies to the factorial clock $C_n=n!$.
\end{corollary}
\begin{proof}
Clear the denominators of the rational exponents. If the monomial were
algebraic, a Laurent polynomial of the form $X^m-b$ with algebraic $b$
and support rank $1$ would vanish. Apply Theorem~\ref{f13:bdsup:main} with
$\ell=1$. For $C_n=n!$ the multipliers are $c_{n+1}=n+1$.
\end{proof}

\begin{corollary}[Homogeneous relations and ratios]
\label{f13:bdsup:ratios}
If $r\ge2$ and
$\limsup_n c_{n+1}/C_n^{r-2}=\infty$, then for any distinct
$A_1,\ldots,A_r\in L(C)$ the $r-1$ ratios
$A_2/A_1,\ldots,A_r/A_1$ are algebraically independent.
If instead $\limsup_n c_{n+1}/C_n^{r-1}=\infty$, then the $r$ constants
themselves are algebraically independent.
\end{corollary}
\begin{proof}
A polynomial relation among the ratios homogenizes to a nonzero
polynomial in the $A_i$, all of whose exponent vectors lie in a single
affine hyperplane, so its support rank is at most $r-1$.
Apply Theorem~\ref{f13:bdsup:main} with $\ell=r-1$.
An arbitrary polynomial in $r$ variables has support rank at most $r$,
which proves the second assertion.
\end{proof}

\begin{example}[Long square stretches with an independent entire left boundary]
Take $C_1=2$ and recursively define
\[
 c_{n+1}=\begin{cases}
 C_n^n,&n\text{ is a power of }2,\\
 2,&\text{otherwise},
 \end{cases}
 \qquad C_{n+1}=c_{n+1}C_n.
\]
For every fixed $\ell$, the limsup condition of
Theorem~\ref{f13:bdsup:main} holds along the powers of $2$. Therefore the
entire set $L(C)$ is algebraically independent. It is countable and
dense in $W(C)$, as the right endpoints of the complementary intervals
of a closed nowhere-dense perfect set, together with its minimum.
The square stretches are unbounded in length, and $c_{n+1}=2$
infinitely often; the growth limit in Saito--Takeda's Theorem~1.3
therefore does not hold. This example illustrates the sparse-burst
limsup theorem, which is stronger than that published growth-limit
hypothesis.
\end{example}

\begin{remark}[Scope of the improvement]
For the factorial clock, Saito--Takeda's Proposition~8.2 gives rational
linear independence of a pair with the stated boundary orientations.
Corollary~\ref{f13:bdsup:monomials} proves the stronger transcendence of every
pairwise quotient when both endpoints are local left minima, and also
handles all finite nontrivial multiplicative combinations.
It does not prove algebraic independence of arbitrary factorial-clock
boundary tuples: a general polynomial may have support rank greater
than one. Nor does any statement here apply to all arbitrary admissible
constants in $W(C)$. The short approximation is a consequence of the
left-gap property.
\end{remark}


\section{Density-one answers to the cubic Pisot problem}
\label{f13:sec:r2-cubic}

Saito's Problem~1.1 in arXiv:2504.14968v1 asks whether, for every
cubic Pisot number $\beta$, the numbers $\lfloor\beta^{3^n}\rfloor$
are composite infinitely often. The results below give uniform
infinite-family and density statements for this exact exponent
sequence. They do not settle the word ``every'' in that problem,
or Question~1.7 of arXiv:2404.19461v2 about the least Mills constant.

For a totally real cubic Pisot number write
\[
 f(X)=X^3-aX^2+bX-c,
 \qquad T_r=\operatorname{Tr}(\beta^{3^r}).
\]
The two small conjugates have unequal moduli: equality would make
them opposites and hence make their remaining conjugate $\beta$
an integer. Therefore there is a fixed $s\in\{0,1\}$ such that
\[
 \lfloor\beta^{3^r}\rfloor=T_r-s
\]
for all sufficiently large $r$. Here $s=0$ when the dominant small
conjugate is negative and $s=1$ when it is positive.

\begin{lemma}[A uniform small-prime return condition]
\label{f13:lem:r2-return}
Let $p\equiv2\pmod3$ be a prime with $p\nmid c\operatorname{disc}(f)$.
Suppose $a\equiv s\pmod p$ and that $f$ either splits completely
or is irreducible over $\mathbb F_p$. Then the sequence
$\lfloor\beta^{3^r}\rfloor$ has a positive lower density of
composite terms.
\end{lemma}
\begin{proof}
The roots of $f$ lie in $\mathbb F_{p^e}$ with $e=1$ or $3$.
Since $p\equiv-1\pmod3$ and $e$ is odd,
$\gcd(3,p^e-1)=1$. Choose $L\ge1$ with
$3^L\equiv1\pmod{p^e-1}$. All roots are nonzero, so their
$3^{jL}$th powers equal the original roots, for every $j\ge0$.
It follows that $T_{jL}\equiv a\equiv s\pmod p$.
The corresponding floors eventually exceed $p$, and hence are
composite. The progression has density $1/L$.
\end{proof}

\begin{theorem}[Relative density one among totally real cubic Pisot numbers]
\label{f13:thm:r2-density-all}
Let $\mathcal T(B)$ be the finite set of totally real cubic Pisot
numbers $\beta<B$. Let $\mathcal E(B)$ be the subset for which the
composite indices of $\lfloor\beta^{3^r}\rfloor$ have lower density
zero. Then
\[
 \frac{|\mathcal E(B)|}{|\mathcal T(B)|}\longrightarrow0
 \qquad(B\longrightarrow\infty).
\]
In particular, Saito's Problem~1.1 has an affirmative answer for
a relative density-one subset of totally real cubic Pisot numbers
under this ordering.
\end{theorem}
\begin{proof}
Fix $s\in\{0,1\}$ and a prime $p\equiv2\pmod3$ with $p>6$.
Choose ordered triples $(x,y,z)$ of distinct nonzero elements of
$\mathbb F_p$ with $x+y+z=s$. There are $p^2$ ordered choices
before the restrictions; the three zero-coordinate conditions and
three equality conditions exclude at most $6p$ choices. Thus at
least $(p^2-6p)/6$ distinct monic cubics with trace $s$ split into
three distinct nonzero roots. Each unordered root set corresponds
to a unique coefficient pair $(b,c)$. Consequently a proportion at
least
\[
 \delta_p=\frac{p^2-6p}{6p^3}
\]
of all triples $(a,b,c)\pmod p$ meet Lemma~\ref{f13:lem:r2-return}
for this fixed $s$.

For any finite set $S$ of such primes, the Chinese remainder theorem
shows that the proportion of integer coefficient triples in a box
of side lengths comparable to $B$ which avoid every corresponding
return condition tends, as $B\to\infty$, to at most
$\prod_{p\in S}(1-\delta_p)$. The classical divergence
$\sum_{p\equiv2\ (3)}1/p=\infty$ implies that these finite products
tend to zero as $S$ increases. The avoiding triples therefore have
count $o(B^3)$. Applying this separately for $s=0$ and $s=1$ and
taking their union preserves the $o(B^3)$ bound.

Every $\beta<B$ under consideration has
\[
 |a|<B+2,\qquad |b|<2B+1,\qquad 0<|c|<B,
\]
so its polynomial belongs to such a coefficient box. Conversely,
there are $\gg B^3$ distinct such Pisot numbers: for each integer
$2\le a<B-1$, all the polynomials
\[
 X^3-aX^2-dX+c,\qquad d,c\ge1,\quad d+c\le a,
\]
have roots in $(-1,0),(0,1),(a,a+1)$, and their negative small
root has greater modulus. These cubics are irreducible, since an
integer large root would leave an integral quadratic constant term
of nonzero modulus less than one. Their number is
$\sum_{a<B-1}a(a-1)/2\gg B^3$, and distinct irreducible polynomials
give distinct Pisot numbers. Dividing the exceptional $o(B^3)$
bound by this lower bound proves the assertion.
\end{proof}

The argument is a finite-residue sieve followed by a density limit.
It makes no distribution assertion about the individual lacunary
angles of complex conjugates. The theorem is deliberately stated
for totally real cubic Pisot numbers, including the entire real
configuration relevant to an algebraic Mills shadow.

\begin{theorem}[Density one at every fixed nonzero norm]
\label{f13:thm:r2-fixednorm}
Fix a nonzero integer $c$. Among totally real cubic Pisot numbers
$\beta<B$ with norm $c$ and negative dominant small conjugate, a
relative proportion tending to one have a positive lower density
of composite terms $\lfloor\beta^{3^r}\rfloor$.
This applies separately to the unit norms $c=1$ and $c=-1$.
\end{theorem}
\begin{proof}
Take $p\equiv2\pmod3$ with $p\nmid6c$, and impose $a\equiv0\pmod p$.
The residual discriminant is
\[
 \operatorname{disc}(f)\equiv-4b^3-27c^2\pmod p.
\]
Cubing permutes $\mathbb F_p$, so exactly $(p-1)/2$ choices of $b$
make this discriminant a nonzero square. For a separable cubic this
means that Frobenius is an even permutation of the roots; its
factorization type is therefore $1+1+1$ or $3$.
Lemma~\ref{f13:lem:r2-return} applies with $s=0$.
Thus the fraction of good coefficient pairs $(a,b)\pmod p$ is
exactly $(p-1)/(2p^2)$. The same Chinese-remainder and divergent
reciprocal-prime argument gives $o(B^2)$ exceptional pairs in a
box of size $O(B)^2$.

There are $\gg_c B^2$ Pisot numbers in the stated family.
If $c<0$, write $q=-c>0$ and use
$X^3-aX^2-dX+q$ with $1\le d\le a-q$.
If $c>0$, take sufficiently large $a<B-2$ and integer
$3a/4<d<7a/8$. For
$g(X)=X^3-aX^2-dX-c$,
\[
 g(-1)<0,\qquad g(-1/2)>0,\qquad g(0)<0.
\]
Hence its two negative roots are in $(-1,-1/2)$ and $(-1/2,0)$;
the positive root is in $(a,a+2)\subset(1,B)$.
Irreducibility follows as above. Each construction supplies
$\gg_c B^2$ distinct polynomials, which proves the relative-density
claim.
\end{proof}

\begin{theorem}[Density one within Saito's shrinking conjugate locus]
\label{f13:thm:r2-saito-locus}
Fix a nonzero integer $c$. Let $\mathcal S_c(B)$ consist of the totally
real cubic Pisot numbers $\beta<B$ of norm $c$, with conjugates
$\beta,-\rho,\gamma$ satisfying
\[
 0<\sigma=|\gamma|<\rho<1,\qquad
 \rho\le\min(\sigma^{17/23},\beta^{-17/40}).
\]
These are precisely the conjugate inequalities in Saito's
Theorem~1.7 of arXiv:2508.16068v3, restricted to the fixed norm.
Then
\[
 |\mathcal S_c(B)|\gg_c B^{63/40},
\]
and the proportion of members for which
$\lfloor\beta^{3^r}\rfloor$ has a positive lower density of composite
indices tends to one as $B\to\infty$.
\end{theorem}
\begin{proof}
Put $q=|c|\ge1$ and $\alpha=23/40$.
Because $q=\beta\rho\sigma$, the bound
$\rho\le\beta^{-17/40}$ already implies
\[
 \rho^{40}\le\beta^{-17}\le q^{17}\beta^{-17},
 \quad\text{hence}\quad \rho\le\sigma^{17/23}.
\]
For a polynomial $X^3-aX^2+bX-c$ in this locus, the trace and
pairwise-product formulas give
\[
 0<a<B,\qquad |b|\le2\beta\rho+\rho\sigma
                  \le2B^{\alpha}+1.
\]
The lower bound on $a$ follows from the elementary fact that every
totally real cubic Pisot number exceeds $2$; alternatively it is
immediate for the large-$\beta$ portion, and finitely many smaller
polynomials do not affect the density conclusion.

The fixed-norm sieve in Theorem~\ref{f13:thm:r2-fixednorm} works unchanged
in the rectangle of side lengths $B$ and $O(B^\alpha)$: for each
fixed finite set of sieve primes, both side lengths tend to infinity,
so Chinese-remainder residue proportions still have the stated
limits. Increasing the fixed finite set afterward shows that the
number of avoiding pairs in this rectangle is $o(B^{1+\alpha})$.
It remains to match this by a lower bound for actual admissible
polynomials, rather than for the ambient rectangle.

Take large integers $a$, and integers $d$ in the interval
\[
 \frac14a^\alpha\le d\le\frac12a^\alpha,
 \qquad f_{a,d}(X)=X^3-aX^2-dX-c.
\]
Uniformly in these choices, $f_{a,d}(a)<0$ and
$f_{a,d}(a+1)>0$, so there is a root $\beta\in(a,a+1)$.
Write $u=\beta-a$. The equation for this root gives
\[
 u=\frac d\beta+\frac c{\beta^2}
   =\frac da\left(1+O_c(a^{-1})\right).
\]
The other two roots have sum $-u$ and product $c/\beta$, so their
quadratic discriminant is $u^2-4c/\beta$.
As $d^2/a\asymp a^{3/20}\to\infty$, it is positive for all
sufficiently large $a$, for either sign of $c$. Calling the larger
small-root modulus $\rho$ and the smaller one $\sigma$, the
quadratic formula and the norm identity give, uniformly,
\[
 \rho=\frac da\left(1+O_c(a^{-3/20})\right),\qquad
 \sigma=\frac qd\left(1+O_c(a^{-3/20})\right).
\]
In particular the dominant small root is negative,
$\sigma/\rho\to0$, and both small roots have modulus less than one.
The cubic is irreducible by the integer-root argument used above.
Moreover,
\[
 \rho\beta^{17/40}
 =\frac d{a^{23/40}}\left(1+O_c(a^{-3/20})\right)
 \le\frac12+o(1)<1.
\]
Thus all these polynomials belong to the required locus.
Taking $a$ between $B/2$ and $B-2$ gives
$\gg_c B^{1+\alpha}=B^{63/40}$ distinct members. The exceptional
$o(B^{63/40})$ count from the same shrinking rectangle now proves
the relative-density assertion.
\end{proof}

This strengthens the density conclusion inside the actual shrinking
conjugate region allowed by Saito's Mills reduction. It still does not
identify the hypothetical shadow of the distinguished least Mills
constant: a density-zero exceptional family has not been shown empty.

\section{An exact valuation theorem in the surviving unit branch}
\label{f13:sec:r2-valuations}

The preceding density theorems leave a possible thin exceptional
set. The next calculation restricts an important part of that set
without assuming any prime-distribution conjecture.

\begin{theorem}[Exact linear $3$-adic depth]
\label{f13:thm:r2-depth}
Let $\varepsilon\in\{1,-1\}$ and let
\[
 f(X)=X^3-aX^2+bX+\varepsilon
\]
be an irreducible cubic, with
$a\equiv\varepsilon\pmod3$ and $b\equiv-1\pmod3$.
Write $T_r$ and $B_r$ for the trace and second elementary symmetric
function of the $3^r$th powers of its roots, and put
\[
 s=v_3(a+\varepsilon b).
\]
Then $s$ is a finite positive integer and
\[
 \boxed{v_3(T_r-\varepsilon)=r+s\qquad(r\ge s).}
\]
\end{theorem}
\begin{proof}
The norm of the roots is $c=-\varepsilon$, which is unchanged by
odd powering. Put
\[
 x_r=T_r-\varepsilon,\quad y_r=B_r+1,\quad
 z_r=x_r+\varepsilon y_r=T_r+\varepsilon B_r,\quad
 w_r=x_r-\varepsilon y_r=T_r-\varepsilon B_r-2\varepsilon.
\]
Newton's identities, expanded at the fixed state
$(T,B,c)=(\varepsilon,-1,-\varepsilon)$, give the exact factorizations
\[
\begin{split}
 z_{r+1}&=z_r\bigl(3+3\varepsilon w_r
                   +x_r^2-\varepsilon x_ry_r+y_r^2\bigr),\\
 w_{r+1}&=w_r\bigl(9+3\varepsilon w_r
                   +x_r^2+\varepsilon x_ry_r+y_r^2\bigr).
\end{split}
\]
The residue conditions ensure $3\mid x_r,y_r$ at every index.
The first multiplier is $3\pmod9$ and the second is divisible by
$9$. Therefore
\[
 v_3(z_r)=r+s,\qquad v_3(w_r)\ge2r+1.
\]
The initial $z_0=a+\varepsilon b$ cannot vanish, because
$f(-\varepsilon)=-a-\varepsilon b$ and $f$ is irreducible.
Finally, $2x_r=z_r+w_r$. For $r\ge s$ the first summand has
strictly smaller valuation, proving the formula.
\end{proof}

\begin{lemma}[Projective return, recalled with proof]
\label{f13:lem:r2-projective}
Let $p\nmid c\operatorname{disc}(f)$ divide $T_r$, and let
$Y_1,Y_2,Y_3$ be the residual roots raised to $3^r$ in a finite
splitting field. If the subgroup generated by their ratios has
order coprime to $3$, then $p\mid T_{r+jL}$ for every $j\ge0$
and some $L\ge1$.
\end{lemma}
\begin{proof}
Choose $L$ with $3^L\equiv1$ modulo the exponent of that subgroup.
For every $j$, the three $3^{jL}$th powers of the $Y_i$ differ
from the $Y_i$ by a common nonzero scalar. Their sum remains zero.
\end{proof}

\begin{theorem}[Complete splitting forced in a positive-sign unit branch]
\label{f13:thm:r2-splitbranch}
Let $\beta$ be a totally real cubic Pisot number with negative
dominant small conjugate and polynomial
\[
 f(X)=X^3-aX^2+bX+1,
 \qquad a\equiv1\pmod3,\quad b\equiv-1\pmod3,
 \quad v_3(a+b)=1.
\]
If $\lfloor\beta^{3^r}\rfloor$ is eventually prime, then, for every
sufficiently large $r$, $f$ splits completely modulo that prime.
More strongly, one sufficiently late prime trace at which $f$
does not split completely forces a positive-density progression
of composite floors.
\end{theorem}
\begin{proof}
For large $r$, the floor equals $T_r$ and exceeds every fixed
prime divisor of $3\operatorname{disc}(f)$. Let $p=T_r$ be prime.
Theorem~\ref{f13:thm:r2-depth} gives
$p\equiv1\pmod3$ and $v_3(p-1)=r+1$.
If $f$ is irreducible modulo $p$, its roots have the form
$u,u^p,u^{p^2}$. Their ratios have orders dividing $p^2+p+1$,
whose $3$-adic valuation is exactly one. Raising to $3^r$ with
$r\ge1$ kills this $3$-primary part, so
Lemma~\ref{f13:lem:r2-projective} gives a composite progression.

If the factorization type is $1+2$, write the roots as
$u\in\mathbb F_p$ and $v,v^p\in\mathbb F_{p^2}$.
Since $v_3(p^2-1)=r+1$, the $3$-primary components of their
$3^r$th powers are in $\mu_3$. The last two components are equal,
since $p\equiv1\pmod3$. Their product with the first component
is the $3$-primary component of the norm $-1$, namely one.
Thus all three components are equal, and every root ratio has
order coprime to $3$. Projective return again applies.
An eventually prime tail excludes both nonsplit cases.
\end{proof}

\begin{corollary}[A positive-sign discriminant-cycle obstruction]
Under the hypotheses of Theorem~\ref{f13:thm:r2-splitbranch}, suppose
that $\operatorname{disc}(f)$ is not a rational square. Let $D$ be
the fundamental discriminant of its quadratic resolvent field.
Iterate the cubic Newton state $(T_r,B_r,-1)$ modulo $|D|$.
If any state of its eventual cycle has
\[
 \chi_D(T_r)\ne1,
\]
then the floor sequence has a positive lower density of composite
terms. Here $\chi_D$ is the quadratic Dirichlet character.
\end{corollary}
\begin{proof}
A cycle state occurs on an infinite arithmetic progression of
indices. If its character is zero, a fixed prime divisor of $D$
divides every trace on that progression. If the character is $-1$,
then either the progression is eventually composite, or it
contains a sufficiently late prime trace. At such a prime,
Frobenius has odd sign, so the factorization type is $1+2$.
Theorem~\ref{f13:thm:r2-splitbranch} gives a composite progression.
\end{proof}

The exact valuation theorem extends the earlier character sieve
into a positive-sign, norm-$-1$ branch; it does not assert that
the cycle obstruction always occurs. A surviving orbit could
have character $1$ at every cycle state, and completely split
prime traces are not ruled out. Chebotarev's theorem supplies
many primes of other splitting types, but does not force one of
those primes to divide this prescribed sparse trace orbit after
its $3$-primary transient. Proving that such an appropriate
recurrent divisor always exists remains the precise missing
step in this approach to Saito's full Problem~1.1.


\section{Square clocks: rational-function avoidance and interior gap tests}
\label{f13:sec:square-second-round}

Saito's Question 1.6 asks whether the least $A>1$ for which every
$\lfloor A^{2^k}\rfloor$ is prime is irrational.  The results of this
section do not resolve that question.  They give unconditional auxiliary
theorems that restrict a rational candidate and an additional
transcendence test distinct from the previously considered small-gap
condition.  Their deep inputs are established approximation theorems,
so these deductions are not presented as a new transcendence method or
as priority claims.

Write, for any $A>1$,
\[
 p_k=\lfloor A^{2^k}\rfloor,\qquad
 u_k=A^{2^k}-p_k\in[0,1),\qquad
 D_k=p_{k+1}-p_k^2.
\]
The exact relation
\begin{equation}\label{f13:eq:square-defect-identity}
 D_k=2p_ku_k+u_k^2-u_{k+1}
\end{equation}
will be used repeatedly.  Primality is not needed for the first two
theorems.

\subsection{Rational candidates cannot obey a rational-function law}

\begin{theorem}[Finite rational-function coincidences]
\label{f13:thm:square-rational-avoidance}
Let $A>1$ be rational but not an integer.  For every rational function
$R\in\mathbb Q(X)$, there are only finitely many $k$ at which
\[
 p_{k+1}=R(p_k).
\]
Consequently, for every fixed finite collection of rational functions,
all sufficiently late consecutive pairs $(p_k,p_{k+1})$ avoid all
their graphs.
The same statement holds with $2^k$ replaced by $m^k$ for any fixed
integer $m\ge2$.
\end{theorem}

\begin{proof}
Suppose there are infinitely many coincidences.  Since $p_k\to\infty$
and $p_{k+1}/p_k^2\to1$, the rational function must have the expansion
\[
 R(X)=X^2+bX+O(1)\qquad(X\to+\infty),\qquad b\in\mathbb Q.
\]
Equation~\eqref{f13:eq:square-defect-identity} then gives, at the successful
indices,
\[
 u_k=b/2+O(p_k^{-1}).
\]
Choose a positive integer $q$ with $qb/2\in\mathbb Z$.  We obtain
\begin{equation}\label{f13:eq:square-rational-nearest}
 \|qA^{2^k}\|=O(A^{-2^k}).
\end{equation}
The Corvaja--Zannier--Dubickas approximation criterion, in the form
of Saito's Theorem 2.6 in arXiv:2404.19461v2, implies that some power
$A^{2^j}$ at a successful exponent is Pisot.  A rational Pisot number
is an integer, whereas no positive integer power of a nonintegral
rational number is an integer.  This is a contradiction.  Equivalently,
one can use the singleton rational-base consequence of the
Kulkarni--Mavraki--Nguyen theorem.

For $m\ge2$, the asymptotic ratio forces
$R(X)=X^m+bX^{m-1}+O(X^{m-2})$.  Expanding
$(p_k+u_k)^m$ gives $u_k=b/m+O(p_k^{-1})$, and the same proof applies.
The finite-collection assertion follows by taking a finite union of
the exceptional index sets.
\end{proof}

\begin{corollary}[A necessary irregularity for a rational square constant]
If a rational $A>1$ has prime outputs $p_k=\lfloor A^{2^k}\rfloor$
for every $k\ge1$, its output sequence cannot satisfy even infinitely
many instances of any one fixed rational-function recurrence.
In particular it cannot eventually follow a finite menu of rational
polynomial recurrences.
\end{corollary}

\begin{proof}
An integral $A>1$ has composite outputs and is impossible.  Apply
Theorem~\ref{f13:thm:square-rational-avoidance}.
\end{proof}

\subsection{Transcendence from concentration inside the square interval}

\begin{theorem}[Interior algebraic-slope criterion]
\label{f13:thm:square-interior-slope}
Suppose that for some real algebraic number $0<\lambda<2$, some
$\varepsilon>0$, and some $C>0$, infinitely many $k$ satisfy
\begin{equation}\label{f13:eq:square-interior-condition}
 |D_k-\lambda p_k|\le C p_k^{1-\varepsilon}.
\end{equation}
Then $A$ is transcendental.  In particular,
\[
 D_k=p_k+O(p_k^{1-\varepsilon})
 \quad\text{on an infinite subsequence}
 \quad\Longrightarrow\quad A\text{ is transcendental}.
\]
\end{theorem}

\begin{proof}
Put $\eta=\min\{\varepsilon,1\}>0$.  From
\eqref{f13:eq:square-defect-identity}, since $u_k,u_{k+1}\in[0,1)$, we
obtain
\[
 u_k=\lambda/2+O(p_k^{-\eta}).
\]
Thus on the successful subsequence
\begin{equation}\label{f13:eq:square-translated-approx}
 \bigl\|A^{2^k}+1-\lambda/2\bigr\|
 =O(A^{-\eta2^k}).
\end{equation}
Assume that $A$ is algebraic.  If $\lambda$ is irrational, then
$1-\lambda/2$ is an algebraic irrational in $(0,1)$; Kumar's
Theorem 2.1 in arXiv:2207.08614v3 excludes infinitely many such
approximations, with singleton base $A$ and fixed coefficient $1$.
If $\lambda$ is rational, the shift is a nonintegral rational number;
Kumar's Theorem 2.2 applies because the coefficient $1$ is an
algebraic integer.  In both cases select
$A^{-\eta}<\theta<1$ to absorb the implied constant into the
required bound $\theta^{2^k}$.  This contradicts
\eqref{f13:eq:square-translated-approx} and proves transcendence.
\end{proof}

\begin{corollary}[Finite interior schedules]
Let $\Lambda\subset(0,2)$ be a finite nonempty set of real algebraic
numbers.  If, for some $\varepsilon>0$ and $C>0$, infinitely many $k$
satisfy
\[
 \min_{\lambda\in\Lambda}|D_k-\lambda p_k|
 \le C p_k^{1-\varepsilon},
\]
then $A$ is transcendental.
\end{corollary}
\begin{proof}
One member of the finite set is successful infinitely often.  Apply
Theorem~\ref{f13:thm:square-interior-slope}.
\end{proof}

This criterion concerns concentration near an interior fraction of
the interval between consecutive squares.  It does not require
$D_k=o(p_k)$, and therefore is distinct from the previous small-gap
criterion.  The power-saving rate in
\eqref{f13:eq:square-interior-condition} is essential to this proof:
convergence $D_k/p_k\to\lambda$ without a quantitative rate is not
enough to invoke the approximation theorem.

\subsection{The unique algebraic polynomial-law exception}

For this classification we use Saito's Theorem~2.6 and Lemma~2.7
in arXiv:2404.19461v2: exponentially close algebraic powers produce
a Pisot power, and the sum of the smaller conjugate powers has a
lower bound at its largest conjugate modulus, up to a polynomial
factor. The argument below supplies the required quadratic
classification directly.

\begin{theorem}[Classification under an eventual polynomial law]
\label{f13:thm:square-polynomial-exception}
Suppose every $p_k=\lfloor A^{2^k}\rfloor$ is prime and, for some
$P\in\mathbb Q[X]$, one has $p_{k+1}=P(p_k)$ for all sufficiently
large $k$.  Then exactly the following alternatives apply:
\begin{enumerate}
\item $P(X)\ne X^2+2X-2$, and $A$ is transcendental;
\item $P(X)=X^2+2X-2$, and for every sufficiently large starting
index $K$,
\begin{equation}\label{f13:eq:square-polynomial-value}
 A^{2^K}=\frac{p_K+1+\sqrt{(p_K+1)^2-4}}2.
\end{equation}
In particular $A$ is algebraic and irrational in this alternative.
\end{enumerate}
Thus in every case $A$ is irrational.
\end{theorem}

\begin{proof}
The asymptotic ratio $p_{k+1}/p_k^2\to1$ forces
$P(X)=X^2+bX+c$, with $b,c\in\mathbb Q$. By
\eqref{f13:eq:square-defect-identity},
\[
 u_k=b/2+O(p_k^{-1}),
\]
so $0\le b\le2$. If $0<b<2$, the interior-slope theorem already
makes $A$ transcendental. Suppose instead that $A$ is algebraic.
Then $b\in\{0,2\}$, and
\[
 \|A^{2^k}\|=O(A^{-2^k}).
\]
Saito's Theorem~2.6 gives a Pisot number $\beta=A^{2^K}$ for some
$K$. We may increase $K$ freely, since a positive integer power of
a Pisot number is again Pisot. If $\beta$ is an integer, then later
outputs are proper integer powers and are composite, a contradiction.
Thus its degree $d$ is at least two.

Let $\beta_2,\ldots,\beta_d$ be its other conjugates and put
$\rho=\max_{i\ge2}|\beta_i|<1$. The trace of $\beta^m$ is an
integer and differs from $\beta^m$ by a quantity tending to zero.
At the exponents $m=2^j$, this trace is therefore the nearest integer
for all sufficiently large $j$, and the preceding estimate gives
\[
 \left|\sum_{i=2}^d\beta_i^{2^j}\right|=O(\beta^{-2^j}).
\]
Saito's Lemma~2.7 supplies, for some fixed $v>0$, the lower bound
$\rho^{2^j}(2^j)^{-v}$. Taking $2^j$th roots yields
$\rho\le\beta^{-1}$. The nonzero integral norm consequently satisfies
\[
 1\le |N_{\mathbb Q(\beta)/\mathbb Q}(\beta)|
    \le\beta\rho^{d-1}\le\beta^{2-d}.
\]
It follows that $d=2$ and that the norm is $\pm1$. After increasing
$K$ once, the conjugate is positive and equals $\beta^{-1}$.
The squared Pisot number still has degree two: it cannot be rational,
since its conjugate would then equal itself despite having modulus
less than one.

The integer traces now satisfy
\[
 T_j=\beta^{2^j}+\beta^{-2^j}\in\mathbb Z,\qquad
 T_{j+1}=T_j^2-2.
\]
Since $0<\beta^{-2^j}<1$, the exact floors are
$p_{K+j}=T_j-1$. Hence
$p_{K+j+1}=p_{K+j}^2+2p_{K+j}-2$.
The labels tend to infinity, so equality at infinitely many distinct
arguments forces the polynomial identity $P(X)=X^2+2X-2$.

Conversely, suppose this polynomial law holds from $K$ onward.  Put
$t=p_K+1\ge3$ and $\alpha=(t+\sqrt{t^2-4})/2$.
Then $\alpha>1$, $\alpha+\alpha^{-1}=t$, and the same trace
recurrence proves by induction that
$p_{K+n}=\alpha^{2^n}+\alpha^{-2^n}-1$.
Taking growth limits gives \eqref{f13:eq:square-polynomial-value}.
The discriminant $t^2-4$ is not a square for $t\ge3$, so this
$\alpha$ is irrational; $A$ therefore cannot be rational.
\end{proof}

The exceptional polynomial has not been shown to possess an infinite
all-prime orbit here.  The theorem identifies the only possible
algebraic law under the stated hypothesis; it does not establish the
existence of that law for Saito's least square constant.

\subsection{What still prevents a solution of Question 1.6}

For a rational candidate $A=a/b>1$ in lowest terms, $b>1$, the elementary
lower bound
\[
 \{A^{2^k}\}\ge b^{-2^k}
\]
is too small to contradict the floor inequalities.  Neither primality
nor minimality has been proved here to supply the exponentially accurate
translated approximations required above.  Likewise, the Type A/B
theorems in arXiv:2508.16068v3 require a ratio limsup greater than
$40/19$, while the square clock has ratio exactly $2$.
No argument in this section supplies that missing input.  The
unconditional irrationality of the least square prime constant remains
unresolved by the present work.


\section{Square clocks: a dimension gap and the remaining computability obstruction}
\label{f13:sqnew:section}

The global dimension results for unbounded multipliers and iterated
exponential clocks do not imply local dimension statements for a fixed
finite prefix. The present question is a local strengthening of Saito's dimension theorem
and Saito--Takeda's perfectness theorem: what happens when every sufficiently
late multiplier is exactly $2$? The extra hypothesis of infinitely many
multipliers at least $3$ in Saito--Takeda's Theorem~1.2 does not apply.

Fix $E_0>0$, put $E_n=E_0 2^n$ for $n\ge0$, and define
\[
 W_2(E_0)=\{A>1:\lfloor A^{E_n}\rfloor\in\mathbb P
                 \text{ for every }n\ge0\}.
\]
This set is nonempty, closed, and has Hausdorff dimension $1$, by the
published existence, closedness, and dimension theorems. None of those
facts alone implies the full local dimension asserted below outside a
small exceptional set.

A prime label $p$ at exponent $E$ has the half-open root cell
$[p^{1/E},(p+1)^{1/E})$. For an admissible square-clock constant all
powers $A^{E_n}$ are nonintegers: an integral value would have a
composite square at the next index. Thus every admissible point lies
in the interior of each of its cells. A valid later prime cell has
closure strictly inside its parent: if $q$ follows $p$, then
$p^2<q$ and $q+1<(p+1)^2$, since both $p^2$ and $(p+1)^2-1$ are
composite. We use half-open cells for finite prefix conditions and
closed, strictly nested descendant cells for limit constructions.

\subsection{Good square seeds and a sparse exceptional set}

Matom\"aki's prime-interval estimate \cite[Lemma~9]{Matomaki2010} gives
constants $d,D>0$ such that,
for sufficiently large $x$, at most $Dx^{1/6}$ disjoint intervals
$[u,u+u^{1/2}]\subset[x,2x]$ contain at most
$d u^{1/2}/\log u$ primes. Choose a fixed positive integer $K$ sufficiently
large, and set
\[
 H(p)=\left\lceil\frac{p}{K\lceil\log_2 p\rceil}\right\rceil.
\]
For sufficiently large $p$, this is at most $dp/(2\log p)$.
Call a prime $p$ \emph{good} if $[p^2,p^2+p]$ contains at least $H(p)$
primes, and \emph{poor} otherwise. A sufficiently large threshold $P$
will be fixed throughout. Both predicates are decidable by finite integer
arithmetic and primality tests; numerical values of the universal
constants are not evaluated in this note. In every computability or
semidecision statement below, fixed integers $K,P$ with the stated
properties are hardcoded into the algorithm. Their existence proves
ordinary computability, even though this argument does not supply their
numerical values or a practical digit-computation procedure.

\begin{lemma}[Few poor prime seeds]\label{f13:sqnew:poorcount}
As $X\to\infty$,
\[
 \#\{p\in\mathbb P\cap[X,2X]:p\text{ is poor}\}\ll X^{1/3}.
\]
\end{lemma}
\begin{proof}
The intervals $[p^2,p^2+p]$ for distinct primes $p$ are disjoint. For
$p\in[X,2X]$, they lie in $[X^2,4X^2+2X]$. Cover this range by a fixed
number of dyadic intervals. At each dyadic boundary at most one of the
disjoint intervals straddles the boundary; discard these finitely many
intervals. All remaining poor intervals are exceptional intervals in
Matom\"aki's estimate, with the squared scale $x\asymp X^2$.
The bound is consequently $O((X^2)^{1/6})=O(X^{1/3})$.
\end{proof}

\begin{lemma}[A good seed has full-dimensional continuations]
\label{f13:sqnew:goodfull}
For every good prime $p\ge P$ and every $E>0$, the square-continuation set
\[
 \{A\in(p^{1/E},(p+1)^{1/E}):
      \lfloor A^{E2^j}\rfloor\in\mathbb P\ (j\ge0)\}
\]
has Hausdorff dimension $1$. If $E$ is computable, it contains a computable element and a
computably embedded binary Cantor set.
\end{lemma}
\begin{proof}
Among the at least $H(p)$ prime children
$q\in[p^2,p^2+p]$, at most $Dp^{2/3}$ are poor, by the same exceptional
interval argument at the scale $p^4$. Thus, for sufficiently large $p$,
there remain at least $a p/\log p$ good children, with a fixed $a>0$.
Their root cells are strictly inside the parent cell: $q>p^2$ because
$p^2$ is composite, and $q+1\le p^2+p+1<(p+1)^2$.
Iterating gives a tree of good primes satisfying the hypotheses of
Saito's Proposition~5.1 \cite{SaitoDimension2021} with $R=2$ and $L=1$.

At a descendant prime $q$, that proposition bounds the dimension of its
continuation subtree below by
$(1+2/(q\log q))^{-1}$, after the appropriate power change of variable.
Every such subtree lies in the original seed cell, and descendant labels
tend to infinity. Taking these lower bounds to $1$ proves the claim.
Choosing the least good child at each step gives a computable limit.
Choosing the first or second good child according to a binary sequence
gives the stated computable embedding; the separated child cells give a
computable inverse on its image.
\end{proof}

\begin{theorem}[An exceptional set of dimension at most one third]
\label{f13:sqnew:exceptionaldimension}
Let $(F_n)$ be any real exponent sequence with $F_n\ge c2^n$ for a
constant $c>0$. The set
\[
 \mathcal E_F=
 \{A>1:\lfloor A^{F_n}\rfloor\text{ is a poor prime for infinitely
 many }n\}
\]
has Hausdorff dimension at most $1/3$.
In particular, outside a set of dimension at most $1/3$, every
sufficiently late prime output of a square clock is good.
\end{theorem}
\begin{proof}
Fix a compact interval $[a,b]\subset(1,\infty)$ and $s>1/3$.
For large $n$, the prime root cell corresponding to a label $p$ that
meets $[a,b]$ has diameter at most
\[
 \frac{B}{F_n p}
\]
for a constant $B$ depending only on $b$; this follows from the mean value
theorem, since $p^{1/F_n}\le b$. Such labels satisfy
$p\ge a^{F_n}-1$. From Lemma~\ref{f13:sqnew:poorcount} and a dyadic sum,
\[
 \sum_{\substack{p\ge Y\\p\ \text{poor}}}p^{-s}
 \ll_s Y^{1/3-s}.
\]
Therefore the sum of the $s$th powers of diameters of all poor-label
cells meeting $[a,b]$ at level $n$ is
\[
 \ll_{s,a,b} F_n^{-s}a^{F_n(1/3-s)}.
\]
This is summable in $n$, because $F_n\ge c2^n$ and $s>1/3$.
For every $N$, the union of these covers over $n\ge N$ covers the
limsup set. The total $s$-content tends to zero as $N\to\infty$, proving
that its $s$-dimensional Hausdorff measure is zero. Let $s\downarrow1/3$
and use a countable cover of $(1,\infty)$ by compact intervals.
\end{proof}

\subsection{A local dimension gap and a canonical perfect part}

\begin{theorem}[No intermediate local dimensions]\label{f13:sqnew:dichotomy}
For every open interval $U\subset(1,\infty)$,
\[
 \boxed{\dim_H(U\cap W_2(E_0))\le\tfrac13
       \quad\text{or}\quad
       \dim_H(U\cap W_2(E_0))=1.}
\]
More precisely, if $A\in W_2(E_0)$ has good prime outputs at infinitely
many indices, then every neighborhood $U$ of $A$ satisfies
$\dim_H(U\cap W_2(E_0))=1$.
\end{theorem}
\begin{proof}
The half-open root cell of an output at level $n$ respects every earlier
prime condition, because the ratios $E_n/E_j=2^{n-j}$ are integers.
The point $A$ lies in its interior, and these cells shrink to $A$. If good outputs occur at
arbitrarily late indices, a good seed cell is contained in any prescribed
neighborhood of $A$, and Lemma~\ref{f13:sqnew:goodfull} applies.

Suppose $\dim_H(U\cap W_2(E_0))>1/3$. By
Theorem~\ref{f13:sqnew:exceptionaldimension}, the set $U\cap W_2(E_0)$ cannot
be contained in the exceptional limsup set. It thus contains a point
whose outputs are all good from some index onward. The preceding
argument gives a full-dimensional subset inside $U$.
\end{proof}

\begin{theorem}[A perfect part with full local dimension]
\label{f13:sqnew:regularpart}
Define
\[
 R=\{A\in W_2(E_0):
       \dim_H(U\cap W_2(E_0))=1\text{ for every open }U\ni A\}.
\]
Then $R$ is a nonempty closed perfect set, every open interval meeting
$R$ intersects $R$ in dimension $1$, and
\[
 \dim_H(W_2(E_0)\setminus R)\le\tfrac13.
\]
Every point of $W_2(E_0)\setminus R$ has only poor outputs from some
index onward.
If $E_0$ is computable, $R$ has a dense sequence of computable elements,
and every nonempty relative open subset of $R$ contains every Turing degree.
\end{theorem}
\begin{proof}
The complement of $R$ in $W_2(E_0)$ is relatively open. It is a union
of countably many relative intersections with rational open intervals
having dimension less than $1$, hence at most $1/3$ by
Theorem~\ref{f13:sqnew:dichotomy}. Countable stability of dimension proves
the bound. Since $W_2(E_0)$ is closed, $R$ is closed; since the full set
has dimension $1$, $R$ is nonempty and has dimension $1$.
If $U$ meets $R$, then $U\cap W_2(E_0)$ has dimension $1$. Removing the
set of dimension at most $1/3$ leaves $\dim_H(U\cap R)=1$.
In particular $R$ has no isolated point. If a point outside $R$ had
infinitely many good outputs, the preceding theorem would place it in
$R$, which proves the assertion about eventual poor outputs.

Assume $E_0$ is computable. Enumerate all finite valid square-prime
prefixes ending at a good prime at least $P$, and for each prefix choose
the computable continuation that always takes the least good child.
Every enumerated point belongs to $R$. Conversely every open interval
meeting $R$ has dimension $1$ in $W_2(E_0)$, hence contains a point
outside the exceptional set, and consequently contains a complete valid
good seed cell. An enumerated continuation lies in this cell. The
sequence is therefore dense in $R$. The binary construction in
Lemma~\ref{f13:sqnew:goodfull}, inside the same cell, lies in $R$ and
preserves Turing degrees.
\end{proof}

\begin{corollary}[A semidecision procedure for full-dimensional prefixes]
\label{f13:sqnew:semidecision}
Suppose $E_0$ is computable. Given a finite valid square-prime prefix,
there is an algorithm that halts exactly when its completion set has
Hausdorff dimension $1$. It searches for a later valid descendant with
a good prime label at least $P$.
There is also an algorithm that halts exactly when the completion set
is empty. Thus the only unresolved output of the combined search is a
nonempty completion set of dimension at most $1/3$.
\end{corollary}
\begin{proof}
Reachability of a good descendant is a search through a finitely
branching, decidable tree of prime labels. If a good descendant is found,
Lemma~\ref{f13:sqnew:goodfull} proves dimension $1$. If none exists, every
infinite branch is eventually poor and its real limit belongs to the
exceptional set, so the dimension is at most $1/3$.
For emptiness, inspect successive finite levels of the prefix tree.
An empty level certifies emptiness. Conversely, if all levels are
nonempty, K\"onig's lemma gives an infinite branch, whose strictly
nested prime root cells define a genuine continuation. Hence an empty
completion set eventually has an empty finite level.
\end{proof}

\subsection{Isolated points and the least square constant}

\begin{proposition}[Every isolated square constant is computable]
\label{f13:sqnew:isolatedcomputable}
If $E_0$ is computable, every isolated point of $W_2(E_0)$ is computable.
Every such point has only poor prime outputs from some index onward.
\end{proposition}
\begin{proof}
The eventual-poor assertion follows from
Theorem~\ref{f13:sqnew:regularpart}. For computability, choose a finite prime
prefix cell whose intersection with $W_2(E_0)$ is the isolated singleton.
Its tree of finite extensions is decidable and finitely branching, with
exactly one infinite path. For any requested level $n$, search for a
later finite level at which all surviving nodes have the same ancestor
at level $n$. Such a level exists: otherwise a second infinite path
would exist by K\"onig's lemma. The common ancestor determines a root
cell of length tending effectively to zero with $n$, which computes the
singleton. This proves existence of an algorithm for each isolated
point; it does not give an algorithm deciding which prefixes isolate a
point.
\end{proof}

\begin{proposition}[Where noncomputability of the least square constant could hide]
\label{f13:sqnew:minimumcriterion}
Suppose $E_0$ is computable, and let $\xi_2=\min W_2(E_0)$.
If $\xi_2\in R$, then $\xi_2$ is computable. Consequently, if $\xi_2$
is not computable, it must be nonisolated, must have eventually only
poor outputs, and must have a neighborhood $U$ for which
\[
 0\le\dim_H(U\cap W_2(E_0))\le\tfrac13.
\]
No unconditional computability or irrationality claim for $\xi_2$ is
made here.
\end{proposition}
\begin{proof}
The least lower endpoint $\ell_n$ among all finite valid prefix cells
through level $n$ is computable by enumerating final primes and testing
all earlier outputs by exact integer roots. The numbers $\ell_n$
increase and are bounded above by $\xi_2$.
Let $\ell=\lim\ell_n$. We verify $\ell\in W_2(E_0)$, including the
possible endpoint issue. If $\ell^{E_j}=h\in\mathbb N$, then $h\ge2$.
If the approximants eventually equal $\ell$, the later output
$h^2$ is composite, contrary to their finite validity. Otherwise they
approach $\ell$ strictly from below; at level $j+2$ their outputs are
then eventually $h^4-1$, also composite. Thus $\ell^{E_j}$ is never an
integer. Its floor is eventually the corresponding prime floor of
$\ell_n$, so every output at $\ell$ is prime. Hence $\ell=\xi_2$.
This gives lower semicomputability.

The computable points enumerated densely in $R$ in
Theorem~\ref{f13:sqnew:regularpart} all lie above $\xi_2$. If $\xi_2\in R$,
their infimum is $\xi_2$, giving upper semicomputability by a running
minimum of rational upper approximations. The two semicomputabilities
make $\xi_2$ computable. If it is not computable, it lies outside $R$,
and Theorem~\ref{f13:sqnew:regularpart} gives the eventual-poor and local
dimension conclusions. Proposition~\ref{f13:sqnew:isolatedcomputable}
excludes isolation.
\end{proof}

\begin{remark}[What remains unresolved]
These arguments do not prove $R=W_2(E_0)$, do not rule out isolated
square constants, and do not show that every feasible finite square
prefix has an infinite continuation. The possible exceptional region
is smaller in dimension, but can still contain distinguished minima.
The passage from a sparse exceptional set of prime labels to the
absence of an infinite exceptional chain is precisely the missing step.
No new algebraic-independence assertion about the boundary elements in
Saito--Takeda's Question~1.1 follows from this dimension argument.
\end{remark}

\begin{remark}[A general quantitative transfer]
The exponent $1/3$ in the exceptional-dimension and local-gap conclusions
comes only from the counting estimate in Lemma~\ref{f13:sqnew:poorcount}.
If that estimate is replaced by $O_\varepsilon(X^{\tau+\varepsilon})$
for every $\varepsilon>0$, with $\tau<1/2$, the same proofs replace
$1/3$ by $\tau$. The restriction $\tau<1/2$ ensures that among the
$\gg p/\log p$ candidate children of a good seed, only
$O_\varepsilon(p^{2\tau+2\varepsilon})=o(p/\log p)$ are poor.
This is a transfer principle, not a claim of a new prime-distribution
estimate.
\end{remark}


\section{Logarithmic prime-free windows for non-reversible clocks}
\label{f13:sec:nonrev-log-windows}

The previously established return theorem for an exponential polynomial
plus a nonzero linear drift gives composite windows of every fixed
radius.  Here the radius grows logarithmically with the index.  This is
an extension of the quantitative conclusion of Saito's prime-free-window
theorem to this particular non-reversible family; it is not a second
resolution of his already treated Problem~1.7.

We use the following previously proved arithmetic input, with all its
quantifiers stated explicitly.  It is recalled rather than counted as
a new result.
\begin{lemma}[Universal return input]\label{f13:lem:window-return-input}
Let $a\ge2$, $b\ne0$, $e\in\mathbb Z$, and $P\in\mathbb Z[X]$.
For $R(n)=P(a^n)+bn+e$, there is $m_0$ such that, for every $m\ge m_0$
and every positive integer $q$, the congruence
$R(n)\equiv R(m)\pmod q$ has arbitrarily large solutions $n$.
One may take $m_0=1$ when $|b|=1$ and
$m_0=\max(1,\lceil2\log_2|b|\rceil)$ otherwise.
\end{lemma}
The input, including the displayed seed threshold, is proved in
Theorem~\ref{sur:thm:clock-return}. The quantitative argument below
does not require a constructive bound from that proof; eventual
periodicity supplies the needed bound independently.

\begin{lemma}[A uniform bound on a recurrent index]
\label{f13:lem:window-return-bound}
Under the hypotheses of Lemma~\ref{f13:lem:window-return-input}, fix
$m\ge m_0$ and $q\ge2$.  There is a solution of
$R(n)\equiv R(m)\pmod q$ satisfying
\[
 m<n\le \max\{m+1,\lceil\log_2q\rceil\}+q^2.
\]
\end{lemma}
\begin{proof}
Write $q=uv$, with every prime divisor of $u$ dividing $a$ and
$\gcd(v,a)=1$.  For $n\ge\lceil\log_2q\rceil$ one has $a^n\equiv0$
modulo $u$, while $a^n$ is periodic modulo $v$ with period
$t=\operatorname{ord}_v(a)$ (take $t=1$ if $v=1$).  Thus $R(n)$ is
periodic from that point with period
$\operatorname{lcm}(q,t)\le q\varphi(v)\le q^2$.
Universal return ensures that a solution occurs in the eventual cycle.
Every interval of one full period within that cycle contains one.
Apply this to the interval starting at
$\max\{m+1,\lceil\log_2q\rceil\}$.
\end{proof}

\begin{theorem}[Logarithmic windows beyond reversibility]
\label{f13:thm:nonrev-log-windows}
Let $\alpha$ be a Pisot or Salem number of degree $D$.  Let
\[
 R(n)=P(a^n)+bn+e,
\]
where $a\ge2$, $b\ne0$, $e\in\mathbb Z$, and
$P\in\mathbb Z[X]$ has positive leading coefficient and positive degree.
For every $\varepsilon>0$, there are infinitely many indices $n$ such
that all integers in
\[
 \left[
 \lfloor\alpha^{R(n)}\rfloor-
       \left(\frac1{4D}-\varepsilon\right)\log n,
 \lfloor\alpha^{R(n)}\rfloor+
       \left(\frac1{4D}-\varepsilon\right)\log n
 \right]
\]
are composite.  Here $\log$ denotes the natural logarithm; a nonpositive
radius is interpreted trivially.  The constant is independent of the
degree and coefficients of $P$, and of the nonzero coefficient $b$.
\end{theorem}
\begin{proof}
Put $T(k)=\operatorname{Tr}_{\mathbb Q(\alpha)/\mathbb Q}(\alpha^k)$.
It is an integer linear recurrence of order $D$, and
\begin{equation}\label{f13:eq:window-trace-error}
 |T(k)-\alpha^k|\le D-1,\qquad
 |T(k)-\lfloor\alpha^k\rfloor|\le D.
\end{equation}
There is a fixed integer $s_0$ such that, for every prime $p$, the
sequence $T(k)\bmod p$ is periodic for $k\ge s_0$, with some period
$\ell_p\le p^D$.  Indeed, its state vector has at most $p^D$ states.
For primes not dividing the nonzero trailing coefficient its transition
matrix is invertible, hence the sequence is purely periodic.  There are
only finitely many other primes, so their preperiods have a common
bound.  This is also Saito's recurrence-period lemma.

Let $m\to\infty$ through integers satisfying $m\ge m_0$ and
$R(m)\ge s_0$.  Write
\[
 t=T(R(m)),\qquad H=t-2.
\]
For all large $m$, $t$ is a positive integer and $H\ge0$.  Every integer
$t+j$, $|j|\le H$, lies in $[2,2t-2]$, so it has a prime divisor
$p_j\le2t-2$.  Let
\[
 L=\operatorname{lcm}\{\ell_p:p\le2t-2,\ p\text{ prime}\}.
\]
By the prime number theorem in the form
$\vartheta(x)=\sum_{p\le x}\log p=x+o(x)$,
\begin{equation}\label{f13:eq:window-period-size}
 \log L\le D\vartheta(2t-2)=(2D+o(1))t.
\end{equation}
For $L\ge2$, Lemma~\ref{f13:lem:window-return-bound} gives $n>m$ with
$R(n)\equiv R(m)\pmod L$ and
\[
 n\le L^2+\max\{m+1,\lceil\log_2L\rceil\}.
\]
If $L=1$, take $n=m+1$, for which the same congruence is automatic.
Since $R$ is eventually increasing and $R(m+1)-R(m)\to\infty$,
\[
 \frac{T(R(m+1))}{T(R(m))}\longrightarrow\infty.
\]
Thus for all sufficiently large $m$ and every $|j|\le H$,
$T(R(n))+j>2t-2\ge p_j$.  On the other hand, periodicity gives
\[
 T(R(n))+j\equiv T(R(m))+j\equiv0\pmod{p_j}.
\]
Each integer in the trace window $[T(R(n))-H,T(R(n))+H]$ is therefore
composite.  Equation~\eqref{f13:eq:window-trace-error} transfers this to a
window of radius $H-D=t-D-2$ around the required floor.

Finally $\log(m+1)=o(t)$, because $R(m)$ grows exponentially and
$t=\alpha^{R(m)}+O(1)$.  Combining the bound on $n$ with
\eqref{f13:eq:window-period-size} gives
\[
 \log n\le(4D+o(1))t.
\]
Consequently $t-D-2\ge(1/(4D)-\varepsilon)\log n$ for all sufficiently
large $m$.  The selected indices are unbounded since $n>m$, and the
proof is complete.
\end{proof}

For example, every Pisot or Salem number $\alpha$ of degree $D$ has
infinitely many prime-free windows of radius
$(1/(4D)-\varepsilon)\log n$ around
$\lfloor\alpha^{3^n+n}\rfloor$ and
$\lfloor\alpha^{2^{5n}+7n-4}\rfloor$.
The theorem does not include a pure geometric clock $a^n$: its nonzero
linear coefficient is the hypothesis that makes universal return
available.  No improvement of Saito's constant in the reversible case
is asserted.

\section{A finite-cover classification for quadratic norm-one unit clocks}
\label{f13:sec:quadratic-clock-cover}

This section exhibits both fixed-divisor cascades and all-composite
sequences for which every fixed finite set of prime divisors eventually
fails.  The factorization ingredient is classical and underlies the
quadratic-unit results discussed by Cass and Saito.  The exact gcd and
cover classification are supplied explicitly here, with no priority
claim.

Let $V_0(X)=2$, $V_1(X)=X$, and
$V_{r+2}(X)=XV_{r+1}(X)-V_r(X)$.  These integer polynomials satisfy
\[
 V_r(z+z^{-1})=z^r+z^{-r},\qquad
 V_r(V_s(X))=V_{rs}(X).
\]
Their values at $X=1$ are the six-periodic sequence
\begin{equation}\label{f13:eq:chebyshev-at-one}
 V_r(1)=2\cos(r\pi/3)=2,1,-1,-2,-1,1
 \quad(r\equiv0,1,2,3,4,5\pmod6).
\end{equation}

\begin{theorem}[Exact pairwise gcd]
\label{f13:thm:quadratic-clock-gcd}
Let $A\ge3$ be an integer and let $\beta>1$ be the root of
$X^2-AX+1$.  For $a\ge2$ put $F_n=\lfloor\beta^{a^n}\rfloor$, $n\ge1$.
For every $m<n$,
\[
 \gcd(F_m,F_n)=\gcd(F_m,c_a),\qquad
 c_a=
 \begin{cases}
 0,&\gcd(a,6)=1,\\
 2,&2\mid a,\ 3\nmid a,\\
 1,&6\mid a,\\
 3,&3\mid a,\ 2\nmid a.
 \end{cases}
\]
In particular a fixed finite set of primes covers every sufficiently late
$F_n$ if and only if at least one of the following holds:
\begin{enumerate}
\item $\gcd(a,6)=1$;
\item $a$ is even, $3\nmid a$, and $A$ is odd;
\item $a$ is odd, $3\mid a$, and $A\equiv1\pmod3$.
\end{enumerate}
In each listed case one prime suffices.  In every other case the
$F_n$ are pairwise coprime and their least prime factors tend to infinity.
\end{theorem}
\begin{proof}
The conjugate of $\beta$ is $\beta^{-1}\in(0,1)$, so
$F_n=V_{a^n}(A)-1$.  Reducing the composition identity modulo $F_m$ gives
\[
 F_n\equiv V_{a^{n-m}}(1)-1\pmod{F_m}.
\]
Equation~\eqref{f13:eq:chebyshev-at-one} proves the asserted gcd formula.
If $\gcd(a,6)=1$, it shows $F_1\mid F_n$ for every $n$, and any prime
divisor of $F_1$ is proper at all sufficiently late indices.

If $a$ is even and not divisible by three, then $V_{a^n}(A)$ is odd
exactly when $A$ is odd.  For even $A$ this follows directly from the
recurrence modulo two.  For odd $A$, that recurrence has period three
with values $0,1,1$, and $3\nmid a^n$.  Thus in this case either every
$F_n$ is even or every $F_n$ is odd; the gcd formula gives the second
alternative in the statement, or pairwise coprimality, respectively.

If $a$ is odd and divisible by three, reduction modulo three gives
$V_{a^n}(A)\equiv A\pmod3$.  To see this, if $A\equiv0$, the
recurrence makes every odd-index value zero.  If $A\equiv1$, the
six-periodic values in~\eqref{f13:eq:chebyshev-at-one} have all odd-index
values congruent to one modulo three.  Finally
$V_r(-X)=(-1)^rV_r(X)$ handles $A\equiv-1$.
Thus the gcd formula gives a common divisor three exactly when
$A\equiv1\pmod3$, and otherwise pairwise coprimality.  The case
$6\mid a$ is pairwise coprime immediately.

A pairwise coprime sequence of integers greater than one cannot have an
eventual cover by finitely many primes.  Each bounded prime occurs in
at most one term, so its least prime factors tend to infinity.
\end{proof}

\begin{proposition}[All composite, but no finite prime cover]
\label{f13:prop:pairwise-composite-example}
Let $A\ge4$ be even, let $\beta$ be the root greater than one of
$X^2-AX+1$, and let $a$ be even, not divisible by three, and have a
prime divisor $q\ge5$.  Then all the integers
$\lfloor\beta^{a^n}\rfloor$, $n\ge1$, are composite and pairwise coprime.
In particular,
\[
 \boxed{\lfloor(2+\sqrt3)^{10^n}\rfloor\quad(n\ge1)}
\]
is an explicit all-composite sequence whose least prime factors tend to
infinity.  No fixed finite prime-cover argument can certify all its
terms.
\end{proposition}
\begin{proof}
Pairwise coprimality follows from Theorem~\ref{f13:thm:quadratic-clock-gcd}.
Since $q\equiv\pm1\pmod6$, the integer polynomial $V_q(X)-1$ is
divisible by $X-1$.  Write $a^n=qu$.  Then
\[
 \lfloor\beta^{a^n}\rfloor
   =V_q(V_u(A))-1
\]
is divisible by $V_u(A)-1\ge A-1\ge3$.  This divisor is smaller than
the number itself, since $q>1$ and $\beta^{qu}>\beta^u$ imply strict
increase of the traces.  The quotient is an integer greater than one,
proving compositeness.  The explicit example takes $A=4,a=10$.
\end{proof}


\section{A smaller unconditional cutoff in Saito's Type C}
\label{f13:sec:type-c-refined}

All source theorem numbers in this section refer to Saito's
\emph{Transcendency of variants of Mills' constant},
arXiv:2508.16068v3.  His Theorem~1.9(C), proved using Theorem~2.6,
establishes transcendence of $\xi(r3^k-1)$ for
$r\ge4.003\times10^{14}$.  The discussion preceding Corollary~2.7
asks whether the threshold can be replaced by $1$; that corollary
obtains this conclusion assuming a prime in $[x,x+3x^{2/3}]$ for
every $x\ge1$.  We lower the unconditional threshold, but do not
reach $r=1$.

Write $A_r=\xi((r3^k-1)_{k\ge1})$ and
\[
 D_r=\begin{cases}1,&r\text{ even},\\2,&r\text{ odd}.\end{cases}
\]
Saito's Type-B and Type-C reductions give the following facts:
if $A_r$ is algebraic, there is a divisor $g$ of $D_r$ for which
$\beta=A_r^g$ is a cubic Pisot number, and
\begin{equation}\label{f13:eq:tc-source-traces}
 \lfloor A_r^{r3^k-1}\rfloor
   =\operatorname{Tr}\bigl(\beta^{(r3^k-1)/g}\bigr)
 \quad\text{for all sufficiently large }k.
\end{equation}
Here the trace equality is Proposition~3.1(iv); all sufficiently
large indices belong to its set $I$ because the ratios exceed $3$.
With
\[
 x_0=\exp(3\exp(32.76)),\qquad M=2x_0^{1/3}+1,
\]
Theorem~2.6 also gives
\begin{equation}\label{f13:eq:tc-source-bound}
 \beta\le M^{g/(3r-1)}\le M^{D_r/(3r-1)}.
\end{equation}
The explicit prime-interval input here is exactly the one used by
Saito, from Mossinghoff--Trudgian--Yang and Cully-Hugill; no stronger
prime-gap assumption is introduced.

\begin{lemma}[The complete cubic Pisot list below two]
The cubic Pisot numbers below $2$ are the large real roots of
\[
 X^3-X-1,\quad X^3-X^2-X-1,\quad X^3-X^2-1,\quad
 X^3-2X^2+X-1.
\]
If $\alpha$ is the root of the first polynomial, the root of the
last polynomial is $\alpha^2$.
\end{lemma}
\begin{proof}
First, a totally real cubic Pisot number $\beta$ exceeds $2$.
If its other roots are $u,v$, then $|uv|\ge1/\beta$, and its
minimal polynomial satisfies
\[
 1\le|f(1)f(-1)|
 =(\beta^2-1)(1-u^2)(1-v^2)
 \le(\beta^2-1)(1-1/\beta)^2.
\]
The right side is at most $3/4$ when $\beta\le2$.

Thus a cubic Pisot number $1<\beta<2$ has a nonreal conjugate pair.
Its positive integer norm is less than $2$, hence equals $1$.
Write its polynomial as $X^3-aX^2+bX-1$.  The trace bounds give
$a\in\{0,1,2,3\}$.  Since the polynomial is negative at $1$ and
positive at $2$, its integer coefficients satisfy
\[
 2a-3\le b\le a-1.
\]
This leaves the four displayed polynomials and the two possibilities
$X^3-2X-1$ and $X^3-3X-1$.  The former is reducible, and the latter
has three real roots and so cannot be Pisot below $2$.
Each displayed polynomial is irreducible, has one real root in
$(1,2)$, and has a nonreal pair with common modulus $\beta^{-1/2}<1$.
Finally direct substitution, using $\alpha^3=\alpha+1$, shows that
$\alpha^2$ is the root of the fourth polynomial.
\end{proof}

\begin{lemma}[Finite trace certificates for the four small candidates]
\label{f13:lem:type-c-finite-certificate}
Set $R=400300000000000$.  For $1\le r\le R$, none of the four
cubic Pisot numbers $\beta<2$ can satisfy
\eqref{f13:eq:tc-source-traces} with $g\mid D_r$ and all outputs prime.
\end{lemma}
\begin{proof}[Proof with exact finite-ring certificates]
Write $f(X)=X^3-aX^2+bX-1$.  For a prime $p$, work in
$\mathbb F_p[X]/(f)$, and let $t$ be the multiplicative order of
$X$.  If $X^e=v_0+v_1X+v_2X^2$, its trace modulo $p$ is
\[
 \tau(e)=3v_0+av_1+(a^2-2b)v_2.
\]
The powers and traces are computed exactly by the update
\[
 (v_0,v_1,v_2)\longmapsto(v_2,v_0-bv_2,v_1+av_2)\pmod p,
 \qquad (v_0,v_1,v_2)=(1,0,0)\text{ initially}.
\]
For $g=1$, the exponent begins at $3r-1$ and evolves by
$e\mapsto3e+2\pmod t$.  A zero of $\tau$ on its eventual cycle
gives infinitely many composite traces.  The following finite
certificates show that avoiding such a zero forces $r\equiv0\pmod t$
for every listed pair $(p,t)$:
\[
\begin{array}{c|l}
f& (p,t)\\\hline
X^3-X-1&(2,7),(41,1723),(131,17293),(347,173),(397,52669)\\
X^3-X^2-X-1&(5,31),(31,331),(59,3541),(71,5113),(89,8011)\\
X^3-2X^2+X-1&(2,7),(41,1723),(71,5113),(131,17293),\\
&\quad(347,173),(409,55897)
\end{array}
\]
The respective least common multiples of these orders are
\[
 1900442940636401,\qquad1488257702170843,\qquad
 10312502210725960469,
\]
all larger than $R$.  Thus no positive $r\le R$ can avoid every
obstruction.  For $X^3-X^2-1$, modulo $2$ the order is $7$ and
every eventual exponent cycle has a zero trace, for both $g=1$
and $g=2$.

For $g=2$, necessarily $r$ is odd; write $r=2j+1$.
The initial exponent is $3j+1$ and evolves by $e\mapsto3e+1$.
For the first and second polynomials, modulo $7$ the order is
$48$ and every eventual exponent cycle has a zero trace.
For the fourth polynomial $\beta=\alpha^2$, hence $A_r=\alpha$;
its traces are exactly those of the first polynomial in the
already excluded case $g=1$.

These certificates require no approximated roots or prime outputs.
They are verified by enumerating the indicated finite exponent sets,
marking each eventual cycle that contains a zero trace, and testing
all possible initial residues.  The complete standard-library verifier in Appendix~\ref{f13:app:typec-code}
reproduces the calculation; it may be copied verbatim into a Python file.  A zero appearing on a cycle recurs infinitely often,
and the corresponding positive traces tend to infinity, so they
eventually exceed $p$.  They therefore contradict the eventual
prime trace identity \eqref{f13:eq:tc-source-traces}.
\end{proof}

\begin{theorem}[An unconditional Type-C cutoff improvement]
\label{f13:thm:type-c-cutoff-new}
The full minimum $\xi((r3^k-1)_{k\ge1})$ is transcendental in each
of the following ranges:
\begin{enumerate}
\item every integer $r\ge162400000000000$;
\item every even integer $r\ge81200000000000$.
\end{enumerate}
More precisely, it suffices that
\[
 r>\frac{D_r\log(2x_0^{1/3}+1)}{3\log2}+\frac13,
 \qquad x_0=\exp(3\exp(32.76)).
\]
\end{theorem}
\begin{proof}
For $r\ge4.003\times10^{14}$ this is Saito's existing theorem.
For smaller $r$, algebraicity would produce a cubic Pisot number
$\beta=A_r^g$ as above.  Lemma~\ref{f13:lem:type-c-finite-certificate}
excludes $\beta<2$, and $\beta=2$ is impossible because its degree
is three.  Thus $\beta>2$.  On the other hand, the displayed
sufficient condition and \eqref{f13:eq:tc-source-bound} give $\beta<2$,
a contradiction.

The rounded integer cutoffs have rigorous slack.  For example,
\[
 e^{32.76}<168845000000000,\qquad
 \log2>693147/10^6,\qquad
 \log(2e^{e^{32.76}}+1)<e^{32.76}+1
\]
give both cutoffs by rational arithmetic.  The first two inequalities
are certified by, respectively, a Taylor sum with a geometric upper
bound on its tail and the positive series
$\log2=2\sum_{j\ge0}((2j+1)3^{2j+1})^{-1}$; these rational checks
are included in the verifier.
\end{proof}

The gain is a factor of about $2.46$ compared with using the smallest
Pisot number as the lower bound (and the corresponding parity refinement).
The threshold remains large.  No conclusion for $r=1$, or for all
positive $r$, is claimed.  In particular this does not replace the
prime-gap hypothesis of Saito's Corollary~2.7.

\section{Exact obstructions to the proposed universal-return routes}

\begin{proposition}[Every Type-C seed obstructs universal return]
For $C_n=r3^n-1$, every seed $m$ fails condition (B6), even if the
index-set restriction in that condition is discarded.
\end{proposition}
\begin{proof}
For every $n>m$,
\[
 v_3(C_n-C_m)=v_3(r)+m.
\]
But $3\nmid C_m$.  Taking $L=3^{v_3(r)+m+1}$ makes
$C_n\equiv C_m\pmod{LC_m}$ impossible.  Thus deleting any finite
prefix cannot repair this particular return condition.
\end{proof}

\begin{proposition}[A coupled obstruction for $3^n+n^2$]
For the full schedule $C_n=3^n+n^2$, condition (B6) fails at the
first seed.  Specifically,
\[
 C_n\not\equiv4\pmod{9\cdot757}\qquad(n\ge2).
\]
\end{proposition}
\begin{proof}
For $n\ge2$, return modulo $9$ would give
$n^2\equiv4\pmod9$, hence $n\equiv2$ or $7\pmod9$.
The prime $757$ has $\operatorname{ord}_{757}(3)=9$;
indeed $3^9\equiv1$, $3^3\equiv27\not\equiv1$.
Consequently return modulo $757$ would require
\[
 n^2\equiv4-3^2=-5
 \quad\text{or}\quad
 n^2\equiv4-3^7\equiv88\pmod{757}.
\]
Both are quadratic nonresidues.  Quadratic reciprocity gives
$(-5/757)=(5/757)=(2/5)=-1$, while
$(88/757)=(2/757)(11/757)=(-1)(9/11)=-1$.
Thus no return is possible.  Since $C_1=4$, the choice
$L=9\cdot757$ suffices to disprove (B6).
\end{proof}

This last proposition rules out a direct full-schedule application of
the universal-return clause.  It does not rule out eventual return
after a sufficiently large prefix, does not disprove transcendence,
and does not settle the full minimum for $3^n+n^2$.  The previously
proved sufficiently-late-tail transcendence remains the available
positive conclusion for that example.


\section{A density-one version of the prime--totient quotient conjecture}
\label{f13:sec:quotient-density}

Let
\[
 T(N)=\sum_{p\leq N}\varphi\!\left(\left\lfloor\frac Np\right\rfloor\right),
 \qquad \kappa=\frac1{\zeta(2)}.
\]
Conjecture 10 in Saito--Suzuki--Takeda--Yoshida,
\emph{Some remarks on the $[x/n]$-sequence}, arXiv:2312.15642v1,
asks whether $T(x)\sim\kappa x\log\log x$ pointwise as $x\to\infty$.
The paper appeared in Research in Number Theory \textbf{11} (2025),
Article 50; the conjecture number here refers specifically to the arXiv
version. We prove the following density and mean statements. They do not
establish the pointwise conjecture. An earlier announced pointwise
formula and a difficulty with its proof are discussed below.

\begin{theorem}[Bounded mean absolute error and density-one asymptotic]
\label{f13:thm:quotient-density}
For every sufficiently large integer $X$,
\[
 \boxed{\frac1X\sum_{X\leq N<2X}
 \left|\frac{T(N)}{N}-\kappa\log\log N\right|\ll1.}
\]
Consequently, for every fixed $\varepsilon>0$,
\[
 \#\left\{X\leq N<2X:
  \left|\frac{T(N)}{N\log\log N}-\kappa\right|>\varepsilon\right\}
 \ll_\varepsilon\frac{X}{\log\log X}=o(X).
\]
In particular, the conjectured asymptotic holds on a subset of the
positive integers of natural density one.
\end{theorem}

The proof uses only the identity $\varphi(n)/n=\sum_{d\mid n}\mu(d)/d$,
the prime reciprocal estimate $\sum_{p\leq z}p^{-1}=\log\log z+O(1)$,
and the following elementary correlation calculation.

\begin{lemma}[Correlation of quotient divisibility blocks]
\label{f13:lem:quotient-correlation}
For positive integers $p,q,d,e$, put
\[
 I_{p,d}(N)=\mathbf 1_{d\mid\lfloor N/p\rfloor},
 \qquad F_{p,d}(N)=I_{p,d}(N)-\frac1d,
 \qquad g=\gcd(pd,qe).
\]
Then, for every interval of $X$ consecutive integers,
\[
 \left|\frac1X\sum_N F_{p,d}(N)F_{q,e}(N)\right|
 \leq\frac{g^2}{pqde}+O\!\left(\frac{pqde}{X}\right).
\]
The implicit constant is absolute; $p=q$ is permitted.
\end{lemma}

\begin{proof}
The first indicator has period $pd$ and selects the $p$ residue classes
$0,\ldots,p-1$ modulo $pd$, so its mean is $1/d$. Similarly the other
mean is $1/e$. Their common period is $L=pqde/g$.
For each $r\bmod g$, let $A_r$ count the integers in $[0,p)$ congruent
to $r\bmod g$, and define $B_r$ similarly with $q$. The Chinese
remainder theorem shows that the joint indicator selects
$\sum_{r\bmod g}A_rB_r$ classes modulo $L$. Write
$A_r=p/g+a_r$ and $B_r=q/g+b_r$. Then
$|a_r|,|b_r|\leq1$ and $\sum a_r=\sum b_r=0$, giving
\[
 \frac1L\sum_{r\bmod g}A_rB_r
 =\frac1{de}+\frac1L\sum_{r\bmod g}a_rb_r.
\]
Thus the covariance over a full common period has absolute value at
most $g/L=g^2/(pqde)$. The product $F_{p,d}F_{q,e}$ has
absolute value at most one. Its average over an arbitrary interval
differs from its periodic average by $O(L/X)$, proving the claim.
\end{proof}

\begin{lemma}[A summable covariance majorant]\label{f13:lem:quotient-gcd-sum}
Writing $\omega(n)$ for the number of distinct prime factors of $n$,
with $\omega(1)=0$, one has
\[
 \sum_{m,n\geq1}
 \frac{\omega(m)\omega(n)\gcd(m,n)^2}{m^2n^2}<\infty.
\]
\end{lemma}
\begin{proof}
Set $A_j=\sum_{n\geq1}\omega(n)^j/n^2$ for $j=1,2$ and
$A_0=\zeta(2)$. These are finite since $\omega(n)\leq\log_2n$.
Write $m=ga,n=gb$ with $g=\gcd(m,n)$, drop the restriction
$\gcd(a,b)=1$, and use $\omega(ga)\leq\omega(g)+\omega(a)$.
The sum in the statement is bounded by
\[
 \sum_{g,a,b\geq1}
 \frac{(\omega(g)+\omega(a))(\omega(g)+\omega(b))}{g^2a^2b^2}
 = A_2A_0^2+3A_0A_1^2<\infty.
\]
\end{proof}

\begin{proof}[Proof of Theorem~\ref{f13:thm:quotient-density}]
Write $W=\log\log X$, choose $D=\lfloor W\rfloor$ and
$Y=\lfloor X^{1/4}\rfloor$, and put
\[
 f(n)=\frac{\varphi(n)}n,\quad
 f_D(n)=\sum_{\substack{d\mid n\\d\leq D}}\frac{\mu(d)}d,
 \quad \kappa_D=\sum_{d\leq D}\frac{\mu(d)}{d^2}.
\]
Since $0<f(n)\leq1$,
\[
 \frac{T(N)}N
 =\sum_{p\leq N}\frac1p f(\lfloor N/p\rfloor)+O(1).
\]
The primes $Y<p\leq N$ contribute $O(1)$ uniformly for
$X\leq N<2X$, by the prime reciprocal estimate. We retain all primes
$p\leq Y$, including the small ones.

The discarded divisor tail has small mean. For every $p\leq Y$,
each value of $\lfloor N/p\rfloor$ is attained at most $p$ times, and
\[
 \frac1X\#\{X\leq N<2X:d\mid\lfloor N/p\rfloor\}
 \leq \frac1d+O(p/X).
\]
It follows that
\begin{align*}
 &\frac1X\sum_{X\leq N<2X}
 \sum_{p\leq Y}\frac1p
 |f(\lfloor N/p\rfloor)-f_D(\lfloor N/p\rfloor)|\\
 &\qquad\ll
 \frac1D\sum_{p\leq Y}\frac1p
 +\frac{\pi(Y)\log(2X)}X\ll1.
\end{align*}
Here the exact divisor identity bounds the absolute tail by
$\sum_{d>D,d\mid n}1/d$, and $d\leq2X/p$ in this average.

Define the centered truncated sum
\[
 Z(N)=\sum_{p\leq Y}\frac1p
 \sum_{d\leq D}\frac{\mu(d)}d F_{p,d}(N).
\]
Lemma~\ref{f13:lem:quotient-correlation}, including its diagonal case,
gives
\begin{align*}
 \frac1X\sum_{X\leq N<2X}|Z(N)|^2
 &\ll
 \sum_{\substack{p,q\leq Y\\p,q\text{ prime}}}
 \sum_{d,e\leq D}\frac{\gcd(pd,qe)^2}{(pd)^2(qe)^2}
 +\frac{D^2\pi(Y)^2}{X}\\
 &\leq
 \sum_{m,n\geq1}\frac{\omega(m)\omega(n)\gcd(m,n)^2}{m^2n^2}
 +O(D^2Y^2/X)\ll1.
\end{align*}
For the second line, each integer $m=pd$ has at most $\omega(m)$
representations with $p$ prime. Lemma~\ref{f13:lem:quotient-gcd-sum}
therefore applies. Cauchy--Schwarz bounds the mean of $|Z(N)|$ by $O(1)$.
Finally,
\[
 \kappa_D=\kappa+O(1/D),\qquad
 \sum_{p\leq Y}\frac1p=W+O(1),\qquad W/D=O(1).
\]
Combining these estimates proves the asserted bounded mean absolute
error. Replacing $\log\log X$ by $\log\log N$ changes the quantity
by $O(1/\log X)$.
Markov's inequality proves the displayed exceptional-set estimate.
To obtain one density-one subset on which the relative error tends to
zero, use the shrinking tolerance $(\log\log X)^{-1/2}$ on dyadic
intervals. The relative exceptional proportions then tend to zero.
\end{proof}

\begin{corollary}[Ces\`aro form of Conjecture 10]
\label{f13:cor:quotient-cesaro}
As $X\to\infty$ through the positive integers,
\[
 \sum_{N\leq X}T(N)
 =\frac{X^2\log\log X}{2\zeta(2)}+O(X^2).
\]
\end{corollary}

\begin{proof}
Multiply the dyadic mean absolute bound by $2X$ on each interval
$[X,2X)$, and sum over dyadic intervals up to the requested endpoint.
The total error is $O(X^2)$. Partial summation gives
$\sum_{3\leq N\leq X}N\log\log N=(X^2/2)\log\log X+O(X^2/\log X)$.
\end{proof}

\begin{remark}[Earlier announcement and scope]
Carella, \emph{Euler Totient Function And The Largest Integer Function
Over The Shifted Primes}, arXiv:2105.00790v2 (2021), Theorem~2.1,
announces the stronger formula
$T(x)=\kappa x\log\log x+c_1x+O(x/\log x)$.
Its equation~(20), however, uses an incorrect prime-counting estimate.
Already at its fixed parameter $n=1$ that estimate would imply
\[
 \pi(\lfloor x\rfloor)-\pi(\lfloor x/2\rfloor)
 =\frac12\operatorname{li}(x)+O(xe^{-c\sqrt{\log x}}).
\]
The prime number theorem and the expansion of $\operatorname{li}$ give
instead
\[
 \pi(\lfloor x\rfloor)-\pi(\lfloor x/2\rfloor)
 -\tfrac12\operatorname{li}(x)
 =-\frac{(\log2)x}{2\log^2x}+O(x/\log^3x),
\]
contradicting that error bound. This identifies a failure of the supplied
proof, not a disproof of its announced final asymptotic. The later
Saito--Suzuki--Takeda--Yoshida preprint explicitly poses the leading
asymptotic as Conjecture~10; it does not discuss Carella's manuscript.
Our independent averaging argument does not claim bibliographic
priority and does not establish the pointwise conjecture. Its bounds
permit exceptional individual integers with error of order
$N\log\log N$. The density and averaged conclusions concern the same
conjecture and are not two separate resolved open problems.
\end{remark}

\subsection{Explaining the observed half-square inclusion}

Conjecture 8.3 of Kanado--Saito, \emph{A system of certain linear
Diophantine equations on analogs of squares}, arXiv:2205.12226v4,
asks for infinitely many $a,b,c\geq2$ such that
\[
 a^2+b^2-1,\quad a^2+c^2-1,\quad b^2+c^2-1,\quad
 a^2+b^2+c^2-2
\]
are all squares. Let $V(X)$ denote its ordered solutions
$a\leq b\leq c\leq X$. Immediately after the displayed list $V(46300)$,
the authors state that they do not know the reason for the observed
inclusion $V(46300)\subset U_{\leq46300}(1/2)$. Here
$U_{\leq X}(\alpha)$ consists of triples of indices at most $X$ whose
positive values $\lfloor\alpha a^2\rfloor,\lfloor\alpha b^2\rfloor,
\lfloor\alpha c^2\rfloor$, all three pair sums, and their total sum
belong to $S(\alpha)=\{\lfloor\alpha n^2\rfloor:n\geq\alpha^{-1/2}\}$.
The observed inclusion has an exact explanation valid at every cutoff.

\begin{theorem}[Universal transfer of near-brick solutions]
\label{f13:thm:nearbrick-transfer}
For every $X>0$ and every
\[
 \alpha\in\left\{\frac18,\frac28,\ldots,\frac78,
                         \frac13,\frac23\right\},
 \qquad V(X)\subseteq U_{\leq X}(\alpha).
\]
In particular, the authors' observed inclusion at $\alpha=1/2$ holds
for every $X$.
\end{theorem}

\begin{proof}
Let $(a,b,c)\in V(X)$, and write the four companion squares as
\[
 u_{ab}^2=a^2+b^2-1,\quad u_{ac}^2=a^2+c^2-1,\quad
 u_{bc}^2=b^2+c^2-1,\quad v^2=a^2+b^2+c^2-2.
\]
At most one of $a,b,c$ is even: two even coordinates would make
their companion square congruent to $3\pmod4$.
Every odd square is $1\pmod8$. If there is an even coordinate,
write its square residue modulo $8$ as $e\in\{0,4\}$.
A pair of odd coordinates has companion-square residue $1$;
a pair containing the even coordinate has companion-square residue $e$.
The total companion square likewise has residue $1$ if all coordinates
are odd, and residue $e$ otherwise.

Write each coordinate square as $8h+r$. The exact square equations
show that each companion's quotient $h$ is the sum of the corresponding
coordinate quotients: the residue corrections are
$1+1-1=1$, $e+1-1=e$, $1+1+1-2=1$, and $e+1+1-2=e$.
For $1\leq p\leq7$,
\[
 \left\lfloor\frac p8(8h+r)\right\rfloor
 =ph+\left\lfloor\frac{pr}{8}\right\rfloor.
\]
The residue $1$ contributes zero, while the sole possible residue $e$
contributes the same term on both sides of every relevant sum identity.
This proves all four identities for $\alpha=p/8$.

Modulo $3$, squares have residues $0,1$, and at most one of $a,b,c$
is divisible by $3$, since two such coordinates would give a companion
square congruent to $2\pmod3$. The identical argument with modulus $3$
and exceptional residue $e=0$ proves the assertions for $p/3$, $p=1,2$.

Finally no coordinate can equal $2$: for another coordinate $b\geq2$,
$b^2+3$ lies strictly between $b^2$ and $(b+1)^2$.
Thus $a,b,c\geq3$, so all their floor values are positive for the
listed $\alpha\geq1/8$. The companion floor values are their positive
sums. All membership and positivity requirements in $U_{\leq X}$
are therefore satisfied.
\end{proof}

\begin{remark}
The transfer proves the universal inclusion behind an explicitly
unexplained numerical observation in the source. It is an elementary
congruence result; it does not prove that $V(X)$ grows without bound,
does not resolve Conjecture 8.3, and does not answer Question 8.1 about
the existence of an $\alpha\in(0,1)$ with $T(\alpha)$ empty.
\end{remark}

A useful exact search reduction is also available.

\begin{proposition}[Finite fibers and an exact divisor reduction]
Fix $a\geq2$. All possible positive integers $b$ satisfying
$a^2+b^2-1=u^2$ are obtained by
\[
 b=\frac{v-w}{2},\qquad u=\frac{v+w}{2},\qquad
 vw=a^2-1,\quad v>w>0,\quad v\equiv w\pmod2.
\]
In particular $b\leq(a^2-2)/2$. Consequently every fixed smallest
coordinate occurs in only finitely many solutions of Conjecture 8.3,
and any infinite family would have all three coordinates tending to
infinity along a subsequence.
\end{proposition}

\begin{proof}
Factor $(u-b)(u+b)=a^2-1$. Both factors are positive integers of the
same parity, giving the formula and its converse. Since $v\leq a^2-1$
and $w\geq1$, the bound follows. Apply the same argument to $c$.
For $a\leq b\leq c$, the set of solutions with $a$ bounded is therefore
finite.
\end{proof}

\section{What would complete the remaining transfers?}
The results locate explicit missing steps rather than removing them by
renaming auxiliary problems. For Mills' constant, the density sieve must
be upgraded to a recurrent-divisor theorem for each individual permitted
cubic orbit. For the least square constant, one must force an
approximation or recurrence of the kind proved sufficient here; neither
minimality nor primality has yet supplied it. For Type C, the improved
cutoff still leaves small parameters. For the prime--totient sum, an
average estimate cannot exclude large errors at a sparse set of integers.
For near-Euler bricks, the exact rational transfer does not create an
unbounded family of integer solutions.

The boundary theorem illustrates a successful auxiliary-problem route:
first sharpen prime-root separation according to affine support rank,
then prove simultaneous short approximations at every local left minimum,
and finally transfer the separation to those limits. The cube sieve uses
a different connection: finite-field return, Chinese-remainder residue
counts, and a shrinking coefficient region. Both arguments give statements
over infinite families, without replacing their remaining universal
quantifiers by finite computations.

\clearpage
\part{Focused reductions for the three remaining problems}
\label{part:focused}
\section{The nonreal cubic Pisot case}
\label{ft:sec:nonreal-cubic}

Problem~1.1 of \cite{SaitoIntervals} asks whether every cubic Pisot
number $\beta$ satisfies
\[
 \#\{r\geq 0:\lfloor\beta^{3^r}\rfloor\text{ is composite}\}=\infty.
\]
Changing the starting index does not change this assertion. The following
result settles its nonreal-conjugate case. It uses no equidistribution
assertion for the angles $3^r\theta$, either on the whole sequence or
on an arithmetic progression.

\begin{theorem}[Nonreal cubic conjugates]
\label{ft:thm:nonreal-cubic}
Let $\beta$ be a cubic Pisot number whose two other conjugates are
nonreal. Then $\lfloor\beta^{3^r}\rfloor$ is composite for infinitely
many integers $r\geq0$.
\end{theorem}

We give all of the trace, sign, and return arguments needed for the proof.
Write the minimal polynomial and conjugates as
\[
 f(X)=X^3-aX^2+bX-c\in\mathbb Z[X],
 \qquad \beta,\alpha,\overline\alpha,
 \qquad \alpha=\rho e^{i\theta},\quad 0<\rho<1.
\]
For $r\geq0$ set
\[
 T_r=\beta^{3^r}+\alpha^{3^r}+\overline\alpha^{3^r},\qquad
 B_r=(\beta\alpha)^{3^r}+(\beta\overline\alpha)^{3^r}
       +(\alpha\overline\alpha)^{3^r},\qquad E_r=c^{3^r}.
\]
These are integers. Newton's identities give
\begin{equation}\label{ft:eq:cubic-newton}
 \begin{split}
 T_{r+1}&=T_r^3-3T_rB_r+3E_r,\\
 B_{r+1}&=B_r^3-3T_rB_rE_r+3E_r^2,\\
 E_{r+1}&=E_r^3.
 \end{split}
\end{equation}
In particular $T_r\equiv a\pmod3$ for every $r$.

\begin{lemma}[Distinct powered conjugates]
\label{ft:lem:pisot-power-fields}
For a Pisot number $\gamma$ and every positive integer $m$,
$\mathbb Q(\gamma^m)=\mathbb Q(\gamma)$. Consequently the $m$th
powers of its distinct conjugates are distinct. In the present setting,
$\alpha/\overline\alpha$ is not a root of unity, and
$\cos(3^r\theta)\ne0$ for every $r\geq0$.
\end{lemma}
\begin{proof}
Any embedding of $\mathbb Q(\gamma)$ fixing $\gamma^m$ sends
$\gamma$ to a conjugate of modulus $\gamma>1$. The Pisot property
makes $\gamma$ the only such conjugate. Separability therefore gives
$[\mathbb Q(\gamma):\mathbb Q(\gamma^m)]=1$.
The distinctness assertion follows by restricting embeddings to these
equal fields. If $\alpha/\overline\alpha$ were a root of unity,
some powered conjugates would coincide. If $\cos(3^r\theta)=0$,
then $(\alpha/\overline\alpha)^{3^r}=-1$, which again makes this
ratio a root of unity.
\end{proof}

\begin{lemma}[Trace congruence]
\label{ft:lem:cubic-ghost}
For every $r\geq0$,
\begin{equation}\label{ft:eq:cubic-ghost}
 T_{r+1}-T_r\in3^{r+1}\mathbb Z.
\end{equation}
\end{lemma}
\begin{proof}
Let $A$ be an integral companion matrix of $f$, so
$T_r=\operatorname{Tr}(A^{3^r})$. In the matrix expansion of
$\operatorname{Tr}(A^{3^j})$, group closed index words by cyclic
rotation. A word of least period $3^i$ is the repetition of a
primitive cyclic word of length $3^i$, with an integral weight $w$,
and its orbit contributes $3^i w^{3^{j-i}}$. Consequently
\[
 \operatorname{Tr}(A^{3^j})=
 \sum_{i=0}^{j}3^i\sum_{w\in\mathcal V_i}w^{3^{j-i}},
\]
where $\mathcal V_i$ is the finite multiset of weights of primitive
cyclic words of length $3^i$. The new terms with $i=j$ are divisible
by $3^j$. For every integer $w$ and $h\geq1$,
\[
 w^{3^h}-w^{3^{h-1}}\in3^h\mathbb Z;
\]
this follows from Fermat's congruence at $h=1$ and induction using
the factorization of a difference of cubes. Thus every older term
in the difference of the expressions for $j$ and $j-1$ is also
divisible by $3^j$. Take $j=r+1$.
\end{proof}

For all sufficiently large $r$, define
\[
 s_r=\mathbf 1_{\{\cos(3^r\theta)>0\}},\qquad
 q_r=\lfloor\beta^{3^r}\rfloor=T_r-s_r.
\]
The formula for $q_r$ follows from
$T_r-\beta^{3^r}=2\rho^{3^r}\cos(3^r\theta)$, whose absolute
value is strictly between zero and one for all sufficiently large $r$.

\begin{lemma}[Both signs recur]
\label{ft:lem:cubic-both-signs}
Both values $0$ and $1$ occur infinitely often in $(s_r)$.
\end{lemma}
\begin{proof}
If two consecutive cosines are positive, the identity
$\cos(3x)=\cos x(4\cos^2x-3)$ gives $\cos x>\sqrt3/2$.
If both are negative, it gives $\cos x<-\sqrt3/2$.
An eventually constant sign would therefore keep every sufficiently
late angle within $\pi/6$ of a fixed center, either $0$ or $\pi$
modulo $2\pi$. Write the signed displacement from that center as
$\delta_r\in(-\pi/6,\pi/6)$. Tripling the angle gives
$\delta_{r+1}=3\delta_r$: since $3\delta_r\in(-\pi/2,\pi/2)$,
no change by a multiple of $2\pi$ is possible. Boundedness for all
future $r$ forces $\delta_r=0$. It follows that
$(\alpha/\overline\alpha)^{3^r}=1$, contradicting
Lemma~\ref{ft:lem:pisot-power-fields}.
\end{proof}

\begin{proof}[Proof of Theorem~\ref{ft:thm:nonreal-cubic}]
Suppose, for a contradiction, that every sufficiently late $q_r$ is
prime. The floors grow strictly eventually, so these primes eventually
avoid any fixed finite set. In particular, they are eventually odd and
different from $3$. If $a\equiv0\pmod3$, the infinitely many indices
with $s_r=0$ give multiples of $3$. If $a\equiv1\pmod3$, those with
$s_r=1$ do so. Lemma~\ref{ft:lem:cubic-both-signs} therefore forces
\begin{equation}\label{ft:eq:cubic-a-two}
 a\equiv2\pmod3.
\end{equation}

Modulo $2$, the Newton map \eqref{ft:eq:cubic-newton} is
\[
 (T,B,E)\longmapsto(T+BT+E,\ B+EBT+E,\ E).
\]
Its complete eight-state graph is
\[
 \begin{gathered}
 000\mapsto000,\quad001\mapsto111,\quad010\mapsto010,\quad
 011\leftrightarrow101,\\
 100\mapsto100,\quad110\mapsto010,\quad111\mapsto111.
 \end{gathered}
\]
Every cycle except $011\leftrightarrow101$ has constant trace parity.
Since $q_r=T_r-s_r$ is odd eventually, such a cycle would make $s_r$
eventually constant, contrary to Lemma~\ref{ft:lem:cubic-both-signs}.
Thus eventual primality forces the exact sign alternation
\begin{equation}\label{ft:eq:cubic-alternation}
 s_{r+1}=1-s_r
 \qquad\text{for all sufficiently large }r.
\end{equation}

Fix a sufficiently late index $r$ such that the prime $q=q_r$ does
not divide $3c\operatorname{disc}(f)$. Let $d\in\{1,2,3\}$ be the
degree of the splitting field of $f$ over $\mathbb F_q$; these are
the only possibilities for a separable cubic. Suppose that
$v_3(q^d-1)\leq r$. Write
\[
 q^d-1=3^tN,\qquad 3\nmid N,\qquad t\leq r.
\]
Choose a positive even integer $K$ for which $3^K\equiv1\pmod N$;
if $N=1$, take $K=2$. For every $j\geq0$,
\[
 q^d-1\mid 3^r(3^{jK}-1).
\]
All three residual roots are nonzero, so their $3^r$th powers
return at indices $r+jK$. Their traces therefore return modulo $q$.
Since $K$ is even, \eqref{ft:eq:cubic-alternation} returns the shift
as well. It follows that
\[
 q\mid T_{r+jK}-s_{r+jK}=q_{r+jK}\qquad(j\geq0).
\]
For sufficiently large $j$, this contradicts primality because
$q_{r+jK}>q$. We have proved
\[
 v_3(q_r^d-1)\geq r+1.
\]

For $d=1$, this gives $q_r\equiv1\pmod{3^{r+1}}$. For $d=2$,
exactly one of $q_r-1,q_r+1$ is divisible by $3$, so
$q_r\equiv\pm1\pmod{3^{r+1}}$. For $d=3$, it first forces
$q_r\equiv1\pmod3$, and
\[
 v_3(q_r^3-1)=v_3(q_r-1)+1
\]
gives $q_r\equiv1\pmod{3^r}$. Uniformly, therefore,
\[
 q_r\equiv\varepsilon_r\pmod{3^r},\qquad
 \varepsilon_r\in\{-1,1\}.
\]
Equation~\eqref{ft:eq:cubic-a-two} fixes the sign:
$\varepsilon_r=2s_r-1$. Hence
\[
 T_r\equiv3s_r-1\pmod{3^r}.
\]
By \eqref{ft:eq:cubic-alternation}, the representatives on the right
alternate between $-1$ and $2$. Comparing indices $r$ and $r+1$
therefore gives
\[
 T_{r+1}-T_r\equiv\pm3\pmod{3^r}.
\]
For $r\geq2$ this contradicts \eqref{ft:eq:cubic-ghost}. Thus
infinitely many floors are composite, as asserted.
\end{proof}

\begin{remark}[Scope of the angle argument]
\label{ft:rem:cubic-sign-scope}
Lemma~\ref{ft:lem:cubic-both-signs} concerns the full tripling orbit.
Irrationality alone does not give both signs on each progression of
indices. For example, choose a digit sequence $d_j\in\{1,2\}$ that
is not eventually periodic and put $t=\sum_{j\geq1}d_j9^{-j}$.
Every shifted base-$9$ tail lies strictly between $1/8$ and $1/4$.
Thus $t$ is irrational,
$\cos(2\pi9^nt)>0$, and $\cos(2\pi\,3\cdot9^nt)<0$ for every
$n\geq0$. The signs alternate forever. The proof above does not
try to rule out this itinerary by dynamics: it rules it out for an
eventually prime cubic trace sequence by the arithmetic trace congruence.
No assertion that $e^{2\pi it}$ is algebraic is made in this example.
\end{remark}

\section{The three remaining problems and their logical relations}
\label{ft:sec:bridge}

This part concerns the irrationality of $\xi(2^k)$, the unconditional
transcendence of the least Mills constant $A=\xi(3^k)$, and Saito's
Problem~1.1 asking whether every cubic Pisot number $\beta$ has infinitely
many composite values $\lfloor\beta^{3^r}\rfloor$.
The square problem is Question~1.6 of \cite{SaitoMills}; the cubic
problem is posed in \cite{SaitoIntervals}. The nonreal cubic case is
settled by Theorem~\ref{ft:thm:nonreal-cubic}. The results below give
further exact restrictions in the remaining cases, without settling
the three full assertions.

\begin{proposition}[One composite term is the universal target]
\label{ft:prop:one-composite}
Let $\mathcal C$ be a class of cubic Pisot numbers closed under
$\beta\mapsto\beta^{3^j}$ for every $j\ge0$. The following assertions
are equivalent:
\begin{enumerate}
\item every $\beta\in\mathcal C$ has at least one composite value
$\lfloor\beta^{3^r}\rfloor$ with $r\ge1$;
\item every $\beta\in\mathcal C$ has infinitely many such composite values.
\end{enumerate}
This applies both to the full cubic Pisot class and to the class
$\mathcal S$ of totally real cubic Pisot numbers with conjugates
$\beta,-\rho,\gamma$ satisfying
\begin{equation}
\label{ft:eq:saito-locus}
 0<\sigma=|\gamma|<\rho<1,
 \qquad \rho\le\min\{\sigma^{17/23},\beta^{-17/40}\}.
\end{equation}
Either universal assertion for $\mathcal S$ implies the unconditional
transcendence of the least Mills constant.
\end{proposition}
\begin{proof}
Apply the first assertion to $\beta^{3^j}$. It gives a composite
value at an index larger than $j$ for the original $\beta$, for every
$j$. The reverse implication is immediate.
Lemma~\ref{ft:lem:pisot-power-fields}
preserves cubic degree. Odd powers preserve the signs, order of the
small-root moduli, and both inequalities in
\eqref{ft:eq:saito-locus}, proving closure of the stated classes.

If the least Mills constant were algebraic, Saito's
Theorem~1.7 in \cite{SaitoVariants} would supply
$\beta=A^{3^m}\in\mathcal S$. Every value
$\lfloor\beta^{3^r}\rfloor=\lfloor A^{3^{m+r}}\rfloor$ would be
prime, contradicting the first assertion.
\end{proof}

The implication does not identify every algebraic member of $W(3^k)$
with its minimum. An eventually all-prime cubic Pisot orbit, if one
existed, would give an algebraic member after a power change; it would
not determine the arithmetic nature of the least member.

\begin{proposition}[The norm restriction and a nonunit example]
\label{ft:prop:nonunit}
For every $\beta\in\mathcal S$, the positive integer
$q=|\operatorname{Norm}(\beta)|$ satisfies
\[
 1\le q<\beta^{3/20}.
\]
The bound $\rho\le\sigma^{17/23}$ follows already from
$\rho\le\beta^{-17/40}$ and $q\ge1$. These restrictions do not force
$q=1$: the largest root of
\[
 f(X)=X^3-10007X^2-90X+2
\]
belongs to $\mathcal S$ and has norm $-2$.
\end{proposition}
\begin{proof}
The norm identity is $q=\beta\rho\sigma$. Since $\sigma<\rho$,
\[
 q<\beta\rho^2\le\beta^{3/20}.
\]
Furthermore
$\rho^{40}\le\beta^{-17}\le q^{17}\beta^{-17}
=\rho^{17}\sigma^{17}$, giving the claimed redundancy.

For the example, the signs at rational endpoints place its roots in
\[
 (10007,10008),\qquad
 (-39/2000,-19/1000),\qquad(1/100,11/1000).
\]
For the two small-root intervals the exact endpoint values are,
in the displayed order,
\[
 -\frac{401353319}{8000000000},\quad
 \frac{97466141}{1000000000},\quad
 \frac{99301}{1000000},\quad
 -\frac{200845669}{1000000000}.
\]
Also $f(10007)<0<f(10008)$. The three sign-changing intervals account
for all roots and contain no integer; the monic cubic is therefore
irreducible by the rational-root theorem. The dominant small conjugate
is negative and the norm is $-2$. The exact integer inequality
\[
 39^{40}10008^{17}<2000^{40}
\]
implies $\rho^{40}\beta^{17}<1$. The remaining conjugate inequality
follows from the first part of the proof.
\end{proof}

Passing to $\beta^{3^j}$ also raises $q$ to $q^{3^j}$; it does not
improve this logarithmic norm bound. A bound
$\rho\le\beta^{-1/2}$ would instead give
$1\le q<\beta\rho^2\le1$. The exponent $17/40$ does not reach that
contradiction, and the nonunit example shows why an argument restricted
to units cannot exhaust the possible shadows.

\section{A rational square minimum would force a prime deficit}
\label{ft:sec:square-deficit}

Let $A=\xi(2^k)$ and put $p_n=\lfloor A^{2^n}\rfloor$ for $n\ge1$.
Call a prime $q$ \emph{viable} if there is a real number $B>1$ such
that $\lfloor B^{2^j}\rfloor$ is prime for every $j\ge0$ and
$\lfloor B\rfloor=q$. Write $V$ for the set of viable primes.
Equivalently, $q\in V$ if its finitely branching tree of prime children
\[
 p^2<q'<(p+1)^2-1
\]
has an infinite path. Strict nesting of the closed descendant root
cells proves the equivalence.

\begin{lemma}[The least viable child]
\label{ft:lem:viable-child}
At every index $n\ge1$,
\[
 p_{n+1}=
 \min\{q\in V:p_n^2<q<(p_n+1)^2-1\}.
\]
\end{lemma}
\begin{proof}
The actual tail shows $p_{n+1}\in V$. Neither endpoint is prime:
$p_n^2$ is composite and $(p_n+1)^2-1=p_n(p_n+2)$ is composite.
If a smaller viable child $q$ existed, choose its infinite continuation.
After taking the $2^{n+1}$th root, this continuation lies strictly inside
the $p_n$ parent cell and below the $p_{n+1}$ cell. It respects every
earlier prime condition. Indeed, an earlier integer-power boundary
becomes an integer at level $n$, so it cannot cut the interior of a
single level-$n$ prime cell. The resulting admissible constant would
be smaller than $A$, contradicting minimality.
\end{proof}

For the good-prime predicate used in the preceding square-clock section,
fix sufficiently large integers $K,P$ and set
\[
 H(p)=\left\lceil\frac{p}{K\lceil\log_2p\rceil}\right\rceil.
\]
A prime $p\ge P$ is good when $[p^2,p^2+p]$ contains at least $H(p)$
primes, and poor otherwise. Lemmas~\ref{f13:sqnew:poorcount}
and~\ref{f13:sqnew:goodfull}, proved from Matom\"aki's prime-interval
estimate, give the two inputs
\begin{equation}
\label{ft:eq:good-poor-inputs}
 \{p\ge P:p\text{ is good}\}\subset V,
 \qquad
 \#\{p\in[X,2X]:p\text{ is a poor prime}\}\le C X^{1/3}
\end{equation}
for a fixed constant $C$ and all sufficiently large $X$.
Only existence of a continuation, rather than its dimension, is needed
in the next theorem.

\begin{theorem}[Necessary prime deficit for a rational minimum]
\label{ft:thm:square-deficit}
If $A=\xi(2^k)$ is rational, then for every fixed $\epsilon\in(0,1)$
and all sufficiently large $n$,
\begin{equation}
\label{ft:eq:square-deficit}
 \pi(p_n^2+p_n^{1-\epsilon})-\pi(p_n^2)
 \le C p_n^{2/3},
\end{equation}
where $C$ is the poor-prime constant in
\eqref{ft:eq:good-poor-inputs}. For $0<\epsilon<1/3$, this upper bound
is $o(p_n^{1-\epsilon}/\log p_n)$.
\end{theorem}
\begin{proof}
An integer $A>1$ cannot have prime square outputs, so a rational $A$
would be nonintegral. Write $u_n=A^{2^n}-p_n\in(0,1)$.
Mahler's approximation theorem, in the form of Theorem~2.5 of
\cite{SaitoMills}, states that for every $\eta>0$,
\[
 \|A^m\|>\exp(-\eta m)
\]
for all sufficiently large integers $m$. Choose
$\eta=\epsilon\log(A)/2$ and use $p_n\sim A^{2^n}$. It follows that
\[
 u_n\ge\|A^{2^n}\|>p_n^{-\epsilon}
\]
for all sufficiently large $n$.

The exact repeated-squaring relation is
\[
 p_{n+1}-p_n^2=2p_nu_n+u_n^2-u_{n+1}.
\]
Since $0<u_{n+1}<1$, it yields
\[
 p_{n+1}-p_n^2>2p_n^{1-\epsilon}-1>p_n^{1-\epsilon}
\]
eventually. Every prime in
$(p_n^2,p_n^2+p_n^{1-\epsilon}]$ is therefore an earlier child of
$p_n$. Lemma~\ref{ft:lem:viable-child} says it is nonviable, and
\eqref{ft:eq:good-poor-inputs} says it is poor. The whole interval
lies in $[p_n^2,2p_n^2]$ for large $n$. Applying the poor-prime bound
with $X=p_n^2$ proves \eqref{ft:eq:square-deficit}.
Finally,
\[
 \frac{p_n^{2/3}}{p_n^{1-\epsilon}/\log p_n}
 =p_n^{-1/3+\epsilon}\log p_n\longrightarrow0
\]
when $\epsilon<1/3$.
\end{proof}

\begin{corollary}[A viable-prime window sufficient for irrationality]
\label{ft:cor:square-window}
Suppose that, for some fixed $\delta\in(0,1)$ and infinitely many
indices of the actual least square-clock orbit, there is a viable prime
\[
 p_n^2<q\le p_n^2+p_n^{1-\delta}.
\]
Then $\xi(2^k)$ is irrational. In particular, it suffices that such
a good prime occur infinitely often. Alternatively, for some fixed
$\epsilon\in(0,1/3)$, a prime count exceeding the right side of
\eqref{ft:eq:square-deficit} infinitely often also suffices.
\end{corollary}
\begin{proof}
Minimality gives $p_{n+1}\le q$ at the indicated indices. The exact
square relation then gives
\[
 2p_nu_n\le p_n^{1-\delta}+1,
 \qquad u_n=O(p_n^{-\delta}).
\]
This contradicts Mahler's lower bound, for example with exponent
$\delta/2$, if $A$ is nonintegral rational. The integer case is
already excluded. The other assertions follow from
\eqref{ft:eq:good-poor-inputs} and Theorem~\ref{ft:thm:square-deficit}.
\end{proof}

At the scale $x=p_n^2$, the deficit window has length
$x^{(1-\epsilon)/2}$, below the exponent $1/2$. An averaged bound on
exceptional intervals does not exclude this particular orbit: since
$p_{n+1}>p_n^2$, it has at most one sufficiently large label in a
dyadic interval. The theorem supplies a precise necessary deficit,
while the viable-prime window in Corollary~\ref{ft:cor:square-window}
remains an unproved input for the actual minimum.

\section{Exact parity depths in a cubic-unit branch}
\label{ft:sec:final-mills-parity}

In the unit branch whose norm equals the limiting trace sign,
the exact $3$-adic depth can depend on the parity of the index.
The next theorem treats that branch and supplies a terminating
integer procedure for identifying both depths.

We use the Baker--Brumer theorem in the following standard form:
if the $p$-adic logarithms of nonzero algebraic numbers are linearly
independent over $\mathbb Q$, they are linearly independent over
$\overline{\mathbb Q}$. See Brumer~\cite{Brumer1967} and the explicit formulation in
Maksoud~\cite[Theorem~2.1]{Maksoud2023}.

\begin{lemma}[Multiplicative independence of two conjugates]
\label{ft:lem:final-unit-independence}
Let $\beta$ be a totally real cubic Pisot unit. Any two distinct
conjugates of $\beta$ are multiplicatively independent modulo
roots of unity.
\end{lemma}
\begin{proof}
Write the three real conjugates as $\beta_0,\beta_1,\beta_2$,
with $\beta_0=\beta$, and put $\ell_i=\log|\beta_i|$.
Then $\ell_0>0$, $\ell_1,\ell_2<0$, and
$\ell_0+\ell_1+\ell_2=0$ because the norm has absolute value one.
The transitive Galois group of an irreducible cubic contains a
three-cycle on its roots. Applying that cycle twice to a
multiplicative relation, and taking real absolute logarithms,
shows that its exponent vector is orthogonal to all cyclic
permutations of $(\ell_0,\ell_1,\ell_2)$.
These permutations span the plane $x_0+x_1+x_2=0$: the two
nontrivial Fourier eigenvalues of the associated circulant matrix
are nonzero, since a real triple with one positive and two negative
entries cannot be constant. Consequently the exponent vector is a
multiple of $(1,1,1)$. A relation involving only two roots has its
remaining exponent zero and is therefore trivial.
\end{proof}

\begin{theorem}[Exact parity-dependent $3$-adic depth]
\label{ft:thm:final-parity-depth}
Let $\beta$ be a totally real cubic Pisot unit with minimal polynomial
\[
 f(X)=X^3-aX^2+bX-\varepsilon,
 \qquad \varepsilon\in\{1,-1\},\qquad
 a\equiv\varepsilon\pmod3,\quad b\equiv1\pmod3.
\]
Let $T_r=\operatorname{Tr}(\beta^{3^r})$ for $r\ge0$.
There are finite positive integers $s_0,s_1$ such that
\[
 \boxed{v_3(T_r-\varepsilon)=r+s_j
 \quad\text{if }r\equiv j\pmod2\text{ and }r\ge s_j-1.}
\]
Moreover, the following statements hold.
\begin{enumerate}
\item Both integers $s_0,s_1$ are computable by a terminating
algorithm using only modular integer arithmetic.
\item One has
\[
 s_0=1\iff v_3(a-\varepsilon)=1,
 \qquad s_1=1\iff v_3(b-1)=1.
\]
Here the valuation of zero is understood to be infinity.
\item If $s=v_3(a-\varepsilon b)$, then $s$ is finite and
$\min(s_0,s_1)\le s$. In particular, $s=1$ forces at least one
parity offset to equal one.
\end{enumerate}
\end{theorem}
\begin{proof}
Modulo $3$ the polynomial is
$(X-\varepsilon)(X^2+1)$ and is separable. Its $3$-adic splitting
field is the unramified quadratic extension $K=\mathbb Q_3(i)$,
where $i^2=-1$. Hensel lifting labels its roots as
\[
 \alpha=\varepsilon\exp(L),\qquad
 u=i\exp(M),\qquad v=-i\exp(N),
 \qquad L,M,N\in3\mathcal O_K.
\]
Here $\alpha,L\in\mathbb Q_3$, and the nontrivial automorphism
of $K/\mathbb Q_3$ interchanges $u,v$ and $M,N$, and sends
$i$ to $-i$. Since $\alpha uv=\varepsilon$, injectivity of the
exponential on $3\mathcal O_K$ gives $L+M+N=0$.
For $m=3^r$ we obtain
\[
 T_r=\varepsilon\exp(mL)+(-1)^r i
                   \bigl(\exp(mM)-\exp(mN)\bigr).
\]
Define
\[
 \lambda_0=\varepsilon L+i(M-N),\qquad
 \lambda_1=\varepsilon L-i(M-N).
\]
Both belong to $3\mathbb Z_3$. The exponential series gives
\begin{equation}
\label{ft:eq:final-parity-remainder}
 T_r-\varepsilon=3^r\lambda_{r\bmod2}+R_r,
 \qquad v_3(R_r)\ge2r+2.
\end{equation}
Indeed, for $k\ge2$ and $Z\in3\mathcal O_K$, the $k$th
exponential term has valuation at least
$k(r+1)-v_3(k!)\ge2r+2$.

It remains to exclude a vanishing linear coefficient. Fixing an
embedding of the global splitting field into $\overline{\mathbb Q}_3$,
the roots $u,v$ are images of two distinct global conjugates.
The logarithms $M,N$ are linearly independent over $\mathbb Q$:
an integral relation would, by exponentiating in
$1+3\mathcal O_K$, make a product of powers of $u,v$ a root
of unity, contradicting Lemma~\ref{ft:lem:final-unit-independence}.
The Baker--Brumer theorem makes them linearly independent over
$\overline{\mathbb Q}$. Since $L=-M-N$, either equality
$\lambda_j=0$ would give a nontrivial algebraic linear relation
between $M$ and $N$, with coefficients
$-\varepsilon\pm i$ and $-\varepsilon\mp i$.
Thus $\lambda_0,\lambda_1$ are nonzero. Put
$s_j=v_3(\lambda_j)\ge1$.
When $r\ge s_j-1$, we have $r+s_j<2r+2$; the asserted
valuation formula follows from \eqref{ft:eq:final-parity-remainder}.

For effectiveness, let $B_r$ be the second elementary symmetric
function of the $3^r$th powers of the roots. Starting from
$(T_0,B_0)=(a,b)$, use the exact Newton recurrence
\begin{equation}
\label{ft:eq:final-unit-newton}
 T_{r+1}=T_r^3-3T_rB_r+3\varepsilon,
 \qquad B_{r+1}=B_r^3-3\varepsilon T_rB_r+3.
\end{equation}
For each parity, search for an $r$ of that parity for which
$T_r-\varepsilon$ is nonzero modulo $3^{2r+2}$.
For the first such certificate, if
$v=v_3(T_r-\varepsilon)<2r+2$, equation
\eqref{ft:eq:final-parity-remainder} proves $s_j=v-r$.
The nonvanishing of $\lambda_j$ ensures that this search
terminates for both parities. All recurrence calculations may be
performed modulo $3^{2r+2}$ for the current test.

At $r=0$, equation \eqref{ft:eq:final-parity-remainder} gives
$\lambda_0\equiv a-\varepsilon\pmod9$.
Writing $x=a-\varepsilon$ and $y=b-1$, both divisible by $3$,
the first recurrence gives
$T_1-\varepsilon=-3\varepsilon y+3\varepsilon x^2+x^3-3xy$.
Consequently $\lambda_1\equiv-\varepsilon(b-1)\pmod9$,
which proves the two criteria for $s_j=1$.

Finally set $x_r=T_r-\varepsilon$, $y_r=B_r-1$,
$z_r=x_r+\varepsilon y_r$, and
$w_r=x_r-\varepsilon y_r=T_r-\varepsilon B_r$.
The recurrence factors as
\[
 w_{r+1}=w_r
 \bigl(3+3\varepsilon z_r+x_r^2+\varepsilon x_ry_r+y_r^2\bigr).
\]
Since $3\mid x_r,y_r$, its multiplier is $3\pmod9$.
Also $w_0=a-\varepsilon b\ne0$, since otherwise
$f(\varepsilon)=0$. Thus $v_3(w_r)=r+s$ exactly.
Expanding the inverse-root expression for $B_r$ gives
\[
 w_r=2\varepsilon3^r L+O_3(3^{2r+2}).
\]
It follows that $v_3(L)=s$. Since
$\lambda_0+\lambda_1=2\varepsilon L$, the two offsets cannot
both exceed $s$.
\end{proof}

\begin{example}[The offsets need not agree]
For
\[
 f(X)=X^3-13X^2-8X-1
\]
there is one root in each of $(-1,-1/4)$, $(-1/4,0)$, and
$(13,14)$. The rational-root test proves irreducibility, so the
large root is a totally real cubic Pisot unit with negative
dominant small conjugate. Newton iteration gives
\[
 T_0-1=12,\qquad T_3-1\equiv2187=3^7\pmod{3^8}.
\]
These are exact certificates for $s_0=1$ and $s_1=4$.
Therefore
\[
 v_3(T_r-1)=
 \begin{cases}
 r+1,&r\text{ even},\\
 r+4,&r\text{ odd and }r\ge3.
 \end{cases}
\]
This example demonstrates the valuation phenomenon; it is not
claimed to be a prime-generating Pisot number.
\end{example}

\begin{corollary}[Parity-specific complete splitting]
\label{ft:cor:final-parity-splitting}
Assume the hypotheses of Theorem~\ref{ft:thm:final-parity-depth}
with $\varepsilon=1$, and suppose the dominant small conjugate
is negative. Fix a parity $j$ for which $s_j=1$.
If $\lfloor\beta^{3^r}\rfloor$ is eventually prime, then $f$
splits completely modulo every sufficiently late such prime
with $r\equiv j\pmod2$.
In particular this conclusion holds on the even indices when
$a\not\equiv1\pmod9$, and on the odd indices when
$b\not\equiv1\pmod9$.
\end{corollary}
\begin{proof}
For all sufficiently large $r$, the negative dominant small
conjugate gives $\lfloor\beta^{3^r}\rfloor=T_r$.
Let $p=T_r$ be prime, large enough that
$p\nmid3\operatorname{disc}(f)$, with $r\equiv j\pmod2$.
Then $p\equiv1\pmod3$ and $v_3(p-1)=r+1$.

If $f$ is irreducible modulo $p$, the ratios of its three roots
have orders dividing $p^2+p+1$, whose $3$-adic valuation is one.
After raising the roots to $3^r$, their ratios therefore have
orders coprime to $3$. If the factorization type is $1+2$, the
roots lie in $\mathbb F_{p^2}$, and their $3^r$th powers have
$3$-primary components in $\mu_3$, since
$v_3(p^2-1)=r+1$. The two conjugate quadratic roots have the
same component because $p\equiv1\pmod3$. Their product with
the remaining component is one, by the norm condition.
All three components are consequently equal, and their ratios
again have orders coprime to $3$.

In either nonsplit case, choose $L\ge1$ such that $3^L$ is one
modulo the exponent of the group generated by these ratios.
Cubing $L$ times returns all three powered roots up to a common
nonzero scalar. Their vanishing sum modulo $p$ therefore
recurs at $r+L,r+2L,\ldots$. The corresponding traces eventually
exceed $p$, yielding a positive-density progression of composite
floors. This contradicts an eventually prime tail.
\end{proof}


\section{A completely split prime seed need not return}
\label{ft:sec:final-cubic-return}

The splitting conditions above and their opposite-norm counterparts
do not imply recurrence of the seed prime as a divisor. The following
example lies in the precise shrinking-conjugate region and admits a
finite exhaustive certificate of nonreturn.

For an irreducible polynomial
\[
 f(X)=X^3-aX^2+bX-c
\]
with roots $\beta,\gamma_2,\gamma_3$, write
\[
 T_r=\beta^{3^r}+\gamma_2^{3^r}+\gamma_3^{3^r},\qquad
 B_r=(\beta\gamma_2)^{3^r}+(\beta\gamma_3)^{3^r}
       +(\gamma_2\gamma_3)^{3^r}.
\]
If $\beta$ is a totally real cubic Pisot number and the small conjugate
of largest absolute value is negative, then
$\lfloor\beta^{3^r}\rfloor=T_r$ for all sufficiently large $r$.
In the example below this holds for every $r\ge0$.

\begin{proposition}[A nonreturning prime seed in the narrowed cubic locus]
\label{ft:prop:seed-nonreturn}
Let $\beta$ be the largest root of
\[
 f(X)=X^3-211X^2-7X+1.
\]
Then $\beta$ is an irreducible totally real cubic Pisot number, with
conjugates $\beta,-\rho,\sigma$ satisfying
\[
 0<\sigma<\rho<1,
 \qquad \rho\le\min\{\sigma^{17/23},\beta^{-17/40}\}.
\]
Its first two sparse floors are the primes
\[
 T_0=211,\qquad T_1=p=9\,398\,359.
\]
Nevertheless,
\[
 p\nmid T_r\qquad\text{for every }r\ge2.
\]
Moreover $f$ splits completely modulo $p$.  The example satisfies
$a\equiv1$, $b\equiv-1\pmod3$ and $v_3(a+b)=1$, as well as the
odd-trace parity condition.
\end{proposition}
\begin{proof}
Sign changes give roots in $(-1,0)$, $(0,1)$ and $(211,212)$.
A rational root would be $1$ or $-1$, neither of which is a root;
hence the cubic is irreducible.  From the sum and product of its roots,
\[
 \rho-\sigma=\beta-211>0,\qquad \rho\sigma=\beta^{-1}.
\]
Since $f(-1/10)<0<f(0)$, one has $\rho<1/10$.  The exact inequality
$212^{17}<10^{40}$ gives
$1/10<212^{-17/40}<\beta^{-17/40}$.
The norm identity then implies
\[
 \rho^{40}\le\beta^{-17}=(\rho\sigma)^{17},
 \qquad \rho\le\sigma^{17/23}.
\]
For every odd integer $N\ge1$, the trace error is
$-\rho^N+\sigma^N\in(-1,0)$, proving equality of floors and traces.

Newton's identities give $T_1=211^3+3\cdot211\cdot7-3=p$ and
$B_1=-7^3-3\cdot211\cdot7+3=-4771$.
The primality assertions are certified by trial division up to the
respective integer square roots.  Substitution modulo $p$ certifies
the three distinct residual roots
\[
 -1\,584\,786,\qquad -1\,942\,624,\qquad 3\,527\,621.
\]
For norm $-1$, the exact modular Newton map is
\[
 \Phi(t,b)=
 (t^3-3tb-3,\ b^3+3tb+3)\pmod p.
\]
Start at $S=(0,-4771\bmod p)$, the state of index $1$.
Its next state is
\[
 S_1=(9\,398\,356,7\,830\,596).
\]
The integer-only verifier in Appendix~\ref{ft:app:seed-verifier} checks
\[
 \Phi^{85712}(S_1)=S_1,
 \qquad
 \text{the first coordinate of }\Phi^j(S_1)\ne0
 \quad(0\le j<85712).
\]
These checks are a finite exhaustive certificate of the asserted
nondivisibility at every later index.  They involve no numerical
approximation and no unproved primality test.
Finally, $a+b=204$ has $3$-adic valuation one; modulo $2$ the state
$(T_r,B_r,c)=(1,1,1)$ is fixed.
\end{proof}

This is not a counterexample to infinite compositeness.  Indeed,
\[
 T_2=830\,149\,079\,789\,343\,894\,643
     =7\cdot31\cdot3\,825\,571\,796\,264\,257\,579.
\]
It rules out deriving the desired theorem merely by asserting that
an observed prime trace must recur as a divisor, even after imposing
the known shrinking-conjugate, parity and $3$-adic restrictions.

For a fixed prime not dividing the norm, the finite Newton state is
eventually periodic. A zero trace in the transient part need not reappear
on its eventual cycle. Proposition~\ref{ft:prop:seed-nonreturn} therefore
excludes a universal proof based solely on return of an arbitrarily
chosen prime seed; it does not exclude a different recurrent divisor.

\section{The norm cascade in the positive-dominant real branch}
\label{ft:sec:positive-cascade}

For the remaining real sign configuration, the actual floor is
eventually one less than the trace. Let $\beta$ be a totally real cubic
Pisot number whose dominant small conjugate is positive, and write its
conjugates as $\beta,\rho,\delta$, where
$0<|\delta|<\rho<1$. Put $M=3^r$ and
\[
 \begin{split}
 T_r&=\beta^M+\rho^M+\delta^M,\\
 B_r&=(\beta\rho)^M+(\beta\delta)^M+(\rho\delta)^M,\\
 E_r&=(\beta\rho\delta)^M,\qquad F_r=T_r-1.
 \end{split}
\]
These are integers, and $F_r=\lfloor\beta^{3^r}\rfloor$ for all
sufficiently large $r$.

\begin{proposition}[A sufficient cascade condition and its eventual exclusion]
\label{ft:prop:positive-cascade}
If a prime $p=F_r$ divides
$\operatorname{Norm}(\beta^{3^r}-1)$, then $p\mid F_{r+j}$ for every
$j\ge0$. However, for every fixed $\beta$ in this real sign
configuration, this divisibility hypothesis fails at every sufficiently
large prime seed $F_r$. More precisely,
\[
 0<B_r-E_r<F_r
\]
for all sufficiently large $r$. Also
\[
 \gcd(F_r,F_{r+1})=\gcd(F_r,3(B_r-E_r)),
\]
and a sufficiently large prime $F_r$ never divides $F_{r+1}$.
\end{proposition}
\begin{proof}
The exact norm identity is
\[
 \operatorname{Norm}(\beta^M-1)=E_r-B_r+T_r-1
 =E_r-B_r+F_r.
\]
If $p=F_r$ divides this norm, then modulo $p$ the polynomial whose
roots are the three $M$th powers becomes
\[
 X^3-X^2+B_rX-E_r=(X-1)(X^2+E_r).
\]
Its roots in an algebraic closure of the residue field are
$1,w,-w$, counted with multiplicity. Their trace after every odd
power is $1$. Since all powers $3^j$ are odd, $p\mid T_{r+j}-1$.

On the other hand, an exact division gives
\[
 \frac{B_r-E_r}{(\beta\rho)^M}
 =1+\left(\frac\delta\rho\right)^M
   +\left(\frac\delta\beta\right)^M-\delta^M
 \longrightarrow1.
\]
Also $F_r/\beta^M\to1$, and consequently
\[
 \frac{B_r-E_r}{F_r}\sim\rho^M\longrightarrow0.
\]
This proves the displayed strict inequalities. The positive integer
$B_r-E_r$ is eventually smaller than the prime $F_r$, so it cannot
be divisible by that prime; the norm identity excludes the cascade
hypothesis.

Finally Newton's identity
$T_{r+1}=T_r^3-3T_rB_r+3E_r$ gives
$F_{r+1}\equiv3(E_r-B_r)\pmod{F_r}$, proving the gcd formula.
The same asymptotic gives $0<3(B_r-E_r)<F_r$ eventually, so a prime
$F_r$ cannot divide its next floor.
\end{proof}

This proposition concerns the particular norm
$\operatorname{Norm}(\beta^{3^r}-1)$ and propagation of the entire
prime seed. It leaves open whether another proper divisor or a later
finite-field return supplies infinitely many composite terms.

\section{Exact remaining requirements}
\label{ft:sec:remaining}

The square-clock reduction asks for one viable prime in a sublinear
window infinitely often along the actual least orbit, or for a violation
of its necessary prime deficit. The cubic reduction asks for a composite
term for every member of a power-closed real Pisot class; the nonreal
case is already covered. For the least Mills constant, Saito's shrinking
real class suffices. The parity-depth theorem and the seed-return
counterexample distinguish an exact local restriction from a universal
divisor assertion: complete splitting can occur, and a prime divisor can
lie only in the transient portion of a modular orbit. These results
leave the irrationality of $\xi(2^k)$, unconditional Mills transcendence,
and the full cubic infinite-compositeness problem unresolved.

\clearpage
\section{Scope of the consolidated conclusions}
The least constant, an arbitrary prime-representing base, a selected
branch, a local boundary point, and a minimum after deleting initial
conditions remain different objects throughout. None of the results
permits these distinctions to be dropped. The precise final status of
the three principal arithmetic targets is recorded in
Part~\ref{part:focused}. Saito's unconditional irrationality of Mills'
constant and his conditional transcendence theorem are external inputs,
not conclusions first established here.

The twin-output window equivalence and certified finite twin pairs do
not prove that their windows are nonempty infinitely often. Extremely
close approximations to a prime-generating base do not remove the
precision or primality-certification costs of evaluating its large outputs.
Failure of a fixed finite prime cover does not imply infinitely many
prime values, as the pairwise-coprime composite examples demonstrate.

For the polynomial-pattern results, the balanced-equation bound does
not determine the exact exceptional dimension. The waiting-time
parameter is fixed. The normality criterion does not establish its
analytic cancellation condition for a named algebraic irrational.
Density-one and averaged asymptotics do not establish a pointwise
conjecture. The general lattice dimension-profile result concerns
arbitrary subsets and leaves room for additional restrictions on groups
or homogeneous spaces.

\paragraph{Coefficient fields.}
In Theorem~\ref{fol:thm:relative-exponential-field}, the prescribed
field $F_0$ may contain $e$, $\pi$, or other named constants. This does
not put them into $F=\overline{\mathbb Q}(K)$. The corollary makes $e$
transcendental over $F$ and implies $\pi\notin F$: otherwise
$e^{i\pi}+e^0=0$ contradicts the stated exponential independence.
This last assertion is nonmembership, not transcendence over $F$.
No general form of Schanuel's conjecture is asserted.

\paragraph{Provenance and assistance.}
This edition extends the earlier collected manuscript by incorporating
the subsequent work on twelve classification and endpoint problems,
the thirteen further research routes, and the final focused reductions.
The counts in those working documents are not counts of solved named
conjectures. Repeated front matter and bibliographic entries have been
removed, and repeated arguments are cross-referenced where appropriate.
The preparation used AI assistance and independent internal checks.
The manuscript remains a proof draft requiring authorial and external
review before any publication or priority claim.

\clearpage\appendix
\part*{Appendices}
\addcontentsline{toc}{part}{Appendices}
\lstset{language=Python,basicstyle=\ttfamily\fontsize{8.5}{10}\selectfont,
frame=none,numbers=left,xleftmargin=14pt,xrightmargin=0pt,
breaklines=true,columns=fullflexible,keepspaces=true,showstringspaces=false}

\lstset{basicstyle=\ttfamily\fontsize{8.5}{10}\selectfont}
\section{Executable exact verification}
\label{enc:app:verification-code}
The following is the complete verifier. Copy this single listing to
\texttt{verify\_certificates.py} and run it with Python~3.8 or later.
It uses only the standard library, no internet access, and no probabilistic
primality tests. The seed and all seven later primes are proved by the
criteria in Section~\ref{enc:sec:certificates}. The final blocks verify the
finite floor table, the integer inequalities used for precision, and the
561 certified composite offsets. They do not attempt to construct the
astronomical next greedy prime or print the asserted common decimal block.

The listing is embedded directly in this single \LaTeX{} source.
It has no external data dependencies.
\begin{lstlisting}
"""Exact certificates for the sparse-prime-coding examples.
Python 3 standard library only. No probabilistic primality test is used.
Run: python verify_certificates.py
"""
from math import factorial, gcd, isqrt, prod

P0 = 5132842708382182842735812571653
# n: (witness, complete factorization of n-1).
CERT = {
    3: (2, {2: 1}),
    5: (2, {2: 2}),
    7: (3, {2: 1, 3: 1}),
    11: (2, {2: 1, 5: 1}),
    13: (2, {2: 2, 3: 1}),
    19: (2, {2: 1, 3: 2}),
    23: (5, {2: 1, 11: 1}),
    29: (2, {2: 2, 7: 1}),
    37: (2, {2: 2, 3: 2}),
    61: (2, {2: 2, 3: 1, 5: 1}),
    131: (2, {2: 1, 5: 1, 13: 1}),
    139: (2, {2: 1, 3: 1, 23: 1}),
    149: (2, {2: 2, 37: 1}),
    157: (5, {2: 2, 3: 1, 13: 1}),
    181: (2, {2: 2, 3: 2, 5: 1}),
    241: (7, {2: 4, 3: 1, 5: 1}),
    557: (2, {2: 2, 139: 1}),
    3259: (3, {2: 1, 3: 2, 181: 1}),
    4457: (3, {2: 3, 557: 1}),
    12541: (14, {2: 2, 3: 1, 5: 1, 11: 1, 19: 1}),
    14149: (6, {2: 2, 3: 3, 131: 1}),
    56597: (2, {2: 2, 14149: 1}),
    59242103: (5, {2: 1, 61: 1, 149: 1, 3259: 1}),
    355452619: (2, {2: 1, 3: 1, 59242103: 1}),
    853482761: (3, {2: 3, 5: 1, 13: 1, 29: 1,
                    56597: 1}),
    411378690803: (2, {2: 1, 241: 1, 853482761: 1}),
    P0: (2, {2: 2, 157: 1, 4457: 1, 12541: 1,
             355452619: 1, 411378690803: 1}),
}
proved = {2}


def require(condition):
    if not condition:
        raise ArithmeticError("Certificate verification failed")


def prove_seed(n):
    if n in proved:
        return
    a, factors = CERT[n]
    for q in factors:
        prove_seed(q)
    require(prod(q ** e for q, e in factors.items()) == n - 1)
    require(pow(a, n - 1, n) == 1)
    for q in factors:
        require(gcd(pow(a, (n - 1) // q, n) - 1, n) == 1)
    proved.add(n)


def lucas_u(k, n, P, Q):
    """U_k(P,Q) modulo n, by exact binary doubling."""
    u, v = 0, 1
    for bit in bin(k)[2:]:
        u2 = u * (2 * v - P * u) % n
        v2 = (v * v - Q * u * u) % n
        if bit == "0":
            u, v = u2, v2
        else:
            u, v = v2, (P * v2 - Q * u2) % n
    return u


def check_cell(s, n):
    require(s ** 4 < n < n + 1 < (s + 1) ** 4)


def prove_plus(s, j, a):
    require(s in proved)
    n, F = s ** 4 + j * s ** 2 + 1, 2 * s ** 2
    require((n - 1) % F == 0 and F * F > n)
    require(pow(a, n - 1, n) == 1)
    for q in (2, s):
        require(gcd(pow(a, (n - 1) // q, n) - 1, n) == 1)
    check_cell(s, n)
    proved.add(n)
    return n


def prove_minus(s, j, P, Q):
    require(s in proved)
    n, F = s ** 4 + j * s ** 2 - 1, 2 * s ** 2
    D = P * P - 4 * Q
    require(gcd(n, 2 * Q * D) == 1)
    require((n + 1) % F == 0 and (F - 1) ** 2 > n)
    require(lucas_u(n + 1, n, P, Q) == 0)
    for q in (2, s):
        require(gcd(lucas_u((n + 1) // q, n, P, Q), n) == 1)
    check_cell(s, n)
    proved.add(n)
    return n


prove_seed(P0)
require(len(proved) == 28)
require(P0 == 29 ** 21 + 24)
# Each entry is a proper divisor of 29**21+i, for i=1,...,23.
DIVISORS = [2, 31, 2, 3, 2, 5, 2, 571, 2, 3, 2, 13,
            2, 17, 2, 3, 2, 72673, 2, 7, 2, 3, 2]
for i, d in enumerate(DIVISORS, 1):
    n = 29 ** 21 + i
    require(1 < d < n and n % d == 0)
print("Seed: 28 certificate nodes; all 23 intervening integers composite.")

p1 = prove_plus(P0, 33, 2)
p2 = prove_plus(p1, 765, 5)
p3 = prove_plus(p2, 4583, 13)
require([len(str(n)) for n in (p1, p2, p3)] == [123, 492, 1966])
print("Prime tower: 123, 492, 1966 digits; all certificates passed.")

t_plus = prove_plus(P0, 20303, 17)
t_minus = prove_minus(P0, 20303, 1, -1)
w_plus = prove_plus(t_plus, 447479, 7)
w_minus = prove_minus(t_plus, 447479, 1, -3)
for lo, hi, digits in ((t_minus, t_plus, 123),
                      (w_minus, w_plus, 492)):
    require(hi - lo == 2)
    require(len(str(lo)) == len(str(hi)) == digits)
print("Twin pairs: 123 and 492 digits; all certificates passed.")
print("All seven selected child cells are strictly nested.")


# All finite root comparisons use exact integer powers.
EXPONENTS = (4, 6, 7, 12)
ROWS = [
    (2, 4, (2, 2, 3, 8)),
    (11, 12, (2, 3, 4, 11)),
    (5, 7, (2, 3, 5, 15)),
    (17, 12, (2, 4, 5, 17)),
    (5, 6, (2, 5, 6, 25)),
    (7, 7, (3, 5, 7, 28)),
    (29, 12, (3, 5, 7, 29)),
]
for p, d, outputs in ROWS:
    for c, m in zip(EXPONENTS, outputs):
        require(m ** d <= p ** c < (m + 1) ** d)
for c, m in zip(EXPONENTS[:3], (3, 5, 7)):
    require(m ** 12 <= 29 ** c < 30 ** c < (m + 1) ** 12)
require(P0 + 1 < 2 * 29 ** 21 < 30 ** 21)
require(30 * 3 ** 12 < 4 ** 12)
require(29 * 10 ** 12 > 13 ** 12)
for c, q in zip(range(7, 15), (5, 5, 7, 7, 11, 11, 11, 13)):
    require(3 ** c <= q ** 6 < 4 ** c)
require(4 ** 15 > 2 ** 6 * 3 ** 15)
print("Finite minima and all displayed root comparisons passed.")

R = factorial(252)
E1 = 252 * R
require(2 * 10 ** 497 < R < 21 * 10 ** 496)
require(len(str(R)) == 498 and len(str(E1)) == 500)
require(E1 > 5 * 10 ** 499)
require(10 ** 30 < P0 < 10 ** 31)
# Bounds for the number of decimal digits of the next greedy prime.
require(30 * R > 6 * 10 ** 498)
require(31 * R + 2 < 7 * 10 ** 498)
require(252 * 2 ** 251 < 10 ** 79)
require(252 * 10 ** 495 + 79 < 10 ** 498)
require(252 ** 252 > 10 ** 605)
require(252 * 10 ** 601 + 79 < 10 ** 604)
print("Factorial sizes and decimal-boundary inequalities passed.")

def small_prime(n):
    return n >= 2 and all(n % d for d in range(2, isqrt(n) + 1))

small_primes = [q for q in range(2, 563) if small_prime(q)]
for q in small_primes:
    require(R % (q - 1) == 0)
    require(pow(P0, R, q) == 1)
for h in range(1, 562):
    q = next(q for q in small_primes if (h + 1) % q == 0)
    require((pow(P0, R, q) + h) % q == 0)
print("All 561 offsets in the composite block are certified.")

\end{lstlisting}


\clearpage
\lstset{language=Python,basicstyle=\ttfamily\fontsize{8.5}{10}\selectfont,
frame=none,numbers=left,xleftmargin=14pt,xrightmargin=0pt,
breaklines=true,columns=fullflexible,keepspaces=true,showstringspaces=false}

\section{Reproducing the finite cubic verification}\label{tw:app:cubic-verification}
The following Python 3 code uses only the standard library. It regenerates
the entire finite family and finds periodic modular obstructions; no
stored certificate file, floating-point root approximation, or external
primality package is required. The reported count is the size of the exact coefficient catalogue
proved in the main text; the verification uses its integer inequalities
directly, rather than numerically approximating its roots.

\begingroup
\lstset{language=Python,basicstyle=\ttfamily\fontsize{8.5}{10}\selectfont,
breaklines=true,columns=fullflexible,keepspaces=true,showstringspaces=false,
frame=single,rulecolor=\color{black!20},xleftmargin=4pt,xrightmargin=4pt}

\begin{lstlisting}
"""Self-contained exact verifier: cubic Pisot shadows beta < 100.
Python 3 standard library only; no primality or real-root oracle is used.
"""
from collections import Counter

BOUND = 100
PRIMES = (2, 3, 5, 7, 11, 13, 17, 19, 23, 29, 31, 37, 41, 43, 47)
EXPECTED = {2:134584, 3:27162, 5:27780, 7:14024, 11:4856,
            13:2501, 17:1163, 19:544, 23:492, 29:182,
            31:142, 37:86, 41:35, 43:14, 47:4}

def polynomials(B):
    # Coefficients of X^3 - a X^2 + b X - c.
    for a in range(1, B):
        # Two negative small roots.
        for d in range(1, 2*B):
            for c in range(1, B):
                D = a*a*d*d + 4*d**3 - 4*a**3*c - 27*c*c - 18*a*d*c
                if (d <= 2*a+2 and d <= a+c and D > 0
                        and B**3 - a*B*B - d*B - c > 0):
                    yield (a, -d, c)
        # Opposite signs, negative root dominant in modulus.
        for d in range(1, a):
            for c in range(1, a-d+1):
                yield (a, -d, -c)

def step(s, p):
    a, b, c = s
    return ((a**3-3*a*b+3*c) % p,
            (b**3-3*a*b*c+3*c*c) % p, pow(c, 3, p))

def certificate(f, p):
    s = tuple(x % p for x in f)
    seen, states = {}, []
    while s not in seen:
        seen[s] = len(states)
        states.append(s)
        s = step(s, p)
    mu, period = seen[s], len(states)-seen[s]
    for j in range(mu, len(states)):
        if states[j][0] == 0:
            return mu, period, j
    return None

def run():
    remaining = list(polynomials(BOUND))
    assert len(remaining) == 213569
    assert len(set(remaining)) == len(remaining)
    counts = Counter()
    max_period = max_mu = 0
    for p in PRIMES:
        assert all(p % d for d in range(2, p) if d*d <= p)
        survivors = []
        for f in remaining:
            cert = certificate(f, p)
            if cert is None:
                survivors.append(f)
            else:
                mu, period, zero = cert
                counts[p] += 1
                max_period, max_mu = max(max_period, period), max(max_mu, mu)
        remaining = survivors
    assert not remaining
    assert dict(counts) == EXPECTED
    assert (max_period, max_mu) == (1148, 3)
    small = list(polynomials(4))
    assert small == [(2,-1,-1), (3,-1,-1), (3,-1,-2), (3,-2,-1)]
    print('PASS: 213569 polynomials; all certified by primes <= 47.')
    print('First-witness counts:', dict(counts))
    print('Maximum period =', max_period, '; maximum preperiod =', max_mu)

if __name__ == '__main__':
    run()

\end{lstlisting}
\endgroup

\clearpage
\lstset{language=Python,basicstyle=\ttfamily\fontsize{8.5}{10}\selectfont,
frame=none,numbers=left,xleftmargin=14pt,xrightmargin=0pt,
breaklines=true,columns=fullflexible,keepspaces=true,showstringspaces=false}

\section{Exact certificate for the Type-C improvement}
\label{f13:app:typec-code}
The following Python 3 code uses only the standard library. It enumerates
complete finite residue cycles, checks every survivor claimed in the
proof, and verifies the rounded cutoffs with rational inequalities.
The assertion that the four small Pisot candidates are exhaustive is
proved analytically in the main text, not inferred from this computation.
\begingroup
\lstset{language=Python,basicstyle=\ttfamily\footnotesize,
breaklines=true,columns=fullflexible,keepspaces=true,
showstringspaces=false,frame=single,rulecolor=\color{gray!40},
aboveskip=8pt,belowskip=8pt}
\begin{lstlisting}
"""Exact finite-field certificate for the strengthened Type-C cutoff.
Python standard library only. No numerical root or primality oracle.
"""
from math import lcm
from fractions import Fraction as Q

TABLE = {
 (0,-1,1): [(2,7),(41,1723),(131,17293),(347,173),(397,52669)],
 (1,-1,1): [(5,31),(31,331),(59,3541),(71,5113),(89,8011)],
 (2,1,1): [(2,7),(41,1723),(71,5113),(131,17293),(347,173),(409,55897)]
}
EXPECTED = [1900442940636401,1488257702170843,10312502210725960469]

def trace_period(f,p):
 a,b,c=f
 assert c==1
 assert all(p%d for d in range(2,p) if d*d<=p)
 v=(1,0,0); traces=[]
 while True:
  traces.append((3*v[0]+a*v[1]+(a*a-2*b)*v[2])%p)
  v=(c*v[2]%p,(v[0]-b*v[2])%p,(v[1]+a*v[2])%p)
  if v==(1,0,0):return traces

def surviving_parameters(f,p,g):
 tr=trace_period(f,p); t=len(tr); status={}
 for start in range(t):
  if start in status:continue
  z=start; seen={}; seq=[]
  while z not in seen and z not in status:
   seen[z]=len(seq);seq.append(z);z=(3*z+2//g)%t
  good=status[z] if z in status else all(tr[w]!=0 for w in seq[seen[z]:])
  for w in seq:status[w]=good
 # g=1: parameter j=r, initial exponent 3r-1.
 # g=2: parameter j=(r-1)/2, initial exponent 3j+1.
 good=[j for j in range(t) if status[(3*j-1 if g==1 else 3*j+1)%t]]
 return t,good

def verify():
 moduli=[]
 for f,rows in TABLE.items():
  for p,t in rows:
   assert surviving_parameters(f,p,1)==(t,[0]),(f,p)
  M=lcm(*(t for p,t in rows));moduli.append(M)
  assert M>400300000000000
  print('g=1:',f,'forces r divisible by',M)
 assert moduli==EXPECTED
 for f,p,g,t in [((0,-1,1),7,2,48),((1,-1,1),7,2,48),
                 ((1,0,1),2,1,7),((1,0,1),2,2,7)]:
  assert surviving_parameters(f,p,g)==(t,[]),(f,p,g)
  print('No surviving parameter:',f,'p=',p,'g=',g)
 # Rational enclosures: exp(32.76) and log(2).
 x=Q(819,25);term=Q(1);s=term
 for j in range(1,201):term*=x/j;s+=term
 first=term*x/201
 upper=s+first/(1-x/202)
 assert upper<168845000000000
 lower=2*sum((Q(1,(2*j+1)*3**(2*j+1)) for j in range(12)),Q(0))
 assert lower>Q(693147,1000000)
 for D,R in [(1,81200000000000),(2,162400000000000)]:
  assert D*(168845000000000+1)/(3*Q(693147,1000000))+Q(1,3)<R
 print('PASS: all exact certificates verified.')

if __name__=='__main__':verify()

\end{lstlisting}
\endgroup

\clearpage
\lstset{language=Python,basicstyle=\ttfamily\fontsize{8.5}{10}\selectfont,
frame=none,numbers=left,xleftmargin=14pt,xrightmargin=0pt,
breaklines=true,columns=fullflexible,keepspaces=true,showstringspaces=false}

\section{Exact certificate for the nonreturning prime seed}
\label{ft:app:seed-verifier}

The following standalone Python program uses only the standard library.
It proves primality by exhaustive trial division, checks the residual
roots by substitution, and traverses every state in the exact modular
cycle used in Proposition~\ref{ft:prop:seed-nonreturn}. Thus the
nonreturn assertion concerns all subsequent indices, rather than a
finite tested prefix. The final additional assertion checks the integer
inequality in Proposition~\ref{ft:prop:nonunit}.


\begin{lstlisting}
"""Exact certificate for a nonreturning prime trace in Saito's cubic locus.

This refutes a proposed divisor-return lemma; it does not refute Saito's
infinite-compositeness problem.  No floating-point calculations or packages.
"""
from math import isqrt

def prime(n):
    return n >= 2 and all(n % d for d in range(2, isqrt(n) + 1))

def phi(t, b, p):
    return ((t**3 - 3*t*b - 3) % p,
            (b**3 + 3*t*b + 3) % p)

a, d, p = 211, 7, 9_398_359
assert prime(a) and prime(p)
assert a**3 + 3*a*d - 3 == p
B1 = -d**3 - 3*a*d + 3
assert B1 == -4771

# Three distinct roots of f mod p, all certified by substitution.
roots = [-1_584_786, -1_942_624, 3_527_621]
assert len({r % p for r in roots}) == 3
assert all((r**3-a*r*r-d*r+1) % p == 0 for r in roots)

# An explicit period after the seed.  The full Newton state, not merely
# its trace, returns, so this is a finite proof about every later index.
L = 85_712
s = (0, B1 % p)
first = phi(*s, p)
s = first
cycle = set()
for _ in range(L):
    assert s[0] != 0
    assert s not in cycle
    cycle.add(s)
    s = phi(*s, p)
assert s == first
assert len(cycle) == L

# The example is not a counterexample to infinite compositeness.
T2 = p**3 - 3*p*B1 - 3
assert T2 == 7*31*3_825_571_796_264_257_579

# Rational sign bounds used to certify Saito's conjugate inequalities.
# f(-1/10)<0<f(0); beta<212; and 212^17<10^40.
assert -1 - 10*a + 100*d + 1000 < 0
assert -1-a+d+1 < 0 < 1
assert 1-a-d+1 < 0
assert -d*a+1 < 0 < (a+1)**3-a*(a+1)**2-d*(a+1)+1
assert 212**17 < 10**40

print("PASS: prime seed p occurs exactly once, at index 1;")
print("the 85,712-state certified tail has no zero trace modulo p.")

# Additional exact inequality for the nonunit example.
assert 39**40 * 10008**17 < 2000**40
\end{lstlisting}

\clearpage
\phantomsection\addcontentsline{toc}{part}{References}
\begingroup\footnotesize
\begin{thebibliography}{99}

\bibitem{AlkauskasDubickas2004}
G. Alkauskas and A. Dubickas,
\emph{Prime and composite numbers as integer parts of powers},
Acta Math. Hungar. \textbf{105} (2004), no.~3, 249--256.
The bibliographic data and the cited prime-constant result are also
recorded in \cite{SaitoMills,SaitoTakeda2022}.

\bibitem{Baker}
A. Baker, \emph{Transcendental Number Theory}, Cambridge University Press, 1975.

\bibitem{BellDranishnikov}
G. Bell and A. Dranishnikov, \emph{Asymptotic dimension in B\k{e}dlewo},
arXiv:math/0507570.
\url{https://arxiv.org/abs/math/0507570}.

\bibitem{Bertrand}
D.~Galvin,
\emph{Erd\H{o}s's proof of Bertrand's postulate},
lecture notes, University of Notre Dame, May 1, 2015.
\url{https://academicweb.nd.edu/~dgalvin1/pdf/bertrand.pdf}.

\bibitem{BezivinRobba1989}
J.-P. B\'ezivin and P. Robba,
\emph{A new $p$-adic method for proving irrationality and transcendence results},
Ann. of Math. (2) \textbf{129} (1989), no.~1, 151--160.
\href{https://doi.org/10.2307/1971488}{doi:10.2307/1971488}.

\bibitem{BHP}
R. C. Baker, G. Harman, and J. Pintz,
\emph{The difference between consecutive primes. II},
Proc. London Math. Soc. (3) \textbf{83} (2001), 532--562.
\href{https://doi.org/10.1112/plms/83.3.532}{doi:10.1112/plms/83.3.532}.

\bibitem{Brumer1967}
A. Brumer, \emph{On the units of algebraic number fields},
Mathematika \textbf{14} (1967), 121--124.
\href{https://doi.org/10.1112/S0025579300003703}{doi:10.1112/S0025579300003703}.

\bibitem{Carella}
N. A. Carella, \emph{Euler Totient Function And The Largest Integer
Function Over The Shifted Primes}, arXiv:2105.00790v2 (2021).
\url{https://arxiv.org/abs/2105.00790}.
The claim is discussed critically in Section~\ref{f13:sec:quotient-density};
it is not used as an input.

\bibitem{CarrOSullivan2009}
R. Carr and C. O'Sullivan,
\emph{On the linear independence of roots},
Int. J. Number Theory \textbf{5} (2009), no.~1, 161--171.
\href{https://doi.org/10.1142/S1793042109002018}{doi:10.1142/S1793042109002018}.
\href{https://arxiv.org/abs/math/0703522}{arXiv:math/0703522}.

\bibitem{CorvajaZannier2004}
P. Corvaja and U. Zannier,
\emph{On the rational approximations to the powers of an algebraic
number: solution of two problems of Mahler and Mend\`es France},
Acta Math. \textbf{193} (2004), 175--191.

\bibitem{CullyHugill2023}
M. Cully-Hugill,
\emph{Primes between consecutive powers},
J. Number Theory \textbf{247} (2023), 100--117.
\href{https://doi.org/10.1016/j.jnt.2022.12.002}{doi:10.1016/j.jnt.2022.12.002}.
\href{https://arxiv.org/abs/2107.14468v3}{arXiv:2107.14468v3}.

\bibitem{Dubickas}
A. Dubickas, \emph{Transcendency of some constants related to integer
sequences of polynomial iterations}, Ramanujan J. \textbf{57} (2022),
569--581.
\href{https://doi.org/10.1007/s11139-021-00428-5}{doi:10.1007/s11139-021-00428-5}.

\bibitem{Dubickas2006}
A. Dubickas, \emph{Arithmetical properties of powers of algebraic numbers},
Bull. London Math. Soc. \textbf{38} (2006), no.~1, 70--80.
\href{https://doi.org/10.1112/S0024609305017728}{doi:10.1112/S0024609305017728}.

\bibitem{EdgarMiller2003}
G. A. Edgar and C. Miller,
\emph{Borel subrings of the reals},
Proc. Amer. Math. Soc. \textbf{131} (2003), no.~4, 1121--1129.
\href{https://doi.org/10.1090/S0002-9939-02-06653-4}{doi:10.1090/S0002-9939-02-06653-4}.

\bibitem{Elsholtz2020}
C. Elsholtz, \emph{Unconditional prime-representing functions, following Mills},
2020, arXiv:2004.01285v1.
\url{https://arxiv.org/abs/2004.01285}.

\bibitem{HIKK2026}
S. Hahn, D. Ismailescu, G. Kim, and M. Kim,
\emph{Explicit algebraic numbers all whose integer parts of powers are
always composite}, arXiv:2608.05309v1, 5 August 2026.
\url{https://arxiv.org/pdf/2608.05309}.

\bibitem{KanadoSaito}
Y. Kanado and K. Saito, \emph{A system of certain linear Diophantine
equations on analogs of squares}, arXiv:2205.12226v4 (2023).
\url{https://arxiv.org/pdf/2205.12226}.

\bibitem{KanadoSaitoNormality}
Y. Kanado and K. Saito, \emph{Normality of algebraic numbers and the Riemann zeta function},
arXiv:2412.02337v2.
\url{https://arxiv.org/html/2412.02337v2}.

\bibitem{KelleyMeka}
Z. Kelley and R. Meka, \emph{Strong bounds for 3-progressions},
arXiv:2302.05537v6.
\url{https://arxiv.org/abs/2302.05537v6}.

\bibitem{KMN2019}
A. Kulkarni, N. M. Mavraki, and K. D. Nguyen,
\emph{Algebraic approximations to linear combinations of powers:
An extension of results by Mahler and Corvaja--Zannier},
Trans. Amer. Math. Soc. \textbf{371} (2019), 3787--3804.
The precise form used here is reproduced in \cite[Theorem~1.2]{Kumar2023}.

\bibitem{KobayashiSaitoTakeda2023}
H. Kobayashi, K. Saito, and W. Takeda,
\emph{Transcendence of values of the iterated exponential function at algebraic points},
J. Integer Seq. \textbf{26} (2023), Article 23.3.3.
\url{https://cs.uwaterloo.ca/journals/JIS/VOL26/Takeda/tak6.pdf}.

\bibitem{Kumar2023}
V. Kumar, \emph{On the transcendence of growth constants associated with
polynomial recursions}, arXiv:2207.08614v3, 19 December 2023.
\url{https://arxiv.org/html/2207.08614v3}.
All theorem and equation numbers in this note refer to this version.

\bibitem{Mahler1957}
K. Mahler, \emph{On the fractional parts of the powers of a rational number. II},
Mathematika \textbf{4} (1957), 122--124.
The precise lower bound used here is reproduced in \cite[Theorem~2.5]{SaitoMills}.

\bibitem{Maksoud2023}
A. Maksoud, \emph{On the rank of Leopoldt's and Gross's regulator maps},
Documenta Mathematica \textbf{28} (2023), 1441--1471.
\href{https://doi.org/10.4171/DM/935}{doi:10.4171/DM/935}.
Theorem~2.1 states the Baker--Brumer theorem used here.

\bibitem{Matomaki2010}
K. Matom\"aki, \emph{Prime-representing functions},
Acta Math. Hungar. \textbf{128} (2010), no.~4, 307--314.
\url{https://users.utu.fi/ksmato/wp-content/uploads/sites/1432/2025/04/Primerepfunc.pdf}.

\bibitem{Mills1947}
W. H. Mills, \emph{A prime-representing function},
Bull. Amer. Math. Soc. \textbf{53} (1947), 604.
\href{https://doi.org/10.1090/S0002-9904-1947-08849-2}{doi:10.1090/S0002-9904-1947-08849-2}.

\bibitem{MontgomeryVaughan1973}
H. L. Montgomery and R. C. Vaughan, \emph{The large sieve},
Mathematika \textbf{20} (1973), no.~2, 119--134.
\href{https://doi.org/10.1112/S0025579300004708}{doi:10.1112/S0025579300004708}.

\bibitem{Mycielski1964}
J. Mycielski,
\emph{Independent sets in topological algebras},
Fund. Math. \textbf{55} (1964), no.~2, 139--147.
\href{https://doi.org/10.4064/fm-55-2-139-147}{doi:10.4064/fm-55-2-139-147}.

\bibitem{Polymath2012}
D. H. J. Polymath, \emph{A new proof of the density Hales--Jewett theorem},
Ann. of Math. \textbf{175} (2012), 1283--1327.
\url{https://annals.math.princeton.edu/2012/175-3/p06}.

\bibitem{PongsriiamSaito2023}
P. Pongsriiam and K. Saito, \emph{Palindromes and Antipalindromes in Short Intervals},
J. Integer Seq. \textbf{26} (2023), Article 23.9.5.
\url{https://cs.uwaterloo.ca/journals/JIS/VOL26/Saito/saito1.pdf}.

\bibitem{SaitoAvoidance}
K. Saito, \emph{New bounds for dimensions of a set uniformly avoiding multi-dimensional arithmetic progressions},
arXiv:1910.13071v1.
\url{https://arxiv.org/abs/1910.13071v1}.

\bibitem{SaitoBalanced}
K. Saito, \emph{Finiteness of solutions to linear Diophantine equations on Piatetski--Shapiro sequences},
arXiv:2306.17813v2.
\url{https://arxiv.org/abs/2306.17813v2}.

\bibitem{SaitoDimension2021}
K. Saito, \emph{Prime-representing functions and Hausdorff dimension},
Acta Math. Hungar. \textbf{165} (2021), no.~1, 203--217.
\href{https://doi.org/10.1007/s10474-021-01170-6}{doi:10.1007/s10474-021-01170-6}.
Public text and numbering used: arXiv:2102.04038v1.
\url{https://arxiv.org/html/2102.04038v1}.

\bibitem{SaitoFractalProfiles}
K. Saito, \emph{New fractal dimensions and some applications to arithmetic patches},
arXiv:1802.08839v4.
\url{https://arxiv.org/abs/1802.08839v4}.

\bibitem{SaitoIntervals}
K. Saito, \emph{Intervals without primes near an iterated linear recurrence
sequence}, arXiv:2504.14968v1 (2025).
\url{https://arxiv.org/html/2504.14968v1}.

\bibitem{SaitoMills}
K. Saito, \emph{Mills' constant is irrational},
Mathematika \textbf{71} (2025), e70027.
\href{https://doi.org/10.1112/mtk.70027}{doi:10.1112/mtk.70027}.
Public text used: arXiv:2404.19461v2.
\url{https://arxiv.org/html/2404.19461v2}.

\bibitem{SaitoQuotients}
K. Saito, Y. Suzuki, W. Takeda, and Y. Yoshida,
\emph{Some remarks on the $[x/n]$-sequence},
Research in Number Theory \textbf{11} (2025), Article 50.
Question numbering and public text: arXiv:2312.15642v1.
\url{https://arxiv.org/pdf/2312.15642}.

\bibitem{SaitoTakeda2022}
K. Saito and W. Takeda,
\emph{Topological properties and algebraic independence of sets of
prime-representing constants},
Mathematika \textbf{68} (2022), no.~2, 429--453.
\href{https://doi.org/10.1112/mtk.12142}{doi:10.1112/mtk.12142}.
\url{https://arxiv.org/pdf/2112.14383v2}.

\bibitem{SaitoVariants}
K. Saito, \emph{Transcendency of variants of Mills' constant},
Ramanujan J. \textbf{70} (2026), article 69.
\href{https://doi.org/10.1007/s11139-026-01443-0}{doi:10.1007/s11139-026-01443-0}.
Public text and numbering used: arXiv:2508.16068v3, 7 December 2025.
\url{https://arxiv.org/html/2508.16068v3}.

\bibitem{SaitoYoshidaCurved2019}
K. Saito and Y. Yoshida, \emph{Arithmetic progressions in the graphs of slightly curved sequences},
J. Integer Seq. \textbf{22} (2019), Article 19.2.1.
\url{https://cs.uwaterloo.ca/journals/JIS/VOL22/Yoshida/yoshida4.pdf}.

\bibitem{SaitoYoshidaPatterns2021}
K. Saito and Y. Yoshida, \emph{Distributions of finite sequences represented by polynomials in Piatetski--Shapiro sequences},
J. Number Theory (2021); source version arXiv:2006.13930v2.
\url{https://arxiv.org/abs/2006.13930v2}.

\bibitem{TaoZiegler2014}
T. Tao and T. Ziegler, \emph{Narrow progressions in the primes}, arXiv:1409.1327v2.
\url{https://arxiv.org/abs/1409.1327v2}.

\bibitem{Toth2017}
L. T\'oth, \emph{A variation on Mills-like prime-representing functions},
J. Integer Seq. \textbf{20} (2017), Article 17.9.8.
\url{https://cs.uwaterloo.ca/journals/JIS/VOL20/Toth2/toth32.pdf}.

\bibitem{WagnerZiegler2021}
S. Wagner and V. Ziegler,
\emph{Irrationality of growth constants associated with polynomial recursions},
J. Integer Seq. \textbf{24} (2021), Article 21.1.6.
\url{https://cs.uwaterloo.ca/journals/JIS/VOL24/Wagner/wagner4.pdf}.
\end{thebibliography}
\endgroup
\end{document}
